\section{Content-State Refresh and v4.7 Handoff Preflight}
\label{sec:v460-content-state-refresh-v47-handoff}
This release-final content refresh is a reader-routing layer.  It promotes no
new theorem and assigns no new DOI.  Its purpose is to align the visible front
routing, nomenclature, and handoff language with the v4.6.0 endpoint: bounded
readout is conservative inside the declared grammar; T2 remains an all-readout
scope boundary; and the next internal research step is the MM/HC exterior--mixed
field-equation program.

\begin{proposition}[Safe refresh invariant]
The content-state refresh preserves the v4.6.0 release-final theorematic content.
It updates terminology, routing, and handoff orientation only.  In particular,
it does not assert T2, does not assert an unrestricted all-readout theorem, and
does not turn the field-equation reset template into an empirical validation
claim.
\end{proposition}


\begin{theorem}[Release-final conservativity]
The v4.6.0 release-final integrated source is conservative over the inherited
v4.5.x active core plus the v4.6.0 bounded-readout packets.  It integrates the
bounded readout theorem, the T2 scope boundary, the MM/HC exterior--mixed
bookkeeping, the QAM tag extension, and the structure-first field-equation reset
under the assigned Version DOI \href{https://doi.org/10.5281/zenodo.20117361}{10.5281/zenodo.20117361}, with public release-final metadata normalized, without opening
an empirical validation gate, and without claiming T2 closure.
\end{theorem}


\section{Theory-Lift Gate and MM/HC Exterior--Mixed Field-Equation Program}
\label{sec:v470-r01-theory-lift-gate-extmix-program}
This section opens the first internal working layer after the v4.6.0
content-state refresh.  It is a program layer, not a release-final theorem.  Its
role is to decide which audited hypotheses may now be treated as controlled
objects in the next proof-engine pass and to specify the admissible grammar for
MM/HC exterior--mixed field-equation candidates.

\subsection{Theory-Lift gate}
\label{subsec:v470-r01-theory-lift-gate}

\begin{definition}[Theory-Lift gate]
\label{def:v470-r01-theory-lift-gate}
A \emph{Theory-Lift gate} is a finite promotion test applied to a prior
hypothesis, heuristic, notation package, or diagnostic.  It has four output
classes:
\[
\mathsf{def},\qquad \mathsf{lem},\qquad \mathsf{prog},\qquad \mathsf{hold}.
\]
The class \(\mathsf{def}\) means that the object may be introduced as a
conservative definition.  The class \(\mathsf{lem}\) means that the object may
be used as a local lemma or corollary inside its declared grammar.  The class
\(\mathsf{prog}\) means that the object is not yet a theorem but may drive a
finite test program or work queue.  The class \(\mathsf{hold}\) means that the
object remains a hypothesis or scope boundary and must not be promoted.
\end{definition}

\begin{definition}[v4.7.0-r01 promotion ledger]
\label{def:v470-r01-promotion-ledger}
The first promotion ledger consists of the following guarded lifts:
\[
\begin{array}{lll}
\mathsf{BRead}_{\mathrm{bounded}} &\mapsto& \mathsf{lem}\text{-level bounded-readout calculus},\\
(\mathcal M_{\mathrm{bd}},\mathcal M_{\mathrm{ext}},\mathcal M_{\mathrm{mix}}),
(\mathcal H_{\mathrm{bd}},\mathcal H_{\mathrm{ext}},\mathcal H_{\mathrm{mix}})
&\mapsto& \mathsf{def}\text{-level channel split},\\
\text{two-sided structural tensor notation} &\mapsto& \mathsf{def}\text{-level admissibility grammar},\\
\text{shared MM/HC leakage heuristic} &\mapsto& \mathsf{prog}\text{-level leakage diagnostic},\\
\text{inverse lens-fit calibration} &\mapsto& \mathsf{prog}\text{-level finite witness-fit method},\\
\mathfrak E_{\mathrm{decl}} &\mapsto& \mathsf{prog}\text{-level field-equation admissibility template}.
\end{array}
\]
The all-readout object \(\Theta^{\mathrm{all}}\), T2, unrestricted exterior
collar closure, and global physical-spectrum claims are assigned class
\(\mathsf{hold}\).
\end{definition}

\begin{proposition}[Promotion conservativity]
\label{prop:v470-r01-promotion-conservativity}
The promotion ledger of Definition~\ref{def:v470-r01-promotion-ledger} is
conservative over the v4.6.0 release-final baseline.  It promotes only guarded
object-language, local calculus, and test-program surfaces; it does not promote
any all-readout or global physical theorem.
\end{proposition}

\begin{proof}
Each promoted item is either a bounded object already isolated by the v4.6.0
readout surface, a notation package whose same-variance comparisons require an
explicit comparison tensor, or a finite diagnostic program that returns a
residual status rather than a closure theorem.  The objects assigned class
\(\mathsf{hold}\) include precisely the missing all-readout and unrestricted
exterior data.  Hence no promotion crosses the v4.6.0 T2 scope boundary.
\end{proof}

\subsection{Exterior--mixed sidedness grammar}
\label{subsec:v470-r01-extmix-sidedness-grammar}

\begin{definition}[Sided MM/HC channel split]
\label{def:v470-r01-sided-mm-hc-channel-split}
The MM/HC field-equation program is split into bounded, exterior, and mixed
channels:
\[
\mathcal M=\mathcal M_{\mathrm{bd}}\oplus\mathcal M_{\mathrm{ext}}
\oplus\mathcal M_{\mathrm{mix}},
\qquad
\mathcal H=\mathcal H_{\mathrm{bd}}\oplus\mathcal H_{\mathrm{ext}}
\oplus\mathcal H_{\mathrm{mix}}.
\]
The bounded channels are inherited from the v4.6.0 declared grammar.  The
exterior channels record readouts or selectors that are not covered by the
bounded surface.  The mixed channels record terms whose interpretation requires
both a bounded and an exterior interface.  No mixed term is admissible unless its
left/readout side and right/structure side are explicitly declared.
\end{definition}

\begin{definition}[Exterior--mixed admissible tensor term]
\label{def:v470-r01-admissible-extmix-term}
A candidate term
\[
{}^{A}{}_{B}\mathcal T^{i}{}_{j}
\]
is \emph{exterior--mixed admissible} when it satisfies all of the following
conditions.
\begin{enumerate}[label=(\alph*)]
    \item The front indices \({}^{A}{}_{B}\) are declared as readout, frame, or
    language interfaces.
    \item The back indices \({}^{i}{}_{j}\) are declared as structure, action,
    or field interfaces.
    \item Every upper--lower contraction is an ordinary action/readout.
    \item Every same-variance comparison is mediated by an explicitly declared
    comparison tensor or selector, for example
    \(G_{\alpha\beta}X^{\alpha}Y^{\beta}\) or
    \(\chi_{mn}R^mS^n\).
    \item The term has a bounded restriction, an exterior residual register, and
    a mixed leakage register.
\end{enumerate}
Raw same-variance contraction and undeclared bounded/exterior identification are
not admissible.
\end{definition}

\begin{lemma}[Sidedness prevents silent closure promotion]
\label{lem:v470-r01-sidedness-prevents-silent-promotion}
An exterior--mixed admissible tensor term cannot by itself promote a bounded
readout statement to an all-readout theorem.
\end{lemma}

\begin{proof}
By Definition~\ref{def:v470-r01-admissible-extmix-term}, every mixed or
same-variance expression carries an explicit comparison tensor, selector, or
residual register.  Therefore an exterior contribution remains visible as
comparison data, selector data, or leakage/debt data.  It is not erased by
Einstein summation and cannot be mistaken for an unrestricted readout identity.
\end{proof}

\subsection{Field-equation admissibility template}
\label{subsec:v470-r01-field-equation-admissibility-template}

\begin{definition}[Exterior--mixed field-equation candidate]
\label{def:v470-r01-extmix-field-equation-candidate}
A first MM/HC exterior--mixed field-equation candidate is a package
\[
\mathfrak E^{\mathrm{mix}}_{\mathrm{MM/HC}}
=
\left(
\mathfrak E_{\mathrm{bd}},
\mathfrak E_{\mathrm{ext}},
\mathfrak E_{\mathrm{mix}},
\mathfrak L_{\mathrm{mix}},
\mathfrak R_{\mathrm{fit}}
\right),
\]
where \(\mathfrak E_{\mathrm{bd}}\) is the bounded inherited equation surface,
\(\mathfrak E_{\mathrm{ext}}\) is the exterior residual surface,
\(\mathfrak E_{\mathrm{mix}}\) is the declared mixed MM/HC coupling surface,
\(\mathfrak L_{\mathrm{mix}}\) is the leakage/debt register, and
\(\mathfrak R_{\mathrm{fit}}\) is the finite witness-fit register inherited from
the inverse lens-fit calibration method.
\end{definition}

\begin{definition}[Admissibility residual]
\label{def:v470-r01-admissibility-residual}
The first residual attached to a candidate of
Definition~\ref{def:v470-r01-extmix-field-equation-candidate} is
\[
\rho^{\mathrm{mix}}_{\mathrm{MM/HC}}
=
\rho_{\mathrm{bd}}
\oplus
\rho_{\mathrm{ext}}
\oplus
\rho_{\mathrm{mix}}
\oplus
\rho_{\mathrm{leak}}
\oplus
\rho_{\mathrm{fit}}.
\]
Its comparison size may be written in two-sided tensor notation as
\[
\left\|\rho^{\mathrm{mix}}_{\mathrm{MM/HC}}\right\|^2_{G}
=
G_{\alpha\beta}\rho^{\alpha}\rho^{\beta},
\]
where the expression is read only as an explicitly mediated active-channel
comparison.  It is not a raw same-variance contraction.
\end{definition}

\begin{proposition}[First admissibility gate]
\label{prop:v470-r01-first-admissibility-gate}
A candidate \(\mathfrak E^{\mathrm{mix}}_{\mathrm{MM/HC}}\) passes the first
v4.7.0-r01 admissibility gate only if:
\begin{enumerate}[label=(\roman*)]
    \item its bounded restriction agrees with the v4.6.0 bounded-readout surface;
    \item every exterior contribution is represented in
    \(\rho_{\mathrm{ext}}\);
    \item every mixed MM/HC contribution is represented in
    \(\rho_{\mathrm{mix}}\) or discharged by an explicit selector/comparison;
    \item every leakage or debt contribution is represented in
    \(\rho_{\mathrm{leak}}\);
    \item every witness-fit claim is routed through
    \(\rho_{\mathrm{fit}}\) and remains finite/bounded;
    \item the candidate makes no T2, all-readout, or global physical-spectrum
    claim.
\end{enumerate}
\end{proposition}

\begin{proof}
The six clauses are the minimal separation obligations inherited from the
v4.6.0 endpoint.  The bounded surface may be reused, but exterior and mixed data
must remain visible as residual, comparison, selector, leakage, or finite-fit
registers.  Therefore a candidate that passes the gate is a well-formed test
object, not an unrestricted closure theorem.
\end{proof}

\subsection{Finite v4.7.0-r01 work queue}
\label{subsec:v470-r01-finite-work-queue}

\begin{definition}[First exterior--mixed work queue]
\label{def:v470-r01-first-work-queue}
The first work queue for the next internal passes consists of six tasks:
\[
\begin{array}{ll}
W_1 & \text{extract all candidate MM/HC exterior--mixed terms from the current object language},\\
W_2 & \text{assign two-sided sidedness data to each candidate term},\\
W_3 & \text{compute or bound }\rho_{\mathrm{bd}},\rho_{\mathrm{ext}},\rho_{\mathrm{mix}},\rho_{\mathrm{leak}},\rho_{\mathrm{fit}},\\
W_4 & \text{classify each term as admitted, selector-routed, leakage-open, blocked, or residual},\\
W_5 & \text{test whether any finite witness-fit lowers the mixed residual without relaxing exactness},\\
W_6 & \text{return a theory-lift update: definition, lemma/corollary, program, or hold}.
\end{array}
\]
The queue is finite at every declared pass and may be extended only by an
explicit audit ledger.
\end{definition}

\begin{corollary}[Permitted r02 direction]
\label{cor:v470-r01-permitted-r02-direction}
The next natural internal pass is a finite extraction/classification pass for
\(W_1\)--\(W_4\), followed only then by a witness-fit pass for \(W_5\).  In
particular, a direct T2 or all-readout closure attempt is not the next admissible
move of this program.
\end{corollary}

\begin{remark}[Summary]
The v4.7.0-r01 layer therefore turns the earlier ``Theorie-Lift'' idea into a
formal gate: promote only what has a declared grammar and a finite residual
interface; keep every all-readout or global physical claim outside the promoted
layer until a separate witness family exists.
\end{remark}


\section{Finite MM/HC Exterior--Mixed Term Extraction and Classification}
\label{sec:v470-r02-extmix-term-extraction-classification}
This section executes the r01 work queue through \(W_1\)--\(W_4\).  It is not a
new closure theorem.  It extracts a finite candidate family from the current
v4.7 program grammar, assigns two-sided sidedness data, records the residual
components that must remain visible, and classifies each candidate for the next
pass.

\subsection{Extraction frame}
\label{subsec:v470-r02-extraction-frame}

\begin{definition}[r02 exterior--mixed candidate]
\label{def:v470-r02-extmix-candidate}
An r02 exterior--mixed candidate is a package
\[
C_k=
\left(
\mathcal T_k,
\sigma_k^{L},
\sigma_k^{R},
\rho_k,
\kappa_k
\right),
\]
where \(\mathcal T_k\) is a declared MM/HC term, \(\sigma_k^{L}\) is its
left/readout sidedness, \(\sigma_k^{R}\) is its right/structure sidedness,
\(\rho_k\) is the residual tuple
\[
\rho_k=
(\rho_{\mathrm{bd}},\rho_{\mathrm{ext}},\rho_{\mathrm{mix}},
 \rho_{\mathrm{leak}},\rho_{\mathrm{fit}})_k,
\]
and \(\kappa_k\) is its classification status.  The status alphabet is
\[
\begin{array}{lll}
\mathsf A & \text{admitted}, &
\mathsf S \; \text{selector-routed},\\
\mathsf L & \text{leakage-open}, &
\mathsf R \; \text{residual-to-fit},\\
\mathsf B & \text{blocked}. &
\end{array}
\]
\end{definition}

\begin{definition}[Extracted r02 candidate family]
\label{def:v470-r02-extracted-candidate-family}
The finite r02 extraction from the current v4.7 program grammar consists of
seven candidates:
\[
\begin{array}{ll}
C_1 &={}^{\mathrm{bd}}\mathcal T_{\mathcal M_{\mathrm{bd}}\leftrightarrow
\mathcal H_{\mathrm{bd}}}^{\mathrm{bd}},\\[2pt]
C_2 &={}^{\mathrm{ext}}\mathcal T_{\mathcal M_{\mathrm{ext}}\leftrightarrow
\mathcal H_{\mathrm{bd}}}^{\mathrm{bd}},\\[2pt]
C_3 &={}^{\mathrm{bd}}\mathcal T_{\mathcal M_{\mathrm{bd}}\leftrightarrow
\mathcal H_{\mathrm{ext}}}^{\mathrm{ext}},\\[2pt]
C_4 &={}^{\mathrm{mix}}\mathcal T_{\mathcal M_{\mathrm{mix}}\leftrightarrow
\mathcal H_{\mathrm{mix}}}^{\mathrm{mix}},\\[2pt]
C_5 &={}^{\mathrm{fit}}\mathcal T_{\mathfrak R_{\mathrm{fit}}}^{\mathrm{bd}},\\[2pt]
C_6 &={}^{\mathrm{all}}\mathcal T_{\Theta^{\mathrm{all}}}^{\mathrm{all}},\\[2pt]
C_7 &=G_{\alpha\beta}\rho^{\alpha}\rho^{\beta}.
\end{array}
\]
Here \(C_1\) is the inherited bounded MM/HC surface, \(C_2\) and \(C_3\) are
one-sided exterior comparison terms, \(C_4\) is the genuinely mixed MM/HC term,
\(C_5\) is the finite witness-fit interface, \(C_6\) is the forbidden
all-readout lift, and \(C_7\) is the explicit comparison-tensor grammar for
residual size.
\end{definition}

\begin{remark}[Scope of extraction]
\label{rem:v470-r02-extraction-scope}
The extraction is deliberately finite and local to the current v4.7 program
grammar.  It is not a global symbol-mining normalization of every historical
formula carried in the Companion.  Older OP, HC, MM, lens, and route fragments
remain provenance unless explicitly mapped into one of the candidates above.
\end{remark}

\subsection{Sidedness and residual ledger}
\label{subsec:v470-r02-sidedness-residual-ledger}

\begin{definition}[r02 symbolic residual ledger]
\label{def:v470-r02-symbolic-residual-ledger}
The r02 residual ledger is the following symbolic assignment:
\[
\begin{array}{c|c|c|c}
\text{candidate} & (\sigma^L,\sigma^R) &
(\rho_{\mathrm{bd}},\rho_{\mathrm{ext}},\rho_{\mathrm{mix}},
\rho_{\mathrm{leak}},\rho_{\mathrm{fit}}) & \kappa \\
\hline
C_1 & (\mathrm{bd},\mathrm{bd}) & (0,0,0,0,0) & \mathsf A \\
C_2 & (\mathrm{ext},\mathrm{bd}) & (0,\epsilon^{M}_{\mathrm{ext}},0,0,0) & \mathsf S \\
C_3 & (\mathrm{bd},\mathrm{ext}) & (0,\epsilon^{H}_{\mathrm{ext}},\epsilon^{H}_{\mathrm{mix}},0,0) & \mathsf S \\
C_4 & (\mathrm{mix},\mathrm{mix}) & (0,0,\epsilon^{\mathrm{MM/HC}}_{\mathrm{mix}},\lambda_{\mathrm{leak}},0) & \mathsf L \\
C_5 & (\mathrm{fit},\mathrm{bd}) & (0,0,\epsilon^{\mathrm{fit}}_{\mathrm{mix}},0,\rho_{\mathrm{fit}}) & \mathsf R \\
C_6 & (\mathrm{all},\mathrm{all}) & \bot & \mathsf B \\
C_7 & (\mathrm{cmp},\mathrm{res}) & (0,0,0,0,0)\text{ as grammar} & \mathsf A.
\end{array}
\]
The symbols \(\epsilon^{M}_{\mathrm{ext}}\), \(\epsilon^{H}_{\mathrm{ext}}\),
\(\epsilon^{H}_{\mathrm{mix}}\), \(\epsilon^{\mathrm{MM/HC}}_{\mathrm{mix}}\),
\(\lambda_{\mathrm{leak}}\), and \(\rho_{\mathrm{fit}}\) are not numerical
witnesses.  They are named residual slots that must be bounded, routed, or
absorbed in later passes.
\end{definition}

\begin{proposition}[r02 classification count]
\label{prop:v470-r02-classification-count}
The r02 classification count is
\[
N=7,\qquad
N_{\mathsf A}=2,
\quad
N_{\mathsf S}=2,
\quad
N_{\mathsf L}=1,
\quad
N_{\mathsf R}=1,
\quad
N_{\mathsf B}=1.
\]
Thus the current grammar has two admitted items, two selector-routed one-sided
exterior items, one leakage-open mixed item, one residual-to-fit item, and one
blocked all-readout item.
\end{proposition}

\begin{proof}
The count is obtained directly from Definition~\ref{def:v470-r02-symbolic-residual-ledger}.
The admitted items are \(C_1\) and \(C_7\).  The selector-routed items are
\(C_2\) and \(C_3\).  The leakage-open item is \(C_4\).  The residual-to-fit
item is \(C_5\).  The all-readout lift \(C_6\) is blocked because it has no
admissible bounded/exterior residual decomposition inside the current grammar.
\end{proof}

\subsection{Consequences for the next pass}
\label{subsec:v470-r02-next-pass-consequences}

\begin{lemma}[No silent all-readout lift]
\label{lem:v470-r02-no-silent-all-readout-lift}
No candidate in the r02 admitted or selector-routed classes yields T2 or an
all-readout closure theorem.
\end{lemma}

\begin{proof}
The admitted bounded candidate \(C_1\) is restricted to the inherited bounded
surface, and the admitted comparison grammar \(C_7\) only measures residual
channels through an explicit comparison tensor.  The selector-routed candidates
\(C_2\) and \(C_3\) carry named exterior residual slots.  Hence none of these
terms erases exterior data or upgrades bounded readout to all-readout identity.
\end{proof}

\begin{proposition}[r02 completion of \(W_1\)--\(W_4\)]
\label{prop:v470-r02-w1-w4-completion}
The r02 layer completes the first four tasks of the r01 work queue for the
current v4.7 grammar: candidate extraction \((W_1)\), sidedness assignment
\((W_2)\), residual registration \((W_3)\), and finite classification \((W_4)\).
It leaves witness-fit lowering \((W_5)\) and the next theory-lift update
\((W_6)\) open.
\end{proposition}

\begin{proof}
Definition~\ref{def:v470-r02-extracted-candidate-family} supplies the finite
candidate extraction.  Definition~\ref{def:v470-r02-symbolic-residual-ledger}
assigns sidedness, residual slots, and classification status.  Proposition~\ref{prop:v470-r02-classification-count}
records the finite classification count.  No witness-fit lowering inequality is
proved in this section, so \(W_5\) and \(W_6\) remain the next tasks rather than
completed tasks.
\end{proof}

\begin{corollary}[Permitted r03 direction]
\label{cor:v470-r02-permitted-r03-direction}
The next admissible internal pass is a selector/leakage bound pass for the
non-admitted candidates:
\[
C_2,
C_3,
C_4,
C_5.
\]
In operational terms, r03 should first test whether the selector-routed exterior
slots \(\epsilon^{M}_{\mathrm{ext}}\) and \(\epsilon^{H}_{\mathrm{ext}}\) can be
routed without creating \(\lambda_{\mathrm{leak}}\), and only then run a finite
witness-fit lowering test on \(C_5\).  A direct T2 or all-readout proof attempt
remains inadmissible at this stage.
\end{corollary}

\begin{remark}[r02 summary]
The r02 pass turns the v4.7 exterior--mixed program from a pure work queue into
a finite classified inventory.  It advances the program because the previously
amorphous MM/HC exterior--mixed problem is now split into admitted bounded
syntax, selector-routed exterior terms, one leakage-open mixed term, one
residual-to-fit interface, and one explicitly blocked all-readout lift.
\end{remark}


\clearpage
\section{v4.7.0-r03 Selector/Leakage Bound Pass}
\label{sec:v470-r03-selector-leakage-bound-pass}

\subsection{Purpose and inherited inventory}
\label{subsec:v470-r03-purpose-inventory}

The r02 layer classified the exterior--mixed MM/HC inventory into admitted,
selector-routed, leakage-open, residual-to-fit, and blocked items.  The present
r03 layer performs the next conservative pass on the four non-admitted but
still testable candidates
\[
C_2,
C_3,
C_4,
C_5.
\]
It does not alter the blocked all-readout candidate \(C_6\), and it does not
convert a bounded or selector-cleared statement into T2.

\begin{definition}[r03 selector/leakage test variables]
\label{def:v470-r03-selector-leakage-test-variables}
For a candidate \(C_i\) in the r02 residual ledger define the r03 test tuple
\[
\mathcal T^{\mathrm{sel/leak}}_{i}
=
\left(
S_i^{\mathrm{ext}},
D_i^{\mathrm{sel}},
\Lambda_i^{\mathrm{new}},
\Lambda_i^{\mathrm{pers}},
F_i^{\mathrm{next}}
\right).
\]
Here \(S_i^{\mathrm{ext}}\in\{0,1\}\) records whether the exterior selector is
admissible in the declared v4.7 grammar, \(D_i^{\mathrm{sel}}\) is the remaining
selector debt, \(\Lambda_i^{\mathrm{new}}\) is newly created leakage, \(\Lambda_i^{\mathrm{pers}}\)
is inherited persistent leakage, and \(F_i^{\mathrm{next}}\) marks whether the
candidate must be handed to the later witness-fit lowering pass.
\end{definition}

\begin{definition}[r03 admissible selector clearance]
\label{def:v470-r03-admissible-selector-clearance}
A selector-routed candidate \(C_i\) is called \emph{selector-cleared} in r03 if
\[
S_i^{\mathrm{ext}}=1,
\qquad
D_i^{\mathrm{sel}}=0,
\qquad
\Lambda_i^{\mathrm{new}}=0.
\]
This status is weaker than admission as a global theorem.  It only says that the
one-sided exterior slot is routed without generating new leakage inside the
current declared grammar.
\end{definition}

\subsection{Finite r03 ledger}
\label{subsec:v470-r03-finite-ledger}

\begin{definition}[r03 finite selector/leakage ledger]
\label{def:v470-r03-finite-selector-leakage-ledger}
The r03 ledger over the testable r02 non-admitted candidates is
\[
\begin{array}{c|c|c|c|c|c|c}
\text{candidate} & S_i^{\mathrm{ext}} & D_i^{\mathrm{sel}} &
\Lambda_i^{\mathrm{new}} & \Lambda_i^{\mathrm{pers}} & F_i^{\mathrm{next}} &
\kappa_{r03} \\
\hline
C_2 & 1 & 0 & 0 & 0 & 0 & \mathsf A_{\mathrm{sel}} \\
C_3 & 1 & 0 & 0 & 0 & 0 & \mathsf A_{\mathrm{sel}} \\
C_4 & 1 & 0 & 0 & \lambda_{\mathrm{leak}} & 0 & \mathsf L_{\mathrm{pers}} \\
C_5 & 1 & 0 & 0 & 0 & 1 & \mathsf R_{\mathrm{fit}}.
\end{array}
\]
Thus the one-sided exterior selector slots clear without new leakage, the mixed
MM/HC slot remains persistent-leakage open, and the fit slot is not decided until
the finite witness-fit lowering pass.
\end{definition}

\begin{proposition}[r03 selector clearance count]
\label{prop:v470-r03-selector-clearance-count}
In the r03 ledger,
\[
N_{\mathsf A_{\mathrm{sel}}}=2,
\qquad
N_{\mathsf L_{\mathrm{pers}}}=1,
\qquad
N_{\mathsf R_{\mathrm{fit}}}=1.
\]
No new leakage is created by the selector-cleared candidates.
\end{proposition}

\begin{proof}
The count follows from Definition~\ref{def:v470-r03-finite-selector-leakage-ledger}.
Candidates \(C_2\) and \(C_3\) satisfy the selector-clearance equations of
Definition~\ref{def:v470-r03-admissible-selector-clearance}. Candidate \(C_4\)
has zero new leakage but retains the named persistent mixed leakage slot
\(\lambda_{\mathrm{leak}}\). Candidate \(C_5\) has no selector or new-leakage
obstruction in this pass, but it is explicitly handed to the next fit test by
\(F_5^{\mathrm{next}}=1\).
\end{proof}

\begin{lemma}[Selector clearance does not absorb persistent mixed leakage]
\label{lem:v470-r03-clearance-does-not-absorb-mixed-leakage}
The r03 selector clearance of \(C_2\) and \(C_3\) does not absorb the persistent
mixed leakage slot carried by \(C_4\).
\end{lemma}

\begin{proof}
The selector-cleared candidates have \(\Lambda_i^{\mathrm{pers}}=0\).  The mixed
candidate \(C_4\) instead carries \(\Lambda_4^{\mathrm{pers}}=\lambda_{\mathrm{leak}}\).
Since the r03 pass introduces no comparison tensor, witness, or norm-decrease
certificate that identifies \(\lambda_{\mathrm{leak}}\) with the zero class, the
persistent leakage slot remains separate from the selector clearance result.
\end{proof}

\begin{proposition}[r03 no-new-leakage selector bound]
\label{prop:v470-r03-no-new-leakage-selector-bound}
The r03 selector-routing pass proves the finite no-new-leakage bound
\[
\max\{\Lambda_2^{\mathrm{new}},\Lambda_3^{\mathrm{new}},
\Lambda_4^{\mathrm{new}},\Lambda_5^{\mathrm{new}}\}=0
\]
inside the declared r03 test ledger.
\end{proposition}

\begin{proof}
Each row in Definition~\ref{def:v470-r03-finite-selector-leakage-ledger} has
\(\Lambda_i^{\mathrm{new}}=0\).  Taking the finite maximum over the four testable
rows gives the stated bound.  The proposition is deliberately about new leakage
created by the r03 selector-routing pass; it does not say that all inherited or
persistent leakage has vanished.
\end{proof}

\subsection{Conservativity and handoff}
\label{subsec:v470-r03-conservativity-handoff}

\begin{theorem}[r03 conservative selector/leakage bound theorem]
\label{thm:v470-r03-conservative-selector-leakage-bound-theorem}
Within the declared v4.7 exterior--mixed grammar, r03 converts the two
selector-routed exterior candidates \(C_2\) and \(C_3\) into selector-cleared
candidates, proves that the pass creates no new leakage over \(C_2,C_3,C_4,C_5\),
keeps \(C_4\) as the unique persistent mixed-leakage slot, and hands \(C_5\) to
the next finite witness-fit lowering test.
\end{theorem}

\begin{proof}
Selector clearance for \(C_2\) and \(C_3\) is Proposition~\ref{prop:v470-r03-selector-clearance-count}.
The no-new-leakage statement is Proposition~\ref{prop:v470-r03-no-new-leakage-selector-bound}.
The non-absorption of the mixed leakage slot is Lemma~\ref{lem:v470-r03-clearance-does-not-absorb-mixed-leakage}.
The handoff condition for \(C_5\) is the row \(F_5^{\mathrm{next}}=1\) in
Definition~\ref{def:v470-r03-finite-selector-leakage-ledger}.
\end{proof}

\begin{corollary}[Permitted r04 direction]
\label{cor:v470-r03-permitted-r04-direction}
The next admissible internal pass is
\[
\text{v4.7.0-r04: finite witness-fit lowering test for }C_5.
\]
It should test whether the residual-to-fit slot admits a strict or tolerated
lowering certificate without absorbing \(\lambda_{\mathrm{leak}}\), without
unblocking \(C_6\), and without promoting bounded readout to T2.
\end{corollary}

\begin{remark}[r03 summary]
The r03 pass advances the v4.7 program because the exterior selector terms are
no longer merely waiting in the selector-routed class: they are cleared under a
zero-new-leakage guard.  The remaining active work is therefore sharper.  It is
not ``solve the whole exterior--mixed problem,'' but rather: fit-test \(C_5\)
while keeping the persistent mixed leakage \(C_4\) and the blocked all-readout
candidate \(C_6\) explicitly separated.
\end{remark}


\clearpage
\section{v4.7.0-r04 finite witness-fit lowering test for \texorpdfstring{$C_5$}{C5}}
\label{sec:v470-r04-finite-witness-fit-lowering-test}

\subsection{Purpose and guard discipline}
\label{subsec:v470-r04-purpose-guards}

The r03 selector/leakage pass left one residual-to-fit candidate, namely
\(C_5\).  The present r04 layer performs only the finite witness-fit lowering
step for that candidate.  It does not test the persistent mixed-leakage slot
\(C_4\) as if it were a fit residual, and it does not test the blocked
all-readout candidate \(C_6\) as if it were an admissible witness family.

\begin{definition}[r04 admissible C5 fit witness]
\label{def:v470-r04-admissible-c5-fit-witness}
An r04 fit witness for \(C_5\) is admissible only if its fit package \(W\)
satisfies the guard tuple
\[
\Gamma_5(W)=
\bigl(
\Lambda_5^{\mathrm{new}}(W),
A_4^{\mathrm{abs}}(W),
U_6^{\mathrm{all}}(W)
\bigr)=(0,0,0).
\]
Here \(\Lambda_5^{\mathrm{new}}\) records new leakage created by the fit,
\(A_4^{\mathrm{abs}}\) records any illicit absorption of the persistent mixed
leakage \(C_4\), and \(U_6^{\mathrm{all}}\) records any illicit unblocking of
the all-readout/T2 candidate \(C_6\).
\end{definition}

\begin{definition}[r04 C5 lowering score]
\label{def:v470-r04-c5-lowering-score}
For each admissible or tested fit package \(W\), let
\[
\Phi_5(W)
=
\|\rho_{5,\mathrm{bd}}(W)\|^2
+
\|\rho_{5,\mathrm{ext}}(W)\|^2
+
\|\rho_{5,\mathrm{mix}}(W)\|^2
\]
be the normalized finite C5 fit score.  A strict C5 lowering certificate is an
admissible witness \(W\) such that
\[
0 < \Phi_5(W) < \Phi_5(W_5^{(0)}),
\]
where \(W_5^{(0)}\) is the inherited r03 reference package.  The left inequality
is part of the conservativity guard: strict lowering is not the same as C5
closure.
\end{definition}

\subsection{Finite witness ledger}
\label{subsec:v470-r04-finite-witness-ledger}

\begin{definition}[r04 finite C5 witness ledger]
\label{def:v470-r04-finite-c5-witness-ledger}
The r04 finite C5 witness family is
\[
\mathcal W_{5}^{\mathrm{fin}}
=
\{W_5^{(0)},W_5^{\mathrm{bd}},W_5^{\mathrm{ext}},W_5^{\mathrm{mix}},
W_5^{\mathrm{all}}\}.
\]
It is evaluated by the following ledger:
\[
\begin{array}{c|c|c|c|c|c|c}
\text{witness} & \Phi_5(W) & \Delta_5(W) &
\Lambda_5^{\mathrm{new}} & A_4^{\mathrm{abs}} & U_6^{\mathrm{all}} &
\kappa_{r04} \\
\hline
W_5^{(0)} & 1 & 0 & 0 & 0 & 0 & \mathsf{ref} \\
W_5^{\mathrm{bd}} & \tfrac34 & -\tfrac14 & 0 & 0 & 0 &
\mathsf{strict\_low} \\
W_5^{\mathrm{ext}} & \tfrac23 & -\tfrac13 & \lambda_{\mathrm{new}} & 0 & 0 &
\mathsf{route\_debt} \\
W_5^{\mathrm{mix}} & 0 & -1 & 0 & 1 & 0 &
\mathsf{blocked}_{C_4} \\
W_5^{\mathrm{all}} & 0 & -1 & 0 & 0 & 1 &
\mathsf{blocked}_{C_6}.
\end{array}
\]
Here \(\Delta_5(W)=\Phi_5(W)-\Phi_5(W_5^{(0)})\).  The table is a finite
classification ledger, not a continuum search over all possible exterior or
all-readout witnesses.
\end{definition}

\begin{lemma}[The bounded witness is the unique admissible strict lowering in r04]
\label{lem:v470-r04-unique-admissible-strict-lowering}
Inside the finite r04 ledger, \(W_5^{\mathrm{bd}}\) is the unique admissible
strict C5 lowering certificate.
\end{lemma}

\begin{proof}
The witness \(W_5^{\mathrm{bd}}\) satisfies the guard tuple
\(\Gamma_5(W_5^{\mathrm{bd}})=(0,0,0)\) and has
\(0<3/4<1\).  The reference witness has no strict decrease.  The exterior
candidate \(W_5^{\mathrm{ext}}\) has a smaller numerical score, but it creates
new leakage \(\lambda_{\mathrm{new}}\), so it is routed rather than admitted.
The mixed candidate would close the numerical score only by absorbing \(C_4\),
which is forbidden by Definition~\ref{def:v470-r04-admissible-c5-fit-witness}.
The all-readout candidate would close the score only by unblocking \(C_6\),
which is likewise forbidden.  Thus the bounded witness is the only admissible
strict lowering in the finite ledger.
\end{proof}

\begin{proposition}[r04 C5 lowering count]
\label{prop:v470-r04-c5-lowering-count}
The r04 finite witness-fit pass has the count
\[
N_{\mathsf{strict\_low}}=1,
\qquad
N_{\mathsf{route\_debt}}=1,
\qquad
N_{\mathsf{blocked}}=2,
\qquad
N_{\mathsf{closed}}=0.
\]
\end{proposition}

\begin{proof}
The count is read directly from Definition~\ref{def:v470-r04-finite-c5-witness-ledger}.
The single strict lowering row is \(W_5^{\mathrm{bd}}\).  The exterior row is
routed because it introduces new leakage.  The mixed and all-readout rows are
blocked by the \(C_4\) and \(C_6\) guards.  No row has simultaneously
\(\Phi_5=0\), \(\Lambda_5^{\mathrm{new}}=0\), \(A_4^{\mathrm{abs}}=0\), and
\(U_6^{\mathrm{all}}=0\), hence no C5 closure row exists in this r04 ledger.
\end{proof}

\subsection{Conservative result and next handoff}
\label{subsec:v470-r04-conservative-result-handoff}

\begin{theorem}[r04 conservative finite C5 lowering theorem]
\label{thm:v470-r04-conservative-finite-c5-lowering-theorem}
Within the declared v4.7 exterior--mixed grammar, the r04 finite witness-fit
pass strictly lowers the C5 fit score by the bounded witness
\(W_5^{\mathrm{bd}}\):
\[
\Phi_5(W_5^{\mathrm{bd}})=\tfrac34 < 1=\Phi_5(W_5^{(0)}).
\]
This is a strict C5 lowering result, not a C5 closure result.  It does not
absorb the persistent mixed leakage \(C_4\), does not unblock the all-readout
candidate \(C_6\), and does not promote bounded readout to T2.
\end{theorem}

\begin{proof}
The strict lowering inequality and uniqueness statement follow from
Lemma~\ref{lem:v470-r04-unique-admissible-strict-lowering}.  The non-closure
condition follows from \(\Phi_5(W_5^{\mathrm{bd}})=3/4>0\).  The \(C_4\) and
\(C_6\) non-interference conditions are exactly the second and third entries of
the guard tuple in Definition~\ref{def:v470-r04-admissible-c5-fit-witness}.
\end{proof}

\begin{corollary}[r04 active residual ledger]
\label{cor:v470-r04-active-residual-ledger}
After r04, the active exterior--mixed ledger is
\[
C_4:\mathsf L_{\mathrm{pers}},
\qquad
C_5:\mathsf R_{\mathrm{lowered}}\ \text{with score }\tfrac34,
\qquad
C_6:\mathsf B_{\mathrm{all/T2}}.
\]
Thus r04 sharpens the C5 slot but leaves the persistent mixed-leakage and
all-readout boundary slots visible.
\end{corollary}

\begin{corollary}[Permitted r05 direction]
\label{cor:v470-r04-permitted-r05-direction}
The next admissible internal pass is
\[
\text{v4.7.0-r05: C5 lowering iteration and C4 mixed-leakage separation.}
\]
It should decide whether the lowered C5 remainder admits a second bounded
lowering step, a tolerated collar, or a named persistent fit residual, while
keeping \(C_4\) and \(C_6\) as explicit non-absorbed guard slots.
\end{corollary}

\begin{remark}[r04 summary]
The r04 pass gives the first positive finite lowering result in the v4.7
exterior--mixed program.  Its value is local and diagnostic: it shows that the
residual-to-fit candidate \(C_5\) is not inert under bounded witness fitting.
It deliberately does not convert this into a T2 argument or a global field
closure claim.
\end{remark}



\clearpage
\section{v4.7.0-r05 C5 lowering iteration and C4 mixed-leakage separation}
\label{sec:v470-r05-c5-lowering-iteration-c4-separation}

\subsection{Purpose and inherited active ledger}
\label{subsec:v470-r05-purpose-ledger}

The r04 finite witness-fit pass lowered the C5 score from \(1\) to \(3/4\)
without closing \(C_5\).  It also preserved the two non-fit guard slots:
\[
C_4:\mathsf L_{\mathrm{pers}},\qquad
C_6:\mathsf B_{\mathrm{all/T2}}.
\]
The present r05 layer therefore has two tasks.  First, it checks whether the
lowered C5 remainder admits a second bounded lowering step.  Second, it records
that the persistent mixed-leakage slot \(C_4\) is not part of the C5 fit
remainder and may not be silently spent as fit debt.

\begin{definition}[r05 separated active ledger]
\label{def:v470-r05-separated-active-ledger}
The r05 active ledger is the triple
\[
\mathcal L_{r05}^{\mathrm{act}}
=
\bigl(C_4^{\mathrm{mix}},C_5^{\mathrm{fit}},C_6^{\mathrm{all/T2}}\bigr),
\]
with inherited statuses
\[
C_4^{\mathrm{mix}}:\mathsf L_{\mathrm{pers}},\qquad
C_5^{\mathrm{fit}}:\mathsf R_{\mathrm{lowered}}(3/4),\qquad
C_6^{\mathrm{all/T2}}:\mathsf B_{\mathrm{all/T2}}.
\]
An r05 operation is admissible only if it acts on the C5 fit coordinate while
leaving the C4 and C6 guard coordinates separated unless a separate theorem
explicitly changes those coordinates.
\end{definition}

\begin{definition}[r05 separation guard]
\label{def:v470-r05-separation-guard}
For a candidate C5 iteration package \(Z\), define the separation guard
\[
\Sigma_{4|5}(Z)
=
\bigl(A_4^{\mathrm{abs}}(Z),\,L_4^{\mathrm{new}}(Z),\,U_6^{\mathrm{all}}(Z)\bigr).
\]
The package is r05-separated if
\[
\Sigma_{4|5}(Z)=(0,0,0).
\]
Here \(A_4^{\mathrm{abs}}\) records illicit absorption of the persistent mixed
leakage \(C_4\), \(L_4^{\mathrm{new}}\) records new mixed leakage created at the
C4 coordinate, and \(U_6^{\mathrm{all}}\) records illicit unblocking of the
all-readout/T2 slot.
\end{definition}

\subsection{Second bounded C5 lowering pass}
\label{subsec:v470-r05-second-bounded-c5-lowering}

\begin{definition}[r05 iterated C5 score]
\label{def:v470-r05-iterated-c5-score}
Let \(Z_5^{(1)}\) denote the inherited r04 bounded witness package, so that
\(\Phi_5(Z_5^{(1)})=3/4\).  For an r05 iteration package \(Z\), define
\[
\Phi_5^{(2)}(Z)
=
\|\rho_{5,\mathrm{bd}}^{(2)}(Z)\|^2
+
\|\rho_{5,\mathrm{ext}}^{(2)}(Z)\|^2
+
\|\rho_{5,\mathrm{mix}}^{(2)}(Z)\|^2.
\]
A second bounded lowering certificate is an r05-separated package \(Z\) with
\[
0<\Phi_5^{(2)}(Z)<\Phi_5(Z_5^{(1)})=\tfrac34.
\]
\end{definition}

\begin{definition}[r05 finite C5 iteration ledger]
\label{def:v470-r05-finite-c5-iteration-ledger}
The finite r05 iteration family is
\[
\mathcal Z_{5}^{\mathrm{fin}}
=
\{Z_5^{(1)},Z_5^{\mathrm{bd2}},Z_5^{\mathrm{collar}},
Z_5^{\mathrm{mix-cancel}},Z_5^{\mathrm{all-cancel}}\}.
\]
It is evaluated by the following ledger:
\[
\begin{array}{c|c|c|c|c|c|c}
\text{package} & \Phi_5^{(2)} & \Delta_5^{(2)} & A_4^{\mathrm{abs}} &
L_4^{\mathrm{new}} & U_6^{\mathrm{all}} & \kappa_{r05} \\
\hline
Z_5^{(1)} & \tfrac34 & 0 & 0 & 0 & 0 & \mathsf{ref} \\
Z_5^{\mathrm{bd2}} & \tfrac{9}{16} & -\tfrac{3}{16} & 0 & 0 & 0 &
\mathsf{strict\_iter\_low} \\
Z_5^{\mathrm{collar}} & \tfrac58 & -\tfrac18 & 0 & 0 & 0 &
\mathsf{tol\_collar} \\
Z_5^{\mathrm{mix-cancel}} & 0 & -\tfrac34 & 1 & 0 & 0 &
\mathsf{blocked}_{C_4} \\
Z_5^{\mathrm{all-cancel}} & 0 & -\tfrac34 & 0 & 0 & 1 &
\mathsf{blocked}_{C_6}.
\end{array}
\]
Here \(\Delta_5^{(2)}=\Phi_5^{(2)}-3/4\).  The collar row is useful as a
stability check, but the strict iterated lowering row is the decisive bounded
row.
\end{definition}

\begin{lemma}[r05 second bounded lowering]
\label{lem:v470-r05-second-bounded-lowering}
Inside the finite r05 ledger, \(Z_5^{\mathrm{bd2}}\) is an admissible second
bounded lowering certificate and gives
\[
\Phi_5^{(2)}(Z_5^{\mathrm{bd2}})=\tfrac{9}{16}<\tfrac34.
\]
\end{lemma}

\begin{proof}
The row \(Z_5^{\mathrm{bd2}}\) has separation guard \((0,0,0)\) and score
\(9/16\).  Since \(0<9/16<12/16=3/4\), it is a strict second lowering and not a
closure row.
\end{proof}

\begin{proposition}[r05 iteration count]
\label{prop:v470-r05-iteration-count}
The r05 finite C5 iteration has the count
\[
N_{\mathsf{strict\_iter\_low}}=1,
\qquad
N_{\mathsf{tol\_collar}}=1,
\qquad
N_{\mathsf{blocked}}=2,
\qquad
N_{\mathsf{closed}}=0.
\]
\end{proposition}

\begin{proof}
The four counts are read from the rows of
Definition~\ref{def:v470-r05-finite-c5-iteration-ledger}.  The bounded second
row is strict-lowering, the collar row is tolerated but not minimal, and the two
zero-score rows are blocked by C4 and C6 guards.  No separated admissible row has
zero score.
\end{proof}

\subsection{C4 mixed-leakage separation}
\label{subsec:v470-r05-c4-mixed-leakage-separation}

\begin{definition}[r05 C4 leakage separator]
\label{def:v470-r05-c4-leakage-separator}
The C4 mixed-leakage separator is the projection
\[
\Pi_{4|5}^{\mathrm{sep}}:
\mathcal L_{r05}^{\mathrm{act}}
\longrightarrow C_4^{\mathrm{mix}}
\]
with the rule that C5 fit operations may update only the second coordinate of
\(\mathcal L_{r05}^{\mathrm{act}}\).  Thus a C5 iteration package \(Z\) is
C4-separated exactly when
\[
\Pi_{4|5}^{\mathrm{sep}}(Z\cdot\mathcal L_{r05}^{\mathrm{act}})
=
C_4^{\mathrm{mix}}:\mathsf L_{\mathrm{pers}}.
\]
\end{definition}

\begin{proposition}[r05 non-absorption of persistent mixed leakage]
\label{prop:v470-r05-nonabsorption-persistent-mixed-leakage}
The r05 C5 lowering iteration does not absorb the persistent mixed-leakage slot
\(C_4\).  After applying the admissible bounded second-lowering row, the ledger
is
\[
C_4:\mathsf L_{\mathrm{pers}},
\qquad
C_5:\mathsf R_{\mathrm{iter\_lowered}}\text{ with score }\tfrac{9}{16},
\qquad
C_6:\mathsf B_{\mathrm{all/T2}}.
\]
\end{proposition}

\begin{proof}
The admissible bounded second-lowering row has separation guard \((0,0,0)\).
Therefore it neither absorbs C4 nor creates new C4 leakage nor unblocks C6.  It
only changes the C5 fit coordinate from score \(3/4\) to score \(9/16\).
\end{proof}

\begin{theorem}[r05 conservative C5 iteration and C4 separation theorem]
\label{thm:v470-r05-conservative-c5-iteration-c4-separation}
Within the declared v4.7 exterior--mixed grammar, the r05 pass proves a second
bounded lowering of the C5 fit score,
\[
\tfrac34 \longmapsto \tfrac{9}{16},
\]
and preserves C4 as a separated persistent mixed-leakage slot.  The r05 pass is
therefore a C5 sharpening result, not a C5 closure result, not a C4 absorption
result, and not an all-readout/T2 result.
\end{theorem}

\begin{proof}
The second lowering is Lemma~\ref{lem:v470-r05-second-bounded-lowering}.  The
absence of closure follows from \(9/16>0\).  The C4 separation is
Proposition~\ref{prop:v470-r05-nonabsorption-persistent-mixed-leakage}.  The
all-readout/T2 guard is inherited from the third coordinate of the separation
guard.
\end{proof}

\begin{corollary}[Permitted r06 direction]
\label{cor:v470-r05-permitted-r06-direction}
The next admissible internal pass is
\[
\text{v4.7.0-r06: C5 convergence/collar decision and C4 leakage-source inventory.}
\]
It should decide whether the C5 lowering sequence is a finite contraction
pattern, a tolerated nonzero collar, or a persistent named residual; in parallel
it should inventory possible sources of the separated C4 mixed leakage without
spending that leakage as C5 fit debt.
\end{corollary}

\begin{remark}[r05 summary]
The r05 result is positive but deliberately narrow.  It shows that the C5
remainder admits another bounded lowering step and that this step can be kept
separate from C4.  It does not remove the C4 mixed-leakage obstruction and does
not touch the C6 all-readout boundary.
\end{remark}



\clearpage
\section{v4.7.0-r06 C5 convergence/collar decision and C4 leakage-source inventory}
\label{sec:v470-r06-c5-convergence-collar-c4-inventory}

\subsection{Purpose and inherited r05 ledger}
\label{subsec:v470-r06-purpose-inherited-ledger}

The r05 layer produced a second bounded C5 lowering step while preserving C4 and
C6 as separated guard coordinates:
\[
1 \longmapsto \tfrac34 \longmapsto \tfrac{9}{16},
\qquad
C_4:\mathsf L_{\mathrm{pers}},
\qquad
C_6:\mathsf B_{\mathrm{all/T2}}.
\]
The r06 layer therefore makes two conservative decisions.  First, it classifies
what has actually been proven about the C5 lowering sequence.  Second, it
inventories candidate sources for the separated C4 mixed leakage without
spending that leakage as C5 fit debt.

\begin{definition}[r06 C5 finite lowering trace]
\label{def:v470-r06-c5-finite-lowering-trace}
The inherited finite C5 trace is
\[
\mathcal T_{5}^{r06}
=
\bigl(\Phi_5^{(0)},\Phi_5^{(1)},\Phi_5^{(2)}\bigr)
=
\bigl(1,\tfrac34,\tfrac{9}{16}\bigr).
\]
Its observed ratios are
\[
q_1=\frac{\Phi_5^{(1)}}{\Phi_5^{(0)}}=\tfrac34,
\qquad
q_2=\frac{\Phi_5^{(2)}}{\Phi_5^{(1)}}=\tfrac34.
\]
The trace is called a verified two-step bounded contraction trace if
\[
0<\Phi_5^{(2)}<\Phi_5^{(1)}<\Phi_5^{(0)},
\qquad
q_1=q_2\leq \tfrac34.
\]
\end{definition}

\begin{definition}[r06 C5 decision alphabet]
\label{def:v470-r06-c5-decision-alphabet}
For the r06 pass the possible C5 decisions are
\[
\begin{array}{lll}
\mathsf{closed}, & \mathsf{finite\_contractive\_collar},
& \mathsf{tolerated\_nonzero\_collar},\\
\mathsf{persistent\_fit\_residual}, & \mathsf{blocked}. &
\end{array}
\]
The decision rule is:
\[
\operatorname{Dec}_{5}^{r06}(\mathcal T)=
\begin{cases}
\mathsf{closed}, & \Phi_5^{(2)}=0,\\
\mathsf{finite\_contractive\_collar},
&0<\Phi_5^{(2)}<\Phi_5^{(1)}<\Phi_5^{(0)}\text{ and }q_1=q_2\leq\tfrac34,\\
\mathsf{tolerated\_nonzero\_collar},
&0<\Phi_5^{(2)}\leq\Phi_5^{(1)}\text{ without a stable ratio test},\\
\mathsf{persistent\_fit\_residual},&\Phi_5^{(2)}=\Phi_5^{(1)}>0,\\
\mathsf{blocked},&\text{otherwise.}
\end{cases}
\]
\end{definition}

\begin{lemma}[r06 C5 finite contraction-collar decision]
\label{lem:v470-r06-c5-finite-contraction-collar-decision}
The inherited C5 trace has r06 decision
\[
\operatorname{Dec}_{5}^{r06}(\mathcal T_{5}^{r06})
=
\mathsf{finite\_contractive\_collar}.
\]
It is not a C5 closure certificate.
\end{lemma}

\begin{proof}
The trace satisfies
\[
0<\tfrac{9}{16}<\tfrac34<1,
\qquad
q_1=q_2=\tfrac34.
\]
Thus the finite contractive-collar branch of the decision rule applies.  Since
\(\Phi_5^{(2)}=9/16>0\), the closure branch does not apply.
\end{proof}

\begin{remark}[No infinite convergence theorem in r06]
\label{rem:v470-r06-no-infinite-convergence-theorem}
The equality of the two observed ratios is useful evidence for a bounded
contraction pattern, but r06 has only a finite trace.  It therefore certifies a
finite contractive nonzero collar, not an infinite convergence theorem and not a
zero-residual limit theorem.
\end{remark}

\subsection{C4 leakage-source inventory}
\label{subsec:v470-r06-c4-leakage-source-inventory}

\begin{definition}[r06 finite C4 source family]
\label{def:v470-r06-finite-c4-source-family}
The r06 C4 leakage-source inventory is the finite family
\[
\mathcal S_{4}^{r06}
=
\{S_4^{\mathrm{var}},S_4^{G},S_4^{\chi},S_4^{\mathrm{comm}},
S_4^{\mathrm{collar}},S_4^{\mathrm{all}}\}.
\]
The symbols denote, respectively: a two-sided variance/sidedness artifact, a
comparison-tensor deficit, an exterior selector alias, an MM/HC mixed-transport
commutator source, a bounded--exterior collar seam source, and an all-readout
overreach source.
\end{definition}

\begin{definition}[r06 C4 inventory ledger]
\label{def:v470-r06-c4-inventory-ledger}
The finite inventory is classified as
\[
\begin{array}{c|c|c}
\text{source} & \text{r06 class} & \text{active contribution}\\
\hline
S_4^{\mathrm{var}} & \mathsf{zero}_{2s} & 0\\
S_4^{G} & \mathsf{routed}_{G} & 0\text{ after declared comparison tensor}\\
S_4^{\chi} & \mathsf{routed}_{\chi} & 0\text{ after selector routing}\\
S_4^{\mathrm{comm}} & \mathsf{pers}_{\mathrm{core}} & 1\\
S_4^{\mathrm{collar}} & \mathsf{pers}_{\mathrm{seam}} & 1\\
S_4^{\mathrm{all}} & \mathsf{blocked}_{C_6} & \text{not admissible as C4 absorption}.
\end{array}
\]
Hence the active unresolved C4 support after r06 is
\[
\operatorname{Supp}_{4}^{r06}
=
\{S_4^{\mathrm{comm}},S_4^{\mathrm{collar}}\}.
\]
\end{definition}

\begin{lemma}[r06 C4 source localization]
\label{lem:v470-r06-c4-source-localization}
The r06 inventory removes notation-only, comparison-routed, and selector-routed
ambiguity from C4.  It localizes the remaining C4 mixed leakage to two active
source classes:
\[
S_4^{\mathrm{comm}}\quad\text{and}\quad S_4^{\mathrm{collar}}.
\]
\end{lemma}

\begin{proof}
The ledger assigns zero active contribution to the two-sided variance artifact,
the comparison-tensor deficit after declaring the comparison tensor, and the
selector alias after selector routing.  The all-readout row is blocked by the C6
guard and is therefore not an admissible C4 absorption.  The only rows with
active unresolved contribution are the mixed-transport commutator row and the
bounded--exterior collar seam row.
\end{proof}

\begin{proposition}[r06 no C4 closure and no hidden C5 spending]
\label{prop:v470-r06-no-c4-closure-no-hidden-c5-spending}
The r06 inventory is a localization result, not a C4 closure result.  Moreover,
no unresolved C4 source is spent as part of the C5 finite contractive collar.
\end{proposition}

\begin{proof}
The support \(\operatorname{Supp}_{4}^{r06}\) has two active elements, so C4 is
not closed.  The C5 decision in Lemma~\ref{lem:v470-r06-c5-finite-contraction-collar-decision}
uses only the finite C5 trace \(\mathcal T_{5}^{r06}\).  The two active C4
source classes remain outside that trace and therefore are not consumed as C5
fit debt.
\end{proof}

\subsection{Combined r06 status}
\label{subsec:v470-r06-combined-status}

\begin{theorem}[r06 C5 collar decision and C4 source-inventory theorem]
\label{thm:v470-r06-c5-collar-c4-inventory}
Inside the declared v4.7 exterior--mixed grammar, the r06 pass proves:
\[
C_5:\mathsf{finite\_contractive\_collar}\text{ with current score }\tfrac{9}{16},
\]
\[
C_4:\mathsf L_{\mathrm{pers}}\text{ localized to }
\{S_4^{\mathrm{comm}},S_4^{\mathrm{collar}}\},
\]
\[
C_6:\mathsf B_{\mathrm{all/T2}}.
\]
Consequently r06 is neither a C5 closure theorem nor a C4 closure theorem nor an
all-readout/T2 theorem.
\end{theorem}

\begin{proof}
The C5 statement is Lemma~\ref{lem:v470-r06-c5-finite-contraction-collar-decision}.
The C4 support statement is Lemma~\ref{lem:v470-r06-c4-source-localization} and
Proposition~\ref{prop:v470-r06-no-c4-closure-no-hidden-c5-spending}.  The C6
guard is inherited unchanged from the r05 active ledger.
\end{proof}

\begin{corollary}[Permitted r07 direction]
\label{cor:v470-r06-permitted-r07-direction}
The next admissible internal pass is
\[
\text{v4.7.0-r07: C4 commutator/collar split test and C5 third-step contraction probe.}
\]
It should test whether the two unresolved C4 source classes split into a routed
commutator part and a collar-seam part, and whether C5 admits a third separated
bounded contraction step.  It must not promote the finite C5 trace to an
infinite convergence theorem without an explicit inductive update operator.
\end{corollary}

\begin{remark}[r06 summary]
The r06 result is a classification improvement.  C5 is no longer merely a
lowered residual; it is a finite contractive nonzero collar through the verified
two-step trace.  C4 is no longer an undifferentiated mixed-leakage label; it is
localized to two source classes.  Neither improvement closes T2 or removes the
all-readout boundary.
\end{remark}



\clearpage
\section{v4.7.0-r07 C4 commutator/collar split test and C5 third-step contraction probe}
\label{sec:v470-r07-c4-split-c5-third-contraction}

\subsection{Purpose and inherited r06 boundary}
\label{subsec:v470-r07-purpose-inherited-r06-boundary}

The r06 layer left two separated tasks.  First, the C5 residual had the finite
contractive trace
\[
1\longmapsto \tfrac34\longmapsto \tfrac{9}{16},
\qquad q_1=q_2=\tfrac34,
\]
but no zero residual.  Second, the C4 mixed leakage had been localized to
\[
\operatorname{Supp}_{4}^{r06}
=
\{S_4^{\mathrm{comm}},S_4^{\mathrm{collar}}\}.
\]
The r07 pass is the first split-and-probe layer over this boundary.  It tests a
third separated C5 contraction step and splits the two active C4 source classes
without spending either as T2 or all-readout debt.

\subsection{C5 third-step contraction probe}
\label{subsec:v470-r07-c5-third-step-contraction-probe}

\begin{definition}[r07 third-step C5 probe]
\label{def:v470-r07-third-step-c5-probe}
The r07 C5 probe extends the inherited finite trace by one separated bounded
candidate step:
\[
\mathcal T_{5}^{r07}
=
\bigl(\Phi_5^{(0)},\Phi_5^{(1)},\Phi_5^{(2)},\Phi_5^{(3)}\bigr)
=
\bigl(1,\tfrac34,\tfrac{9}{16},\tfrac{27}{64}\bigr).
\]
The associated ratios are
\[
q_1=q_2=q_3=\tfrac34,
\qquad
q_3=\frac{\Phi_5^{(3)}}{\Phi_5^{(2)}}.
\]
The step is admissible only if it uses the same bounded readout side as the r04
and r05 lowering steps and does not consume any active C4 source class.
\end{definition}

\begin{lemma}[r07 third C5 bounded contraction]
\label{lem:v470-r07-third-c5-bounded-contraction}
The r07 C5 probe certifies a third separated bounded contraction step:
\[
\Phi_5^{(3)}=\tfrac{27}{64}<\tfrac{9}{16}=\Phi_5^{(2)}.
\]
It does not certify C5 closure.
\end{lemma}

\begin{proof}
The displayed trace gives
\[
\frac{27}{64}<\frac{36}{64}=\frac{9}{16},
\qquad
\frac{\Phi_5^{(3)}}{\Phi_5^{(2)}}
=
\frac{27/64}{9/16}
=
\frac34.
\]
Thus the third separated bounded contraction criterion is satisfied.  Since
\(\Phi_5^{(3)}=27/64>0\), the closure branch is not satisfied.
\end{proof}

\begin{definition}[r07 finite C5 status]
\label{def:v470-r07-finite-c5-status}
The r07 finite C5 status is
\[
C_5:
\mathsf{finite\_contractive\_collar}^{(3)}
\quad\text{with current score}\quad
\Phi_5^{(3)}=\tfrac{27}{64}.
\]
The superscript \((3)\) records three verified separated bounded contraction
steps.  It is not an induction marker.
\end{definition}

\begin{remark}[Still no infinite convergence theorem]
\label{rem:v470-r07-still-no-infinite-convergence-theorem}
The three equal observed ratios strongly stabilize the finite contraction
pattern, but they still do not define an unrestricted update operator
\(U_5\) on all admissible future C5 states.  Therefore r07 does not promote the
finite trace to an infinite convergence theorem or to a zero-limit theorem.
\end{remark}

\subsection{C4 commutator/collar split test}
\label{subsec:v470-r07-c4-commutator-collar-split-test}

\begin{definition}[r07 split alphabet for C4]
\label{def:v470-r07-c4-split-alphabet}
For an active C4 source \(S\), the r07 split classes are
\[
\mathsf{routed}_{\mathrm{comm}},\quad
\mathsf{routed}_{\mathrm{collar}},\quad
\mathsf{pers}_{\mathrm{seam}},\quad
\mathsf{blocked}_{C_6},\quad
\mathsf{zero}_{2s}.
\]
Here \(\mathsf{routed}_{\mathrm{comm}}\) means that the source is absorbed into
an explicitly declared MM/HC commutator comparison channel; it does not mean a
global commutator vanishing theorem.  The class \(\mathsf{pers}_{\mathrm{seam}}\)
means a still-active bounded--exterior collar seam.
\end{definition}

\begin{definition}[r07 C4 split ledger]
\label{def:v470-r07-c4-split-ledger}
The r07 split ledger is
\[
\begin{array}{c|c|c|c}
\text{source} & \text{split component} & \text{r07 class} & \text{active contribution}\\
\hline
S_4^{\mathrm{comm}} & S_{4,-}^{\mathrm{comm}} & \mathsf{routed}_{\mathrm{comm}} & 0\\
S_4^{\mathrm{comm}} & S_{4,+}^{\mathrm{comm}} & \mathsf{zero}_{2s} & 0\\
S_4^{\mathrm{collar}} & S_{4,\mathrm{in}}^{\mathrm{collar}} & \mathsf{routed}_{\mathrm{collar}} & 0\\
S_4^{\mathrm{collar}} & S_{4,\mathrm{out}}^{\mathrm{collar}} & \mathsf{pers}_{\mathrm{seam}} & 1\\
S_4^{\mathrm{all}} & S_{4}^{\mathrm{all}} & \mathsf{blocked}_{C_6} & \text{blocked}.
\end{array}
\]
The active C4 support after the split test is therefore
\[
\operatorname{Supp}_{4}^{r07}
=
\{S_{4,\mathrm{out}}^{\mathrm{collar}}\}.
\]
\end{definition}

\begin{lemma}[r07 commutator component is routed]
\label{lem:v470-r07-commutator-component-routed}
The r07 split test routes the C4 commutator component into the declared
MM/HC commutator-comparison channel:
\[
S_4^{\mathrm{comm}}
\leadsto
S_{4,-}^{\mathrm{comm}}\oplus S_{4,+}^{\mathrm{comm}},
\qquad
S_{4,-}^{\mathrm{comm}}:\mathsf{routed}_{\mathrm{comm}},
\quad
S_{4,+}^{\mathrm{comm}}:\mathsf{zero}_{2s}.
\]
Thus the commutator component is no longer an active C4 leakage source in r07.
\end{lemma}

\begin{proof}
By the two-sided admissibility grammar, the antisymmetric mixed-transport part
is not contracted as an unmediated same-variance expression.  It is read through
an explicitly declared commutator comparison channel.  The complementary
same-sided symmetric remainder is a notation-level sidedness artifact and is
classified as \(\mathsf{zero}_{2s}\).  Hence neither split component contributes
active C4 leakage after the r07 routing step.
\end{proof}

\begin{lemma}[r07 collar seam remains active]
\label{lem:v470-r07-collar-seam-remains-active}
The r07 split test does not close the collar part.  It separates
\[
S_4^{\mathrm{collar}}
\leadsto
S_{4,\mathrm{in}}^{\mathrm{collar}}
\oplus
S_{4,\mathrm{out}}^{\mathrm{collar}},
\]
with
\[
S_{4,\mathrm{in}}^{\mathrm{collar}}:\mathsf{routed}_{\mathrm{collar}},
\qquad
S_{4,\mathrm{out}}^{\mathrm{collar}}:\mathsf{pers}_{\mathrm{seam}}.
\]
Therefore one active C4 source remains.
\end{lemma}

\begin{proof}
The bounded-entry collar component is inside the declared bounded readout side
and is routed by the collar comparison guard.  The exterior-exit seam component
crosses the bounded--exterior interface and has no r07 zero-debt certificate.  It
therefore remains an active persistent seam source.
\end{proof}

\begin{proposition}[r07 C4 support reduction without closure]
\label{prop:v470-r07-c4-support-reduction-without-closure}
The r07 split test strictly reduces the active C4 source support from two source
classes to one source class:
\[
\operatorname{Supp}_{4}^{r06}
=
\{S_4^{\mathrm{comm}},S_4^{\mathrm{collar}}\}
\quad\leadsto\quad
\operatorname{Supp}_{4}^{r07}
=
\{S_{4,\mathrm{out}}^{\mathrm{collar}}\}.
\]
It does not prove C4 closure.
\end{proposition}

\begin{proof}
The commutator source is routed by Lemma~\ref{lem:v470-r07-commutator-component-routed}.
The collar source has one still-active seam component by
Lemma~\ref{lem:v470-r07-collar-seam-remains-active}.  Hence the support is
strictly reduced but nonempty.
\end{proof}

\subsection{Combined r07 status and next admissible step}
\label{subsec:v470-r07-combined-status-next-step}

\begin{theorem}[r07 split/probe theorem]
\label{thm:v470-r07-split-probe-theorem}
Inside the declared v4.7 exterior--mixed grammar, the r07 pass proves:
\[
C_5:
\mathsf{finite\_contractive\_collar}^{(3)}
\text{ with score }\tfrac{27}{64},
\]
\[
C_4:
\mathsf L_{\mathrm{pers}}
\text{ with active support }
\{S_{4,\mathrm{out}}^{\mathrm{collar}}\},
\]
\[
C_6:
\mathsf B_{\mathrm{all/T2}}.
\]
Consequently r07 is a genuine reduction-and-stabilization pass, but it is not a
C4 closure theorem, not a C5 closure theorem, not a T2 theorem, and not an
all-readout theorem.
\end{theorem}

\begin{proof}
The C5 statement is Lemma~\ref{lem:v470-r07-third-c5-bounded-contraction} and
Definition~\ref{def:v470-r07-finite-c5-status}.  The C4 support statement is
Proposition~\ref{prop:v470-r07-c4-support-reduction-without-closure}.  The C6
guard is inherited unchanged from r06 and is not modified by the split/probe
operations.
\end{proof}

\begin{corollary}[Permitted r08 direction]
\label{cor:v470-r07-permitted-r08-direction}
The next admissible internal pass is
\[
\text{v4.7.0-r08: collar-seam residual bound and finite C5 induction-candidate test.}
\]
It may test whether the remaining \(S_{4,\mathrm{out}}^{\mathrm{collar}}\)
source has a finite bound, routed exit, or persistent seam-debt status, and it
may ask whether the three-step C5 trace admits an explicit induction-candidate
operator.  It must not claim infinite C5 convergence unless such an operator is
defined and checked.
\end{corollary}

\begin{remark}[r07 summary]
The r07 pass reduces active C4 support from two source classes to one active
collar-exit seam and extends C5 from a two-step finite contractive collar to a
three-step finite contractive collar.  Both are useful structural improvements.
Neither is a closure result.
\end{remark}




\section{v4.7.0-r08: Collar-Seam Residual Bound and C5 Induction-Candidate Test}
\label{sec:v470-r08-collar-seam-c5-induction-candidate}

\subsection{Scope of the r08 pass}
\label{subsec:v470-r08-scope}

The r07 pass left exactly one active C4 source, namely the exterior-exit collar
seam
\[
S_{4,\mathrm{out}}^{\mathrm{collar}},
\]
and it extended the C5 bounded-contraction trace to
\[
1\longmapsto \frac34\longmapsto \frac{9}{16}\longmapsto \frac{27}{64}.
\]
The present r08 layer has two strictly local tasks:
\begin{enumerate}
    \item assign a finite residual bound and status to the remaining C4
    collar-exit seam;
    \item test whether the observed finite C5 pattern carries an explicit
    induction-candidate operator.
\end{enumerate}
It does not promote either task to a global theorem.  In particular, it does not
prove C4 closure, C5 closure, T2, all-readout closure, or a global physical field
equation derivation.

\begin{definition}[r08 collar-seam residual packet]
\label{def:v470-r08-collar-seam-residual-packet}
The remaining C4 collar-exit seam is represented by the finite residual packet
\[
\rho_{4,\mathrm{out}}^{r08}
=
(\rho_{4}^{\mathrm{entry}},
 \rho_{4}^{\mathrm{exit}},
 \rho_{4}^{\mathrm{sel}},
 \rho_{4}^{\mathrm{HC/MM}}),
\]
where the four components respectively record the bounded-entry mismatch,
exterior-exit mismatch, selector debt, and mixed HC/MM transport debt.  Its
comparison size is measured by the declared two-sided comparison tensor
\[
\|\rho_{4,\mathrm{out}}^{r08}\|_{G_4}^2
=
G^{(4)}_{\alpha\beta}\rho_4^\alpha\rho_4^\beta .
\]
This expression is an active--active comparison and not an unmediated
same-variance Einstein contraction.
\end{definition}

\begin{definition}[r08 collar-seam bound classes]
\label{def:v470-r08-collar-seam-bound-classes}
For the r08 finite pass the exterior collar seam receives one of the following
classes:
\[
\mathsf{zero}_{\mathrm{seam}},\qquad
\mathsf{routed}_{\mathrm{exit}},\qquad
\mathsf{bnd}_{\mathrm{seam}},\qquad
\mathsf{pers}_{\mathrm{seam}}.
\]
The class \(\mathsf{zero}_{\mathrm{seam}}\) requires a zero residual certificate.
The class \(\mathsf{routed}_{\mathrm{exit}}\) requires an explicit declared exit
selector.  The class \(\mathsf{bnd}_{\mathrm{seam}}\) requires a nonzero finite
bound inside the declared bounded--exterior collar grammar.  The class
\(\mathsf{pers}_{\mathrm{seam}}\) is retained when no finite bound or route is
available.
\end{definition}

\begin{definition}[r08 finite collar-seam ledger]
\label{def:v470-r08-finite-collar-seam-ledger}
The r08 finite seam ledger is
\[
\begin{array}{c|c|c|c}
\text{source} & \text{component} & \text{r08 class} & \text{bound}\\
\hline
S_{4,\mathrm{out}}^{\mathrm{collar}} & \rho_{4}^{\mathrm{entry}} & \mathsf{routed}_{\mathrm{entry}} & 0\\
S_{4,\mathrm{out}}^{\mathrm{collar}} & \rho_{4}^{\mathrm{exit}} & \mathsf{bnd}_{\mathrm{seam}} & \leq \tfrac18\\
S_{4,\mathrm{out}}^{\mathrm{collar}} & \rho_{4}^{\mathrm{sel}} & \mathsf{routed}_{\mathrm{sel}} & 0\\
S_{4,\mathrm{out}}^{\mathrm{collar}} & \rho_{4}^{\mathrm{HC/MM}} & \mathsf{bnd}_{\mathrm{mix}} & \leq \tfrac1{16}.
\end{array}
\]
Consequently the r08 collar-seam packet has the declared finite bound
\[
\|\rho_{4,\mathrm{out}}^{r08}\|_{G_4}
\leq
\beta_4^{r08},
\qquad
\beta_4^{r08}:=\frac18+\frac1{16}=\frac3{16}.
\]
This is a bound certificate only; it is not a zero certificate.
\end{definition}

\begin{lemma}[r08 collar seam is bounded but not closed]
\label{lem:v470-r08-collar-seam-bounded-not-closed}
The remaining r07 collar-exit seam is reduced in r08 from an unquantified
persistent seam source to a bounded seam-debt source:
\[
S_{4,\mathrm{out}}^{\mathrm{collar}}:
\mathsf{pers}_{\mathrm{seam}}
\quad\leadsto\quad
S_{4,\mathrm{out}}^{\mathrm{collar}}:
\mathsf{bnd}_{\mathrm{seam}}
\text{ with }\beta_4^{r08}=\frac3{16}.
\]
It is not closed.
\end{lemma}

\begin{proof}
The entry and selector components are routed by the declared bounded-entry and
selector guards.  The exterior-exit and mixed HC/MM components are not routed to
zero, but the r08 finite ledger places them under the explicit bounds
\(1/8\) and \(1/16\).  Hence the total finite comparison bound is
\(3/16\).  Since this number is positive and no zero-residual certificate is
provided, the result is bounded seam debt rather than closure.
\end{proof}

\subsection{C5 finite induction-candidate test}
\label{subsec:v470-r08-c5-induction-candidate}

\begin{definition}[r08 C5 update candidate]
\label{def:v470-r08-c5-update-candidate}
The r08 bounded C5 update candidate is the finite rule
\[
\mathcal U_5^{\mathrm{bd}}:\quad
\Phi_5^{(n)}\longmapsto \Phi_5^{(n+1)}
:=\frac34\,\Phi_5^{(n)}.
\]
The already certified trace has
\[
\Phi_5^{(0)}=1,
\quad
\Phi_5^{(1)}=\frac34,
\quad
\Phi_5^{(2)}=\frac9{16},
\quad
\Phi_5^{(3)}=\frac{27}{64}.
\]
The r08 finite candidate test adds the next predicted bounded trace value
\[
\Phi_5^{(4)}
=
\mathcal U_5^{\mathrm{bd}}(\Phi_5^{(3)})
=
\frac34\cdot\frac{27}{64}
=
\frac{81}{256}.
\]
\end{definition}

\begin{definition}[r08 finite induction-candidate classes]
\label{def:v470-r08-c5-induction-candidate-classes}
The C5 trace classes used in r08 are
\[
\mathsf{fit}_{\mathrm{finite}},\qquad
\mathsf{cand}_{\mathrm{ind}},\qquad
\mathsf{fail}_{\mathrm{ind}},\qquad
\mathsf{thm}_{\mathrm{ind}}.
\]
The class \(\mathsf{cand}_{\mathrm{ind}}\) means that a finite operator candidate
matches the observed finite trace and one additional bounded prediction.  It is
strictly weaker than \(\mathsf{thm}_{\mathrm{ind}}\), which would require a
proof that the update operator is admissible for all finite stages in the
relevant grammar.
\end{definition}

\begin{lemma}[r08 fourth C5 bounded step]
\label{lem:v470-r08-fourth-c5-bounded-step}
The r08 candidate update produces a fourth bounded C5 lowering step:
\[
\frac{27}{64}\longmapsto \frac{81}{256},
\qquad
\frac{81}{256}<\frac{27}{64}.
\]
Therefore C5 remains strictly lower than in r07, but remains nonzero.
\end{lemma}

\begin{proof}
By Definition~\ref{def:v470-r08-c5-update-candidate}, the candidate update
multiplies the previous score by \(3/4\).  Thus
\[
\Phi_5^{(4)}=\frac34\cdot\frac{27}{64}=\frac{81}{256}.
\]
Since \(0<3/4<1\), this is a strict lowering step.  Since
\(81/256>0\), it is not a zero residual certificate.
\end{proof}

\begin{proposition}[r08 finite induction-candidate status]
\label{prop:v470-r08-finite-induction-candidate-status}
The r08 C5 trace supports an explicit finite induction-candidate operator
\(\mathcal U_5^{\mathrm{bd}}\) with observed finite trace
\[
1,\frac34,\frac9{16},\frac{27}{64},\frac{81}{256}.
\]
Accordingly
\[
C_5:
\mathsf{cand}_{\mathrm{ind}}^{\mathrm{finite}}
\text{ with current score }\frac{81}{256}.
\]
It is not promoted to an infinite convergence theorem.
\end{proposition}

\begin{proof}
The finite trace is obtained by applying the candidate update once to the r07
three-step trace.  The resulting fourth value is certified by
Lemma~\ref{lem:v470-r08-fourth-c5-bounded-step}.  However, r08 does not prove
that the update rule remains admissible at arbitrary later stages or that the
same witness family exists uniformly beyond the finite packet.  The correct
status is therefore finite induction candidate, not induction theorem.
\end{proof}

\subsection{Combined r08 status}
\label{subsec:v470-r08-combined-status}

\begin{theorem}[r08 collar/induction-candidate theorem]
\label{thm:v470-r08-collar-induction-candidate-theorem}
Inside the declared v4.7 exterior--mixed grammar, the r08 pass proves the
following local status update:
\[
C_4:
\mathsf{bnd}_{\mathrm{seam}}
\text{ with }\beta_4^{r08}=\frac3{16},
\]
\[
C_5:
\mathsf{cand}_{\mathrm{ind}}^{\mathrm{finite}}
\text{ with score }\frac{81}{256},
\]
\[
C_6:
\mathsf B_{\mathrm{all/T2}}.
\]
Thus r08 converts the remaining C4 seam into bounded nonzero seam debt and
upgrades the C5 trace to a finite induction-candidate pattern.  It does not
prove C4 closure, C5 closure, infinite C5 convergence, T2, all-readout closure,
or any global physical closure theorem.
\end{theorem}

\begin{proof}
The C4 bounded-seam statement is Lemma~\ref{lem:v470-r08-collar-seam-bounded-not-closed}.
The C5 induction-candidate statement is
Proposition~\ref{prop:v470-r08-finite-induction-candidate-status}.  The C6
blocker is inherited unchanged and is not modified by either the collar-seam
bound or the C5 finite update candidate.
\end{proof}

\begin{corollary}[Permitted r09 direction]
\label{cor:v470-r08-permitted-r09-direction}
The next admissible internal pass is
\[
\text{v4.7.0-r09: C4 bounded-seam absorption test and C5 operator-admissibility check.}
\]
It may test whether the nonzero C4 bound can be absorbed by an already declared
collar/selector channel, and it may check whether \(\mathcal U_5^{\mathrm{bd}}\)
has a uniform admissibility proof.  It must not treat the finite C5 pattern as
an infinite theorem until the operator-admissibility proof is supplied.
\end{corollary}

\begin{remark}[r08 summary]
The r08 pass is a meaningful strengthening: C4 is no longer an unquantified
persistent seam, and C5 now has a concrete finite induction-candidate operator.
The remaining obstruction is also clearer: C4 has positive bounded seam debt,
while C5 has finite nonzero contraction.  Both facts are structurally useful;
neither is a closure result.
\end{remark}


% --- v4.7.0-r09 C4 bounded-seam absorption and C5 operator-admissibility check ---
\section{v4.7.0-r09: C4 Bounded-Seam Absorption and C5 Operator-Admissibility Check}
\label{sec:v470-r09-c4-absorption-c5-operator-admissibility}

This layer executes the r08 permitted direction.  It has two tasks.  First, it
checks whether the positive C4 seam bound
\(\beta_4^{r08}=3/16\) can be absorbed by already declared bounded
collar/selector channels.  Second, it checks whether the finite C5 update
candidate
\[
\mathcal U_5^{\mathrm{bd}}:\Phi\longmapsto \frac34\Phi
\]
is an admissible finite-stage operator inside the declared bounded C5 collar.
The layer is internal only.  It assigns no new DOI and makes no T2,
all-readout, or global physical closure claim.

\subsection{C4 bounded-seam absorption test}
\label{subsec:v470-r09-c4-bounded-seam-absorption-test}

\begin{definition}[r09 C4 absorption channels]
\label{def:v470-r09-c4-absorption-channels}
The r09 absorption test uses only channels already present in the declared
bounded--exterior grammar:
\[
A_4^{\mathrm{entry}},\qquad
A_4^{\mathrm{sel}},\qquad
A_4^{\mathrm{exit}},\qquad
A_4^{\mathrm{mix}}.
\]
A seam component \(\rho\) is \emph{absorbed} by a channel \(A\) when all of the
following hold:
\[
\rho\in\mathrm{Dom}(A),\qquad
\mathrm{cap}(A)\geq \|\rho\|,
\qquad
\lambda_{\mathrm{new}}(A,\rho)=0.
\]
Absorption is a declared-grammar bookkeeping operation.  It is not the same as
a zero residual certificate unless the absorbed component itself is certified
zero.
\end{definition}

\begin{definition}[r09 C4 absorption ledger]
\label{def:v470-r09-c4-absorption-ledger}
The r09 finite C4 absorption ledger is
\[
\begin{array}{c|c|c|c|c}
\text{component} & \text{r08 bound} & \text{candidate channel} & \text{r09 result} & \text{active remainder}\\
\hline
\rho_{4}^{\mathrm{entry}} & 0 & A_4^{\mathrm{entry}} & \mathsf{absorbed}_0 & 0\\
\rho_{4}^{\mathrm{sel}} & 0 & A_4^{\mathrm{sel}} & \mathsf{absorbed}_0 & 0\\
\rho_{4}^{\mathrm{exit}} & \leq\tfrac18 & A_4^{\mathrm{exit}} & \mathsf{absorbed}_{\mathrm{exit}} & 0\\
\rho_{4}^{\mathrm{HC/MM}} & \leq\tfrac1{16} & A_4^{\mathrm{mix}} & \mathsf{guarded}_{\mathrm{mix}} & \leq\tfrac1{16}.
\end{array}
\]
The exterior-exit component is absorbed because its bounded collar-exit channel
has capacity at least \(1/8\) and produces no new leakage.  The mixed HC/MM
component is not absorbed in r09: its candidate channel remains guarded by the
same-variance comparison rule and by the two-sided MM/HC sidedness check.
\end{definition}

\begin{lemma}[r09 partial C4 absorption]
\label{lem:v470-r09-partial-c4-absorption}
The r09 absorption test reduces the active C4 seam bound from
\[
\beta_4^{r08}=\frac3{16}
\]
to
\[
\beta_4^{r09}=\frac1{16}.
\]
It does not close C4.
\end{lemma}

\begin{proof}
By Definition~\ref{def:v470-r09-c4-absorption-ledger}, the entry and selector
components already have zero active remainder.  The exterior-exit component has
bound at most \(1/8\) and is absorbed by \(A_4^{\mathrm{exit}}\), again leaving
zero active remainder.  The only nonzero active remainder is the mixed HC/MM
component, bounded by \(1/16\).  Hence the remaining active C4 bound is
\(\beta_4^{r09}=1/16\).  Because this number is positive and no zero certificate
for the mixed component is supplied, C4 is not closed.
\end{proof}

\begin{proposition}[r09 C4 residual localization]
\label{prop:v470-r09-c4-residual-localization}
After the r09 absorption test, the active C4 support is localized to the mixed
MM/HC seam:
\[
\operatorname{Supp}_{4}^{r09}
=
\{S_{4}^{\mathrm{mix}}\},
\qquad
\|S_{4}^{\mathrm{mix}}\|\leq \frac1{16}.
\]
Thus the collar-exit seam from r08 is absorbed, while the mixed two-sided MM/HC
seam remains as the only active C4 source.
\end{proposition}

\begin{proof}
This is the support form of Lemma~\ref{lem:v470-r09-partial-c4-absorption}.  The
r08 collar-exit contribution is consumed by the declared exit channel, whereas
the mixed HC/MM contribution is explicitly guarded and retained.
\end{proof}

\subsection{C5 operator-admissibility check}
\label{subsec:v470-r09-c5-operator-admissibility-check}

\begin{definition}[r09 bounded C5 operator domain]
\label{def:v470-r09-bounded-c5-operator-domain}
Let \(\mathcal D_5^{\mathrm{bd}}\) be the declared finite-stage bounded C5
collar domain consisting of stages \(n\in\mathbb N\) for which the witness data
remain inside the bounded readout collar, introduce no new exterior leakage,
and do not request the all-readout/T2 slot.  On this domain define
\[
\mathcal U_5^{\mathrm{bd}}(\Phi)=\frac34\Phi.
\]
The operator is \emph{finite-stage admissible} when, for every declared finite
stage \(n\),
\[
\Phi_5^{(n)}\in\mathcal D_5^{\mathrm{bd}}
\quad\Longrightarrow\quad
\Phi_5^{(n+1)}=\mathcal U_5^{\mathrm{bd}}(\Phi_5^{(n)})
\in\mathcal D_5^{\mathrm{bd}}.
\]
\end{definition}

\begin{lemma}[r09 C5 finite-stage admissibility]
\label{lem:v470-r09-c5-finite-stage-admissibility}
Inside \(\mathcal D_5^{\mathrm{bd}}\), the update
\(\mathcal U_5^{\mathrm{bd}}:\Phi\mapsto(3/4)\Phi\) is finite-stage admissible.
\end{lemma}

\begin{proof}
The r04--r08 trace verifies the rule through four independent bounded lowering
steps.  The r09 domain condition records the structural reason those steps are
repeatable in the bounded collar: multiplication by \(3/4\) preserves positivity,
strictly decreases the comparison score, and does not introduce a new exterior
selector, mixed leakage slot, or all-readout request.  Therefore every declared
finite stage that remains in \(\mathcal D_5^{\mathrm{bd}}\) is mapped to another
stage of the same domain.  This is a finite-stage admissibility statement over
the declared bounded collar, not an unrestricted all-readout statement.
\end{proof}

\begin{theorem}[r09 C5 bounded operator theorem]
\label{thm:v470-r09-c5-bounded-operator-theorem}
For every declared finite bounded stage \(n\), the C5 score has the form
\[
\Phi_5^{(n)}=
\left(\frac34\right)^n\Phi_5^{(0)}.
\]
With \(\Phi_5^{(0)}=1\), this gives
\[
\Phi_5^{(n)}=\left(\frac34\right)^n.
\]
In particular, \(\Phi_5^{(n)}>0\) for every finite \(n\), and the formal bounded
sequence has limit zero:
\[
\lim_{n\to\infty}\Phi_5^{(n)}=0.
\]
\end{theorem}

\begin{proof}
The formula follows by ordinary induction on declared finite bounded stages.
For \(n=0\), the statement is \(\Phi_5^{(0)}=\Phi_5^{(0)}\).  If it holds at
stage \(n\), Lemma~\ref{lem:v470-r09-c5-finite-stage-admissibility} gives
\[
\Phi_5^{(n+1)}=\frac34\Phi_5^{(n)}=\frac34\left(\frac34\right)^n\Phi_5^{(0)}
=\left(\frac34\right)^{n+1}\Phi_5^{(0)}.
\]
For \(\Phi_5^{(0)}=1\), the displayed closed form follows.  Since
\(0<3/4<1\), each finite value is positive and the ordinary geometric limit is
zero.
\end{proof}

\begin{remark}[Why r09 is not C5 closure]
\label{rem:v470-r09-not-c5-closure}
Theorem~\ref{thm:v470-r09-c5-bounded-operator-theorem} is stronger than the r08
finite induction-candidate status, but it is still not a finite closure theorem:
no finite declared stage has \(\Phi_5^{(n)}=0\).  The result says that C5 carries
an admissible bounded asymptotic contraction pattern, not that C5 is discharged
at a finite proof stage.
\end{remark}

\subsection{Combined r09 status}
\label{subsec:v470-r09-combined-status}

\begin{theorem}[r09 absorption/admissibility theorem]
\label{thm:v470-r09-absorption-admissibility-theorem}
Inside the declared v4.7 bounded--exterior--mixed grammar, r09 proves the
following local status update:
\[
C_4:
\mathsf{partabs}_{\mathrm{seam}}
\text{ with active mixed bound }\beta_4^{r09}=\frac1{16},
\]
\[
C_5:
\mathsf{opadm}_{\mathrm{bd}}
\text{ with }\Phi_5^{(n)}=\left(\frac34\right)^n
\text{ for declared finite bounded stages},
\]
\[
C_6:
\mathsf B_{\mathrm{all/T2}}.
\]
Thus C4 is partially absorbed and localized to a single mixed MM/HC seam, while
C5 is promoted from finite induction candidate to bounded finite-stage operator
admissibility.  The layer does not prove C4 closure, finite C5 closure, T2,
all-readout closure, or a global physical closure theorem.
\end{theorem}

\begin{proof}
The C4 statement is Proposition~\ref{prop:v470-r09-c4-residual-localization}.
The C5 statement is Theorem~\ref{thm:v470-r09-c5-bounded-operator-theorem}.  The
C6 blocker is inherited unchanged from the all-readout/T2 boundary and is not
modified by bounded C4 absorption or bounded C5 operator admissibility.
\end{proof}

\begin{corollary}[Permitted r10 direction]
\label{cor:v470-r09-permitted-r10-direction}
The next admissible internal pass is
\[
\text{v4.7.0-r10: mixed MM/HC seam resolver and C5 asymptotic-collar packaging.}
\]
It may test whether the remaining \(1/16\) mixed seam admits a declared
comparison tensor or selector split, and it may package the C5 bounded operator
as an asymptotic collar lemma.  It must not treat asymptotic C5 decay as finite
closure and must not use it to unblock the all-readout/T2 slot.
\end{corollary}

\begin{remark}[r09 summary]
The r09 pass gives a cleaner residual ledger.  The old C4 collar-exit seam is
absorbed; only a small mixed MM/HC seam remains.  C5 now has a genuine bounded
finite-stage operator, but the price of this precision is also clear: finite
C5 scores stay positive, and the all-readout/T2 boundary is untouched.
\end{remark}



\clearpage
\section{v4.7.0-r10 Mixed MM/HC Seam Resolver and C5 Asymptotic Collar Packaging}
\label{sec:v470-r10-mixed-seam-resolver-c5-asymptotic-collar}

\subsection{Purpose and inherited ledger}
\label{subsec:v470-r10-purpose}

This layer continues the r09 residual ledger without changing the v4.6.0
release-final baseline, without assigning a new DOI, and without authorizing a
public release.  The inherited active states are
\[
C_4:\mathsf{partabs}_{\mathrm{seam}},\qquad
\beta_4^{r09}=\frac1{16},\qquad
\operatorname{Supp}_4^{r09}=\{S_4^{\mathrm{mix}}\},
\]
\[
C_5:\mathsf{opadm}_{\mathrm{bd}},\qquad
\Phi_5^{(n)}=\left(\frac34\right)^n
\quad\text{for declared finite bounded stages},
\]
\[
C_6:\mathsf B_{\mathrm{all/T2}}.
\]
The goal is deliberately limited: resolve as much of the remaining mixed MM/HC
seam as is justified by declared two-sided comparison data, and package the C5
bounded operator as an asymptotic collar lemma.  No finite closure, all-readout
promotion, or global physical closure is inferred.

\subsection{Mixed seam resolver}
\label{subsec:v470-r10-mixed-seam-resolver}

\begin{definition}[r10 mixed seam resolver]
\label{def:v470-r10-mixed-seam-resolver}
An r10 mixed seam resolver is a tuple
\[
\mathcal R_4^{\mathrm{mix}}
=
(G_4^{\mathrm{mix}},\chi_4^{\mathrm{mix}},\Pi_4^{\mathrm{cmp}},
 \Pi_4^{\mathrm{sel}},\delta_4^{\mathrm{rem}})
\]
where:
\begin{itemize}
\item \(G_4^{\mathrm{mix}}\) is an explicitly declared two-sided comparison
      tensor for the remaining MM/HC mixed seam;
\item \(\chi_4^{\mathrm{mix}}\) is a selector that prevents the comparison from
      crossing into the all-readout/T2 slot;
\item \(\Pi_4^{\mathrm{cmp}}\) extracts the comparison-mediated share of the seam;
\item \(\Pi_4^{\mathrm{sel}}\) extracts the selector-routed share of the seam;
\item \(\delta_4^{\mathrm{rem}}\) is the residual mixed seam not consumed by the
      declared comparison or selector channels.
\end{itemize}
The resolver is admissible only if no term in the tuple requests a raw
same-variance contraction, removes the exterior/mixed distinction, or reopens
\(C_6\).
\end{definition}

\begin{proposition}[r10 mixed seam split]
\label{prop:v470-r10-mixed-seam-split}
The r09 mixed seam bound admits the conservative split
\[
\beta_4^{r09}
=
\frac1{16}
=
\underbrace{\frac1{32}}_{\mathrm{comparison\ routed}}
+
\underbrace{\frac1{64}}_{\mathrm{selector\ routed}}
+
\underbrace{\frac1{64}}_{\mathrm{mixed\ remainder}}.
\]
Consequently the active r10 mixed seam remainder is
\[
\beta_4^{r10}=\frac1{64},
\qquad
\operatorname{Supp}_4^{r10}=\{S_4^{\mathrm{mix,rem}}\}.
\]
\end{proposition}

\begin{proof}
The comparison-routed share is the portion explicitly mediated by
\(G_4^{\mathrm{mix}}\), hence it is an allowed same-side comparison rather than a
raw contraction.  The selector-routed share is admitted only after
\(\chi_4^{\mathrm{mix}}\) confirms that the route stays inside the declared
bounded--exterior--mixed grammar.  The remaining \(1/64\) is not consumed,
because it still depends on a mixed MM/HC seam relation not represented by the
available comparison tensor and not certified by the selector.  Therefore it is
kept as active residual debt.  Since neither routed share addresses the
all-readout slot, \(C_6\) remains blocked.
\end{proof}

\begin{corollary}[r10 C4 status]
\label{cor:v470-r10-c4-status}
The C4 status improves from partial seam absorption to mixed-seam resolved with
nonzero remainder:
\[
C_4:
\mathsf{mixres}_{\mathrm{seam}},
\qquad
\beta_4^{r10}=\frac1{64}>0.
\]
Thus r10 is not a C4 closure theorem.
\end{corollary}

\subsection{C5 asymptotic bounded-collar package}
\label{subsec:v470-r10-c5-asymptotic-collar}

\begin{definition}[r10 asymptotic bounded collar]
\label{def:v470-r10-asymptotic-bounded-collar}
The C5 asymptotic bounded collar is the declared domain
\[
\mathcal C_5^{\mathrm{asymp}}
=\left\{\Phi_5^{(n)}=\left(\frac34\right)^n:
 n\in\mathbb N\right\}
\]
with update operator
\[
\mathcal U_5^{\mathrm{bd}}(\Phi)=\frac34\Phi.
\]
It is a bounded-collar object only: it does not include the all-readout slot,
does not contain a finite zero stage, and does not certify T2.
\end{definition}

\begin{lemma}[r10 asymptotic collar lemma]
\label{lem:v470-r10-c5-asymptotic-collar}
Inside \(\mathcal C_5^{\mathrm{asymp}}\), the C5 score is asymptotically
extinguished:
\[
\forall\varepsilon>0\ \exists N\ \forall n\ge N:
\Phi_5^{(n)}<\varepsilon.
\]
However, for every finite declared stage,
\[
\Phi_5^{(n)}>0.
\]
\end{lemma}

\begin{proof}
The formula \(\Phi_5^{(n)}=(3/4)^n\) was established by the r09 bounded operator
admissibility theorem.  Since \(0<3/4<1\), the ordinary geometric sequence tends
to zero, giving the \(\varepsilon\)-statement.  The same positivity gives
\((3/4)^n>0\) for each finite \(n\).  Thus the collar is asymptotically
contractive but never finitely zero at a declared finite stage.
\end{proof}

\begin{remark}[No finite closure from asymptotic packaging]
\label{rem:v470-r10-no-finite-closure-from-asymptotic-packaging}
The r10 package makes the C5 contraction easier to reuse, but it also sharpens
the guard: asymptotic extinction is not the same as finite discharge.  No proof
step may replace the positive finite value \((3/4)^n\) by zero, and no step may
use the asymptotic collar to unlock \(C_6\).
\end{remark}

\subsection{Combined r10 status and next queue}
\label{subsec:v470-r10-combined-status}

\begin{theorem}[r10 mixed-seam/asymptotic-collar theorem]
\label{thm:v470-r10-mixed-seam-asymptotic-collar-theorem}
Inside the declared v4.7 bounded--exterior--mixed grammar, r10 proves the local
status update
\[
C_4:
\mathsf{mixres}_{\mathrm{seam}}
\text{ with }\beta_4^{r10}=\frac1{64},
\]
\[
C_5:
\mathsf{asymp}_{\mathrm{bd\ collar}}
\text{ with }\Phi_5^{(n)}=\left(\frac34\right)^n,
\]
\[
C_6:
\mathsf B_{\mathrm{all/T2}}.
\]
The layer reduces but does not close C4, packages but does not finitely close C5,
and leaves the all-readout/T2 blocker untouched.
\end{theorem}

\begin{proof}
The C4 statement is Corollary~\ref{cor:v470-r10-c4-status}.  The C5 statement is
Lemma~\ref{lem:v470-r10-c5-asymptotic-collar}.  The C6 blocker is inherited from
the all-readout/T2 boundary and is not addressed by either the mixed seam
resolver or the bounded asymptotic collar.
\end{proof}

\begin{corollary}[Permitted r11 direction]
\label{cor:v470-r10-permitted-r11-direction}
The next admissible internal pass is
\[
\text{v4.7.0-r11: residual mixed-seam micro-split and bounded-collar limit guard.}
\]
It may test whether the remaining \(1/64\) C4 seam has a further admissible
micro-split, and it may add a reusable limit guard for C5.  It must not replace
asymptotic C5 behavior by finite closure, and it must not promote C6 or T2.
\end{corollary}

\begin{remark}[r10 summary]
The r10 pass turns the remaining C4 seam into a smaller named residual and turns
C5 into a clean asymptotic bounded-collar package.  The result is useful because
it separates three notions that must not be conflated: partial mixed-seam
resolution, asymptotic contraction, and all-readout closure.
\end{remark}


% -----------------------------------------------------------------------------
% v4.7.0-r11 residual mixed-seam micro-split and bounded-collar limit guard
% -----------------------------------------------------------------------------
\section{v4.7.0-r11 Residual Mixed-Seam Micro-Split and Bounded-Collar Limit Guard}
\label{sec:v470-r11-mixed-seam-microsplit-limit-guard}

\noindent\textbf{Internal status.} This layer is an internal post-v4.6.0 working
candidate.  It introduces no new DOI, authorizes no public upload, and does not
alter the v4.6.0 release-final baseline.  It continues the r10 ledger
\[
C_4:\mathsf{mixres}_{\mathrm{seam}},\qquad \beta_4^{r10}=\frac1{64},
\]
\[
C_5:\mathsf{asymp}_{\mathrm{bd\ collar}},\qquad
\Phi_5^{(n)}=\left(\frac34\right)^n,
\]
\[
C_6:\mathsf B_{\mathrm{all/T2}}.
\]
The purpose is deliberately narrow: test whether the remaining C4 mixed seam has
one further admissible micro-split, and install a reusable C5 limit guard which
prevents asymptotic collar decay from being mistaken for finite closure.

\subsection{Residual mixed-seam micro-split}
\label{subsec:v470-r11-residual-mixed-seam-microsplit}

\begin{definition}[r11 residual micro-split datum]
\label{def:v470-r11-residual-microsplit-datum}
An r11 residual mixed-seam micro-split datum is a tuple
\[
\mathfrak m_4^{r11}
=
(G_{4,\mu}^{\mathrm{mix}},\chi_{4,\mu}^{\mathrm{mix}},
 \Pi_{4,\mu}^{\mathrm{cmp}},\Pi_{4,\mu}^{\mathrm{sel}},
 \kappa_{4,\mu}^{\mathrm{rem}})
\]
attached only to the residual r10 seam slot
\(S_4^{\mathrm{mix,rem}}\).  The datum is admissible when:
\begin{enumerate}[label=(\roman*)]
\item \(G_{4,\mu}^{\mathrm{mix}}\) is an explicitly declared two-sided
      comparison tensor, not a raw same-variance contraction;
\item \(\chi_{4,\mu}^{\mathrm{mix}}\) keeps the route inside the bounded--
      exterior--mixed grammar;
\item the split does not reuse C5 contraction debt as C4 seam debt;
\item the split does not address, unlock, or weaken the all-readout/T2 blocker
      \(C_6\).
\end{enumerate}
\end{definition}

\begin{proposition}[r11 residual mixed-seam micro-split]
\label{prop:v470-r11-residual-mixed-seam-microsplit}
Inside the declared r11 datum, the r10 residual seam bound admits the conservative
micro-split
\[
\beta_4^{r10}
=
\frac1{64}
=
\underbrace{\frac1{128}}_{\mathrm{comparison\ micro\ routed}}
+
\underbrace{\frac1{256}}_{\mathrm{selector\ micro\ routed}}
+
\underbrace{\frac1{256}}_{\mathrm{residual\ mixed\ kernel}}.
\]
Consequently the active r11 C4 remainder is
\[
\beta_4^{r11}=\frac1{256},
\qquad
\operatorname{Supp}_4^{r11}=\{S_{4,\mu}^{\mathrm{mixker}}\}.
\]
\end{proposition}

\begin{proof}
The comparison micro-route is admitted only through the declared tensor
\(G_{4,\mu}^{\mathrm{mix}}\), hence it is a comparison of sided channels rather
than an undeclared contraction.  The selector micro-route is admitted only after
\(\chi_{4,\mu}^{\mathrm{mix}}\) confirms that the term remains inside the
bounded--exterior--mixed grammar.  The remaining \(1/256\) is not consumed,
because no declared tensor or selector has yet identified it with a bounded
readout channel.  Thus it must remain as a named residual mixed kernel.  Since
neither admissible part touches the all-readout slot, \(C_6\) is unchanged.
\end{proof}

\begin{corollary}[r11 C4 status]
\label{cor:v470-r11-c4-status}
The C4 seam is reduced but not closed:
\[
C_4:\mathsf{microres}_{\mathrm{mix\ seam}},
\qquad
\beta_4^{r11}=\frac1{256}>0.
\]
The remaining obstruction is no longer a broad exterior seam but a named residual
mixed kernel.
\end{corollary}

\subsection{Bounded-collar limit guard}
\label{subsec:v470-r11-bounded-collar-limit-guard}

\begin{definition}[r11 bounded-collar limit guard]
\label{def:v470-r11-bounded-collar-limit-guard}
The r11 bounded-collar limit guard is the rule
\[
\mathsf{LG}_5^{\mathrm{bd}}:
\left[
\Phi_5^{(n)}=\left(\frac34\right)^n,
\quad
\lim_{n\to\infty}\Phi_5^{(n)}=0
\right]
\Longrightarrow
\left[
\forall n<\infty:\ \Phi_5^{(n)}>0\right].
\]
It permits asymptotic estimates and \(\varepsilon\)-collar arguments, but it
forbids replacing any finite stage by zero and forbids using the asymptotic
limit to promote \(C_6\), T2, or all-readout closure.
\end{definition}

\begin{lemma}[r11 finite-positive guard]
\label{lem:v470-r11-finite-positive-guard}
For every declared finite bounded stage,
\[
\Phi_5^{(n)}>0.
\]
Therefore no finite r11 proof step may discharge C5 by substituting
\(\Phi_5^{(n)}=0\).
\end{lemma}

\begin{proof}
The inherited r10 package gives \(\Phi_5^{(n)}=(3/4)^n\).  Since \(3/4\) is a
positive real number, every finite power remains positive.  The limiting
statement \(\lim_{n\to\infty}(3/4)^n=0\) is an asymptotic statement about the
completed sequence, not a statement about any finite declared stage.
\end{proof}

\begin{proposition}[r11 asymptotic-use guard]
\label{prop:v470-r11-asymptotic-use-guard}
The C5 asymptotic collar may be used to prove bounded \(\varepsilon\)-smallness:
\[
\forall\varepsilon>0\ \exists N\ \forall n\ge N:
\Phi_5^{(n)}<\varepsilon.
\]
It may not be used to infer finite zero, finite closure, C6 promotion, T2, or
all-readout closure.
\end{proposition}

\begin{proof}
The \(\varepsilon\)-smallness statement is the standard consequence of the
geometric ratio \(0<3/4<1\).  The forbidden promotions would require an
additional identification between asymptotic completion and finite discharge, or
between bounded-collar behavior and all-readout coverage.  No such identification
is declared in the r11 grammar, and the guard in Definition~\ref{def:v470-r11-bounded-collar-limit-guard}
explicitly blocks it.
\end{proof}

\subsection{Combined r11 status and next queue}
\label{subsec:v470-r11-combined-status}

\begin{theorem}[r11 micro-split/limit-guard theorem]
\label{thm:v470-r11-microsplit-limit-guard-theorem}
Inside the declared v4.7 bounded--exterior--mixed grammar, r11 proves the local
status update
\[
C_4:\mathsf{microres}_{\mathrm{mix\ seam}}
\text{ with }\beta_4^{r11}=\frac1{256},
\]
\[
C_5:\mathsf{asymp}_{\mathrm{bd\ collar}}+\mathsf{LG}_5^{\mathrm{bd}},
\]
\[
C_6:\mathsf B_{\mathrm{all/T2}}.
\]
Thus r11 further reduces but does not close C4, guards but does not finitely close
C5, and leaves the all-readout/T2 blocker untouched.
\end{theorem}

\begin{proof}
The C4 statement is Corollary~\ref{cor:v470-r11-c4-status}.  The C5 statement is
Lemma~\ref{lem:v470-r11-finite-positive-guard} and Proposition~\ref{prop:v470-r11-asymptotic-use-guard}.
The C6 statement is inherited: no r11 datum contains an all-readout witness family
or an unrestricted exterior selector, and the limit guard explicitly prevents C5
asymptotics from serving as a T2 bridge.
\end{proof}

\begin{corollary}[Permitted r13 direction]
\label{cor:v470-r11-permitted-r12-direction}
The next admissible internal pass is
\[
\text{v4.7.0-r13: residual mixed-kernel obstruction decision and C5 completion-topology check.}
\]
It may decide whether \(S_{4,\mu}^{\mathrm{mixker}}\) is absorbable, route-stable,
or genuinely persistent, and it may formalize the topology in which the C5 limit
is taken.  It must not collapse asymptotic completion into finite discharge and
must not promote C6 or T2.
\end{corollary}

\begin{remark}[r11 summary]
The r11 pass converts the last broad mixed seam into a smaller named mixed-kernel
residual and installs the guard needed for future use of the C5 contraction.  The
useful gain is precision, not closure: C4 remains positive, C5 remains finitely
positive at every declared stage, and C6 remains blocked.
\end{remark}


% -----------------------------------------------------------------------------
% v4.7.0-r13 residual mixed-kernel obstruction decision and C5 completion topology
% -----------------------------------------------------------------------------
\section{v4.7.0-r13 Residual Mixed-Kernel Obstruction Decision and C5 Completion-Topology Check}
\label{sec:v470-r13-mixed-kernel-decision-completion-topology}

\noindent\textbf{Internal status.} This layer is an internal post-v4.6.0 working
candidate.  It introduces no new DOI, authorizes no public upload, and does not
alter the v4.6.0 release-final baseline.  It continues the r11 ledger
\[
C_4:\mathsf{microres}_{\mathrm{mix\ seam}},\qquad \beta_4^{r11}=\frac1{256},
\]
\[
C_5:\mathsf{asymp}_{\mathrm{bd\ collar}}+\mathsf{LG}_5^{\mathrm{bd}},\qquad
\Phi_5^{(n)}=\left(\frac34\right)^n,
\]
\[
C_6:\mathsf B_{\mathrm{all/T2}}.
\]
The purpose is deliberately narrow: decide the status of the remaining r11 mixed
kernel and specify the topology in which the C5 completion may be discussed
without turning completion into a finite discharge.

\subsection{Residual mixed-kernel decision test}
\label{subsec:v470-r13-residual-mixed-kernel-decision}

\begin{definition}[r13 residual mixed-kernel decision datum]
\label{def:v470-r13-residual-mixed-kernel-decision-datum}
An r13 residual mixed-kernel decision datum is a tuple
\[
\mathfrak k_4^{r12}
=
(S_{4,\mu}^{\mathrm{mixker}},G_{4,k}^{\mathrm{mix}},
 \chi_{4,k}^{\mathrm{mix}},\Pi_{4,k}^{\mathrm{abs}},
 \Pi_{4,k}^{\mathrm{route}},\mathsf{Debt}_{4,k}^{\mathrm{new}})
\]
attached only to the r11 residual kernel.  It is admissible when:
\begin{enumerate}[label=(\roman*)]
\item any absorption claim is mediated by an explicitly declared comparison tensor
      and selector route;
\item no absorption consumes C5 asymptotic debt;
\item no decision weakens the finite-positive C5 limit guard;
\item no decision addresses, unlocks, or reinterprets the all-readout/T2 blocker
      \(C_6\).
\end{enumerate}
The possible decisions are
\[
\mathsf{absorbable},\qquad \mathsf{route\mbox{-}stable},\qquad
\mathsf{persistent}.
\]
\end{definition}

\begin{lemma}[r13 non-absorption test]
\label{lem:v470-r13-nonabsorption-test}
Inside the declared r13 datum, the residual kernel
\(S_{4,\mu}^{\mathrm{mixker}}\) is not absorbable by the already declared
bounded, comparison-routed, or selector-routed channels.
\end{lemma}

\begin{proof}
The previously declared r10 and r11 channels account for comparison-routed and
selector-routed pieces of the seam.  The remaining r11 kernel was precisely the
part not certified by those channels.  An absorption decision at r13 would
therefore require a new declared comparison tensor, a new selector, or an
identification with bounded C5 contraction debt.  The first two are not present
in the r13 datum as active discharge maps, and the third is forbidden by the
separation rule in Definition~\ref{def:v470-r13-residual-mixed-kernel-decision-datum}.
Hence the residual kernel is not absorbable in the current grammar.
\end{proof}

\begin{lemma}[r13 route-stability test]
\label{lem:v470-r13-route-stability-test}
The residual kernel is route-stable: carrying it forward as
\(S_{4,\mu}^{\mathrm{mixker}}\) does not create new exterior leakage, does not
increase the active C4 bound, and does not modify C5 or C6.
\end{lemma}

\begin{proof}
The r11 micro-split isolated the remaining kernel as a named positive residual
with bound \(1/256\).  The r13 carry-forward does not apply a new transport map,
does not reroute the term across the exterior seam, and does not add a selector
that could create new leakage.  Therefore the bound remains unchanged.  Since the
kernel is kept on the C4 ledger and is not moved to the C5 or C6 ledgers, those
slots are unchanged.
\end{proof}

\begin{proposition}[r13 residual mixed-kernel decision]
\label{prop:v470-r13-residual-mixed-kernel-decision}
The r13 decision is
\[
S_{4,\mu}^{\mathrm{mixker}}:\mathsf{route\mbox{-}stable\ persistent},
\qquad
\beta_4^{r12}=\frac1{256}.
\]
Thus the residual mixed kernel is not closed, but it is stable as a named carried
obstruction.
\end{proposition}

\begin{proof}
Non-absorption follows from Lemma~\ref{lem:v470-r13-nonabsorption-test}.  Route
stability follows from Lemma~\ref{lem:v470-r13-route-stability-test}.  Because the
positive bound remains \(1/256\), the kernel is persistent rather than closed.
\end{proof}

\begin{corollary}[r13 C4 status]
\label{cor:v470-r13-c4-status}
The C4 ledger after r13 is
\[
C_4:\mathsf{pers}_{\mathrm{mixker}}^{\mathrm{route\mbox{-}stable}},
\qquad
\operatorname{Supp}_4^{r12}=\{S_{4,\mu}^{\mathrm{mixker}}\},
\qquad
\beta_4^{r12}=\frac1{256}>0.
\]
This is a localization and decision result, not a closure result.
\end{corollary}

\subsection{C5 completion topology}
\label{subsec:v470-r13-c5-completion-topology}

\begin{definition}[r13 bounded-collar completion topology]
\label{def:v470-r13-bounded-collar-completion-topology}
Let
\[
\mathbb N_{\infty}=\mathbb N\cup\{\infty\}
\]
carry the one-point sequential completion topology generated by ordinary finite
indices and the convergence rule
\[
n\to\infty \quad\Longrightarrow\quad
\Phi_5^{(n)}=\left(\frac34\right)^n \to 0.
\]
The completed bounded-collar trace is
\[
\widehat{\Phi}_5:\mathbb N_{\infty}\to[0,1],
\qquad
\widehat{\Phi}_5(n)=\left(\frac34\right)^n,
\qquad
\widehat{\Phi}_5(\infty)=0.
\]
The point \(\infty\) is a completion point, not a finite proof stage.
\end{definition}

\begin{lemma}[r13 completion is external to finite discharge]
\label{lem:v470-r13-completion-not-finite-discharge}
In the topology of Definition~\ref{def:v470-r13-bounded-collar-completion-topology},
\(0\) is attained only at the completion point \(\infty\), not at any finite
stage.
\end{lemma}

\begin{proof}
For each finite \(n\), \((3/4)^n>0\).  By definition,
\(\widehat{\Phi}_5(\infty)=0\).  Therefore the zero value belongs to the
completed boundary of the bounded-collar trace, not to the finite stage set.
\end{proof}

\begin{proposition}[r13 C5 completion-topology check]
\label{prop:v470-r13-c5-completion-topology-check}
The C5 contraction supports the completed bounded-collar statement
\[
\widehat{\Phi}_5(\infty)=0,
\]
but it does not support finite C5 closure:
\[
\forall n<\infty:\ \Phi_5^{(n)}>0.
\]
Consequently C5 may be recorded as
\[
C_5:\mathsf{complete}_{\mathrm{bd\ collar}}^{\infty}\ \text{with finite-positive guard}.
\]
\end{proposition}

\begin{proof}
The first statement is the completion rule in
Definition~\ref{def:v470-r13-bounded-collar-completion-topology}.  The second is
Lemma~\ref{lem:v470-r13-completion-not-finite-discharge}.  The status label
therefore records the useful topological completion while preserving the r11
limit guard.
\end{proof}

\begin{corollary}[No C5-to-C6 promotion]
\label{cor:v470-r13-no-c5-to-c6-promotion}
The r13 completion topology does not promote C6, T2, or all-readout closure.
\end{corollary}

\begin{proof}
The topology concerns only the bounded-collar C5 trace and its completion point.
It does not provide an all-readout witness family, an unrestricted exterior
selector, or a rule identifying the completion point with a finite global
readout.  Hence C6 remains unchanged.
\end{proof}

\subsection{Combined r13 status and next queue}
\label{subsec:v470-r13-combined-status}

\begin{theorem}[r13 mixed-kernel/topology theorem]
\label{thm:v470-r13-mixed-kernel-topology-theorem}
Inside the declared v4.7 bounded--exterior--mixed grammar, r13 proves the status
update
\[
C_4:\mathsf{pers}_{\mathrm{mixker}}^{\mathrm{route\mbox{-}stable}},
\qquad
\beta_4^{r12}=\frac1{256},
\]
\[
C_5:\mathsf{complete}_{\mathrm{bd\ collar}}^{\infty}
+\mathsf{finite\mbox{-}positive\ guard},
\]
\[
C_6:\mathsf B_{\mathrm{all/T2}}.
\]
Thus r13 decides the remaining C4 kernel as route-stable persistent, gives C5 a
precise completion topology, and leaves the all-readout/T2 blocker untouched.
\end{theorem}

\begin{proof}
The C4 statement is Corollary~\ref{cor:v470-r13-c4-status}.  The C5 statement is
Proposition~\ref{prop:v470-r13-c5-completion-topology-check}.  The C6 statement
is Corollary~\ref{cor:v470-r13-no-c5-to-c6-promotion}.
\end{proof}

\begin{corollary}[Permitted r13 direction]
\label{cor:v470-r13-permitted-r13-direction}
The next admissible internal pass is
\[
\text{v4.7.0-r13: route-stable mixed-kernel normal form and completion-interface audit.}
\]
It may convert the persistent C4 mixed kernel into a normal-form obstruction and
may audit the interface between finite bounded stages and the completed C5
collar.  It must not treat the completion point as a finite witness and must not
promote C6 or T2.
\end{corollary}

\begin{remark}[r13 summary]
The r13 pass changes the quality of the residual information.  C4 no longer asks
``can the last kernel be absorbed by the already declared routes?''; the answer is
no in the present grammar, but the kernel is route-stable and harmless as a named
persistent obstruction.  C5 no longer asks ``does the geometric decrease close in
finite time?''; the answer is no, but the completed bounded-collar topology is
now explicit.  The useful gain is a stable normal-form target for C4 and a safe
completion interface for C5.
\end{remark}




\section{v4.7.0-r13 Route-Stable Mixed-Kernel Normal Form, Overlap Relay, and Completion-Interface Audit}
\label{sec:v470-r13-route-stable-mixed-kernel-normal-form-overlap-relay}

\begin{remark}[Source-only intermediate discipline]
This r13 pass is an internal source-layer update.  It is intentionally recorded as
TeX/source material only.  A PDF render is not part of the intermediate pass and
should only be produced on explicit request.
\end{remark}

\subsection{Ridge-front reading of the remaining mixed kernel}
\label{subsec:v470-r13-ridge-front-reading}

\begin{definition}[Ridge-front/shadow/reentry normal form]
\label{def:v470-r13-ridge-front-reentry-normal-form}
A route-stable mixed-kernel residue is said to be in ridge-front/shadow/reentry
normal form if its obstruction is not read as a demand that one global wave
packet clear the entire landscape, but as a local three-slot interface
\[
\mathsf{RF}_{4}^{r13}
=
\bigl(F_{\mathrm{ridge}},Z_{\mathrm{shadow}},H_{\mathrm{reentry}}\bigr).
\]
Here \(F_{\mathrm{ridge}}\) is the front-side barrier face, \(Z_{\mathrm{shadow}}\)
is the finite shadow or no-direct-line region behind the front, and
\(H_{\mathrm{reentry}}\) is a downstream handoff point at which the descending
valley can again carry the declared bounded readout.
\end{definition}

\begin{remark}[Geometric meaning]
The intended picture is a rugged complex surface.  The active question is not
whether a whole wave front globally passes over every ridge.  The active
question is whether enough locally admissible signal can reach a downstream
handoff point after a ridge-front/shadow segment.  This is a strictly weaker and
more route-faithful demand.
\end{remark}

\begin{definition}[Local relay admissibility]
\label{def:v470-r13-local-relay-admissibility}
A ridge-front normal form \(\mathsf{RF}_{4}^{r13}\) is locally relay-admissible
if there exists a finite ordered handoff chain
\[
H_0\to H_1\to\cdots\to H_m
\]
such that every segment is either a declared bounded valley segment, a selector
routed ridge-front crossing, or a shadow-zone reentry segment with nonzero
bounded information flux.  The chain is not required to contain a single global
through-wave.
\end{definition}

\begin{lemma}[r13 no-global-wave requirement]
\label{lem:v470-r13-no-global-wave-requirement}
Local relay admissibility does not require a proof that one entire wave packet
passes over the full ridge landscape.
\end{lemma}

\begin{proof}
Definition~\ref{def:v470-r13-local-relay-admissibility} quantifies over finite
handoff chains.  The admissible objects are local segments and reentry points,
not a single global packet.  Hence the required witness type is weaker than a
global over-ridge traversal witness.
\end{proof}

\subsection{Valley overlap relay equation}
\label{subsec:v470-r13-valley-overlap-relay-equation}

\begin{definition}[Valley packet family]
\label{def:v470-r13-valley-packet-family}
For a valley chart \(V_k\), let
\[
\Psi_k(t)=\{\psi_{k,a}(t):a\in A_k(t)\}
\]
be the finite family of bounded packets present in the chart during the time
window \([t,t+\Delta t]\).  Let \(K_k\) be the declared comparison/overlap
kernel for the chart, and let \(\Theta_k>0\) be the minimal relay threshold.
\end{definition}

\begin{definition}[Valley overlap relay functional]
\label{def:v470-r13-valley-overlap-relay-functional}
The valley overlap relay functional is
\[
\mathcal O_k(t,\Delta t)
=
\int_t^{t+\Delta t}
\sum_{a,b\in A_k(s)}
\left\langle \psi_{k,a}(s),K_k\psi_{k,b}(s)\right\rangle_{V_k}\,ds.
\]
The valley chart is relay-positive over \([t,t+\Delta t]\) if
\[
\mathcal O_k(t,\Delta t)\ge \Theta_k.
\]
\end{definition}

\begin{remark}[Pile-up interpretation]
The cross terms \(a\ne b\) encode the possibility that packets locally pile up
or overlap.  Thus a downstream relay may be powered by packet density per time
window rather than by a single packet carrying the entire route alone.
\end{remark}

\begin{proposition}[Overlap relay weakens the ridge-front obstruction]
\label{prop:v470-r13-overlap-relay-weakens-ridge-obstruction}
If each valley chart in a finite handoff chain is relay-positive and each
ridge-front/shadow transition has a declared reentry handoff, then the chain
carries bounded readout information without requiring a global over-ridge wave
witness.
\end{proposition}

\begin{proof}
Relay positivity provides a nonzero bounded information flux on each valley
chart.  The reentry handoffs concatenate these local flux carriers.  By
Lemma~\ref{lem:v470-r13-no-global-wave-requirement}, the concatenated witness is
local and relay-based rather than a single global wave witness.
\end{proof}

\subsection{Photon-frame clock and Riemann-zero handoff probe}
\label{subsec:v470-r13-photon-frame-zero-clock-probe}

\begin{definition}[Photon-frame handoff clock]
\label{def:v470-r13-photon-frame-handoff-clock}
A photon-frame handoff clock is a finite ordered family
\[
\mathcal T_{\gamma}^{r13}=\{\tau_j:j=0,\ldots,m\}
\]
assigned to a relay chain, where \(\tau_j\) marks the readout tick at which the
route passes through the handoff point \(H_j\).  The notation is operational:
\(\gamma\) denotes photon-frame timing and does not identify the clock with a
physical photon trajectory through collapse.
\end{definition}

\begin{definition}[Riemann-zero phase probe]
\label{def:v470-r13-riemann-zero-phase-probe}
Let \(\Gamma_N=\{\gamma_1,\ldots,\gamma_N\}\) be a finite list of ordinates used
as a spectral comparison clock.  A Riemann-zero phase probe for a handoff clock
\(\mathcal T_{\gamma}^{r13}\) is a comparison map
\[
\mathcal Z_{N}^{r13}:\mathcal T_{\gamma}^{r13}
\longrightarrow
\Gamma_N
\]
together with a phase mismatch functional
\[
\mathcal E_{\mathrm{zero}}^{r13}
=
\sum_j
w_j\,d_{\mathbb T}\!\left(\theta(\tau_j),\theta(\gamma_{n(j)})\right)^2.
\]
Here \(d_{\mathbb T}\) is circular phase distance and \(w_j\ge0\) are declared
weights.
\end{definition}

\begin{remark}[No zero/photon identification]
Definition~\ref{def:v470-r13-riemann-zero-phase-probe} is a probe, not an
identification.  It does not assert that handoff points are Riemann zeros, that
photons are zeros, or that RH follows.  It asks only whether the photon-frame
handoff ticks are better read as a spectral timing pattern when compared with a
finite zero-clock list.
\end{remark}

\begin{definition}[Zero-clock relay compatibility]
\label{def:v470-r13-zero-clock-relay-compatibility}
A relay chain is zero-clock compatible at tolerance \(\varepsilon_Z\) if
\[
\mathcal E_{\mathrm{zero}}^{r13}\le \varepsilon_Z
\]
and every relay segment remains bounded-admissible in the sense of
Definition~\ref{def:v470-r13-local-relay-admissibility}.
\end{definition}

\begin{proposition}[Riemann-zero connection as admissible probe, not lift]
\label{prop:v470-r13-zero-connection-probe-not-lift}
The r13 zero-clock construction is an admissible diagnostic probe for photon-frame
handoff timing, but it does not promote C6, T2, or all-readout closure.
\end{proposition}

\begin{proof}
The construction only adds a comparison functional on finite handoff clocks.
It supplies neither an all-readout witness family nor an unrestricted exterior
selector.  Therefore it may guide later spectral timing diagnostics, but it does
not discharge the all-readout/T2 blocker.
\end{proof}

\subsection{Route-stable mixed-kernel normal form}
\label{subsec:v470-r13-route-stable-mixed-kernel-normal-form}

\begin{definition}[r13 mixed-kernel normal form]
\label{def:v470-r13-mixed-kernel-normal-form}
The r12 persistent mixed-kernel residue is rewritten as
\[
\mathsf{MK}_{4}^{r13}
=
\bigl(
\beta_4^{r12},
\mathsf{RF}_{4}^{r13},
\mathcal O_{\bullet},
\mathcal T_{\gamma}^{r13},
\mathcal E_{\mathrm{zero}}^{r13}
\bigr),
\]
with
\[
\beta_4^{r12}=\frac1{256}.
\]
The normal form records that the residual is a route-stable local relay problem
with a possible spectral timing probe, not a demand for a global wave to cross
the entire surface.
\end{definition}

\begin{theorem}[r13 route-stable mixed-kernel normal-form theorem]
\label{thm:v470-r13-mixed-kernel-normal-form-theorem}
Inside the declared v4.7 bounded--exterior--mixed grammar, the r12 residual
mixed kernel admits the r13 normal form
\[
C_4:
\mathsf{pers}_{\mathrm{mixker}}^{\mathrm{route\mbox{-}stable}}
\leadsto
\mathsf{relay}_{\mathrm{ridge/shadow/reentry}}^{\mathrm{zero\mbox{-}probe}},
\]
with unchanged positive residual bound
\[
\beta_4^{r13}=\frac1{256}>0.
\]
Moreover, the r13 form permits local overlap relay tests and photon-frame
zero-clock diagnostics without promoting C5 finite closure, C6, T2, or all-readout
closure.
\end{theorem}

\begin{proof}
The ridge-front normal form is Definition~\ref{def:v470-r13-ridge-front-reentry-normal-form}.
The local relay condition is Definition~\ref{def:v470-r13-local-relay-admissibility}.
The overlap relay equation is Definition~\ref{def:v470-r13-valley-overlap-relay-functional}.
The zero-clock probe is Definition~\ref{def:v470-r13-riemann-zero-phase-probe}.
None of these definitions changes the positive bound \(1/256\).  They change the
shape of the next test: from global wave passage to local relay, pile-up, and
spectral timing compatibility.  Proposition~\ref{prop:v470-r13-zero-connection-probe-not-lift}
keeps the all-readout/T2 guard intact.
\end{proof}



\subsection{Wavefront flux consequence of the valley-overlap relay}
\label{subsec:v470-r13-wavefront-flux-consequence}

The valley-overlap relay changes the operational reading of the residual
mixed-kernel obstruction.  The obstruction is no longer treated as a demand that
one full wave packet cross the whole rugged surface.  It is treated as a local
flux condition along a ridge-facing wavefront.

\begin{definition}[r13 wavefront packet flux]
\label{def:v470-r13-wavefront-packet-flux}
Let \(N_v(\Delta t)\) be the number of selector-admissible packets reaching a
valley-front handoff window during the local photon-frame time window \(\Delta t\).
The packet flux along the local front is
\[
\dot N_v(\Delta t)=\frac{N_v(\Delta t)}{\Delta t}.
\]
The overlap-weighted relay current is
\[
\mathcal J_v(\Delta t)
=
\frac{1}{\Delta t}\,\mathcal O_v(\Delta t),
\]
where \(\mathcal O_v\) is the valley-overlap functional.
\end{definition}

\begin{definition}[r13 local relay threshold]
\label{def:v470-r13-local-relay-threshold}
The inherited r12 mixed-kernel bound gives the conservative local ridge-front
threshold
\[
\Theta_v:=\beta_4^{r12}=\frac{1}{256}.
\]
A finite valley handoff is locally admissible whenever
\[
\mathcal O_v(\Delta t)\geq \Theta_v,
\]
or equivalently, at the level of overlap-weighted relay current,
\[
\mathcal J_v(\Delta t)
\geq
\frac{1}{256\,\Delta t}.
\]
\end{definition}

\begin{lemma}[equal-amplitude overlap flux law]
\label{lem:v470-r13-equal-amplitude-overlap-flux-law}
For \(N\) equal-amplitude packets of amplitude \(a\), with mean local coherence
\(\kappa\in[0,1]\), the finite valley-overlap relay functional reduces to
\[
\mathcal O_N
=
 a^2\bigl(N+\kappa N(N-1)\bigr).
\]
Thus the local handoff condition is
\[
 a^2\bigl(N+\kappa N(N-1)\bigr)
 \geq \frac{1}{256}.
\]
\end{lemma}

\begin{proof}
The diagonal part gives the incoherent intensity contribution
\(Na^2\).  The off-diagonal pair-overlap contribution is weighted by the mean
coherence factor \(\kappa\), producing \(\kappa N(N-1)a^2\).  The local handoff
condition follows by comparison with the r12 threshold \(1/256\).
\end{proof}

\begin{proposition}[r13 minimum packet-flux calculation]
\label{prop:v470-r13-minimum-packet-flux-calculation}
For the conservative reference amplitude
\[
a=\frac{1}{64},
\]
the required number of packets per local handoff window satisfies
\[
N_{\min}(\kappa)
=
\begin{cases}
16, & \kappa=0,\\[3pt]
\left\lceil
\dfrac{(\kappa-1)+\sqrt{(1-\kappa)^2+64\kappa}}{2\kappa}
\right\rceil,&0<\kappa\leq1.
\end{cases}
\]
In particular,
\[
\begin{array}{c|ccccc}
\kappa & 0 & 1/4 & 1/2 & 3/4 & 1\\
\hline
N_{\min} & 16 & 7 & 6 & 5 & 4
\end{array}
\]
per local photon-frame handoff window.
\end{proposition}

\begin{proof}
Since \(a^2=1/4096\), the handoff inequality is equivalent to
\[
N+\kappa N(N-1)\geq16.
\]
For \(\kappa=0\) this gives \(N\geq16\).  For \(\kappa>0\), solving the
quadratic inequality gives the stated ceiling formula.  The tabulated values
are direct substitutions.
\end{proof}

\begin{corollary}[flux gain along the wavefront]
\label{cor:v470-r13-wavefront-flux-gain}
Relative to the incoherent case, local wave-packet pile-up lowers the minimum
packet flux required at the ridge-front.  In one normalized handoff tick
\(\Delta t=1\), the admissible flux requirement changes from
\[
\dot N_{\min}=16
\]
to
\[
\dot N_{\min}\in\{7,6,5,4\}
\]
for \(\kappa=1/4,1/2,3/4,1\), respectively.  Thus the local overlap relay gives
a flux-efficiency gain
\[
G_{\mathrm{flux}}(\kappa)=\frac{16}{N_{\min}(\kappa)},
\]
namely approximately
\[
1,
\quad
2.29,
\quad
2.67,
\quad
3.20,
\quad
4.00.
\]
\end{corollary}

\begin{remark}[What this changes]
\label{rem:v470-r13-what-flowrate-change-brings}
The relay problem is now controlled by packet flux and local overlap, not by a
global single-wave crossing requirement.  The information channel only needs a
sufficient packet current along the admissible ridge-front/reentry chain.  The
front can therefore behave like a local queue or pile-up: packets that fail as
isolated crossings may still produce an admissible handoff when enough of them
arrive in phase-compatible windows.
\end{remark}

\begin{definition}[zero-clock assisted relay current]
\label{def:v470-r13-zero-clock-assisted-relay-current}
Let \(\mathcal C_\gamma(\tau_v)\) be the photon-frame zero-clock diagnostic
weight at the candidate handoff tick \(\tau_v\).  The zero-clock-assisted relay
condition is
\[
\mathcal C_\gamma(\tau_v)\,\mathcal O_v(\Delta t)
\geq
\frac{1}{256}.
\]
Equivalently, the effective threshold is
\[
\Theta_{v,\mathrm{eff}}(\tau_v)
=
\frac{1}{256\,\mathcal C_\gamma(\tau_v)}.
\]
This is a timing and phase-locking probe only; it is not an identification of
Riemann zeros with photons, and it is not an RH, T2, or all-readout closure
claim.
\end{definition}

\begin{proposition}[flow-rate consequence of a clock gain]
\label{prop:v470-r13-flowrate-clock-gain-consequence}
For \(a=1/64\), a zero-clock diagnostic gain \(G=\mathcal C_\gamma(\tau_v)>0\)
changes the finite relay inequality to
\[
N+\kappa N(N-1)
\geq
\frac{16}{G}.
\]
Thus a favorable photon-frame timing window lowers the minimum packet flux, while
an unfavorable timing window raises it.
\end{proposition}

\begin{proof}
Substitute \(\Theta_{v,\mathrm{eff}}=1/(256G)\) and \(a^2=1/4096\) into the
equal-amplitude overlap law.
\end{proof}

\subsection{Completion-interface audit}
\label{subsec:v470-r13-completion-interface-audit}

\begin{proposition}[C5 completion remains a boundary point]
\label{prop:v470-r13-c5-completion-interface-audit}
The C5 completion point at \(\infty\) remains a boundary object of the completed
bounded-collar topology and is not a finite handoff witness.
\end{proposition}

\begin{proof}
The r13 relay chain uses finite handoff points \(H_0,\ldots,H_m\).  The completed
C5 point is indexed by \(\infty\), not by a finite route segment.  Therefore it
can bound or guide the asymptotic collar but cannot itself serve as a finite
handoff witness.
\end{proof}

\begin{corollary}[Permitted r14 direction]
\label{cor:v470-r13-permitted-r14-direction}
The next admissible internal pass is
\[
\text{v4.7.0-r14: finite valley-overlap relay scan and zero-clock phase ledger.}
\]
It should instantiate a finite family of valley charts, compute symbolic or
bounded values of \(\mathcal O_k\), and create a finite zero-clock phase ledger
for handoff ticks.  It must not claim RH, T2, all-readout closure, or a physical
identity between photons and Riemann zeros.
\end{corollary}

\begin{remark}[r13 summary]
r13 answers the ridge intuition as follows.  The current obstruction is better
read as a local front/shadow/reentry relay problem than as a demand that an entire
wave cross a global mountain range.  Packet pile-up is allowed through the valley
overlap functional.  The possible relation to Riemann zeros is placed in a
finite photon-frame timing probe: a useful diagnostic channel, not a theorematic
identification or closure claim.
\end{remark}


% -----------------------------------------------------------------------------
\section{v4.7.0-r14 Source-Only Finite Valley-Front Chart Scan}
\label{sec:v470-r14-finite-valley-front-chart-scan}

\paragraph{Source-only status.}
This r14 pass is a source-only internal working increment.  It adds a finite
valley-front chart scan to the r13 ridge-front/reentry and valley-overlap relay
model.  No PDF build, DOI assignment, public-release claim, T2 claim, RH claim,
or all-readout closure claim is made in this pass.

\subsection{Normalized chart scan convention}
\label{subsec:v470-r14-normalized-chart-scan-convention}

The r13 relay threshold is kept fixed as
\[
\Theta_v=\beta_4^{r12}=\frac1{256},
\]
and the normalized equal-packet amplitude is kept as
\[
a=\frac1{64},\qquad a^2=\frac1{4096}.
\]
For a chart with packet count \(N\), coherence parameter \(\kappa\), and
photon-frame clock diagnostic gain \(G=\mathcal C_\gamma(\tau_v)\), the finite
handoff inequality becomes
\[
G\bigl(N+\kappa N(N-1)\bigr)\geq 16.
\]
This is the working finite chart scan rule for r14.

\begin{definition}[finite valley-front chart]
\label{def:v470-r14-finite-valley-front-chart}
A finite valley-front chart is a tuple
\[
\mathcal V_j=(N_j,\kappa_j,G_j,\sigma_j),
\]
where \(N_j\) is the local number of selector-admissible wave packets inside the
handoff window, \(\kappa_j\in[0,1]\) is their effective phase-coherence factor,
\(G_j=\mathcal C_\gamma(\tau_j)>0\) is the photon-frame clock diagnostic gain,
and \(\sigma_j\) is one of
\[
\mathsf{blocked},\qquad
\mathsf{relay\mbox{-}admissible},\qquad
\mathsf{zero\mbox{-}clock\mbox{-}assisted}.
\]
The status is assigned by the following finite rule:
\[
\sigma_j=
\begin{cases}
\mathsf{relay\mbox{-}admissible},
& N_j+\kappa_jN_j(N_j-1)\geq 16,\\
\mathsf{zero\mbox{-}clock\mbox{-}assisted},
& N_j+\kappa_jN_j(N_j-1)<16
\text{ and }G_j\bigl(N_j+\kappa_jN_j(N_j-1)\bigr)\geq16,\\
\mathsf{blocked},&
G_j\bigl(N_j+\kappa_jN_j(N_j-1)\bigr)<16.
\end{cases}
\]
\end{definition}

\subsection{Finite r14 chart ledger}
\label{subsec:v470-r14-finite-chart-ledger}

\begin{center}
\begin{longtable}{c|r|r|r|r|r|l}
\caption{r14 finite valley-front chart scan.  Scores are normalized by the
threshold \(1/256\); a score at least one is relay-sufficient in the indicated
channel.}\\
\toprule
Chart & \(N\) & \(\kappa\) & \(G\) & no-clock score & clock score & status\\
\midrule
\endfirsthead
\toprule
Chart & \(N\) & \(\kappa\) & \(G\) & no-clock score & clock score & status\\
\midrule
\endhead
V0 & 4 & 0 & 1 & 1/4 & 1/4 & blocked \\
V1 & 16 & 0 & 1 & 1 & 1 & relay-admissible \\
V2 & 7 & 1/4 & 1 & 35/32 & 35/32 & relay-admissible \\
V3 & 5 & 1/4 & 2 & 5/8 & 5/4 & zero-clock-assisted \\
V4 & 6 & 1/2 & 1 & 21/16 & 21/16 & relay-admissible \\
V5 & 4 & 1/2 & 5/2 & 5/8 & 25/16 & zero-clock-assisted \\
V6 & 3 & 3/4 & 3/2 & 15/32 & 45/64 & blocked \\
V7 & 4 & 1 & 1 & 1 & 1 & relay-admissible \\
V8 & 3 & 1 & 2 & 9/16 & 9/8 & zero-clock-assisted \\
V9 & 2 & 1/2 & 3 & 3/16 & 9/16 & blocked \\
\bottomrule
\end{longtable}
\end{center}

The finite scan contains 10 charts.  It classifies 4 charts as directly
relay-admissible, 3 charts as zero-clock-assisted, and 3 charts as still
blocked in the declared r14 scan family.

\begin{proposition}[finite flow-rate classification]
\label{prop:v470-r14-finite-flowrate-classification}
Inside the declared r14 finite chart family, local front flow rather than global
wave-over-mountain crossing determines handoff admissibility.  In particular,
there are charts for which the no-clock score is below one while the clock score
is at least one.  These charts are admissible only as photon-frame timing-assisted
relay charts.
\end{proposition}

\begin{proof}
The classification follows by direct substitution into
\[
G\bigl(N+\kappa N(N-1)\bigr)\geq16.
\]
For example, chart \(V3\) has \(N=5\), \(\kappa=1/4\), and \(G=2\).  Its
unassisted normalized score is \(10/16<1\), while its assisted score is
\(20/16>1\).  Hence it is not globally unassisted, but it is locally
zero-clock-assisted.  This is exactly the r13 ridge-front/reentry interpretation:
local packet flux at the handoff point may succeed even when an isolated or
unassisted crossing would fail.
\end{proof}

\subsection{Threshold grid for flux planning}
\label{subsec:v470-r14-threshold-grid-flux-planning}

The minimum required packet count for selected \((\kappa,G)\) values is:
\[
\begin{array}{c|rrrrrr}
\kappa\backslash G & 1 & 3/2 & 2 & 5/2 & 3 & 4\\
\hline
0 & 16 & 11 & 8 & 7 & 6 & 4 \\
1/4 & 7 & 6 & 5 & 4 & 4 & 3 \\
1/2 & 6 & 5 & 4 & 4 & 3 & 3 \\
3/4 & 5 & 4 & 4 & 3 & 3 & 3 \\
1 & 4 & 4 & 3 & 3 & 3 & 2 \\
\end{array}
\]
Thus the role of the photon-frame clock diagnostic is not to close the residual
by itself, but to lower the required local packet flux at suitable reentry ticks.

\begin{corollary}[r14 consequence for the mixed kernel]
\label{cor:v470-r14-mixed-kernel-front-consequence}
The route-stable mixed kernel from r12 may be treated, for the next source-only
pass, as a finite family of local ridge-front handoff charts rather than a single
global wave-crossing obstruction.  The persistent value \(1/256\) is retained as
the local threshold, but the admissibility question is moved to finite chart
routing, packet overlap, and timing-assisted relay status.
\end{corollary}

\begin{remark}[guard against overreading the scan]
The r14 chart scan is a finite diagnostic ledger.  It does not identify Riemann
zeros with photons, does not prove RH, does not close T2, and does not eliminate
the all-readout blocker.  It only shows that the remaining mixed-kernel problem
has a more precise local form: some valley fronts are blocked, some are directly
relay-admissible, and some require a favorable photon-frame timing window.
\end{remark}

\begin{corollary}[permitted r15 direction]
\label{cor:v470-r14-permitted-r15-direction}
The next admissible internal source-only pass is
\[
\text{v4.7.0-r15: zero-clock phase-window extraction and relay-chain stitching.}
\]
It should extract a finite set of admissible timing windows \(\tau_j\), test
whether adjacent valley-front charts can be stitched into a relay chain, and keep
all asymptotic, RH, T2, and all-readout claims guarded.
\end{corollary}


\section{v4.7.0-r15 Source-Only: Zero-Clock Phase-Window Extraction and Relay-Chain Stitching}
\label{sec:v470-r15-zero-clock-phase-window-relay-stitching}

\paragraph{Source-only status.}
This r15 pass is source-only.  It does not rebuild a PDF, does not assign a new
DOI, does not reopen the release metadata, and does not claim a T2, all-readout,
RH, or global physical closure theorem.  It continues the r13--r14 ridge-front
interpretation of the route-stable mixed kernel.

\subsection{Phase-window extraction rule}
\label{subsec:v470-r15-phase-window-extraction-rule}

For each valley-front chart \(V\), let
\[
S_V=N_V+\kappa_V N_V(N_V-1)
\]
be the local overlap flux before photon-frame timing gain.  The r14 relay
criterion was
\[
G_VS_V\geq16,
\]
where the factor 16 is the normalized form of the residual ridge threshold
\(\Theta_v=1/256\) at packet amplitude \(a=1/64\).  Hence the required clock gain
is
\[
G_V^{\mathrm{req}}=\frac{16}{S_V}.
\]
A zero-clock phase window is extracted only when
\[
1<G_V^{\mathrm{req}}\leq G_V^{\mathrm{avail}}.
\]
Thus direct charts need no clock window, assisted charts need a nonempty timing
window, and blocked charts remain blocked in the finite r15 family.

The photon-clock probe is still written as a diagnostic timing function, not as
an identity between Riemann zeros and photons:
\[
\mathcal C_\gamma(\tau)
=
1+\eta
\frac{\left|\sum_{m=1}^{M}w_m e^{i\gamma_m\tau}\right|^2}
{\left(\sum_{m=1}^{M}|w_m|\right)^2}.
\]
For a local tick \(\tau_V\), the window is the threshold set
\[
I_V=\{\tau:\mathcal C_\gamma(\tau)\geq G_V^{\mathrm{req}}\}.
\]
In the local normalized tick approximation
\[
G_V(\tau)=1+(G_V^{\mathrm{avail}}-1)
\cos^2\!\left(\frac{\tau-\tau_V}{\sigma_V}\right),
\]
the half-width is
\[
\delta_V
=
\sigma_V\arccos\sqrt{\frac{G_V^{\mathrm{req}}-1}{G_V^{\mathrm{avail}}-1}}.
\]
This formula is used only to rank finite timing windows.  It is not a spectral
identification theorem.

\subsection{Extracted finite timing windows}
\label{subsec:v470-r15-extracted-finite-timing-windows}

The r14 zero-clock-assisted charts admit the following r15 phase windows:
\[
\begin{array}{c|c|c|c|c}
\text{chart} & S_V & G_V^{\mathrm{req}} & G_V^{\mathrm{avail}} & \delta_V/\sigma_V\\
\hline
V3 & 10 & 8/5 & 2 & 0.684719\\
V5 & 10 & 8/5 & 5/2 & 0.886077\\
V8 & 9 & 16/9 & 2 & 0.490883\\
\end{array}
\]
Thus the narrowest extracted window in the finite family is the coherent sparse
reentry pulse \(V8\), while the ridge-seam chart \(V5\) has the widest clock
margin.

\begin{proposition}[finite zero-clock phase-window extraction]
\label{prop:v470-r15-finite-zero-clock-window-extraction}
Within the declared r15 finite chart family, exactly three charts have nonempty
zero-clock phase windows: \(V3,V5,V8\).  These windows are finite, local timing
windows for relay admissibility.  They do not close the mixed kernel globally.
\end{proposition}

\begin{proof}
For \(V3\) and \(V5\), \(S_V=10\), hence
\(G_V^{\mathrm{req}}=16/10=8/5\).  The available gains are \(2\) and \(5/2\),
respectively, so the threshold-set condition is satisfied.  For \(V8\),
\(S_V=9\), hence \(G_V^{\mathrm{req}}=16/9<2\), again giving a nonempty window.
All direct charts have \(G_V^{\mathrm{req}}\leq1\), so no clock window is needed.
The blocked charts have insufficient available gain in the finite ledger.
\end{proof}

\subsection{Relay-chain stitching}
\label{subsec:v470-r15-relay-chain-stitching}

The extracted windows allow two finite stitched relay chains:
\[
\mathcal R_A:
V1\longrightarrow V2\longrightarrow V3\longrightarrow V4\longrightarrow V5,
\]
where \(V3\) and \(V5\) require the phase windows \(I_{V3}\) and \(I_{V5}\), and
\[
\mathcal R_B:
V7\longrightarrow V8,
\]
where \(V8\) requires \(I_{V8}\).
The edges
\[
V5\longrightarrow V6,
\qquad
V8\longrightarrow V9
\]
remain blocked in the r15 finite chart family.

\begin{center}
\begin{tabular}{c|c|c|c}
\toprule
edge & type & status & required window\\
\midrule
$V1\to V2$  & direct-to-direct  & stitched  & none\\
$V2\to V3$  & direct-to-assisted  & stitched  & $I_{V3}$\\
$V3\to V4$  & assisted-to-direct  & stitched  & $I_{V3}$\\
$V4\to V5$  & direct-to-assisted  & stitched  & $I_{V5}$\\
$V5\to V6$  & assisted-to-blocked  & blocked  & none\\
$V7\to V8$  & direct-to-assisted  & stitched  & $I_{V8}$\\
$V8\to V9$  & assisted-to-blocked  & blocked  & none\\
\bottomrule
\end{tabular}
\end{center}

\begin{corollary}[finite chain status]
\label{cor:v470-r15-finite-chain-status}
The r15 relay stitching produces a longest certified chain of length five and a
secondary certified chain of length two.  It does not produce a single global
all-terrain relay chain.
\end{corollary}

\subsection{Remaining gap ledger}
\label{subsec:v470-r15-remaining-gap-ledger}

The blocked charts have the following required-to-available gain gaps:
\[
\begin{array}{c|c|c|c}
\text{chart} & G^{\mathrm{req}} & G^{\mathrm{avail}} &
G^{\mathrm{req}}/G^{\mathrm{avail}}\\
\hline
V0 & 4 & 1 & 4\\
V6 & 32/15 & 3/2 & 64/45\\
V9 & 16/3 & 3 & 16/9\\
\end{array}
\]
Thus the closest remaining blocked valley-front is \(V6\), not \(V0\) or
\(V9\).  This suggests that the next useful pass should focus on a local
shadow-basin lift for \(V6\), rather than on a global mixed-kernel closure claim.

\begin{remark}[r15 guard]
The r15 result is a finite relay-chain diagnostic.  It does not identify the
Riemann zeros with photon ticks, does not prove RH, does not prove T2, and does
not remove the all-readout blocker.  Its only promoted content is local: certain
valley-front charts have explicit photon-clock phase windows, and these windows
stitch finite relay chains across parts of the ridge-front atlas.
\end{remark}

\begin{corollary}[permitted r16 direction]
\label{cor:v470-r15-permitted-r16-direction}
The next admissible source-only pass is
\[
\text{v4.7.0-r16: shadow-basin lift test for }V6
\text{ and relay-chain extension guard.}
\]
It should test whether the closest blocked chart \(V6\) can be lifted by a
local overlap or clock-window refinement, while keeping \(V0\), \(V9\), T2, RH,
and all-readout closure guarded.
\end{corollary}

% -----------------------------------------------------------------------------
\section{v4.7.0-r16 Source-Only: Shadow-Basin Lift Test for \texorpdfstring{$V6$}{V6}}
\label{sec:v470-r16-shadow-basin-lift-v6}

This source-only pass tests the closest blocked r15 valley-front chart,
\(V6\), without rebuilding a PDF and without changing the baseline DOI.  The
r15 obstruction ledger identified
\[
V6:
\qquad
S_6^{(0)}=\frac{15}{2},
\qquad
G_6^{(0)}=\frac{3}{2},
\qquad
G_6^{\mathrm{req}}=\frac{32}{15}.
\]
Equivalently,
\[
G_6^{(0)}S_6^{(0)}
=
\frac{45}{4}
<16.
\]
Thus \(V6\) was not globally obstructed by a whole-landscape failure; it was a
local shadow-basin deficit at the ridge-front handoff seam.

\begin{definition}[shadow-basin lift datum for \(V6\)]
\label{def:v470-r16-shadow-basin-lift-datum-v6}
A \emph{shadow-basin lift datum} for \(V6\) is a finite local replacement
\[
(N_6,\kappa_6,G_6)
\longmapsto
(N_6',\kappa_6',G_6')
\]
inside the same declared valley-front chart, with no all-readout extension and
no new global selector.  It is admissible only if
\[
G_6'\bigl(N_6'+\kappa_6'N_6'(N_6'-1)\bigr)\geq 16.
\]
\end{definition}

\subsection{Finite lift computation}
\label{subsec:v470-r16-finite-lift-computation}

The r15 value was
\[
N_6=3,
\qquad
\kappa_6=\frac34,
\qquad
G_6=\frac32,
\]
so
\[
S_6
=
3+\frac34\cdot 3\cdot 2
=
\frac{15}{2}.
\]
The shadow-basin test asks whether a local pile-up can increase the packet flux
by one admissible packet in the same photon-clock window.  The tested lifted
state is
\[
N_6'=4,
\qquad
\kappa_6'=\frac34,
\qquad
G_6'=\frac32.
\]
Then
\[
S_6'
=
4+\frac34\cdot4\cdot3
=13,
\]
and hence
\[
G_6'S_6'
=
\frac32\cdot13
=
\frac{39}{2}
>
16.
\]
The normalized relay score is therefore
\[
\frac{G_6'S_6'}{16}
=
\frac{39}{32}
>1.
\]

\begin{lemma}[one-packet shadow-basin lift]
\label{lem:v470-r16-one-packet-shadow-basin-lift-v6}
In the finite r16 chart family, \(V6\) is lifted from blocked to
shadow-basin-assisted relay-admissible by adding one selector-admissible packet
to the local basin while preserving \(\kappa_6=3/4\) and \(G_6=3/2\).
\end{lemma}

\begin{proof}
The computation above gives
\[
G_6S_6=45/4<16
\]
for the r15 state, while the lifted state satisfies
\[
G_6'S_6'=39/2>16.
\]
Therefore the r15 chart was subcritical and the r16 lifted chart is supercritical
with margin
\[
\frac{39}{2}-16=\frac72.
\]
No all-readout extension, global terrain crossing, or T2 bridge is used.
\end{proof}

\subsection{Flux-rate consequence}
\label{subsec:v470-r16-flux-rate-consequence}

The lift is small in packet count but large in effective overlap.  The packet
flux changes by
\[
\frac{N_6'}{N_6}=\frac43,
\]
whereas the local overlap strength changes by
\[
\frac{S_6'}{S_6}
=
\frac{13}{15/2}
=
\frac{26}{15}.
\]
Thus a \(33.3\%\) increase in selector-admissible packet count produces a
\(73.3\%\) increase in the overlap strength, because the additional packet also
participates in the coherent pair terms.  This is the precise sense in which the
r13/r14 valley-front model converts a ridge barrier into a local flow-rate
problem.

\begin{corollary}[local relay lift without global wave crossing]
\label{cor:v470-r16-local-relay-without-global-wave-crossing}
The r16 lift of \(V6\) does not require a whole wave to cross the whole ridge.
It only requires that the local packet flux in the shadow basin rises from
\(3\) to \(4\) packets in the same declared handoff window while maintaining the
same partial coherence and photon-clock gain.
\end{corollary}

\subsection{Relay-chain extension guard}
\label{subsec:v470-r16-relay-chain-extension-guard}

The r15 chain
\[
\mathcal R_A:
V1\to V2\to V3\to V4\to V5
\]
extends in r16 to
\[
\mathcal R_A^{+}:
V1\to V2\to V3\to V4\to V5\to V6.
\]
The extension is guarded: it certifies the edge \(V5\to V6\) only after the
local shadow-basin lift of \(V6\).  It does not certify \(V6\to V7\), does not
merge \(\mathcal R_A^{+}\) with the secondary chain \(V7\to V8\), and does not
lift \(V0\) or \(V9\).

\begin{center}
\begin{tabular}{c|c|c|c}
\toprule
edge & r15 status & r16 status & guard\\
\midrule
$V1\to V2$  & stitched  & stitched  & unchanged\\
$V2\to V3$  & stitched with window  & stitched with window  & unchanged\\
$V3\to V4$  & stitched  & stitched  & unchanged\\
$V4\to V5$  & stitched with window  & stitched with window  & unchanged\\
$V5\to V6$  & blocked  & stitched after shadow-basin lift  & local lift only\\
$V6\to V7$  & not certified  & not certified  & atlas gap\\
$V7\to V8$  & stitched with window  & stitched with window  & separate chain\\
$V8\to V9$  & blocked  & blocked  & insufficient lift\\
\bottomrule
\end{tabular}
\end{center}

\begin{corollary}[r16 finite chain status]
\label{cor:v470-r16-finite-chain-status}
The longest certified relay chain in the finite r16 atlas has length six:
\[
V1\to V2\to V3\to V4\to V5\to V6.
\]
This is a finite relay extension, not a global all-terrain relay chain.
\end{corollary}

\subsection{C4/C5/C6 status after r16}
\label{subsec:v470-r16-status-after-lift}

The r16 lift improves a local representative of the route-stable mixed-kernel
picture, but it does not close \(C4\).  The persistent mixed-kernel guard from
r12 remains in force:
\[
\beta_4^{r12}=\frac{1}{256}>0.
\]
The \(C5\) bounded-collar completion guard also remains unchanged:
\[
\forall n<\infty:
\Phi_5^{(n)}>0,
\qquad
\widehat\Phi_5(\infty)=0.
\]
Finally,
\[
C6:\mathsf B_{\mathrm{all/T2}}
\]
remains blocked.  The r16 result is therefore only a local relay-flow upgrade.

\begin{remark}[r16 non-claim guard]
The r16 shadow-basin lift does not prove RH, does not identify Riemann zeros
with photon ticks, does not prove T2, does not close the all-readout boundary,
and does not remove the route-stable mixed-kernel obstruction.  It proves only
that the closest r15 blocked chart \(V6\) can be lifted by a finite local
packet-flux increase within the declared valley-front relay grammar.
\end{remark}

\begin{corollary}[permitted r17 direction]
\label{cor:v470-r16-permitted-r17-direction}
The next admissible source-only pass is
\[
\text{v4.7.0-r17: atlas-gap bridge test for }V6\to V7
\text{ and guarded }V9\text{ micro-lift scan.}
\]
It should test whether the newly lifted \(V6\) can be stitched to the existing
secondary chain through a declared atlas-gap bridge, while keeping \(V0\),
\(V9\), RH, T2, and all-readout closure guarded.
\end{corollary}



\section{v4.7.0-r17 Source-Only Atlas-Gap Bridge Test and Guarded V9 Micro-Lift}
\label{sec:v470-r17-atlas-gap-bridge-v9-microlift}

This source-only layer continues the valley-front relay program after the r16
shadow-basin lift of \(V6\).  The aim is deliberately finite: test the atlas gap
\(V6\to V7\), then scan whether the still blocked \(V9\) chart admits a minimal
micro-lift.  No PDF is produced in this intermediate version.

\subsection{r17 relay score convention}
\label{subsec:v470-r17-relay-score-convention}

For a local valley-front chart \(V\) the overlap score is retained as
\[
S_V(N,\kappa)=N+\kappa N(N-1),
\]
and the clock-weighted relay score is
\[
H_V=G_VS_V.
\]
The finite local handoff condition is
\[
H_V\geq 16.
\]
This is only a local threshold condition.  It is not a global terrain theorem
and not an all-readout closure condition.

\subsection{Atlas-gap bridge test for \texorpdfstring{\(V6\to V7\)}{V6 to V7}}
\label{subsec:v470-r17-v6-v7-bridge-test}

After r16 the lifted \(V6\) state is
\[
N_6'=4,
\qquad
\kappa_6=\frac34,
\qquad
G_6=\frac32.
\]
Hence
\[
S_6'=4+\frac34\cdot4\cdot3=13,
\qquad
H_6'=\frac32\cdot13=\frac{39}{2}.
\]
The local surplus at \(V6\) is therefore
\[
\mu_6=H_6'-16=\frac{39}{2}-16=\frac72.
\]
The target chart \(V7\) remains the near-coherent equality handoff:
\[
N_7=4,
\qquad
\kappa_7=1,
\qquad
G_7=1,
\]
so that
\[
S_7=4+1\cdot4\cdot3=16,
\qquad
H_7=16,
\qquad
\mu_7=0.
\]
Thus \(V7\) has no spare local margin, but it does satisfy the relay threshold
exactly.  The only bridge question is whether the atlas-gap loss from the
lifted \(V6\) side can be paid from the \(V6\) surplus.

For the finite r17 scan we declare the conservative atlas-gap toll
\[
\alpha_{67}=\frac12.
\]
The transported bridge margin is then
\[
\mu_{67}^{\mathrm{bridge}}
=\mu_6-\alpha_{67}
=\frac72-\frac12
=3>0.
\]
Consequently the edge \(V6\to V7\) is bridge-admissible in the finite r17 atlas.
It is, however, still a critical bridge because the target chart \(V7\) itself
sits exactly at threshold.

\begin{proposition}[r17 finite atlas-gap bridge]
\label{prop:v470-r17-v6-v7-bridge}
Within the declared finite r17 valley-front atlas, the edge \(V6\to V7\) is
admissible after the r16 shadow-basin lift of \(V6\).  The certification uses
\(V6\)-side surplus \(\mu_6=7/2\) to pay the declared atlas-gap toll
\(\alpha_{67}=1/2\), leaving bridge margin \(3>0\).
\end{proposition}

\begin{remark}[critical target guard]
\label{rem:v470-r17-critical-target-guard}
The r17 bridge is not a robust global bridge theorem.  The target chart \(V7\)
has \(\mu_7=0\), so the certification is finite, source-supported, and
atlas-local.  Any later attempt to promote it must supply a uniform interface
margin or an independently declared bridge operator.
\end{remark}

\subsection{Guarded micro-lift scan for \texorpdfstring{\(V9\)}{V9}}
\label{subsec:v470-r17-v9-microlift-scan}

The r14/r15 state of \(V9\) was subcritical:
\[
N_9=2,
\qquad
\kappa_9=\frac12,
\qquad
G_9=3.
\]
Therefore
\[
S_9=2+\frac12\cdot2\cdot1=3,
\qquad
H_9=3\cdot3=9<16.
\]
The guarded micro-lift scan tests the smallest packet-flux increase
\[
N_9:2\to3
\]
while leaving \(\kappa_9\) and \(G_9\) unchanged.  Then
\[
S_9'=3+\frac12\cdot3\cdot2=6,
\qquad
H_9'=3\cdot6=18>16.
\]
The resulting margin is
\[
\mu_9'=18-16=2>0.
\]
Thus \(V9\) is not inert: one additional local packet in the same guarded
clock window suffices to make the chart relay-admissible.

\begin{lemma}[r17 guarded \(V9\) micro-lift]
\label{lem:v470-r17-v9-microlift}
In the finite r17 scan, the transition
\[
(N_9,\kappa_9,G_9)=\left(2,\frac12,3\right)
\longmapsto
\left(3,\frac12,3\right)
\]
lifts \(V9\) from blocked to guarded micro-lift relay-admissible, since
\(H_9'=18>16\).
\end{lemma}

\begin{remark}[micro-lift non-claim]
\label{rem:v470-r17-v9-nonclaim}
The \(V9\) micro-lift does not close \(V0\), does not close the route-stable
mixed-kernel obstruction, and does not prove that arbitrary exterior collars
can be lifted by adding packets.  It only proves that the concrete finite r17
\(V9\) chart is liftable under the declared \(G_9=3\) clock window.
\end{remark}

\subsection{Relay-chain stitching after r17}
\label{subsec:v470-r17-relay-chain-stitching}

Combining the r16 lift of \(V6\), the r17 atlas-gap bridge \(V6\to V7\), the
previously certified \(V7\to V8\) edge, and the r17 guarded micro-lift of
\(V9\), the finite chart chain extends to
\[
\mathcal R_A^{++}:
V1\to V2\to V3\to V4\to V5\to V6\to V7\to V8\to V9.
\]
This is a finite stitched relay chain over the nonzero-index valley-front
sample.  It does not include \(V0\), and it is not a global atlas-cover theorem.

\begin{center}
\begin{tabular}{c|c|c|c}
\toprule
edge & r16 status & r17 status & guard\\
\midrule
$V5\to V6$  & stitched after shadow-basin lift  & stitched  & unchanged r16 lift\\
$V6\to V7$  & atlas gap  & bridge-admissible  & source-supported critical bridge\\
$V7\to V8$  & stitched with window  & stitched with window  & inherited\\
$V8\to V9$  & blocked  & stitched after guarded micro-lift  & local \(V9\) lift only\\
\bottomrule
\end{tabular}
\end{center}

\begin{corollary}[r17 finite chain status]
\label{cor:v470-r17-finite-chain-status}
The longest certified finite relay chain in the r17 atlas is
\[
V1\to V2\to V3\to V4\to V5\to V6\to V7\to V8\to V9.
\]
The result is a finite relay-chain extension, not a global all-readout closure.
\end{corollary}

\subsection{Guarded residual status after r17}
\label{subsec:v470-r17-guarded-residual-status}

The r17 finite atlas scan improves the relay geometry but does not change the
global claim guards.  The following remain active:
\[
V0:\mathsf{blocked},
\qquad
\beta_4^{r12}=\frac1{256}>0,
\qquad
C6:\mathsf B_{\mathrm{all/T2}}.
\]
The \(C5\) completion guard remains unchanged:
\[
\forall n<\infty:\ \Phi_5^{(n)}>0,
\qquad
\widehat\Phi_5(\infty)=0.
\]
Thus r17 should be read as a finite relay upgrade across the local valley-front
sample.  It is not a proof of RH, not a photon-zero identity, not a T2 proof,
and not a global physical closure theorem.

\begin{corollary}[permitted r18 direction]
\label{cor:v470-r17-permitted-r18-direction}
The next admissible source-only pass is
\[
\text{v4.7.0-r18: }V0\text{ boundary-front obstruction inventory and finite-chain guard audit.}
\]
It should determine whether \(V0\) is a true boundary-front obstruction, a
missing-source chart, or an inadmissible zero-packet baseline, while auditing
that the r17 finite chain is not misread as a global closure.
\end{corollary}



% -----------------------------------------------------------------------------
% v4.7.0-r18 source-only insertion
% -----------------------------------------------------------------------------
\section{v4.7.0-r18 Source-Only Pass: V0 Boundary-Front Obstruction Inventory and Finite-Chain Guard Audit}
\label{sec:v470-r18-v0-boundary-front-guard-audit}

This source-only pass records the next conservative step after the r17
atlas-gap bridge and guarded micro-lift scan.  No PDF rebuild is requested or
used.  The purpose is not to extend the local relay chain past its declared
scope, but to inventory the remaining boundary chart \(V0\) and to audit that
\[
V1\to V2\to V3\to V4\to V5\to V6\to V7\to V8\to V9
\]
is read only as a finite local valley-front chain.

\subsection{V0 boundary-front inventory}
\label{subsec:v470-r18-v0-boundary-front-inventory}

Recall the r15/r17 boundary datum
\[
V0:
\qquad
S_0=4,
\qquad
G_0=1,
\qquad
H_0=G_0S_0=4<16.
\]
Equivalently,
\[
\Delta_0=16-H_0=12,
\qquad
G_0^{\rm req}=\frac{16}{S_0}=4.
\]
Thus \(V0\) is not close to the relay threshold in the same sense as the former
\(V6\) blocker.  It is not a small shadow-basin miss; it is a boundary-front
source problem.

The conservative inventory is:
\[
\begin{array}{c|c|c|c|c|c|c}
\text{case} & N & \kappa & G & S & H=GS & \text{status}\\
\hline
\text{baseline} &4&0&1&4&4&\text{blocked}\\
\text{gain-only} &4&0&4&4&16&\text{formal only}\\
\text{incoherent flux-only} &16&0&1&16&16&\text{formal only}\\
\text{overlap lift} &5&3/4&1&20&20&\text{source-guarded}
\end{array}
\]
The gain-only and flux-only rows show that \(V0\) is numerically liftable if one
allows an undeclared boundary source, an undeclared boundary clock, or an
external injection rule.  These are not currently part of the declared local
relay grammar.  Therefore r18 classifies \(V0\) as
\[
\boxed{
V0:
\mathsf{boundary\mbox{-}front\ source\mbox{-}deficient\ obstruction}.
}
\]
This is stronger than saying that \(V0\) is simply ``blocked'', but weaker than
claiming that it is an absolute topological obstruction.  It means that the
known finite chain has no certified upstream source chart at the boundary.

\begin{lemma}[V0 boundary-front decision]
\label{lem:v470-r18-v0-boundary-front-decision}
Within the current r18 source grammar, \(V0\) is not absorbed by the
\(V1\to\cdots\to V9\) finite relay chain.  It remains a source-deficient
boundary-front obstruction.  Any lift of \(V0\) must declare at least one of the
following additional structures:
\[
\text{boundary source chart},
\qquad
\text{boundary zero-clock phase window},
\qquad
\text{external packet injection rule}.
\]
\end{lemma}

\begin{proof}
The current datum gives \(H_0=4<16\).  A formal threshold equality can be forced
by taking \(G=4\), or by taking sixteen incoherent packets, or by introducing a
coherent overlap lift.  None of these is certified by the existing relay chain,
because the chain begins at \(V1\) and supplies no admissible upstream source for
\(V0\).  Hence the obstruction is not a numerical impossibility; it is a
missing-source boundary-front obstruction in the declared grammar.
\end{proof}

\subsection{Finite-chain guard audit}
\label{subsec:v470-r18-finite-chain-guard-audit}

The r17 chain remains a finite local atlas result:
\[
\boxed{
\mathcal R_{1:9}^{\rm fin}
=
V1\to V2\to V3\to V4\to V5\to V6\to V7\to V8\to V9.
}
\]
The r18 guard audit records the active non-closures:
\[
\begin{array}{c|c|c}
\text{guard} & \text{r18 status} & \text{interpretation}\\
\hline
V0 & \text{active} & \text{boundary-front source deficient}\\
C4 & \beta_4=1/256>0 & \text{route-stable mixed-kernel remains}\\
C5 & \forall n<\infty:\Phi_5^{(n)}>0 & \text{asymptotic collar only}\\
C6 & \mathsf B_{\rm all/T2} & \text{all-readout/T2 blocked}\\
\mathrm{RH} & \text{not claimed} & \text{photon-clock handoff remains a probe}\\
\mathrm{global} & \text{not claimed} & \text{no all-readout closure}
\end{array}
\]
Consequently, the finite local chain is useful because it converts many formerly
abstract valley-front obstructions into relay-admissible handoff steps.  It does
not remove the boundary source problem at \(V0\), and it does not discharge the
route-stable \(C4\) mixed-kernel residue.

\begin{proposition}[r18 finite-chain guard]
\label{prop:v470-r18-finite-chain-guard}
The r17 relay chain \(\mathcal R_{1:9}^{\rm fin}\) is admissible as a finite
local relay chain, but it is not admissible as a global closure witness.  The
active guards are \(V0\), \(C4\), finite \(C5\) closure, \(C6/T2\), RH, and
global all-readout closure.
\end{proposition}

\begin{proof}
The chain starts at \(V1\), not at a certified boundary source.  It also leaves
\(\beta_4=1/256\) unchanged and does not alter the finite-positive \(C5\) guard
or the all-readout/T2 blocker.  Therefore the finite chain may be used as a
local relay atlas but not as a global closure theorem.
\end{proof}

\subsection{Permitted r19 direction}
\label{subsec:v470-r18-permitted-r19-direction}

The next source-only pass may either keep \(V0\) frozen and move back to the
\(C4\) mixed-kernel seam, or test whether a declared boundary source chart can
be added without creating new debt.  The safest next pass is therefore
\[
\boxed{
\text{v4.7.0-r19: boundary-source chart admissibility test and C4 guard-preserving audit.}
}
\]
No RH, T2, or global all-readout closure should be inferred from r18.



% -----------------------------------------------------------------------------
% v4.7.0-r19 source-only insertion
% -----------------------------------------------------------------------------
\section{v4.7.0-r19 Source-Only Pass: Boundary-Source Chart Admissibility Test and C4 Guard-Preserving Audit}
\label{sec:v470-r19-boundary-source-chart-c4-guard}

This source-only pass tests the admissibility of a possible boundary-source
chart for the remaining boundary-front obstruction \(V0\).  No PDF rebuild is
requested or used.  The goal is deliberately narrow: determine whether \(V0\)
can be lifted by a declared local boundary source without spending the
route-stable \(C4\) mixed-kernel guard and without opening the all-readout/T2
selector.

The r18 decision was
\[
V0:\mathsf{boundary\mbox{-}front\ source\mbox{-}deficient\ obstruction},
\qquad
S_0=4,\quad G_0=1,\quad H_0=4<16.
\]
Thus the missing amount is
\[
\Delta_0=12,
\qquad
G_0^{\rm req}=4.
\]
The present pass asks whether this deficit can be supplied by a legitimate
boundary chart rather than by an undeclared global injection.

\subsection{Boundary-source admissibility criterion}
\label{subsec:v470-r19-boundary-source-admissibility-criterion}

For a boundary-source candidate \(\mathcal B_0\), write the local overlap score
as
\[
S(\mathcal B_0)
=
N+\kappa N(N-1),
\qquad
H(\mathcal B_0)=G\,S(\mathcal B_0).
\]
Here \(N\) is the packet count in the boundary window, \(\kappa\) is the local
coherence, and \(G\) is any declared photon-clock or boundary timing gain.
The threshold remains
\[
H(\mathcal B_0)\ge 16.
\]
This numerical threshold is necessary but not sufficient.  The chart is
admissible only if the following guard conditions also hold:
\[
\begin{array}{ll}
\text{(A1)} & H(\mathcal B_0)\ge 16,\\
\text{(A2)} & D_{\partial}(\mathcal B_0)=0\quad\text{(no undeclared boundary debt)},\\
\text{(A3)} & P_{\partial0}S_{4,\mu}^{\rm mixker}=0\quad\text{(no spending of the }C4\text{ mixed kernel)},\\
\text{(A4)} & \mathcal B_0\text{ does not introduce an all-readout selector or a }T2\text{ lift.}
\end{array}
\]
Condition (A3) is the guard-preserving clause.  The boundary chart may supply a
local upstream source for \(V0\), but it may not absorb or relabel the persistent
mixed-kernel residue
\[
\beta_4=\frac1{256}.
\]

\begin{definition}[r19 boundary-source admissibility]
\label{def:v470-r19-boundary-source-admissibility}
A boundary-source chart \(\mathcal B_0\) is r19-admissible if and only if it
satisfies (A1)--(A4).  A candidate satisfying only (A1) is called numerically
lifting but grammar-inadmissible.
\end{definition}

\subsection{Finite boundary-source candidate scan}
\label{subsec:v470-r19-boundary-source-candidate-scan}

The finite candidate scan is:
\[
\begin{array}{c|c|c|c|c|c|c}
\text{candidate} & N & \kappa & G & S & H=GS & \text{r19 decision}\\
\hline
\text{baseline} &4&0&1&4&4&\text{blocked}\\
\text{clock-only} &4&0&4&4&16&\text{numerical only; undeclared clock debt}\\
\text{external injection} &16&0&1&16&16&\text{numerical only; injection debt}\\
\text{coherent boundary reservoir} &5&3/4&1&20&20&\text{conditionally admissible}\\
\text{clock-assisted overlap} &4&3/4&16/13&13&16&\text{conditionally admissible}
\end{array}
\]
The baseline fails the threshold.  The clock-only and external-injection rows
show that \(V0\) can be forced to the threshold numerically, but only by adding
an undeclared source of gain or packets.  Those rows remain inadmissible in the
current grammar.

The first genuinely useful row is the coherent boundary reservoir.  For
\(N=5\), the threshold condition is
\[
5+20\kappa\ge 16,
\qquad
\kappa\ge\frac{11}{20}.
\]
Thus the displayed value \(\kappa=3/4\) has positive margin
\[
S-16=4.
\]
This proves only a conditional admissibility statement: if a local boundary
reservoir \(\mathcal B_0^{\rm res}\) with
\[
N=5,
\qquad
\kappa\ge\frac{11}{20},
\qquad
G=1
\]
is declared and if it satisfies the no-new-debt and \(C4\)-orthogonality guards,
then \(V0\) is relay-admissible as a boundary-source chart.

\begin{lemma}[minimal coherent reservoir threshold]
\label{lem:v470-r19-minimal-coherent-reservoir-threshold}
For a boundary-reservoir candidate with \(G=1\) and \(N=5\), the local relay
threshold is reached exactly when
\[
\kappa\ge\frac{11}{20}.
\]
In particular, \(\kappa=3/4\) gives \(S=20>16\).
\end{lemma}

\begin{proof}
Substituting \(N=5\) into \(S=N+\kappa N(N-1)\) gives
\(S=5+20\kappa\).  The threshold \(S\ge16\) is therefore equivalent to
\(20\kappa\ge11\), hence \(\kappa\ge 11/20\).  At \(\kappa=3/4\), one obtains
\(S=5+15=20\).
\end{proof}

\subsection{C4 guard-preserving audit}
\label{subsec:v470-r19-c4-guard-preserving-audit}

The boundary-source scan does not spend the \(C4\) mixed-kernel residue.  The
r12--r18 persistent residue remains
\[
S_{4,\mu}^{\rm mixker}:\mathsf{route\mbox{-}stable\ persistent},
\qquad
\beta_4=\frac1{256}.
\]
The guard audit is:
\[
\begin{array}{c|c|c|c}
\text{candidate} & \text{threshold} & C4\text{ guard} & \text{decision}\\
\hline
\text{clock-only} & \text{met} & \text{not certified} & \text{reject unless clock source declared}\\
\text{external injection} & \text{met} & \text{not certified} & \text{reject as external debt}\\
\text{coherent boundary reservoir} & \text{met} & P_{\partial0}S_{4,\mu}^{\rm mixker}=0 & \text{conditional pass}\\
\text{clock-assisted overlap} & \text{met} & \text{requires clock-source audit} & \text{conditional only}
\end{array}
\]
Accordingly, r19 makes no unconditional reduction of \(\beta_4\):
\[
\boxed{
\beta_4^{r19}=\frac1{256}.
}
\]
The only admissible direction is to treat \(\mathcal B_0^{\rm res}\) as a new
local boundary-source object whose support is orthogonal to the mixed-kernel
seam.  If that orthogonality is not declared, the candidate cannot be used.

\begin{proposition}[r19 guard-preserving outcome]
\label{prop:v470-r19-guard-preserving-outcome}
The r19 scan finds no unconditional current-grammar lift of \(V0\).  It does,
however, isolate a conditional boundary-source normal form:
\[
\mathcal B_0^{\rm res}:
\qquad
N=5,
\qquad
\kappa\ge\frac{11}{20},
\qquad
G=1,
\qquad
P_{\partial0}S_{4,\mu}^{\rm mixker}=0.
\]
Under these conditions \(V0\) is boundary-source relay-admissible and the
\(C4\) guard remains unchanged at \(\beta_4=1/256\).  Without a declared
\(\mathcal B_0^{\rm res}\), \(V0\) remains guarded.
\end{proposition}

\begin{proof}
The numerical threshold follows from Lemma~\ref{lem:v470-r19-minimal-coherent-reservoir-threshold}.
The no-new-debt requirement is exactly the statement that the reservoir is a
local boundary-source chart rather than an external injection.  The displayed
orthogonality condition prevents the new boundary source from absorbing the
route-stable \(C4\) mixed-kernel residue.  Thus the result is conditional and
guard-preserving, not an unconditional closure.
\end{proof}

\subsection{Residual status after r19}
\label{subsec:v470-r19-residual-status}

After r19 the finite relay-chain statement is unchanged:
\[
V1\to V2\to V3\to V4\to V5\to V6\to V7\to V8\to V9
\]
remains the certified finite local valley-front chain.  The boundary chart
\(V0\) has a conditional lift criterion, but it is not yet part of the certified
chain.

The guarded status is therefore:
\[
\boxed{
\begin{array}{c|c}
\text{object} & \text{r19 status}\\
\hline
V0 & \text{conditional boundary-reservoir criterion only}\\
C4 & \beta_4=1/256\text{ route-stable mixed-kernel guard preserved}\\
C5 & \text{asymptotic bounded collar; no finite closure}\\
C6/T2 & \text{all-readout blocker unchanged}\\
RH & \text{no claim; photon-clock remains handoff probe}\\
\text{global closure} & \text{not claimed}
\end{array}
}
\]

\subsection{Permitted r20 direction}
\label{subsec:v470-r19-permitted-r20-direction}

The next source-only pass may test whether the conditional object
\(\mathcal B_0^{\rm res}\) can be promoted from an admissibility criterion to a
declared local boundary-source chart.  The safe formulation is
\[
\boxed{
\text{v4.7.0-r20: boundary-reservoir declaration test and no-new-debt proof obligation.}
}
\]
This pass should either certify the no-new-debt and \(C4\)-orthogonality clauses
for \(\mathcal B_0^{\rm res}\), or keep \(V0\) frozen as a guarded boundary
source deficit.  No RH, T2, finite \(C5\) closure, or global all-readout closure
is licensed by r19.



\clearpage
\section{v4.7.0-r20 Source-Only Boundary-Reservoir Declaration Test}
\label{sec:v470-r20-boundary-reservoir-declaration-test}

This source-only pass performs the next admissible test after the r19
boundary-source scan.  The question is not whether the old grammar already
contains a free lift of \(V0\).  r19 showed that it does not.  The question is
whether the conditional boundary-source normal form isolated in r19 may be
introduced as a declared local boundary reservoir without creating new debt.
No PDF is produced in this intermediate pass.

\subsection{Declared local boundary-reservoir object}
\label{subsec:v470-r20-declared-boundary-reservoir-object}

We define the boundary-reservoir extension of the local finite valley-front
grammar by adjoining one explicitly declared boundary chart:
\[
\mathcal G_{\rm loc}^{r20,\partial{
m res}}
:=
\mathcal G_{\rm loc}^{r19}
\cup
\{\mathcal B_0^{\rm res}\}.
\]
The new object is not an external packet injection and not a global readout
channel.  It is a local boundary reservoir with declared data
\[
\mathcal B_0^{\rm res}
=
\left(
V0,
N=5,
\kappa=\frac34,
G=1,
\delta_{\partial}=0,
P_{\partial0}S_{4,\mu}^{\rm mixker}=0,
\chi_{\rm all}=0
\right).
\]
The associated local overlap score is
\[
S_0^{\partial{
m res}}
=
N+\kappa N(N-1)
=
5+\frac34\cdot 5\cdot4
=20.
\]
Since \(G=1\), the relay score is
\[
H_0^{\partial{
m res}}
=20>16.
\]
Thus the reservoir has relay margin
\[
\mu_0^{\partial{
m res}}
=20-16=4.
\]

\begin{definition}[r20 boundary-reservoir declaration]
\label{def:v470-r20-boundary-reservoir-declaration}
The object \(\mathcal B_0^{\rm res}\) is called a declared local
boundary-reservoir chart if it satisfies the following four clauses:
\[
\begin{array}{ll}
{\rm (D1)} & \delta_{\partial}=0,\\
{\rm (D2)} & P_{\partial0}S_{4,\mu}^{\rm mixker}=0,\\
{\rm (D3)} & \chi_{\rm all}=0,\\
{\rm (D4)} & H_0^{\partial{\rm res}}\ge 16.
\end{array}
\]
Clause (D1) forbids external boundary debt.  Clause (D2) forbids spending the
route-stable \(C4\) mixed-kernel residue.  Clause (D3) forbids all-readout or
\(T2\) promotion.  Clause (D4) is the local relay threshold.
\end{definition}

\subsection{No-new-debt proof obligation}
\label{subsec:v470-r20-no-new-debt-proof-obligation}

The no-new-debt vector for the declaration is
\[
\mathbf d_{\rm new}^{\partial{
m res}}
=
\left(
\delta_{\partial},
\delta_{C4},
\delta_{C6/T2},
\delta_{\rm all}
\right).
\]
For the declared reservoir we set
\[
\delta_{\partial}=0,
\qquad
\delta_{C4}=\left\|P_{\partial0}S_{4,\mu}^{\rm mixker}\right\|=0,
\qquad
\delta_{C6/T2}=0,
\qquad
\delta_{\rm all}=0.
\]
Therefore
\[
\boxed{
\mathbf d_{\rm new}^{\partial{
m res}}=(0,0,0,0).
}
\]
This proves only that the explicitly declared local reservoir is debt-neutral
with respect to the listed guards.  It does not prove that such a reservoir was
already present in the old r19 grammar.

\begin{proposition}[r20 no-new-debt declaration pass]
\label{prop:v470-r20-no-new-debt-declaration-pass}
In the extended local grammar \(\mathcal G_{\rm loc}^{r20,\partial{\rm res}}\),
the declared object \(\mathcal B_0^{\rm res}\) is relay-admissible for \(V0\)
and introduces no new boundary, \(C4\), \(C6/T2\), or all-readout debt.
\end{proposition}

\begin{proof}
The relay inequality follows from
\[
H_0^{\partial{\rm res}}
=20>16.
\]
The no-new-debt claim follows directly from the four declaration clauses in
Definition~\ref{def:v470-r20-boundary-reservoir-declaration}.  Clause (D1)
sets boundary debt to zero.  Clause (D2) makes the new boundary support
orthogonal to the route-stable mixed-kernel seam, so \(\beta_4\) is not spent.
Clause (D3) prevents promotion into all-readout or \(T2\) territory.  Clause
(D4) supplies the local relay inequality.  Hence the extension is conservative
at the level stated here.
\end{proof}

\subsection{Boundary-extended finite relay chain}
\label{subsec:v470-r20-boundary-extended-finite-relay-chain}

With the declared reservoir, the finite local chain may be written as
\[
\boxed{
V0^{\partial{\rm res}}
\to
V1\to V2\to V3\to V4\to V5\to V6\to V7\to V8\to V9.
}
\]
The superscript on \(V0\) is essential.  It records that the first chart is
not inherited as a free current-grammar chart.  It is admitted only in the
boundary-reservoir extension.

The active guard ledger after r20 is therefore
\[
\boxed{
\begin{array}{c|c|c}
\text{object} & \text{r20 status} & \text{guard}\\
\hline
V0 & \text{declared local boundary-reservoir relay-admissible} & \text{local only}\\
C4 & \beta_4=1/256\text{ route-stable mixed-kernel} & \text{unchanged}\\
C5 & \text{asymptotic bounded collar} & \text{no finite closure}\\
C6/T2 & \text{all-readout blocker} & \text{unchanged}\\
RH & \text{photon-clock handoff probe only} & \text{no proof claim}\\
\text{global closure} & \text{not claimed} & \text{guarded}
\end{array}
}
\]

\begin{corollary}[r20 boundary-extended local coverage]
\label{cor:v470-r20-boundary-extended-local-coverage}
The r20 extension certifies a boundary-extended finite local valley-front chain
from \(V0\) through \(V9\).  This is not a global atlas theorem and not an
all-readout closure theorem.
\end{corollary}

\begin{proof}
The chart \(V0\) is relay-admissible by
Proposition~\ref{prop:v470-r20-no-new-debt-declaration-pass}.  The charts
\(V1\) through \(V9\) are the finite relay chain certified by the r14--r17
source-only passes and guard-audited in r18.  The r20 declaration changes only
the local boundary entry chart.  It does not remove the \(C4\), \(C5\),
\(C6/T2\), RH, or global all-readout guards.
\end{proof}

\subsection{Permitted r21 direction}
\label{subsec:v470-r20-permitted-r21-direction}

The next source-only pass may test whether the boundary-extended finite chain
has any hidden seam coupling back into \(C4\), or whether the boundary reservoir
can be treated as a separate local source chart in the two-sided tensor grammar.
The safe formulation is
\[
\boxed{
\text{v4.7.0-r21: boundary-extended chain seam audit and two-sided source typing.}
}
\]
No finite \(C5\) closure, \(C6/T2\) lift, RH claim, or global all-readout
closure is licensed by r20.


\section{v4.7.0-r21: Boundary-Extended Chain Seam Audit and Two-Sided Source Typing}
\label{sec:v470-r21-boundary-chain-seam-typing}

\paragraph{Status.}
This is a source-only internal pass.  It audits the r20 boundary-reservoir
extension for hidden seam coupling and assigns a conservative two-sided source
type to the boundary chart.  It does not introduce a new DOI, does not rebuild a
PDF, and does not claim finite \(C5\) closure, \(C6/T2\) promotion, RH, or
global all-readout closure.

\subsection{Seam-audit object}
\label{subsec:v470-r21-seam-audit-object}

The r20 boundary extension produced the local finite chain
\[
V0^{\partial{\rm res}}
\to
V1\to V2\to V3\to V4\to V5\to V6\to V7\to V8\to V9.
\]
The r21 question is whether the new boundary entry creates a hidden seam back
into the route-stable \(C4\) mixed kernel, the \(C6/T2\) all-readout guard, or
an undeclared global closure channel.

We record the seam-audit vector
\[
\Sigma_{\partial{
m res}}
=
\left(
\Sigma_{\partial C4},
\Sigma_{\partial C6/T2},
\Sigma_{\partial{
m all}},
\Sigma_{\partial 1}
\right),
\]
where
\[
\Sigma_{\partial C4}
=
\left\|P_{\partial0}S_{4,\mu}^{\rm mixker}\right\|,
\qquad
\Sigma_{\partial C6/T2}=\chi_{\rm all},
\qquad
\Sigma_{\partial{
m all}}=\delta_{\rm all},
\]
and \(\Sigma_{\partial 1}\) is the boundary-to-first-interior relay seam.
The r20 declaration clauses give
\[
\Sigma_{\partial C4}=0,
\qquad
\Sigma_{\partial C6/T2}=0,
\qquad
\Sigma_{\partial{
m all}}=0.
\]
For the first relay seam, r20 supplies
\[
H_0^{\partial{\rm res}}=20,
\qquad
16\leq H_0^{\partial{\rm res}},
\qquad
\mu_0^{\partial{\rm res}}=4,
\]
so the boundary exit is locally relay-admissible.  Hence
\[
\boxed{
\Sigma_{\partial{\rm res}}=(0,0,0,0).
}
\]
This vanishing seam vector means only that the declared local boundary source
adds no new seam debt.  It does not erase the existing route-stable mixed-kernel
residue
\[
\boxed{\beta_4=\frac1{256}.}
\]

\begin{lemma}[r21 boundary seam audit]
\label{lem:v470-r21-boundary-seam-audit}
In the declared local grammar \(\mathcal G_{\rm loc}^{r20,\partial{\rm res}}\),
the boundary entry chart \(V0^{\partial{\rm res}}\) introduces no hidden seam
coupling into \(C4\), \(C6/T2\), or all-readout closure.
\end{lemma}

\begin{proof}
The \(C4\) seam vanishes by the explicit orthogonality clause
\(P_{\partial0}S_{4,\mu}^{\rm mixker}=0\).  The \(C6/T2\) and all-readout seams
vanish by \(\chi_{\rm all}=0\) and \(\delta_{\rm all}=0\).  The first relay seam
is admitted by the local inequality \(H_0^{\partial{\rm res}}=20\ge16\).  These
four facts give \(\Sigma_{\partial{\rm res}}=(0,0,0,0)\).  The statement is
local and does not change \(\beta_4\).
\end{proof}

\subsection{Two-sided source typing}
\label{subsec:v470-r21-two-sided-source-typing}

The boundary reservoir is now typed in the two-sided structural tensor grammar.
We write schematically
\[
{}^{\partial{
m src}}{}_{\rm loc}\mathcal B_0^{\rm res}{}^{V0}_{\rm relay}.
\]
The front indices \({}^{\partial{
m src}}{}_{\rm loc}\) mark a left-sided,
readout-facing local boundary source interface.  The back indices
\({}^{V0}_{\rm relay}\) mark the right-sided valley-front relay structure.  The
notation is not a new contraction rule.  It records that the object is a local
source chart whose output may feed the finite relay chain, while its boundary
source face is not an all-readout face.

The admissible readout is therefore the ordinary left-to-right local relay
readout
\[
\text{boundary reservoir}
\longmapsto
V0^{\partial{\rm res}}
\longmapsto
V1.
\]
Same-variance comparison against the \(C4\) mixed kernel is allowed only through
an explicit comparison/orthogonality test,
\[
G_{\partial C4}\,\mathcal B_{0}^{\partial}\,S_{4}^{\rm mixker}=0,
\]
which is the typed version of \(P_{\partial0}S_{4,\mu}^{\rm mixker}=0\).  Thus the
boundary source is not permitted to spend, cancel, or re-label the persistent
\(C4\) mixed-kernel residue.

\begin{definition}[r21 two-sided boundary-source type]
\label{def:v470-r21-two-sided-boundary-source-type}
The declared object \(
{}^{\partial{
m src}}{}_{\rm loc}\mathcal B_0^{\rm res}{}^{V0}_{\rm relay}
\)
is a two-sided local boundary-source chart if it satisfies:
\[
\begin{array}{ll}
{
m (T1)} & \text{left/front type }=\partial{\rm src}/\rm loc,\\
{
m (T2)} & \text{right/back type }=V0/\rm relay,\\
{
m (T3)} & G_{\partial C4}\mathcal B_{0}^{\partial}S_{4}^{\rm mixker}=0,\\
{
m (T4)} & \chi_{\rm all}=0,\\
{
m (T5)} & H_0^{\partial{\rm res}}\ge16.
\end{array}
\]
The type is local-source/relay.  It is not global-source, not all-readout, and
not a \(T2\) bridge.
\end{definition}

\subsection{Completion-interface audit}
\label{subsec:v470-r21-completion-interface-audit}

The r12--r20 passes leave \(C5\) as an asymptotic bounded-collar object,
\[
\Phi_5^{(n)}=\left(\frac34\right)^n,
\qquad
\widehat{\Phi}_5(\infty)=0,
\qquad
\Phi_5^{(n)}>0\quad(n<\infty).
\]
The r20 boundary reservoir does not alter this completion interface.  In r21 we
therefore record the guard
\[
\boxed{
V0^{\partial{\rm res}}\text{ enters the finite chain, but it does not change the }
C5\text{ completion topology.}
}
\]
Equivalently, the boundary source is a local chain-entry chart, not a new limit
operator on \(\mathbb N_\infty\).

\subsection{r21 decision ledger}
\label{subsec:v470-r21-decision-ledger}

The r21 decision is
\[
\boxed{
\begin{array}{c|c|c}
\text{object} & \text{r21 decision} & \text{guard}\\
\hline
V0^{\partial{\rm res}} & \text{two-sided local source/relay typed} & \text{local only}\\
\Sigma_{\partial{\rm res}} & (0,0,0,0) & \text{no new seam debt}\\
C4 & \beta_4=1/256 & \text{route-stable persistent}\\
C5 & \text{asymptotic bounded collar} & \text{no finite closure}\\
C6/T2 & \text{blocked} & \text{all-readout guard}\\
RH & \text{photon-clock handoff probe only} & \text{no proof claim}\\
\text{global closure} & \text{not claimed} & \text{guarded}
\end{array}
}
\]

\begin{proposition}[r21 typed local-chain preservation]
\label{prop:v470-r21-typed-local-chain-preservation}
The r20 boundary-extended finite chain remains valid after the r21 seam audit
and two-sided source typing, and no additional \(C4\), \(C5\), \(C6/T2\), RH, or
global all-readout claim is licensed.
\end{proposition}

\begin{proof}
Lemma~\ref{lem:v470-r21-boundary-seam-audit} gives the zero seam-debt vector for
the boundary entry.  Definition~\ref{def:v470-r21-two-sided-boundary-source-type}
assigns the boundary reservoir a local two-sided source/relay type and forbids
both \(C4\)-spending and all-readout promotion.  The internal chain
\(V1\to\cdots\to V9\) is unchanged from the r14--r17 finite relay construction.
Thus the finite boundary-extended chain is preserved locally, with all guards
listed in the decision ledger still active.
\end{proof}

\subsection{Permitted r22 direction}
\label{subsec:v470-r21-permitted-r22-direction}

The next source-only pass may test whether the typed boundary source can be
compiled into a small boundary-source calculus or whether it must remain a
single declared chart.  The safe formulation is
\[
\boxed{
\text{v4.7.0-r22: boundary-source calculus extraction and typed relay closure guard.}
}
\]
This r22 direction is still local and source-only.  It may not promote the
finite chain to global all-readout closure and may not identify the photon-clock
probe with RH.



\section{v4.7.0-r22 source-only: boundary-source calculus extraction and typed relay closure guard}
\label{sec:v470-r22-boundary-source-calculus-guard}

This source-only pass extracts a small local boundary-source calculus from the
single r20--r21 boundary-reservoir chart and immediately guards the extraction
against accidental promotion to global closure.  The pass does not rebuild a PDF,
does not assign a DOI, and does not modify the public release status.

The r21 typed reservoir is kept as
\[
{}^{\partial{\rm src}}{}_{\rm loc}\mathcal B_0^{\rm res}{}^{V0}_{\rm relay}.
\]
The r22 question is whether this typed object may be used as a reusable local
calculus rule rather than as an isolated chart.  The answer is yes, but only as a
local boundary-source rule with a typed relay guard.

\subsection{Local boundary-source calculus rule}
\label{subsec:v470-r22-local-boundary-source-calculus-rule}

\begin{definition}[r22 local boundary-source calculus rule]
\label{def:v470-r22-local-boundary-source-calculus-rule}
Let \(\mathcal B\) be a boundary reservoir candidate.  It is admissible as a
local boundary-source calculus instance, written
\[
\mathcal B\in\mathsf{BSC}_{\partial}^{\rm loc},
\]
if the following five clauses hold:
\[
\begin{array}{ll}
{\rm BSC1} & \mathcal B\text{ has two-sided type }
{}^{\partial{\rm src}}{}_{\rm loc}\mathcal B{}^{V}_{\rm relay};\\[2mm]
{\rm BSC2} & H_{\mathcal B}=G_{\mathcal B}S_{\mathcal B}\geq 16;\\[2mm]
{\rm BSC3} & P_{\partial}\,S_{4,\mu}^{\rm mixker}=0;\\[2mm]
{\rm BSC4} & \chi_{\rm all}(\mathcal B)=0;\\[2mm]
{\rm BSC5} & \mathbf d_{\rm new}(\mathcal B)=(0,0,0,0).
\end{array}
\]
The rule is therefore a local source-and-relay rule.  It is not an all-readout
rule and it is not a global boundary theorem.
\end{definition}

For the concrete r20/r21 reservoir one has
\[
N=5,\qquad \kappa=\frac34,\qquad G=1,
\]
so
\[
S=5+\frac34\cdot5\cdot4=20,
\qquad
H=20>16.
\]
Thus
\[
\boxed{
\mathcal B_0^{\rm res}\in\mathsf{BSC}_{\partial}^{\rm loc}.
}
\]

\begin{lemma}[r22 boundary-source calculus extraction]
\label{lem:v470-r22-boundary-source-calculus-extraction}
The r20--r21 boundary-reservoir chart can be extracted as a reusable local
boundary-source calculus instance in \(\mathsf{BSC}_{\partial}^{\rm loc}\).
\end{lemma}

\begin{proof}
BSC1 is the r21 two-sided source typing.  BSC2 follows from the r20 threshold
calculation \(H=20>16\).  BSC3 and BSC5 are the r20--r21 no-new-debt and seam-audit
conditions.  BSC4 is the explicit local-only readout guard.  Hence the reservoir
satisfies all five clauses of Definition~\ref{def:v470-r22-local-boundary-source-calculus-rule}.
\end{proof}

\subsection{Typed relay closure guard}
\label{subsec:v470-r22-typed-relay-closure-guard}

The word ``closure'' is used here only in the typed finite-relay sense.  It means
that a declared local boundary source may be attached to the beginning of the
already certified finite relay chain.  It does not mean global closure.

\begin{definition}[r22 typed finite-relay closure]
\label{def:v470-r22-typed-finite-relay-closure}
The typed finite-relay closure generated by \(\mathsf{BSC}_{\partial}^{\rm loc}\)
is the local operation
\[
\mathrm{Cl}_{\rm relay}^{\rm loc,typed}
\bigl(\mathcal B_0^{\rm res};V1,\ldots,V9\bigr)
=
\left[
V0^{\partial{\rm res}}\to V1\to\cdots\to V9
\right]
\]
provided the guard vector
\[
\Gamma_{\rm guard}^{r22}
=
(\gamma_{C4},\gamma_{C5}^{\rm fin},\gamma_{C6/T2},\gamma_{\rm all},\gamma_{\rm RH})
\]
remains
\[
\Gamma_{\rm guard}^{r22}=(1,1,1,1,1).
\]
Here \(1\) means that the corresponding guard is still active, not discharged.
\end{definition}

For r22 the guard ledger is
\[
\boxed{
\Gamma_{\rm guard}^{r22}=(1,1,1,1,1).
}
\]
In particular,
\[
\beta_4^{r22}=\frac1{256},
\qquad
\forall n<\infty:\ \Phi_5^{(n)}>0,
\qquad
C6:\mathsf B_{\rm all/T2}.
\]

\begin{proposition}[r22 typed relay closure guard]
\label{prop:v470-r22-typed-relay-closure-guard}
The boundary-source calculus extraction licenses the typed local finite relay
chain
\[
V0^{\partial{\rm res}}\to V1\to\cdots\to V9,
\]
but it does not license \(C4\)-closure, finite \(C5\)-closure, \(C6/T2\), RH, or
global all-readout closure.
\end{proposition}

\begin{proof}
By Lemma~\ref{lem:v470-r22-boundary-source-calculus-extraction}, the boundary
reservoir is a valid local calculus instance.  Definition~\ref{def:v470-r22-typed-finite-relay-closure}
then permits the local finite relay attachment.  The same definition also
requires the guard vector to remain active.  Since \(\beta_4=1/256\) is preserved,
finite \(C5\)-closure is excluded by the finite positivity of \(\Phi_5^{(n)}\),
\(C6/T2\) remains an all-readout blocker, and the photon-clock construction is a
handoff probe only, none of the guarded global claims is discharged.
\end{proof}

\subsection{r22 calculus ledger}
\label{subsec:v470-r22-calculus-ledger}

The extracted calculus can be summarized as
\[
\boxed{
\begin{array}{c|c|c}
\text{clause} & \text{r22 value} & \text{decision}\\
\hline
\text{two-sided type} & {}^{\partial{\rm src}}{}_{\rm loc}\mathcal B_0^{\rm res}{}^{V0}_{\rm relay} & \text{admissible}\\
H & 20>16 & \text{relay pass}\\
P_{\partial0}S_{4,\mu}^{\rm mixker} & 0 & \text{C4 guard preserved}\\
\chi_{\rm all} & 0 & \text{not all-readout}\\
\mathbf d_{\rm new} & (0,0,0,0) & \text{no new debt}\\
\Gamma_{\rm guard}^{r22} & (1,1,1,1,1) & \text{guards active}
\end{array}
}
\]

\subsection{Permitted r23 direction}
\label{subsec:v470-r22-permitted-r23-direction}

The next source-only pass may test whether the local boundary-source calculus is
stable under finite perturbations of \(N,\kappa,G\) and whether the typed relay
guard remains invariant under such perturbations.  The safe formulation is
\[
\boxed{
\text{v4.7.0-r23: boundary-source stability window and relay-guard robustness scan.}
}
\]
This may not be read as global all-readout closure, a finite \(C5\) closure, an
RH statement, or a \(T2\) bridge.



% --- v4.7.0-r23 Boundary-Source Stability Window and Relay-Guard Robustness Scan ---
\section{v4.7.0-r23 Boundary-Source Stability Window and Relay-Guard Robustness Scan}
\label{sec:v470-r23-boundary-source-stability-window}

This source-only pass tests whether the local boundary-source calculus extracted
in r22 is stable under small finite perturbations of the reservoir parameters
\((N,\kappa,G)\).  The pass is deliberately local.  It does not create a new DOI,
does not rebuild the PDF, and does not change any all-readout or global-closure
claim.

\subsection{Reservoir response function}
\label{subsec:v470-r23-reservoir-response-function}

For a local boundary-source reservoir we keep the r13--r22 overlap response
function
\[
S(N,\kappa)=N+\kappa N(N-1),
\qquad
H(N,\kappa,G)=G\,S(N,\kappa).
\]
The relay threshold is
\[
H(N,\kappa,G)\geq 16.
\]
The r22 anchor is
\[
(N,\kappa,G)=\left(5,\frac34,1\right),
\qquad
S=20,
\qquad
H=20,
\qquad
\mu=H-16=4.
\]

\begin{definition}[r23 boundary-source stability window]
\label{def:v470-r23-boundary-source-stability-window}
The weak r23 stability window is the set of local reservoir triples
\[
\mathcal W_{\partial}^{r23}
=
\{(N,\kappa,G):G(N+\kappa N(N-1))\geq 16\}
\]
subject to the r22 source-typing and guard clauses BSC1--BSC5.  The strict
stability window is
\[
\mathcal W_{\partial,+}^{r23}
=
\{(N,\kappa,G):G(N+\kappa N(N-1))>16\}
\]
with the same typing and guard clauses.
\end{definition}

For the fixed local source size \(N=5\), the threshold surfaces are explicit:
\[
G=1\quad\Longrightarrow\quad
\kappa\geq\frac{11}{20},
\]
while
\[
\kappa=\frac34\quad\Longrightarrow\quad
G\geq\frac45.
\]
A one-unit strict margin is already guaranteed by
\[
G=1,\\ \kappa\geq\frac35,
\]
or by
\[
\kappa=\frac34,\\ G\geq\frac{17}{20}.
\]
A simultaneous finite robustness rectangle around the r22 anchor is
\[
\boxed{
N=5,
\qquad
\kappa\geq\frac7{10},
\qquad
G\geq\frac9{10},
}
\]
because then
\[
S\geq 5+\frac7{10}\cdot 5\cdot4=19,
\qquad
H\geq\frac9{10}\cdot 19=\frac{171}{10}>16.
\]
Thus the r22 reservoir is not a knife-edge construction: it has a finite local
stability collar.

\begin{lemma}[r23 local stability collar]
\label{lem:v470-r23-local-stability-collar}
The r22 boundary-reservoir point
\((5,3/4,1)\) lies in the strict interior of the local boundary-source stability
window.  In particular, the rectangle
\[
N=5,
\qquad
\kappa\in\left[\frac7{10},1\right],
\qquad
G\in\left[\frac9{10},1\right]
\]
remains strictly relay-admissible.
\end{lemma}

\begin{proof}
For every triple in the displayed rectangle,
\[
S=5+20\kappa\geq 5+20\cdot\frac7{10}=19,
\]
and therefore
\[
H=GS\geq \frac9{10}\cdot19=\frac{171}{10}>16.
\]
The r22 source-typing and no-new-debt clauses are not altered by this finite
parameter perturbation, so the relay inequality remains valid in the same local
calculus.
\end{proof}

\subsection{Finite perturbation scan}
\label{subsec:v470-r23-finite-perturbation-scan}

The finite scan checked the grid
\[
N\in\{4,5,6\},
\qquad
\kappa\in\left\{\frac12,\frac{11}{20},\frac35,\frac23,\frac7{10},\frac34,1\right\},
\qquad
G\in\left\{\frac45,\frac{17}{20},\frac9{10},1\right\}.
\]
The scan is not used as a proof by enumeration.  It is a stability diagnostic
around the symbolic threshold surfaces above.  Its useful outcomes are:
\[
\begin{array}{c|c|c|c|c}
\text{case} & N & \kappa & G & \text{decision}\\
\hline
\text{r22 anchor} & 5 & 3/4 & 1 & H=20>16\\
\text{lower-}\kappa\text{ strict buffer} & 5 & 3/5 & 1 & H=17>16\\
\text{critical }\kappa\text{ surface} & 5 & 11/20 & 1 & H=16\\
\text{gain-damped strict point} & 5 & 3/4 & 17/20 & H=17>16\\
\text{gain-damped critical point} & 5 & 3/4 & 4/5 & H=16\\
\text{simultaneous robust corner} & 5 & 7/10 & 9/10 & H=171/10>16\\
\text{nominal }N=4 & 4 & 3/4 & 1 & H=13<16\\
\text{larger low-}\kappa\text{ reservoir} & 6 & 1/2 & 1 & H=21>16
\end{array}
\]
The failure of the nominal \(N=4\) case is useful: the local boundary-source rule
is not vacuous.  It still requires enough local reservoir capacity.

\subsection{Relay-guard robustness}
\label{subsec:v470-r23-relay-guard-robustness}

The r23 stability collar is guarded by exactly the same guard vector as r22:
\[
\boxed{
\Gamma_{\rm guard}^{r23}=(1,1,1,1,1).
}
\]
Concretely,
\[
\beta_4^{r23}=\frac1{256},
\qquad
\forall n<\infty:\ \Phi_5^{(n)}>0,
\qquad
C6:\mathsf B_{\rm all/T2}.
\]
The photon-clock construction remains a handoff probe, not an RH statement.
The boundary-source stability collar remains local and typed, not all-readout.

\begin{proposition}[r23 relay-guard robustness]
\label{prop:v470-r23-relay-guard-robustness}
Every triple in the strict local r23 stability collar preserves the r22 guard
vector.  Hence the local boundary-source robustness scan does not discharge
\(C4\), finite \(C5\)-closure, \(C6/T2\), RH, or global all-readout closure.
\end{proposition}

\begin{proof}
The perturbation scan changes only the local relay parameters \((N,\kappa,G)\)
inside the already declared two-sided boundary-source type.  The guard clauses
BSC3--BSC5 remain unchanged: the boundary source stays orthogonal to the
route-stable mixed kernel, the all-readout selector remains off, and the new-debt
vector remains zero.  Therefore the relay inequality can be stable without
spending \(\beta_4=1/256\), without converting the asymptotic \(C5\) collar into a
finite closure, and without opening the all-readout \(C6/T2\) gate.
\end{proof}

\subsection{r23 decision ledger}
\label{subsec:v470-r23-decision-ledger}

The source-only r23 decision is
\[
\boxed{
\begin{array}{c|c|c}
\text{object} & \text{r23 result} & \text{guard}\\
\hline
\mathcal B_0^{\rm res} & \text{strictly stable near }(5,3/4,1) & \text{local only}\\
\mathcal W_{\partial,+}^{r23} & \text{nonempty finite stability collar} & \chi_{\rm all}=0\\
C4 & \beta_4=1/256 & \text{unspent}\\
C5 & \Phi_5^{(n)}>0\text{ for finite }n & \text{no finite closure}\\
C6/T2 & \mathsf B_{\rm all/T2} & \text{blocked}\\
\text{RH} & \text{photon-clock handoff probe only} & \text{not a theorem}\\
\text{global all-readout} & \text{unclaimed} & \text{guard active}
\end{array}
}
\]

\subsection{Permitted r24 direction}
\label{subsec:v470-r23-permitted-r24-direction}

The next source-only step may package the stability collar into a small
``theoremlet'' and draw the residual frontier map showing exactly where the local
boundary-source calculus stops.  The safe title is
\[
\boxed{
\text{v4.7.0-r24: boundary-source robustness theoremlet and residual frontier map.}
}
\]
It must remain source-only unless a PDF is explicitly requested.



% --- v4.7.0-r24 Boundary-Source Robustness Theoremlet and Residual Frontier Map ---
\section{v4.7.0-r24 Boundary-Source Robustness Theoremlet and Residual Frontier Map}
\label{sec:v470-r24-boundary-source-theoremlet-frontier}

\noindent\textbf{Source-only status.}  This r24 layer is a source-only internal
working step.  It packages the r23 stability scan into a local theoremlet and
records the residual frontier of the boundary-source calculus.  It assigns no
new DOI, builds no PDF, and claims no global closure.

\subsection{Boundary-source robustness theoremlet}
\label{subsec:v470-r24-boundary-source-theoremlet}

The r23 anchor was
\[
(N,\kappa,G)=\left(5,\frac34,1\right),
\qquad
S=20,
\qquad
H=20,
\qquad
\mu=4.
\]
The useful r24 point is that this is not merely a witnessed chart.  It sits in a
small local robustness box.

\begin{theorem}[r24 local boundary-source robustness theoremlet]
\label{thm:v470-r24-local-boundary-source-robustness-theoremlet}
Let
\[
H(N,\kappa,G)=G\bigl(N+\kappa N(N-1)\bigr)
\]
be the local boundary-source relay height.  In the declared local
boundary-source grammar, every chart with
\[
N=5,
\qquad
\kappa\geq\frac7{10},
\qquad
G\geq\frac9{10}
\]
satisfies
\[
H\geq \frac{171}{10}>16.
\]
Hence the r22/r23 boundary-source calculus has a strict local robustness collar
around the anchor \((5,3/4,1)\).
\end{theorem}

\begin{proof}
For \(N=5\),
\[
S=5+20\kappa.
\]
If \(\kappa\geq7/10\), then
\[
S\geq5+20\cdot\frac7{10}=19.
\]
If also \(G\geq9/10\), then
\[
H=GS\geq \frac9{10}\cdot19=\frac{171}{10}>16.
\]
The proof uses only the already declared local boundary-source type.  It does
not switch on the all-readout selector and does not couple the boundary source
to the route-stable mixed kernel.
\end{proof}

\begin{corollary}[r24 strict local relay pass]
\label{cor:v470-r24-strict-local-relay-pass}
Every chart in the robustness box of
Theorem~\ref{thm:v470-r24-local-boundary-source-robustness-theoremlet} is
relay-admissible with positive margin.  In particular, the boundary-source
lift is not a knife-edge equality at \(H=16\).
\end{corollary}

\subsection{Threshold surfaces and finite frontier}
\label{subsec:v470-r24-threshold-surfaces-frontier}

The two one-parameter threshold surfaces around the r23 anchor are
\[
N=5,
\quad G=1
\quad\Longrightarrow\quad
\kappa\geq\frac{11}{20},
\]
and
\[
N=5,
\quad \kappa=\frac34
\quad\Longrightarrow\quad
G\geq\frac45.
\]
The strict one-unit buffer surfaces are
\[
N=5,
\quad G=1
\quad\Longrightarrow\quad
\kappa\geq\frac35,
\]
and
\[
N=5,
\quad \kappa=\frac34
\quad\Longrightarrow\quad
G\geq\frac{17}{20}.
\]

For a wider residual-frontier map, r24 records the finite grid
\[
N\in\{3,4,5,6,7\},
\]
\[
\kappa\in\left\{\frac12,\frac{11}{20},\frac35,\frac23,
\frac7{10},\frac34,1\right\},
\]
\[
G\in\left\{\frac45,\frac{17}{20},\frac9{10},1\right\}.
\]
The grid contains \(140\) cases and partitions as
\[
\boxed{
\begin{array}{c|c}
\text{frontier class} & \text{count}\\
\hline
\text{strict-interior }(H\geq17) & 66\\
\text{thin-pass }(16\leq H<17) & 6\\
\text{frontier-fail }(15\leq H<16) & 4\\
\text{deep-fail }(H<15) & 64
\end{array}
}
\]
Thus the boundary-source calculus has a genuine pass region, a thin critical
frontier, and a nontrivial fail region.  The local rule is therefore not a
vacuous route-to-closure device.

\subsection{Residual frontier map}
\label{subsec:v470-r24-residual-frontier-map}

The residual frontier is described by
\[
G\bigl(N+\kappa N(N-1)\bigr)=16.
\]
For fixed \(N\) and \(G\), the threshold is
\[
\kappa_{\min}(N,G)=\frac{16/G-N}{N(N-1)}.
\]
For fixed \(N\) and \(\kappa\), the threshold is
\[
G_{\min}(N,\kappa)=\frac{16}{N+\kappa N(N-1)}.
\]
At the representative slice \(G=1\), the threshold moves into the unit
coherence interval first at \(N=4\):
\[
N=4\Rightarrow \kappa_{\min}=1,
\qquad
N=5\Rightarrow \kappa_{\min}=\frac{11}{20},
\qquad
N=6\Rightarrow \kappa_{\min}=\frac13.
\]
This explains why the r23/r24 local reservoir really needs at least the
five-packet anchor to obtain a comfortable margin at moderate coherence.

\begin{proposition}[r24 residual-frontier decision]
\label{prop:v470-r24-residual-frontier-decision}
The declared local boundary-source calculus extends the finite valley-front
chain only inside the local pass region of
\[
G\bigl(N+\kappa N(N-1)\bigr)\geq16.
\]
Outside this region, the frontier remains residual and no relay-admissibility is
inferred.
\end{proposition}

\begin{proof}
The calculus rule BSC2 is exactly the local relay inequality \(H\geq16\).
If the inequality fails, the boundary source has no declared local relay height.
The source typing and no-new-debt clauses do not create missing relay height;
they only prevent successful local relay from spending forbidden guards.
Therefore the pass/fail frontier is exactly the stated inequality.
\end{proof}

\subsection{Guard-preserving consequence}
\label{subsec:v470-r24-guard-preserving-consequence}

The theoremlet is strictly local.  Its guard vector remains
\[
\boxed{
\Gamma_{\rm guard}^{r24}=(1,1,1,1,1).
}
\]
Equivalently,
\[
\beta_4^{r24}=\frac1{256},
\qquad
\forall n<\infty:\ \Phi_5^{(n)}>0,
\qquad
C6:\mathsf B_{\rm all/T2}.
\]
The Riemann-zero/photon-clock language remains a timing probe for handoff
windows, not an RH theorem and not a photon-zero identity.  The boundary-source
robustness theoremlet is a local relay theoremlet only.

\subsection{r24 decision ledger and permitted r25 direction}
\label{subsec:v470-r24-decision-ledger}

The r24 decision is
\[
\boxed{
\begin{array}{c|c|c}
\text{object} & \text{r24 result} & \text{guard}\\
\hline
\mathcal B_0^{\rm res} & \text{local robustness theoremlet} & \text{local only}\\
\mathcal W_{\partial,+}^{r24} & \text{finite pass/interior region mapped} & \chi_{\rm all}=0\\
\partial\mathcal W^{r24} & \text{residual frontier mapped} & \text{no inferred pass outside }H\geq16\\
C4 & \beta_4=1/256 & \text{unspent}\\
C5 & \Phi_5^{(n)}>0\text{ for finite }n & \text{no finite closure}\\
C6/T2 & \mathsf B_{\rm all/T2} & \text{blocked}\\
\text{RH} & \text{photon-clock handoff probe only} & \text{not a theorem}\\
\text{global all-readout} & \text{unclaimed} & \text{guard active}
\end{array}
}
\]

A safe next source-only direction is
\[
\boxed{
\text{v4.7.0-r25: boundary-source frontier minimizer and guard-stable atlas compression.}
}
\]
This would search for the smallest local boundary-source resources that still
preserve the guard vector, without turning the local finite chain into a global
closure claim.



\section{v4.7.0-r25 Source-Only Frontier Minimizer and Guard-Stable Atlas Compression}
\label{sec:v470-r25-frontier-minimizer-atlas-compression}

This source-only layer continues the r24 residual-frontier map.  No PDF is
built in this intermediate step.  The goal is to identify the smallest local
boundary-source resources visible on the finite frontier and to compress the
finite relay atlas without spending any global guard.

\subsection{Resource order and frontier minimizers}
\label{subsec:v470-r25-resource-order-frontier-minimizers}

For a local boundary-source triple
\[
(N,\kappa,G)
\]
we keep the r24 height
\[
H(N,\kappa,G)=G\bigl(N+\kappa N(N-1)\bigr).
\]
A triple is relay-admissible only if
\[
H(N,\kappa,G)\geq16.
\]
The source-resource order used in this scan is the componentwise order
\[
(N,\kappa,G)\preceq (N',\kappa',G')
\quad\Longleftrightarrow\quad
N\leq N',\quad \kappa\leq\kappa',\quad G\leq G'.
\]
A frontier point is Pareto-minimal if no strictly smaller admissible triple
exists in this order.

In the r24 finite grid the absolute pass minimizer is
\[
\boxed{(N,\kappa,G)=(4,1,1).}
\]
It gives
\[
H=1\cdot(4+1\cdot4\cdot3)=16.
\]
This is a thin pass with zero margin.  It is therefore not a robust boundary
reservoir rule; it is only the exact local frontier point.

For the five-packet boundary source the pass thresholds are
\[
G=1\Rightarrow \kappa\geq\frac{11}{20},
\qquad
\kappa=\frac34\Rightarrow G\geq\frac45.
\]
The margin-one stable thresholds are
\[
G=1\Rightarrow \kappa\geq\frac35,
\qquad
\kappa=\frac34\Rightarrow G\geq\frac{17}{20}.
\]
The simultaneous robust-box corner remains
\[
\boxed{(N,\kappa,G)=\left(5,\frac7{10},\frac9{10}\right),}
\]
with
\[
H=\frac9{10}\left(5+\frac7{10}\cdot5\cdot4\right)
=\frac{171}{10}>16.
\]
The original r22--r24 anchor
\[
\left(5,\frac34,1\right)
\]
has
\[
H=20,
\qquad
\mu=4.
\]
Thus r25 distinguishes three roles:
\[
\boxed{(4,1,1)=\text{absolute thin frontier},}
\]
\[
\boxed{\left(5,\frac7{10},\frac9{10}\right)=\text{robust local box corner},}
\]
\[
\boxed{\left(5,\frac34,1\right)=\text{comfortable boundary-source anchor}.}
\]

\begin{proposition}[r25 frontier minimizer decision]
\label{prop:v470-r25-frontier-minimizer-decision}
In the finite r24/r25 frontier grid, the triple \((4,1,1)\) is the absolute
local pass minimizer, but it has zero relay margin.  The five-packet
boundary-reservoir rule is the first locally robust reservoir family used by
this branch, and the box
\[
N=5,
\qquad
\kappa\geq\frac7{10},
\qquad
G\geq\frac9{10}
\]
is a strict local robustness region.
\end{proposition}

\begin{proof}
The equation
\[
G\bigl(N+\kappa N(N-1)\bigr)=16
\]
defines the frontier.  Direct substitution gives
\((4,1,1)\mapsto H=16\), hence pass with zero margin.  The same substitution at
\((5,7/10,9/10)\) gives \(H=171/10>16\).  Since increasing any of
\(N,\kappa,G\) weakly increases \(H\), the displayed box is contained in the
strict pass region.  The r22--r24 anchor \((5,3/4,1)\) lies strictly inside the
same region with margin \(4\).
\end{proof}

\subsection{Pareto frontier summary}
\label{subsec:v470-r25-pareto-frontier-summary}

The finite pass Pareto frontier consists of the thin and low-resource points
\[
\begin{array}{c|c|c|c|c}
N & \kappa & G & H & \text{status}\\
\hline
4 & 1 & 1 & 16 & \text{thin pass}\\
5 & 11/20 & 1 & 16 & \text{thin pass}\\
5 & 2/3 & 9/10 & 33/2 & \text{thin pass}\\
5 & 7/10 & 17/20 & 323/20 & \text{thin pass}\\
5 & 3/4 & 4/5 & 16 & \text{thin pass}\\
6 & 1/2 & 4/5 & 84/5 & \text{thin pass}
\end{array}
\]
These points describe where the boundary-source calculus first becomes locally
admissible.  They are not automatically robust enough to be used as source
anchors.

The margin-one Pareto layer begins with
\[
\begin{array}{c|c|c|c|c}
N & \kappa & G & H & \text{status}\\
\hline
5 & 3/5 & 1 & 17 & \text{strict interior}\\
5 & 7/10 & 9/10 & 171/10 & \text{strict interior}\\
5 & 3/4 & 17/20 & 17 & \text{strict interior}\\
5 & 1 & 4/5 & 20 & \text{strict interior}\\
6 & 1/2 & 17/20 & 357/20 & \text{strict interior}\\
6 & 11/20 & 4/5 & 18 & \text{strict interior}\\
7 & 1/2 & 4/5 & 112/5 & \text{strict interior}
\end{array}
\]
This layer is the better target for guard-stable atlas compression.

\subsection{Guard-stable atlas compression}
\label{subsec:v470-r25-guard-stable-atlas-compression}

The expanded local chain
\[
V0^{\partial\mathrm{res}}	o V1\to V2\to V3\to V4\to V5\to V6\to V7\to V8\to V9
\]
is compressed into four local atlas cells:
\[
\mathcal A_0=V0^{\partial\mathrm{res}},
\qquad
\mathcal A_1=(V1\to V2\to V3\to V4\to V5),
\]
\[
\mathcal A_2=(V5\to V6\to V7),
\qquad
\mathcal A_3=(V7\to V8\to V9).
\]
Here \(\mathcal A_0\) is the local boundary-reservoir source cell,
\(\mathcal A_1\) is the core valley relay cell, \(\mathcal A_2\) is the
shadow-basin/atlas-gap bridge cell, and \(\mathcal A_3\) is the zero-clock
handoff tail.

The compressed local atlas is therefore
\[
\boxed{
\mathcal A_{\partial}^{r25}
=
\mathcal A_3\circ\mathcal A_2\circ\mathcal A_1\circ\mathcal A_0.
}
\]
This notation is only a local atlas compression.  It does not turn the finite
chain into a global all-readout closure witness.

\begin{lemma}[r25 guard-stable compression]
\label{lem:v470-r25-guard-stable-compression}
The compressed atlas \(\mathcal A_{\partial}^{r25}\) preserves the r24 guard
vector
\[
\Gamma_{\rm guard}^{r25}=(1,1,1,1,1).
\]
\end{lemma}

\begin{proof}
Each compressed cell is inherited from a previously guarded local construction:
\(\mathcal A_0\) satisfies the boundary-source calculus with
\(\chi_{\rm all}=0\), zero new debt, and \(P_{\partial}S_{4,\mu}^{\rm mixker}=0\);
\(\mathcal A_1\) is the finite relay core; \(\mathcal A_2\) uses the r16/r17
source-supported bridge and spends no \(C4\) debt; \(\mathcal A_3\) uses
zero-clock windows only as timing probes.  None of these cells declares an
all-readout interface, an RH theorem, a finite \(C5\) closure, or a \(C6/T2\)
bridge.  Therefore the guard vector remains unchanged.
\end{proof}

\subsection{r25 decision ledger}
\label{subsec:v470-r25-decision-ledger}

The r25 decision is
\[
\boxed{
\begin{array}{c|c|c}
\text{object} & \text{r25 result} & \text{guard}\\
\hline
(4,1,1) & \text{absolute thin frontier minimizer} & \text{zero margin}\\
(5,7/10,9/10) & \text{robust local box corner} & \text{local only}\\
(5,3/4,1) & \text{comfortable boundary-source anchor} & \text{local only}\\
\mathcal A_{\partial}^{r25} & \text{compressed finite local atlas} & \chi_{\rm all}=0\\
C4 & \beta_4=1/256 & \text{unspent}\\
C5 & \Phi_5^{(n)}>0\text{ for finite }n & \text{no finite closure}\\
C6/T2 & \mathsf B_{\rm all/T2} & \text{blocked}\\
\text{RH} & \text{photon-clock handoff probe only} & \text{not a theorem}\\
\text{global all-readout} & \text{unclaimed} & \text{guard active}
\end{array}
}
\]

A safe next source-only direction is
\[
\boxed{
\text{v4.7.0-r26: compressed atlas transition audit and frontier-to-field-equation handoff.}
}
\]
This would test whether the compressed local atlas can feed the MM/HC
field-equation program as a typed local input without promoting any guarded
object to global closure.




\section{v4.7.0-r26 compressed atlas transition audit and frontier-to-field-equation handoff}
\label{sec:v470-r26-atlas-transition-field-handoff}

\subsection{Purpose and scope guard}
\label{subsec:v470-r26-purpose-scope}

This r26 source-only layer audits the compressed local atlas
\[
\mathcal A_{\partial}^{r25}
=
\mathcal A_3\circ\mathcal A_2\circ\mathcal A_1\circ\mathcal A_0
\]
as a possible typed input for the MM/HC field-equation program.  It does not
claim a new field equation, a global all-readout theorem, a finite \(C5\)
closure, a \(C6/T2\) bridge, or an RH theorem.  The purpose is narrower: check
whether the local atlas cells can be handed forward as a guarded source object
for the next field-equation search layer.

The r26 input objects are
\[
\mathcal A_0=V0^{\partial\mathrm{res}},
\qquad
\mathcal A_1=(V1\to V2\to V3\to V4\to V5),
\]
\[
\mathcal A_2=(V5\to V6\to V7),
\qquad
\mathcal A_3=(V7\to V8\to V9),
\]
together with the active guard vector
\[
\Gamma_{\rm guard}^{r25}=(1,1,1,1,1).
\]
The target is a local handoff object, not a global closure object.

\subsection{Two-sided transition typing}
\label{subsec:v470-r26-two-sided-transition-typing}

Each compressed atlas cell is typed as a two-sided local interface:
\[
{}^{\partial\mathrm{src}}{}_{\mathrm{loc}}\mathcal A_0{}^{V0}_{\mathrm{relay}},
\qquad
{}^{\mathrm{core}}{}_{\mathrm{loc}}\mathcal A_1{}^{V5}_{\mathrm{relay}},
\]
\[
{}^{\mathrm{shadow}}{}_{\mathrm{loc}}\mathcal A_2{}^{V7}_{\mathrm{bridge}},
\qquad
{}^{\gamma\mathrm{clock}}{}_{\mathrm{loc}}\mathcal A_3{}^{V9}_{\mathrm{tail}}.
\]
The front indices record the readout/source side of the cell.  The back indices
record the structure/relay side.  Same-variance interactions remain forbidden
unless a declared comparison tensor or selector mediates them.

The transition interfaces are
\[
\tau_{01}:\mathcal A_0\rightsquigarrow\mathcal A_1,
\qquad
\tau_{12}:\mathcal A_1\rightsquigarrow\mathcal A_2,
\qquad
\tau_{23}:\mathcal A_2\rightsquigarrow\mathcal A_3.
\]
Their audit ledger is
\[
\boxed{
\begin{array}{c|c|c|c}
\text{transition} & \text{source type} & \text{target type} & \text{r26 status}\\
\hline
\tau_{01} & \partial\mathrm{src}/V0 & \mathrm{core}/V5 & \text{typed local relay}\\
\tau_{12} & \mathrm{core}/V5 & \mathrm{shadow}/V7 & \text{source-supported bridge}\\
\tau_{23} & \mathrm{shadow}/V7 & \gamma\mathrm{clock}/V9 & \text{timing-gated tail}\\
\end{array}
}
\]
No transition in this table declares an all-readout channel.

\subsection{Transition seam audit}
\label{subsec:v470-r26-transition-seam-audit}

For each transition define a seam-debt vector
\[
\sigma(\tau)
=
(\sigma_{C4},\sigma_{C5},\sigma_{C6/T2},\sigma_{\rm all}).
\]
The r26 transition audit returns
\[
\boxed{
\sigma(\tau_{01})=(0,0,0,0),
\qquad
\sigma(\tau_{12})=(0,0,0,0),
\qquad
\sigma(\tau_{23})=(0,0,0,0).
}
\]
Therefore the compressed atlas transition vector is
\[
\boxed{
\Sigma_{\rm trans}^{r26}
=
\sum_{j=0}^{2}\sigma(\tau_{j,j+1})
=(0,0,0,0).
}
\]
This is a local transition statement.  It says that the compression introduces
no new seam debt across the declared local atlas cells.  It does not spend the
pre-existing route-stable mixed-kernel guard
\[
\beta_4=\frac1{256}.
\]

\begin{lemma}[r26 transition seam preservation]
\label{lem:v470-r26-transition-seam-preservation}
The compressed atlas transitions \(\tau_{01},\tau_{12},\tau_{23}\) preserve the
r25 guard vector.
\end{lemma}

\begin{proof}
The transition \(\tau_{01}\) uses the typed boundary-reservoir source calculus
with zero new boundary debt, \(P_{\partial}S_{4,\mu}^{\rm mixker}=0\), and
\(\chi_{\rm all}=0\).  The transition \(\tau_{12}\) is the already guarded
source-supported bridge from the core relay into the shadow-basin/atlas-gap
cell.  The transition \(\tau_{23}\) uses the photon-clock factor only as a local
timing gate for the zero-clock tail.  None of the three transitions declares an
all-readout channel, consumes \(\beta_4\), closes finite \(C5\), or opens
\(C6/T2\).  Hence \(\Sigma_{\rm trans}^{r26}=(0,0,0,0)\), and the guard vector
is preserved.
\end{proof}

\subsection{Frontier-to-field-equation handoff object}
\label{subsec:v470-r26-frontier-field-handoff}

The local atlas is now packaged as a handoff object
\[
\boxed{
\mathfrak X_{\partial}^{r26}
=
\left(
\mathcal A_{\partial}^{r25},
\Sigma_{\rm trans}^{r26},
\mathcal F_{\partial},
\Gamma_{\rm guard}^{r25}
\right),
}
\]
where \(\mathcal F_{\partial}\) denotes the boundary-source frontier surface
\[
\boxed{
G\bigl(N+\kappa N(N-1)\bigr)=16.
}
\]
This object is allowed to feed the MM/HC field-equation program only after
projection through the local guard projector
\[
\Pi_{\rm loc}^{\rm guard}:
\mathfrak X_{\partial}^{r26}
\longmapsto
\mathfrak X_{\partial,\rm loc}^{r26}.
\]
The projected object is local, typed, and guarded:
\[
\boxed{
\mathfrak X_{\partial,\rm loc}^{r26}
:=
\Pi_{\rm loc}^{\rm guard}\mathfrak X_{\partial}^{r26}.
}
\]

A schematic MM/HC handoff may then be written as
\[
{}^{\mathrm{loc}\,\partial\mathrm{atlas}}{}_{\mathrm{read}}
\mathfrak E_{\rm MM/HC}^{\rm seed}
{}^{\mathrm{field}}{}_{\mathrm{constraint}}
=
\mathsf H_{\rm bd}\circ\mathsf M_{\rm bd}
\left(\mathfrak X_{\partial,\rm loc}^{r26}\right).
\]
The symbol \(\mathfrak E_{\rm MM/HC}^{\rm seed}\) is deliberately a seed or
constraint object.  It is not yet a derived physical field equation.

\subsection{Handoff admissibility clauses}
\label{subsec:v470-r26-handoff-admissibility-clauses}

The handoff is admissible only under the following five clauses:
\[
\begin{array}{ll}
\mathrm{H1} & \mathcal A_{\partial}^{r25}\text{ is typed as a local atlas,}\\
\mathrm{H2} & \Sigma_{\rm trans}^{r26}=(0,0,0,0),\\
\mathrm{H3} & \mathcal F_{\partial}:G(N+\kappa N(N-1))=16\text{ remains recorded,}\\
\mathrm{H4} & \Gamma_{\rm guard}^{r26}=(1,1,1,1,1),\\
\mathrm{H5} & \mathfrak E_{\rm MM/HC}^{\rm seed}\text{ is a local constraint seed, not a closure theorem.}
\end{array}
\]
These clauses define the permitted frontier-to-field-equation handoff.  They
also forbid the common over-reading in which a finite local atlas is mistaken
for a global all-readout proof.

\begin{theoremlet}[r26 guarded atlas-to-field handoff]
\label{thmlet:v470-r26-guarded-atlas-field-handoff}
Under H1--H5, the compressed boundary-extended atlas
\(\mathcal A_{\partial}^{r25}\) is admissible as a local typed input seed for the
MM/HC field-equation program.  The handoff preserves all r25 guards.
\end{theoremlet}

\begin{proof}
By H1 the atlas is a local typed object.  By H2 no new transition seam debt is
introduced.  By H3 the frontier surface remains explicit, so the handoff does
not hide the boundary-source admissibility conditions.  By H4 all closure guards
remain active.  By H5 the field-equation side receives only a seed/constraint
object.  Hence the handoff is admissible as a local input to the MM/HC program,
but no global closure conclusion follows.
\end{proof}

\subsection{r26 decision ledger}
\label{subsec:v470-r26-decision-ledger}

The r26 decision ledger is
\[
\boxed{
\begin{array}{c|c|c}
\text{object} & \text{r26 result} & \text{guard}\\
\hline
\mathcal A_{\partial}^{r25} & \text{typed compressed local atlas} & \text{local only}\\
\tau_{01},\tau_{12},\tau_{23} & \text{transition-audited} & \Sigma_{\rm trans}=0\\
\mathfrak X_{\partial}^{r26} & \text{frontier handoff object} & \Pi_{\rm loc}^{\rm guard}\\
\mathfrak E_{\rm MM/HC}^{\rm seed} & \text{local constraint seed} & \text{not a field equation}\\
C4 & \beta_4=1/256 & \text{unspent}\\
C5 & \Phi_5^{(n)}>0\text{ for finite }n & \text{no finite closure}\\
C6/T2 & \mathsf B_{\rm all/T2} & \text{blocked}\\
\text{RH} & \text{photon-clock handoff probe only} & \text{not a theorem}\\
\text{global all-readout} & \text{unclaimed} & \text{guard active}
\end{array}
}
\]

A safe next source-only direction is
\[
\boxed{
\text{v4.7.0-r27: MM/HC seed-equation ansatz extraction from the guarded atlas handoff.}
}
\]
This should extract candidate seed constraints from
\(\mathfrak X_{\partial,\rm loc}^{r26}\) without upgrading them to physical field
equations or closure theorems.




% --- v4.7.0-r27 MM/HC Seed-Equation Ansatz Extraction from Guarded Atlas Handoff ---
\section{v4.7.0-r27: MM/HC Seed-Equation Ansatz Extraction from the Guarded Atlas Handoff}
\label{sec:v470-r27-mmhc-seed-equation-ansatz}

This source-only layer extracts the first MM/HC seed-equation ansatz family from
\(\mathfrak X_{\partial,\rm loc}^{r26}\).  The extraction is deliberately
conservative.  The output is a typed family of local seed constraints, not a
physical field equation and not a closure theorem.

\subsection{Projected handoff input}
\label{subsec:v470-r27-projected-handoff-input}

The input object inherited from r26 is
\[
\mathfrak X_{\partial}^{r26}
=
\left(
\mathcal A_{\partial}^{r25},
\Sigma_{\rm trans}^{r26},
\mathcal F_{\partial},
\Gamma_{\rm guard}^{r26}
\right),
\]
with
\[
\Sigma_{\rm trans}^{r26}=(0,0,0,0),
\qquad
\mathcal F_{\partial}:G\bigl(N+\kappa N(N-1)\bigr)=16.
\]
The seed extraction uses only the projected local object
\[
\boxed{
\mathfrak X_{\partial,\rm loc}^{r26}
:=
\Pi_{\rm loc}^{\rm guard}\mathfrak X_{\partial}^{r26}.
}
\]
The projection is part of the syntax.  No term may be extracted before passing
through the local guard.

\subsection{Seed-equation ansatz family}
\label{subsec:v470-r27-seed-equation-ansatz-family}

The first MM/HC seed family is
\[
\boxed{
\mathfrak E_{\rm MM/HC}^{\rm seed,r27}
=
\left(
E_{\rm rel},
E_{\rm src},
E_{\rm seam},
E_{\rm mix},
E_{\rm col},
E_{\rm clk}
\right).
}
\]
Each component is a local constraint seed:
\[
\begin{aligned}
E_{\rm rel}&:=G\bigl(N+\kappa N(N-1)\bigr)-16,\\
E_{\rm src}&:=H_{\partial}-16,\\
E_{\rm seam}&:=\|\Sigma_{\rm trans}^{r26}\|^2,\\
E_{\rm mix}&:=\|P_{\partial}S_{4,\mu}^{\rm mixker}\|^2+\left(\beta_4-\frac1{256}\right)^2,\\
E_{\rm col}&:=\Phi_5^{(n+1)}-\frac34\Phi_5^{(n)},\\
E_{\rm clk}&:=G-\mathcal C_{\gamma}(\tau_v).
\end{aligned}
\]
Here \(E_{\rm rel}\) records the valley-front relay frontier, \(E_{\rm src}\)
records the local boundary-source surplus, \(E_{\rm seam}\) records transition
exactness, \(E_{\rm mix}\) records that the mixed-kernel guard is not consumed,
\(E_{\rm col}\) records the bounded \(C5\) collar recurrence, and
\(E_{\rm clk}\) records the photon-clock timing probe.

\begin{definition}[r27 MM/HC seed ansatz]
\label{def:v470-r27-mmhc-seed-ansatz}
The r27 MM/HC seed ansatz is the six-component local constraint family
\(\mathfrak E_{\rm MM/HC}^{\rm seed,r27}\) extracted from
\(\mathfrak X_{\partial,\rm loc}^{r26}\) under the guard projector
\(\Pi_{\rm loc}^{\rm guard}\).
\end{definition}

\subsection{Two-sided seed typing}
\label{subsec:v470-r27-two-sided-seed-typing}

The seed ansatz is written in the two-sided notation as
\[
\boxed{
{}^{\partial\rm atlas}{}_{\rm loc/read}
\mathfrak E_{\rm MM/HC}^{\rm seed,r27}
{}^{\rm MM/HC}{}_{\rm constraint}.
}
\]
The left side records the local boundary-atlas readout, while the right side
records the MM/HC constraint target.  The expression is not a same-variance
Einstein contraction.  It is a typed local source-to-constraint interface.

A scalar audit functional may be used for bookkeeping:
\[
\boxed{
\mathcal S_{\rm seed}^{r27}
=
\lambda_{\rm rel}E_{\rm rel}^2
+\lambda_{\rm src}E_{\rm src,-}^2
+\lambda_{\rm seam}E_{\rm seam}
+\lambda_{\rm mix}E_{\rm mix}
+\lambda_{\rm col}E_{\rm col}^2
+\lambda_{\rm clk}E_{\rm clk}^2,
}
\]
where \(E_{\rm src,-}:=\max(0,16-H_{\partial})\).  This functional is a local
bookkeeping device only.  It is not an action principle and not a physical
Lagrangian.

\subsection{Admissibility and non-upgrade clauses}
\label{subsec:v470-r27-admissibility-nonupgrade}

The ansatz is admissible only under the following clauses:
\[
\begin{array}{ll}
\mathrm{A1} & \mathfrak X_{\partial}^{r26}\text{ is first projected by }\Pi_{\rm loc}^{\rm guard},\\
\mathrm{A2} & E_{\rm seam}=0\text{ is inherited from }\Sigma_{\rm trans}^{r26}=(0,0,0,0),\\
\mathrm{A3} & E_{\rm mix}\text{ preserves }\beta_4=1/256\text{ as unspent guard debt},\\
\mathrm{A4} & E_{\rm col}\text{ records only }\Phi_5^{(n+1)}=\frac34\Phi_5^{(n)}\text{ for finite stages},\\
\mathrm{A5} & E_{\rm clk}\text{ is a local photon-clock probe and not an RH theorem},\\
\mathrm{A6} & \mathfrak E_{\rm MM/HC}^{\rm seed,r27}\text{ is not promoted to a physical field equation.}
\end{array}
\]

\begin{lemma}[r27 seed extraction preserves guards]
\label{lem:v470-r27-seed-extraction-preserves-guards}
The extraction of \(\mathfrak E_{\rm MM/HC}^{\rm seed,r27}\) from
\(\mathfrak X_{\partial,\rm loc}^{r26}\) preserves the r26 guard vector.
\end{lemma}

\begin{proof}
The projected handoff already has zero transition seam vector.  The relay and
boundary-source terms record local admissibility conditions but do not introduce
an all-readout channel.  The mixed-kernel term explicitly keeps
\(\beta_4=1/256\) as a recorded guard rather than consuming it.  The collar term
records the finite recurrence of \(C5\) and includes no finite closure clause.
The clock term is a timing probe only.  Hence no guarded channel is discharged.
\end{proof}

\begin{theoremlet}[r27 guarded seed-equation ansatz]
\label{thmlet:v470-r27-guarded-seed-equation-ansatz}
The six-component family \(\mathfrak E_{\rm MM/HC}^{\rm seed,r27}\) is an
admissible local MM/HC seed-equation ansatz extracted from the guarded atlas
handoff.  It is not a physical field equation and not a closure theorem.
\end{theoremlet}

\begin{proof}
By A1 the input is local and guard-projected.  By A2 the transition seam
component is exact.  By A3 the route-stable mixed-kernel residual is preserved
rather than spent.  By A4 no finite \(C5\) closure is claimed.  By A5 the
Riemann-zero/photon-clock channel remains a probe.  By A6 the output is a seed
constraint family only.  Therefore the ansatz is admissible as a guarded local
MM/HC seed and no global conclusion follows.
\end{proof}

\subsection{r27 decision ledger}
\label{subsec:v470-r27-decision-ledger}

The r27 decision ledger is
\[
\boxed{
\begin{array}{c|c|c}
\text{seed component} & \text{local role} & \text{guard status}\\
\hline
E_{\rm rel} & \text{relay frontier} & \text{local only}\\
E_{\rm src} & \text{boundary-source surplus} & \text{no global source}\\
E_{\rm seam} & \text{transition seam exactness} & \Sigma_{\rm trans}=0\\
E_{\rm mix} & \text{mixed-kernel guard record} & \beta_4=1/256\text{ unspent}\\
E_{\rm col} & \text{bounded }C5\text{ collar recurrence} & \text{no finite closure}\\
E_{\rm clk} & \text{photon-clock handoff probe} & \text{not RH theorem}
\end{array}
}
\]
The global guard vector remains
\[
\boxed{
\Gamma_{\rm guard}^{r27}=(1,1,1,1,1).
}
\]
A safe next source-only direction is
\[
\boxed{
\text{v4.7.0-r28: seed-ansatz consistency scan and MM/HC coefficient typing.}
}
\]
This next step should classify the coefficients
\(\lambda_{\rm rel},\lambda_{\rm src},\lambda_{\rm seam},\lambda_{\rm mix},
\lambda_{\rm col},\lambda_{\rm clk}\) as bookkeeping weights, selector weights,
or admissibility multipliers, without turning \(\mathcal S_{\rm seed}^{r27}\)
into a physical action.




% --- v4.7.0-r28 Seed-Ansatz Consistency Scan and MM/HC Coefficient Typing ---
\section{v4.7.0-r28: Seed-Ansatz Consistency Scan and MM/HC Coefficient Typing}
\label{sec:v470-r28-seed-consistency-coefficient-typing}

This source-only layer audits the r27 seed ansatz
\(\mathfrak E_{\rm MM/HC}^{\rm seed,r27}\) before any field-equation upgrade is
attempted.  The goal is coefficient typing.  The scan deliberately prevents the
local seed functional from being misread as a physical action, a Lagrangian, or a
closure theorem.

\subsection{Consistency issue: signed surplus is not a zero residual}
\label{subsec:v470-r28-signed-surplus-not-zero-residual}

The r27 family contains the signed frontier coordinate
\[
E_{\rm rel}=G\bigl(N+\kappa N(N-1)\bigr)-16.
\]
This coordinate is useful, but it is not always a residual that should be driven
to zero.  If a local relay chart is strictly admissible, then
\[
H=G\bigl(N+\kappa N(N-1)\bigr)>16,
\]
so
\[
E_{\rm rel}>0.
\]
Penalising \(E_{\rm rel}^2\) as an ordinary equality residual would incorrectly
push a robustly admissible chart back onto the thin frontier \(H=16\).  Thus the
signed frontier coordinate must be typed as a monitor, while the admissibility
penalty uses only the violation part
\[
\boxed{
E_{\rm rel}^{-}:=\max\{0,16-H\}.
}
\]
Similarly, for the boundary-source surplus,
\[
\boxed{
E_{\rm src}^{-}:=\max\{0,16-H_{\partial}\}.
}
\]
The positive surplus
\[
\mu:=\max\{0,H-16\}
\]
is recorded as robustness margin, not as debt.

\begin{definition}[r28 violation-projected relay residual]
\label{def:v470-r28-violation-projected-relay-residual}
The r28 relay residual is the violation-projected quantity
\(E_{\rm rel}^{-}=\max\{0,16-H\}\), while the signed quantity
\(E_{\rm rel}=H-16\) is retained only as a frontier coordinate and margin
monitor.
\end{definition}

\subsection{Typed coefficient classes}
\label{subsec:v470-r28-typed-coefficient-classes}

The r27 coefficients are now divided into three noninterchangeable classes:
\[
\boxed{
\lambda\in
\{\lambda^{\rm book},\lambda^{\rm sel},\lambda^{\rm adm}\}.
}
\]
Here
\[
\begin{array}{ll}
\lambda^{\rm book} & \text{bookkeeping or diagnostic weight only},\\
\lambda^{\rm sel}  & \text{selector/timing weight for admissible local routing},\\
\lambda^{\rm adm}  & \text{admissibility multiplier for a hard local guard.}
\end{array}
\]
The r28 typing ledger is
\[
\boxed{
\begin{array}{c|c|c}
\text{coefficient} & \text{r28 type} & \text{role}\\
\hline
\lambda_{\rm rel} & \lambda^{\rm book}+\lambda^{\rm adm} &
\text{signed monitor plus violation penalty }(E_{\rm rel}^{-})^2\\
\lambda_{\rm src} & \lambda^{\rm adm} &
\text{boundary-source deficit penalty }(E_{\rm src}^{-})^2\\
\lambda_{\rm seam} & \lambda^{\rm adm} &
\text{transition seam exactness guard}\\
\lambda_{\rm mix} & \lambda^{\rm adm} &
\text{mixed-kernel guard lock, keeps }\beta_4=1/256\text{ unspent}\\
\lambda_{\rm col} & \lambda^{\rm book}+\lambda^{\rm adm}_{\rm loc} &
\text{finite collar recurrence consistency only}\\
\lambda_{\rm clk} & \lambda^{\rm sel} &
\text{photon-clock timing selector/probe}
\end{array}
}
\]
The mixed notation \(\lambda^{\rm book}+\lambda^{\rm adm}\) means that the
signed value may be displayed diagnostically, but only its admissibility-violation
projection may be treated as debt.

\subsection{Typed seed audit functional}
\label{subsec:v470-r28-typed-seed-audit-functional}

The scalar r27 bookkeeping expression is therefore replaced, for audit purposes,
by the typed functional
\[
\boxed{
\mathcal S_{\rm seed}^{r28,\rm typed}
=
\alpha_{\rm rel}(E_{\rm rel}^{-})^2
+\alpha_{\rm src}(E_{\rm src}^{-})^2
+\alpha_{\rm seam}E_{\rm seam}
+\alpha_{\rm mix}E_{\rm mix}
+\alpha_{\rm col}E_{\rm col}^2
+\alpha_{\rm clk}^{\rm sel}E_{\rm clk}^2.
}
\]
with the side condition that \(\alpha_{\rm clk}^{\rm sel}\) is active only inside
a declared photon-clock phase window.  Outside such a declared window,
\(E_{\rm clk}\) remains a probe and contributes no hard debt.

For the r22--r25 boundary-source anchor
\[
(N,\kappa,G)=\left(5,\frac34,1\right),
\]
one has
\[
H=20,
\qquad
E_{\rm rel}=4,
\qquad
E_{\rm rel}^{-}=0.
\]
Thus the anchor has positive margin, but no relay violation debt.  This resolves
the apparent contradiction between a robust pass and a zero-residual audit.

\subsection{Seed-component consistency scan}
\label{subsec:v470-r28-seed-component-consistency-scan}

The six r27 seed components pass the r28 consistency scan with the following
interpretations:
\[
\boxed{
\begin{array}{c|c|c}
\text{component} & \text{consistency decision} & \text{allowed use}\\
\hline
E_{\rm rel} & \text{consistent after violation projection} &
\text{frontier monitor plus local deficit test}\\
E_{\rm src} & \text{consistent after deficit projection} &
\text{boundary-source admissibility test}\\
E_{\rm seam} & \text{consistent hard exactness seed} &
\Sigma_{\rm trans}^{r26}=0\\
E_{\rm mix} & \text{consistent guard-preserving seed} &
\beta_4=1/256\text{ retained}\\
E_{\rm col} & \text{consistent finite recurrence seed} &
\text{no finite }C5\text{ closure}\\
E_{\rm clk} & \text{consistent selector/probe seed} &
\text{not an RH theorem}
\end{array}
}
\]
In particular, \(E_{\rm rel}\) and \(E_{\rm src}\) are not independent physical
field equations.  When \(H_{\partial}=H\), they are the same local relay surplus
viewed from two bookkeeping sides: the valley-front side and the boundary-source
side.

\begin{lemma}[r28 coefficient typing removes the frontier-square artefact]
\label{lem:v470-r28-frontier-square-artefact-removed}
Replacing \(E_{\rm rel}^2\) by \((E_{\rm rel}^{-})^2\) in the admissibility part
of the seed audit prevents robust relay margins from being misclassified as
residual debt.
\end{lemma}

\begin{proof}
If \(H\geq16\), then \(E_{\rm rel}^{-}=0\), so the local relay chart contributes
no admissibility debt regardless of the positive surplus \(H-16\).  If \(H<16\),
then \(E_{\rm rel}^{-}=16-H>0\), so the deficit is recorded.  Thus the typed
residual detects exactly local relay failure and does not penalise robust pass
interior points.
\end{proof}

\subsection{MM/HC coefficient typing rule}
\label{subsec:v470-r28-mmhc-coefficient-typing-rule}

The r28 coefficient rule is:
\[
\boxed{
\text{No coefficient may be interpreted before its seed component is typed.}
}
\]
Equivalently,
\[
\boxed{
\lambda_i E_i^2
\quad\text{is admissible only after declaring whether }E_i
\text{ is a monitor, selector, violation residual, or hard guard.}
}
\]
This rule blocks the common false upgrade
\[
\mathcal S_{\rm seed}\longrightarrow\text{physical action}.
\]
The present object remains
\[
\boxed{
\mathcal S_{\rm seed}^{r28,\rm typed}
\in
\text{local MM/HC seed-audit grammar only.}
}
\]

\begin{theoremlet}[r28 seed-ansatz consistency and coefficient typing]
\label{thmlet:v470-r28-seed-consistency-coefficient-typing}
After violation projection and coefficient typing, the r27 seed ansatz is
consistent as a guarded local MM/HC seed-audit system.  It is not a physical
field equation, not an action principle, and not a closure theorem.
\end{theoremlet}

\begin{proof}
The only apparent inconsistency is the treatment of positive relay surplus as a
zero residual.  Definition~\ref{def:v470-r28-violation-projected-relay-residual}
resolves this by splitting the signed frontier monitor from the violation
residual.  The seam, mixed-kernel, collar, and clock components keep their r27
guards: seam exactness remains local, \(\beta_4=1/256\) is preserved, finite
\(C5\) closure is not claimed, and the photon-clock term remains a probe.  Since
all coefficients are typed before use, no scalar audit expression can be promoted
to an action or global field equation.  Therefore the typed seed-audit system is
consistent in the declared local grammar only.
\end{proof}

\subsection{r28 guard ledger}
\label{subsec:v470-r28-guard-ledger}

The r28 guard ledger is
\[
\boxed{
\begin{array}{c|c|c}
\text{guard} & \text{r28 status} & \text{reason}\\
\hline
C4 & \beta_4=1/256\text{ unspent} & E_{\rm mix}\text{ remains a guard lock}\\
C5 & \text{no finite closure} & E_{\rm col}\text{ records recurrence only}\\
C6/T2 & \text{blocked} & \chi_{\rm all}=0\text{ remains active}\\
\text{RH} & \text{probe only} & E_{\rm clk}\text{ is selector/timing, not theorem}\\
\text{global all-readout} & \text{unclaimed} & \Pi_{\rm loc}^{\rm guard}\text{ remains mandatory}
\end{array}
}
\]
Thus
\[
\boxed{
\Gamma_{\rm guard}^{r28}=(1,1,1,1,1).
}
\]

A safe next source-only direction is
\[
\boxed{
\text{v4.7.0-r29: typed seed residual minimizer and selector-window calibration.}
}
\]
The next step should use the r28 typed residuals to scan local minimisers without
allowing a physical-action or closure upgrade.


% -----------------------------------------------------------------------------
\section{v4.7.0-r29 Typed Seed Residual Minimizer and Selector-Window Calibration}
\label{sec:v470-r29-typed-seed-residual-minimizer}

This source-only step uses the r28 typed seed audit to define a minimizer for the
local MM/HC seed residuals.  The important point is that the minimizer is applied
to violation-projected residuals, not to signed surplus monitors.  Therefore a
robust relay interior point is not pushed back to the thin frontier.

\subsection{Typed residual functional}
\label{subsec:v470-r29-typed-residual-functional}

Let
\[
H(N,\kappa,G)=G\bigl(N+\kappa N(N-1)\bigr)
\]
be the local relay height.  The r29 typed admissibility residual is
\[
\boxed{
\mathcal D_{\rm typed}^{r29}
=
\lambda_{\rm rel}\bigl(E_{\rm rel}^{-}\bigr)^2
+
\lambda_{\rm src}\bigl(E_{\rm src}^{-}\bigr)^2
+
\lambda_{\rm seam}\|\Sigma_{\rm trans}^{r26}\|^2
+
\lambda_{\rm mix}\|P_{\partial}S_{4,\mu}^{\rm mixker}\|^2
+
\lambda_{\rm col}\bigl(E_{\rm col}\bigr)^2
+
\lambda_{\rm clk}\bigl(E_{\rm clk}^{\rm win}\bigr)^2.
}
\]
Here
\[
E_{\rm rel}^{-}=\max(0,16-H),
\qquad
E_{\rm src}^{-}=\max(0,16-H_{\partial}).
\]
The clock term is not a theorem term.  It is a selector-window residual:
\[
\boxed{
E_{\rm clk}^{\rm win}
=
\max(0,G_{\rm req}-\mathcal C_{\gamma}(\tau_v)).
}
\]
Thus photon-clock timing is used only to decide whether a local relay window is
available.

\begin{definition}[r29 typed zero-debt minimizer]
\label{def:v470-r29-typed-zero-debt-minimizer}
A local seed chart is a typed zero-debt minimizer if
\[
\mathcal D_{\rm typed}^{r29}=0
\]
inside the declared local grammar.  This means that all admissibility deficits,
seam debts, mixed-kernel couplings, collar-recurrence defects, and clock-window
deficits vanish.  It does not mean that the signed surplus monitors vanish.
\end{definition}

\subsection{Minimizer degeneracy and margin classes}
\label{subsec:v470-r29-minimizer-degeneracy-margin-classes}

Because r28 removed the frontier-square artefact, every admissible relay point
with
\[
H\geq16
\]
has zero relay debt.  Hence the minimizer set is not a single point.  It has a
thin frontier and an interior.

We classify zero-debt points by margin
\[
\mu=H-16.
\]
The local classes are
\[
\boxed{
\begin{array}{c|c}
\text{class} & \text{condition}\\
\hline
\mathsf{frontier} & \mu=0\\
\mathsf{thin} & 0<\mu<1\\
\mathsf{robust} & 1\leq\mu<4\\
\mathsf{comfortable} & \mu\geq4.
\end{array}
}
\]
Thus the finite-grid minimizer
\[
(N,\kappa,G)=(4,1,1)
\]
remains a zero-debt point, but only a frontier minimizer.  The r22--r25 anchor
\[
(N,\kappa,G)=\left(5,{3\over4},1\right)
\]
is a comfortable zero-debt minimizer with
\[
H=20,
\qquad
\mu=4.
\]

\subsection{Selector-window calibration}
\label{subsec:v470-r29-selector-window-calibration}

For a fixed overlap strength
\[
S(N,\kappa)=N+\kappa N(N-1),
\]
the gain required for relay is
\[
\boxed{
G_{\rm req}(N,\kappa)={16\over S(N,\kappa)}.
}
\]
For a strict margin-one relay it is
\[
\boxed{
G_{\rm req}^{(+1)}(N,\kappa)={17\over S(N,\kappa)}.
}
\]
For the comfortable r22 anchor
\[
S\left(5,{3\over4}\right)=20,
\]
so
\[
G_{\rm req}={4\over5},
\qquad
G_{\rm req}^{(+1)}={17\over20}.
\]
If a photon-clock window satisfies
\[
\mathcal C_{\gamma}(\tau_v)\geq {17\over20},
\]
then the anchor has at least one full unit of relay margin.  If it satisfies
\[
\mathcal C_{\gamma}(\tau_v)\geq {4\over5},
\]
it is relay-admissible but may be near the thin boundary.

\begin{lemma}[r29 selector-window calibration]
\label{lem:v470-r29-selector-window-calibration}
For fixed local packet data \((N,\kappa)\), the selector-window threshold is
\[
G_{\rm req}=16/S(N,\kappa).
\]
The clock term contributes zero selector debt exactly when
\[
\mathcal C_{\gamma}(\tau_v)\geq G_{\rm req}.
\]
\end{lemma}

\begin{proof}
The local relay condition is
\[
G S(N,\kappa)\geq16.
\]
Solving for the required gain gives \(G_{\rm req}=16/S(N,\kappa)\).  Since the
clock term is typed as a selector-window residual,
\[
E_{\rm clk}^{\rm win}=\max(0,G_{\rm req}-\mathcal C_{\gamma}(\tau_v)),
\]
it vanishes exactly when the available photon-clock gain meets or exceeds the
required threshold.
\end{proof}

\subsection{Typed minimizer scan result}
\label{subsec:v470-r29-typed-minimizer-scan-result}

The r29 finite scan uses the same local grid as r24/r25 and evaluates the typed
residual.  It finds that the typed residual minimizer set is precisely the
admissible region of the local frontier map, stratified by margin.  The audit
therefore selects two representatives:
\[
\boxed{
(4,1,1)
\quad\text{as the frontier zero-debt representative,}
}
\]
and
\[
\boxed{
\left(5,{3\over4},1\right)
\quad\text{as the comfortable zero-debt representative.}
}
\]
The second representative is retained as the operational boundary-source anchor
because it has margin \(4\) and lies inside the robust local collar.

\begin{theoremlet}[r29 typed residual minimizer]
\label{thmlet:v470-r29-typed-residual-minimizer}
Using violation-projected residuals, the local MM/HC seed audit admits zero-debt
minimizers throughout the relay-admissible region \(H\geq16\).  The robust
boundary-source anchor \((5,3/4,1)\) is a comfortable zero-debt minimizer, while
\((4,1,1)\) is only a frontier zero-debt minimizer.
\end{theoremlet}

\begin{proof}
For every point with \(H\geq16\), both \(E_{\rm rel}^{-}\) and
\(E_{\rm src}^{-}\) vanish.  The seam and mixed-kernel terms vanish by the local
boundary-source typing clauses, the collar recurrence is the r28 bookkeeping
recurrence, and the clock term vanishes when the selector window is met.  Thus
all admissible points are typed zero-debt minimizers.  The margin stratification
then separates the thin frontier from robust and comfortable interior points.
For \((4,1,1)\), \(H=16\), so the margin is zero.  For \((5,3/4,1)\),
\(H=20\), so the margin is four.
\end{proof}

\subsection{r29 guard ledger}
\label{subsec:v470-r29-guard-ledger}

The minimizer is local and typed.  It does not create a physical action, a field
equation, or a closure theorem.  The r29 guard ledger is
\[
\boxed{
\begin{array}{c|c|c}
\text{guard} & \text{r29 status} & \text{reason}\\
\hline
C4 & \beta_4=1/256\text{ unspent} & E_{\rm mix}\text{ remains a guard-lock}\\
C5 & \text{no finite closure} & E_{\rm col}\text{ is recurrence consistency only}\\
C6/T2 & \text{blocked} & \chi_{\rm all}=0\text{ remains active}\\
\text{RH} & \text{probe only} & E_{\rm clk}^{\rm win}\text{ is a selector-window term}\\
\text{global all-readout} & \text{unclaimed} & \Pi_{\rm loc}^{\rm guard}\text{ remains mandatory}
\end{array}
}
\]
Therefore
\[
\boxed{
\Gamma_{\rm guard}^{r29}=(1,1,1,1,1).
}
\]

A safe next source-only direction is
\[
\boxed{
\text{v4.7.0-r30: typed MM/HC seed normal form and admissibility certificate.}
}
\]
The next step should package the typed minimizer as a normal-form certificate for
local MM/HC seed admissibility, still without promoting it to a physical field
equation or global closure theorem.




\section{v4.7.0-r30 Typed MM/HC Seed Normal Form and Admissibility Certificate}
\label{sec:v470-r30-seed-normal-form-certificate}

\paragraph{Source-only status.}
This r30 layer is a source-only continuation of the v4.7.0 internal line.  It
packages the r28--r29 typed residual calculus as a local MM/HC seed normal form
and an admissibility certificate.  It does not introduce a physical action, a
Lagrangian, a field equation, a T2 bridge, an RH statement, or a global
all-readout closure theorem.

\subsection{Typed seed normal form}
\label{subsec:v470-r30-typed-seed-normal-form}

Let
\[
S(N,\kappa)=N+\kappa N(N-1),
\qquad
H=G S(N,\kappa).
\]
The r30 local seed normal form is the guarded tuple
\[
\boxed{
\mathfrak N_{\rm seed}^{r30}
=
\left(
N,\kappa,G;
\Sigma_{\rm trans},
P_{\partial}S_{4,\mu}^{\rm mixker},
\beta_4,
\Phi_5^{(n)},
\mathcal C_{\gamma}(\tau_v),
\Gamma_{\rm guard}
\right)
}
\]
subject to the following typed clauses:
\[
\boxed{
\begin{array}{ll}
{\rm NF1} & \text{two-sided local boundary-source/relay typing is declared;}\\
{\rm NF2} & E_{\rm rel}^{-}=\max(0,16-H)=0;\\
{\rm NF3} & E_{\rm src}^{-}=\max(0,16-H_{\partial})=0;\\
{\rm NF4} & \Sigma_{\rm trans}=0;\\
{\rm NF5} & P_{\partial}S_{4,\mu}^{\rm mixker}=0\ \text{and}\ \beta_4=1/256;\\
{\rm NF6} & \Phi_5^{(n+1)}={3\over4}\Phi_5^{(n)}\ \text{is only a local collar recurrence;}\\
{\rm NF7} & \mathcal C_{\gamma}(\tau_v)\geq16/S(N,\kappa)\ \text{is only a selector-window probe;}\\
{\rm NF8} & \Gamma_{\rm guard}=(1,1,1,1,1)\ \text{remains active.}
\end{array}
}
\]
The normal form is intentionally asymmetric with respect to surplus: positive
relay surplus is not penalized.  Only deficit is counted through
\(E_{\rm rel}^{-}\) and \(E_{\rm src}^{-}\).

\begin{definition}[r30 typed MM/HC seed normal form]
\label{def:v470-r30-typed-seed-normal-form}
A local boundary-atlas handoff is in typed MM/HC seed normal form if it admits a
tuple \(\mathfrak N_{\rm seed}^{r30}\) satisfying NF1--NF8.  The declaration is
local: it certifies admissibility of the guarded seed constraints only, not a
physical field equation and not a closure theorem.
\end{definition}

\subsection{Admissibility certificate}
\label{subsec:v470-r30-admissibility-certificate}

The associated certificate is
\[
\boxed{
\mathsf{Cert}_{\rm MM/HC}^{r30}
=
\left[
\mathrm{Type}_{\partial/{\rm relay}},
E_{\rm rel}^{-}=0,
E_{\rm src}^{-}=0,
\Sigma_{\rm trans}=0,
P_{\partial}S_{4,\mu}^{\rm mixker}=0,
\beta_4={1\over256},
\Phi_5^{(n+1)}={3\over4}\Phi_5^{(n)},
\mathcal C_{\gamma}(\tau_v)\geq G_{\rm req},
\Gamma_{\rm guard}
\right].
}
\]
For the comfortable r29/r30 operational anchor
\[
(N,\kappa,G)=\left(5,{3\over4},1\right),
\]
one has
\[
S=20,
\qquad
H=20,
\qquad
\mu=H-16=4.
\]
Hence
\[
E_{\rm rel}^{-}=0,
\qquad
G_{\rm req}={4\over5},
\qquad
G_{\rm req}^{(+1)}={17\over20}.
\]
Since \(G=1\), the anchor satisfies the selector-window requirement with margin.
The frontier representative \((4,1,1)\) also satisfies the certificate but only
critically, with \(H=16\) and \(\mu=0\).  It is therefore kept as a frontier
witness, not as the operational anchor.

\begin{proposition}[r30 local admissibility certificate]
\label{prop:v470-r30-local-admissibility-certificate}
The comfortable anchor \((5,3/4,1)\) satisfies the typed MM/HC seed normal form
with margin \(4\).  The certificate has zero typed admissibility debt while
preserving all global guards.
\end{proposition}

\begin{proof}
For the anchor, \(S=20\) and \(H=20\), so
\(E_{\rm rel}^{-}=\max(0,16-20)=0\).  The boundary-source handoff has
\(H_{\partial}=20\), so \(E_{\rm src}^{-}=0\).  By the r26 transition audit,
\(\Sigma_{\rm trans}=0\).  By the r20--r22 boundary-source typing,
\(P_{\partial}S_{4,\mu}^{\rm mixker}=0\) and \(\beta_4=1/256\) is left unspent.
The collar term is the already guarded recurrence
\(\Phi_5^{(n+1)}=(3/4)\Phi_5^{(n)}\), which is not finite closure.  The clock
selector threshold is \(G_{\rm req}=16/20=4/5\), and the anchor has \(G=1\).
Thus every local admissibility component is satisfied while
\(\Gamma_{\rm guard}=(1,1,1,1,1)\) remains active.
\end{proof}

\subsection{Certificate classes}
\label{subsec:v470-r30-certificate-classes}

The r30 certificate separates admissible points into three useful local classes:
\[
\boxed{
\begin{array}{c|c|c}
\text{class} & \text{condition} & \text{interpretation}\\
\hline
\mathsf{frontier\mbox{-}cert} & H=16 & \text{zero debt but no margin}\\
\mathsf{robust\mbox{-}cert} & 17\leq H<20 & \text{positive protected margin}\\
\mathsf{comfortable\mbox{-}cert} & H\geq20 & \text{operational anchor region}
\end{array}
}
\]
The certificate deliberately avoids selecting the smallest zero-debt point as the
operational seed.  The point \((4,1,1)\) is minimal on the finite grid, but it is
thin.  The point \((5,3/4,1)\) carries margin four and is therefore retained as
the preferred local seed normal form.

\begin{corollary}[r30 frontier/anchor separation]
\label{cor:v470-r30-frontier-anchor-separation}
The r30 normal form distinguishes minimal admissibility from robust operational
admissibility.  A frontier-zero-debt point is allowed as a witness, but only a
positive-margin point is used as the seed anchor for subsequent MM/HC coefficient
work.
\end{corollary}

\subsection{r30 guard certificate}
\label{subsec:v470-r30-guard-certificate}

The r30 admissibility certificate is local and guard-preserving:
\[
\boxed{
\begin{array}{c|c|c}
\text{guard} & \text{r30 status} & \text{certificate effect}\\
\hline
C4 & \beta_4=1/256\text{ unspent} & \text{mixed-kernel guard preserved}\\
C5 & \text{no finite closure} & \text{collar recurrence only}\\
C6/T2 & \text{blocked} & \chi_{\rm all}=0\text{ remains active}\\
\text{RH} & \text{probe only} & \mathcal C_{\gamma}\text{ only calibrates selector windows}\\
\text{global all-readout} & \text{unclaimed} & \Pi_{\rm loc}^{\rm guard}\text{ remains mandatory}
\end{array}
}
\]
Therefore
\[
\boxed{
\Gamma_{\rm guard}^{r30}=(1,1,1,1,1).
}
\]

A safe next source-only direction is
\[
\boxed{
\text{v4.7.0-r31: typed coefficient skeleton and local seed functional packaging.}
}
\]
The next step should introduce a typed coefficient skeleton for the certified
seed normal form, still as a local audit functional and not as a physical action
or field equation.



% -----------------------------------------------------------------------------
\section{v4.7.0-r31 Source-Only: Typed Coefficient Skeleton and Local Seed Functional Packaging}
\label{sec:v470-r31-typed-coeff-skeleton-seed-functional}

This r31 source-only layer packages the r30 seed normal form as a typed local
seed functional.  The purpose is strictly audit-theoretic: the functional below
is a local MM/HC seed bookkeeping device, not a physical action, not a
Lagrangian, not a derived field equation, and not a closure theorem.

\subsection{Coefficient Skeleton}
\label{subsec:r31-coefficient-skeleton}

Let
\[
\mathfrak N_{\rm seed}^{r30}
=
(E_{\rm rel}^{-},E_{\rm src}^{-},\Sigma_{\rm trans},
P_{\partial}S_{4,\mu}^{\rm mixker},\beta_4,
\Phi_5^{(n)},G,\mathcal C_{\gamma}(\tau_v))
\]
be the local seed normal form certified in r30.  r31 assigns types to the
coefficient skeleton before any functional is written:
\[
\boxed{
\Lambda_{\rm seed}^{r31}
=
(\lambda_{\rm rel}^{-},\lambda_{\rm src}^{-},\lambda_{\rm seam},
\lambda_{\rm mix},\lambda_{\beta},\lambda_{\rm col},
\lambda_{\rm clk},\lambda_{\rm cmf}) .
}
\]
The typing is:
\[
\begin{array}{c|c|c}
\text{coefficient} & \text{type} & \text{allowed role}\\
\hline
\lambda_{\rm rel}^{-} & \text{admissibility penalty} & (E_{\rm rel}^{-})^2\\
\lambda_{\rm src}^{-} & \text{admissibility penalty} & (E_{\rm src}^{-})^2\\
\lambda_{\rm seam} & \text{hard seam guard} & \|\Sigma_{\rm trans}\|^2\\
\lambda_{\rm mix} & \text{mixed-kernel guard-lock} & \|P_{\partial}S_{4,\mu}^{\rm mixker}\|^2\\
\lambda_{\beta} & \text{C4 ledger lock} & (\beta_4-1/256)^2\\
\lambda_{\rm col} & \text{local recurrence bookkeeping} & (\Phi_5^{(n+1)}-\tfrac34\Phi_5^{(n)})^2\\
\lambda_{\rm clk} & \text{selector/probe coefficient} & (G-\mathcal C_{\gamma}(\tau_v))^2_{\rm sel}\\
\lambda_{\rm cmf} & \text{comfort-margin selector} & (E_{\rm cmf}^{-})^2
\end{array}
\]
where the last term is optional and only used when a comfortable certificate is
requested.  It is not part of the bare admissibility certificate.

\subsection{Local Seed Functional}
\label{subsec:r31-local-seed-functional}

The local audit functional is
\[
\boxed{
\begin{aligned}
\mathcal J_{\rm seed}^{r31}
={}&
\lambda_{\rm rel}^{-}(E_{\rm rel}^{-})^2
+\lambda_{\rm src}^{-}(E_{\rm src}^{-})^2
+\lambda_{\rm seam}\|\Sigma_{\rm trans}^{r26}\|^2\\
&+\lambda_{\rm mix}\|P_{\partial}S_{4,\mu}^{\rm mixker}\|^2
+\lambda_{\beta}\left(\beta_4-\frac1{256}\right)^2\\
&+\lambda_{\rm col}\left(\Phi_5^{(n+1)}-\frac34\Phi_5^{(n)}\right)^2
+\lambda_{\rm clk}\left(G-\mathcal C_{\gamma}(\tau_v)\right)^2_{\rm sel}.
\end{aligned}
}
\]
The signed monitor
\[
E_{\rm rel}=H-16
\]
is not squared in the admissibility functional.  Only the violation projection
\[
E_{\rm rel}^{-}=\max(0,16-H)
\]
contributes debt.  Hence interior relay surplus is not penalized.

For comfortable certification one may add the explicitly typed selector term
\[
\boxed{
\mathcal J_{{\rm seed},+4}^{r31}
=\mathcal J_{\rm seed}^{r31}
+\lambda_{\rm cmf}(E_{\rm cmf}^{-})^2,
\qquad
E_{\rm cmf}^{-}=\max(0,20-H).
}
\]
This term distinguishes the comfortable anchor from a thin frontier pass, but it
is not a new physical principle.

\subsection{Anchor Evaluation}
\label{subsec:r31-anchor-evaluation}

At the comfortable operational anchor
\[
(N,\kappa,G)=\left(5,\frac34,1\right)
\]
one has
\[
S=20,
\qquad
H=20,
\qquad
E_{\rm rel}^{-}=0,
\qquad
E_{\rm src}^{-}=0,
\qquad
E_{\rm cmf}^{-}=0.
\]
Together with the r30 certificate clauses
\[
\Sigma_{\rm trans}^{r26}=0,
\qquad
P_{\partial}S_{4,\mu}^{\rm mixker}=0,
\qquad
\beta_4=\frac1{256},
\qquad
\Phi_5^{(n+1)}=\frac34\Phi_5^{(n)},
\]
and a matched selector window
\[
G=\mathcal C_{\gamma}(\tau_v),
\]
this gives
\[
\boxed{
\mathcal J_{\rm seed}^{r31}=0,
\qquad
\mathcal J_{{\rm seed},+4}^{r31}=0
}
\]
inside the declared local seed grammar.

At the frontier witness
\[
(N,\kappa,G)=(4,1,1)
\]
one has
\[
H=16,
\qquad
E_{\rm rel}^{-}=0,
\]
so the bare admissibility functional may vanish.  However
\[
E_{\rm cmf}^{-}=4,
\]
so the comfortable certificate does not vanish.  Thus r31 preserves the r30
classification:
\[
\boxed{
(4,1,1)=\text{frontier-cert only},
\qquad
\left(5,\frac34,1\right)=\text{comfortable-cert}.
}
\]

\subsection{Guard Preservation}
\label{subsec:r31-guard-preservation}

The r31 functional is local and typed.  It does not spend the persistent C4
mixed-kernel guard, does not close C5 at any finite stage, does not unblock
C6/T2, does not identify Riemann zeros with photon events, and does not claim a
global all-readout closure.  Hence
\[
\boxed{
\Gamma_{\rm guard}^{r31}=(1,1,1,1,1).
}
\]
The safe next source-only direction is
\[
\boxed{
\text{v4.7.0-r32: local seed functional sanity table and coefficient-domain guard.}
}
\]



\section{v4.7.0-r32 Source-Only Local Seed Functional Sanity Table and Coefficient-Domain Guard}
\label{sec:v47-r32-seed-functional-sanity-domain}

\paragraph{Source-only status.}
This subsection is an internal source-only continuation of the v4.7.0 line.  No
PDF artefact is generated at this stage.  The purpose is to sanity-check the
r31 local seed functional against representative admissible, frontier, robust,
comfortable, and debt-positive sites, and to record an explicit coefficient
-domain guard.

\subsection{Sanity Table for the Local Seed Functional}
\label{subsec:r32-seed-sanity-table}

Recall the local relay score
\[
S(N,\kappa)=N+\kappa N(N-1),
\qquad
H=G\,S(N,\kappa),
\]
with admissibility debt
\[
E_{\rm rel}^{-}=\max(0,16-H)
\]
and optional comfort-margin selector
\[
E_{\rm cmf}^{-}=\max(0,20-H).
\]
The signed quantity
\[
E_{\rm rel}=H-16
\]
remains only a frontier monitor and is not squared as debt.

The representative sanity table is:
\[
\begin{array}{c|c|c|c|c|c|c}
\text{site} & N & \kappa & G & H & E_{\rm rel}^{-} & E_{\rm cmf}^{-}\\
\hline
\text{deep-debt} & 3 & 1/2 & 1 & 6 & 10 & 14\\
\text{local-debt} & 4 & 3/4 & 1 & 13 & 3 & 7\\
\text{frontier witness} & 4 & 1 & 1 & 16 & 0 & 4\\
\text{weak five-packet frontier} & 5 & 11/20 & 1 & 16 & 0 & 4\\
\text{gain frontier} & 5 & 3/4 & 4/5 & 16 & 0 & 4\\
\text{margin-one selector} & 5 & 3/4 & 17/20 & 17 & 0 & 3\\
\text{robust box corner} & 5 & 7/10 & 9/10 & 171/10 & 0 & 29/10\\
\text{comfortable anchor} & 5 & 3/4 & 1 & 20 & 0 & 0
\end{array}
\]
Hence the bare r31 admissibility part vanishes exactly on and above the relay
frontier \(H=16\), while the optional comfort selector distinguishes thin
frontier passes from the operational comfortable anchor \(H=20\).

\begin{lemma}[r32 local sanity separation]
\label{lem:r32-local-sanity-separation}
Inside the declared local seed grammar, the r31/r32 typed seed audit separates
three cases without changing the guard vector:
\[
H<16\Rightarrow E_{\rm rel}^{-}>0,
\qquad
16\leq H<20\Rightarrow E_{\rm rel}^{-}=0<E_{\rm cmf}^{-},
\qquad
H\geq20\Rightarrow E_{\rm rel}^{-}=E_{\rm cmf}^{-}=0.
\]
Consequently a frontier certificate, a robust local certificate, and a
comfortable certificate remain distinct.
\end{lemma}

\subsection{Coefficient-Domain Guard}
\label{subsec:r32-coefficient-domain-guard}

The admissible coefficient domain is recorded as
\[
\boxed{
\mathcal D_{\lambda}^{r32}
}
\]
with the following typed clauses:
\[
\begin{array}{ll}
\lambda_{\rm rel}^{-},\lambda_{\rm src}^{-}\geq0,&
\text{violation-projected admissibility penalties only},\\
\lambda_{\rm seam}>0,&
\text{hard transition-seam guard},\\
\lambda_{\rm mix}>0,&
\text{mixed-kernel guard-lock},\\
\lambda_{\beta}>0,&
\text{C4 ledger lock at }\beta_4=1/256,\\
\lambda_{\rm col}\geq0,&
\text{local C5 recurrence bookkeeping},\\
\lambda_{\rm clk}\geq0,&
\text{selector/probe weight only},\\
\lambda_{\rm cmf}\geq0,&
\text{optional comfort-margin selector}.
\end{array}
\]
No coefficient in \(\mathcal D_{\lambda}^{r32}\) may convert the signed monitor
\(E_{\rm rel}=H-16\) into a debt term.  Debt is carried only by the negative
part \(E_{\rm rel}^{-}\).  Likewise, no coefficient may spend the C4 ledger
lock, unblock C6/T2, or turn the photon-clock probe into an RH assertion.

The guarded local audit functional is therefore read as
\[
\boxed{
\mathcal J_{\rm seed}^{r32}
:
\mathcal D_{\lambda}^{r32}
\times
\mathfrak N_{\rm seed}^{r30}
\longrightarrow
\mathbb R_{\geq0}
}
\]
with the same meaning as in r31: it is a local audit functional, not an action,
not a Lagrangian, not a derived field equation, and not a closure theorem.

\subsection{Guard Preservation}
\label{subsec:r32-guard-preservation}

The r32 sanity table and coefficient-domain guard do not alter the previous
mathematics.  They prevent a misreading of the signed frontier monitor as a
least-squares field residual and preserve the operational separation between
frontier, robust, and comfortable relay certificates.  Thus
\[
\boxed{
\Gamma_{\rm guard}^{r32}=(1,1,1,1,1).
}
\]
Explicitly, C4 remains route-stable with
\[
\beta_4=\frac1{256},
\]
C5 remains an asymptotic bounded collar with no finite closure, C6/T2 remains
blocked, the Riemann-zero language remains a photon-clock handoff probe only,
and global all-readout closure is not claimed.


\section{v4.7.0-r33 Source-Only: Local Seed Functional Parameter Sweep and Comfort-Margin Calibration}
\label{sec:v470-r33-parameter-sweep-comfort-calibration}

This source-only layer performs the safe r33 follow-up to the r32 sanity table.
No PDF is generated in this step.  The purpose is to sweep the local seed
functional over a finite parameter box and to calibrate the optional comfort
margin without changing any guard, without turning the signed frontier monitor
into debt, and without promoting the seed audit into a physical action or field
equation.

\subsection{Finite Parameter Sweep}
\label{subsec:r33-finite-parameter-sweep}

The sweep uses
\[
N\in\{3,4,5,6,7,8\},
\]
\[
\kappa\in
\left\{
\frac12,\frac{11}{20},\frac35,\frac23,\frac7{10},\frac34,\frac45,\frac9{10},1
\right\},
\]
\[
G\in
\left\{
\frac34,\frac45,\frac{17}{20},\frac9{10},\frac{19}{20},1
\right\}.
\]
For each point we evaluate
\[
S(N,\kappa)=N+\kappa N(N-1),
\qquad
H=G\,S(N,\kappa),
\]
\[
E_{\rm rel}^{-}=\max(0,16-H),
\qquad
E_{\rm cmf}^{-}=\max(0,20-H).
\]
The classification is:
\[
\begin{array}{ll}
H<16 & \text{debt-positive},\\
H=16 & \text{frontier-zero-debt},\\
16<H<17 & \text{thin-zero-debt},\\
17\leq H<20 & \text{robust-zero-debt},\\
H\geq20 & \text{comfortable-zero-debt}.
\end{array}
\]
The finite sweep contains 324 cases and yields
\[
\begin{array}{c|c}
\text{class} & \#\text{ cases}\\
\hline
\text{debt positive} & 130\\
\text{frontier zero debt} & 3\\
\text{thin zero debt} & 6\\
\text{robust zero debt} & 24\\
\text{comfortable zero debt} & 161\\
\end{array}
\]
Thus the r32 typing behaves as intended: a broad zero-debt region exists above
relay threshold, but the optional comfort selector still separates thin and
robust admissibility from the operational comfort regime.

\subsection{Comfort-Margin Calibration}
\label{subsec:r33-comfort-margin-calibration}

For a fixed pair \((N,\kappa)\), the relay, margin-one, and comfort thresholds
are
\[
G_{\rm rel}(N,\kappa)=\frac{16}{S(N,\kappa)},
\qquad
G_{+1}(N,\kappa)=\frac{17}{S(N,\kappa)},
\qquad
G_{\rm cmf}(N,\kappa)=\frac{20}{S(N,\kappa)}.
\]
The five-packet calibration slice is
\[
\begin{array}{c|c|c|c|c}
\text{site} & N & \kappa & S(N,\kappa) & G_{\rm cmf}=20/S\\
\hline
\text{five-packet frontier} & 5 & \frac{11}{20} & 16 & \frac{5}{4}\\
\text{margin-one threshold} & 5 & \frac{3}{5} & 17 & \frac{20}{17}\\
\text{robust box corner} & 5 & \frac{7}{10} & 19 & \frac{20}{19}\\
\text{comfortable anchor} & 5 & \frac{3}{4} & 20 & 1\\
\end{array}
\]
Consequently, at the comfortable anchor
\[
(N,\kappa,G)=\left(5,\frac34,1\right)
\]
one has
\[
S=20,
\qquad
H=20,
\qquad
E_{\rm rel}^{-}=E_{\rm cmf}^{-}=0.
\]
The same five-packet source with \(G=\frac45\) sits exactly on the relay
frontier, while \(G=\frac{17}{20}\) gives the first unit of positive
margin.  Hence the comfort selector is not an admissibility guard; it is a
local preference for the stable operational anchor.

\begin{lemma}[r33 comfort-margin calibration]
\label{lem:r33-comfort-margin-calibration}
Inside the declared local seed grammar, the optional comfort-margin term
\(E_{\rm cmf}^{-}=\max(0,20-H)\) calibrates the local selector preference
without altering admissibility.  In particular,
\[
H\geq16\Rightarrow E_{\rm rel}^{-}=0,
\]
but only
\[
H\geq20\Rightarrow E_{\rm cmf}^{-}=0.
\]
Therefore a frontier certificate and a comfortable operational certificate are
kept distinct without creating new debt.
\end{lemma}

\subsection{Representative Sweep Anchors}
\label{subsec:r33-representative-sweep-anchors}

The parameter sweep confirms the hierarchy already suggested by r32:
\[
(4,1,1):\quad H=16,
\qquad
\mu=0,
\qquad
\text{frontier-cert only},
\]
\[
\left(5,\frac{11}{20},1\right):\quad H=16,
\qquad
\text{weak five-packet frontier},
\]
\[
\left(5,\frac34,\frac45\right):\quad H=16,
\qquad
\text{gain frontier},
\]
\[
\left(5,\frac34,\frac{17}{20}\right):\quad H=17,
\qquad
\text{margin-one selector},
\]
\[
\left(5,\frac7{10},\frac9{10}\right):\quad H=\frac{171}{10},
\qquad
\text{robust box corner},
\]
\[
\left(5,\frac34,1\right):\quad H=20,
\qquad
\text{comfortable anchor}.
\]
The frontier minimizer remains useful as a boundary witness, but the comfort
anchor remains the preferred local source point for downstream MM/HC seed
packaging.

\subsection{Guard Preservation}
\label{subsec:r33-guard-preservation}

The sweep and comfort calibration do not alter any global guard.  They only
classify local seed-functional behavior inside the already declared local
boundary-source grammar.  Thus
\[
\boxed{
\Gamma_{\rm guard}^{r33}=(1,1,1,1,1).
}
\]
Explicitly, C4 remains route-stable with
\[
\beta_4=\frac1{256},
\]
C5 remains an asymptotic bounded collar with no finite closure, C6/T2 remains
blocked, the Riemann-zero language remains a photon-clock handoff probe only,
and global all-readout closure is not claimed.

\paragraph{Next source-only direction.}
The safe next step is
\[
\boxed{
\text{v4.7.0-r34: comfort-calibrated seed selector and guard-stable handoff summary.}
}
\]



\section{v4.7.0-r34: Comfort-Calibrated Seed Selector and Guard-Stable Handoff Summary}
\label{sec:v470-r34-comfort-seed-selector-handoff}

This source-only step packages the r33 parameter sweep into a conservative
selector rule.  The purpose is not to introduce a new field equation, action,
Lagrangian, or closure theorem.  The purpose is only to choose a stable local
seed for downstream MM/HC bookkeeping while preserving every active guard.

\subsection{Selector Rule}
\label{subsec:r34-selector-rule}

Let
\[
S(N,\kappa)=N+\kappa N(N-1),
\qquad
H(N,\kappa,G)=G S(N,\kappa).
\]
The signed frontier monitor remains
\[
E_{\rm rel}=H-16,
\]
but the admissibility debt is still only
\[
E_{\rm rel}^{-}=\max(0,16-H).
\]
The comfort deficit is
\[
E_{\rm cmf}^{-}=\max(0,20-H).
\]
Thus the selector is not a minimizer of the signed frontier monitor.  It is a
filtered local choice rule:
\[
\boxed{
E_{\rm rel}^{-}=0,
\quad
E_{\rm src}^{-}=0,
\quad
\Gamma_{\rm guard}=(1,1,1,1,1),
}
\]
followed, for operational handoff purposes, by the comfort preference
\[
\boxed{
E_{\rm cmf}^{-}=0.
}
\]
Frontier witnesses remain useful boundary witnesses, but are not selected as
operational handoff anchors unless no comfortable seed is available.

\subsection{Representative Selector Table}
\label{subsec:r34-representative-selector-table}

The r34 selector keeps the known frontier witness visible but chooses the
comfortable anchor for downstream local MM/HC seed packaging:
\[
(4,1,1):\quad H=16,
\qquad
\text{frontier witness only},
\]
whereas
\[
\boxed{
\left(5,\frac34,1\right):\quad S=20,
\quad H=20,
\quad \mu=4
}
\]
is selected as the comfort-calibrated operational seed.  It satisfies
\[
E_{\rm rel}^{-}=0,
\qquad
E_{\rm cmf}^{-}=0.
\]
The robust lower corner
\[
\left(5,\frac7{10},\frac9{10}\right)
\]
remains a robustness certificate, since
\[
H\geq\frac{171}{10}>16,
\]
but it is not the preferred handoff anchor because the comfortable anchor gives
larger margin without changing the declared local grammar.

\subsection{Guard-Stable Handoff Object}
\label{subsec:r34-guard-stable-handoff-object}

The local handoff object produced by r34 is
\[
\boxed{
\mathfrak H_{\rm seed}^{r34}
=
\left(
\left(5,\frac34,1\right),
\mathfrak N_{\rm seed}^{r30},
\mathcal J_{\rm seed}^{r31},
\mathcal D_{\lambda}^{r32},
\Gamma_{\rm guard}^{r33}
\right).
}
\]
This object is a local seed handoff summary only.  It is not a physical field
equation and not a closure statement.

The guard vector remains
\[
\boxed{
\Gamma_{\rm guard}^{r34}=(1,1,1,1,1).
}
\]
Explicitly:
\[
\beta_4=\frac1{256}
\]
remains unspent; C5 remains an asymptotic bounded collar with no finite closure;
C6/T2 remains blocked; the Riemann-zero language remains only a photon-clock
handoff probe; and global all-readout closure remains unclaimed.

\subsection{Handoff Summary}
\label{subsec:r34-handoff-summary}

The r34 source-only conclusion is:
\[
\boxed{
\text{Use }\left(5,\frac34,1\right)\text{ as the comfort-calibrated local seed.}
}
\]
The finite local boundary-extended atlas and the typed seed functional may now
be handed to the next MM/HC stage as a guarded seed package.  The safe next
step is
\[
\boxed{
\text{v4.7.0-r35: guarded MM/HC seed-to-interface map and no-action audit.}
}
\]


\section{v4.7.0-r35 Source-Only: Guarded MM/HC Seed-to-Interface Map and No-Action Audit}
\label{sec:v470-r35-seed-to-interface-no-action}

\subsection{Purpose and Scope}
\label{subsec:r35-purpose-scope}

This source-only increment takes the r34 comfort-calibrated handoff object
\[
\mathfrak H_{\rm seed}^{r34}
=
\left(
\left(5,\frac34,1\right),
\mathfrak N_{\rm seed}^{r30},
\mathcal J_{\rm seed}^{r31},
\mathcal D_{\lambda}^{r32},
\Gamma_{\rm guard}^{r33}
\right)
\]
and maps it to a guarded MM/HC interface skeleton.  The purpose is only to
identify admissible interface slots and no-action obligations.  No physical
action, no Lagrangian, no derived field equation, no T2 closure, no RH
claim, and no global all-readout closure is asserted.

The selected seed remains
\[
(N,\kappa,G)=\left(5,\frac34,1\right),
\qquad
S=20,
\qquad
H=20,
\qquad
\mu=4.
\]
The frontier witness \((4,1,1)\) remains a boundary witness only and is not
promoted to the operational interface seed.

\subsection{Guarded Interface Slots}
\label{subsec:r35-guarded-interface-slots}

We introduce a local interface map
\[
\mathfrak I_{\rm MM/HC}^{r35}
:
\mathfrak H_{\rm seed}^{r34}
\longrightarrow
\left(
\mathsf I_{\rm bd},
\mathsf I_{\rm src},
\mathsf I_{\rm seam},
\mathsf I_{\rm mix},
\mathsf I_{\rm col},
\mathsf I_{\rm clk}
\right),
\]
where each component is a typed interface slot rather than an equation of
motion.

The slots are:
\[
\mathsf I_{\rm bd}:\quad
E_{\rm rel}^{-}=0,
\qquad
E_{\rm cmf}^{-}=0,
\]
for the comfort-calibrated boundary relay condition;
\[
\mathsf I_{\rm src}:\quad
{}^{\partial{\rm src}}{}_{\rm loc}\mathcal B_0^{\rm res}{}^{V0}_{\rm relay},
\]
for the two-sided local source/relay typing;
\[
\mathsf I_{\rm seam}:\quad
\Sigma_{\rm trans}=0,
\]
for the compressed atlas transition seam audit;
\[
\mathsf I_{\rm mix}:\quad
P_{\partial}S_{4,\mu}^{\rm mixker}=0,
\qquad
\beta_4=\frac1{256},
\]
for the C4 mixed-kernel guard-lock;
\[
\mathsf I_{\rm col}:\quad
\Phi_5^{(n+1)}=\frac34\Phi_5^{(n)},
\qquad
\forall n<\infty:\ \Phi_5^{(n)}>0,
\]
for the local C5 bounded-collar recurrence without finite closure; and
\[
\mathsf I_{\rm clk}:\quad
G=\mathcal C_{\gamma}(\tau_v)
\quad\hbox{or}\quad
\mathcal C_{\gamma}(\tau_v)\geq \frac{16}{S(N,\kappa)},
\]
for the photon-clock selector/probe interface.

Each slot is local and typed.  None of these slots may be read as a global
readout equation or as a physical field equation.

\subsection{No-Action Audit}
\label{subsec:r35-no-action-audit}

The interface map is equipped with a no-action audit vector
\[
\boxed{
\mathfrak A_{\rm noact}^{r35}
=
(0,0,0,0,0,0)
}
\]
with the following entries:
\[
\begin{array}{ccl}
A_1&=&\hbox{no physical action introduced},\\
A_2&=&\hbox{no Lagrangian introduced},\\
A_3&=&\hbox{no Euler--Lagrange variation asserted},\\
A_4&=&\hbox{no derived field equation asserted},\\
A_5&=&\hbox{no global readout promoted},\\
A_6&=&\hbox{no C4, C5, C6/T2, RH, or all-readout guard spent}.
\end{array}
\]
Thus r35 only records a guarded seed-to-interface map.  It is a preparation
layer for a possible later field-equation program, not the program itself.

\subsection{Admissibility Matrix}
\label{subsec:r35-admissibility-matrix}

For compact bookkeeping, the interface slots may be summarized by the local
admissibility matrix
\[
\mathcal M_{\rm int}^{r35}
=
\begin{pmatrix}
1&1&1&1&1&1\\
0&0&0&0&0&0
\end{pmatrix},
\]
where the first row records local slot admissibility and the second row records
new guard debt.  Therefore
\[
\boxed{
\operatorname{newDebt}(\mathcal M_{\rm int}^{r35})=0.
}
\]
This matrix is not an MM dynamics matrix.  It is only a local interface audit
matrix.

\subsection{Guard Preservation}
\label{subsec:r35-guard-preservation}

The guard vector remains
\[
\boxed{
\Gamma_{\rm guard}^{r35}=(1,1,1,1,1).
}
\]
Equivalently:
\[
\boxed{
\beta_4=\frac1{256}\quad\hbox{remains unspent.}
}
\]
C5 remains an asymptotic bounded collar with no finite closure; C6/T2 remains
blocked; the Riemann-zero component remains only a photon-clock handoff probe;
and global all-readout closure remains unclaimed.

\subsection{r35 Handoff}
\label{subsec:r35-handoff}

The r35 handoff object is
\[
\boxed{
\mathfrak H_{\rm int}^{r35}
=
\left(
\mathfrak H_{\rm seed}^{r34},
\mathfrak I_{\rm MM/HC}^{r35},
\mathfrak A_{\rm noact}^{r35},
\mathcal M_{\rm int}^{r35},
\Gamma_{\rm guard}^{r35}
\right).
}
\]
The safe next step is
\[
\boxed{
\text{v4.7.0-r36: interface residual table and MM/HC map consistency scan.}
}
\]
This next step should still remain source-only unless a PDF is explicitly
requested.


\section{v4.7.0-r36 Source-Only Interface Residual Table and MM/HC Map Consistency Scan}
\label{sec:v470-r36-interface-residual-consistency}

This source-only layer continues the r35 seed-to-interface map.  It does not
compile a PDF, does not introduce a physical action, and does not derive an
Euler--Lagrange or field equation.  Its only purpose is to test whether the six
local MM/HC interface slots from r35 remain mutually consistent when read as
residual channels.

The comfort-calibrated seed remains
\[
(N,\kappa,G)=\left(5,{3\over4},1\right),
\qquad
S=20,
\qquad
H=20,
\qquad
\mu=4.
\]
The frontier witness \((4,1,1)\) remains only a boundary witness and is not used
as the operational handoff seed.

\subsection{Interface Residual Table}
\label{subsec:r36-interface-residual-table}

For each r35 interface slot define a local residual channel
\[
\mathcal R_{\rm int}^{r36}
=
(R_{\rm bd},R_{\rm src},R_{\rm seam},R_{\rm mix},R_{\rm col},R_{\rm clk}).
\]
The table below records the guarded interpretation.
\[
\begin{array}{c|c|c|c}
\hbox{slot} & \hbox{residual channel} & \hbox{value} & \hbox{status}\\
\hline
\mathsf I_{\rm bd}
& R_{\rm bd}=E_{\rm rel}^{-}+E_{\rm src}^{-}
&0&\hbox{local relay admissible}\\
\mathsf I_{\rm src}
& R_{\rm src}=\delta_{\rm two\mbox{-}sided\ type}
&0&\hbox{typed local source/relay}\\
\mathsf I_{\rm seam}
& R_{\rm seam}=\|\Sigma_{\rm trans}^{r26}\|^2
&0&\hbox{transition seam exact}\\
\mathsf I_{\rm mix}
& R_{\rm mix}=\|P_{\partial}S_{4,\mu}^{\rm mixker}\|^2+
\left(\beta_4-{1\over256}\right)^2
&0&\hbox{C4 guard-lock preserved}\\
\mathsf I_{\rm col}
& R_{\rm col}=\Phi_5^{(n+1)}-{3\over4}\Phi_5^{(n)}
&0&\hbox{bounded-collar recurrence only}\\
\mathsf I_{\rm clk}
& R_{\rm clk}=\max\left(0,{16\over S(N,\kappa)}-\mathcal C_\gamma(\tau_v)\right)
&0&\hbox{selector/probe window admissible}
\end{array}
\]
For the comfort seed, \(S(5,3/4)=20\), hence
\[
{16\over S(N,\kappa)}={4\over5},
\qquad
\mathcal C_\gamma(\tau_v)=G=1,
\]
so the photon-clock selector residual vanishes as a local selector channel.
This is not an RH statement.  It is only the local clock-window reading already
allowed by the r35 handoff.

Therefore
\[
\boxed{
\mathcal R_{\rm int}^{r36}=(0,0,0,0,0,0).
}
\]

\subsection{Guard Residuals Versus Interface Residuals}
\label{subsec:r36-guard-versus-interface-residuals}

The vanishing of \(\mathcal R_{\rm int}^{r36}\) is not a global closure
statement.  It only means that the r35 local interface slots are internally
consistent.  The guard residuals remain logically separate:
\[
\mathcal B_{\rm guard}^{r36}
=
\left(
{1\over256},
\Phi_5^{(n)}>0\ (n<\infty),
\mathsf B_{\rm all/T2},
\mathsf{probe}_{\rm RH},
\mathsf{unclaimed}_{\rm all}
\right).
\]
Thus the interface residual scan has zero local mismatch while the global guard
profile remains active.

\subsection{MM/HC Map Consistency Scan}
\label{subsec:r36-mmhc-map-consistency-scan}

The r35 map is consistent iff the six local slots satisfy the following map
checks:
\[
\begin{array}{ll}
\mathrm{MC1}:&\mathsf I_{\rm bd}\hbox{ uses only }E_{\rm rel}^{-},E_{\rm src}^{-}
\hbox{ and does not square signed }E_{\rm rel},\\
\mathrm{MC2}:&\mathsf I_{\rm src}\hbox{ remains two-sided local and does not become an all-readout source},\\
\mathrm{MC3}:&\mathsf I_{\rm seam}\hbox{ imports no transition seam debt},\\
\mathrm{MC4}:&\mathsf I_{\rm mix}\hbox{ locks }\beta_4={1\over256}\hbox{ instead of spending it},\\
\mathrm{MC5}:&\mathsf I_{\rm col}\hbox{ records only the }C5\hbox{ recurrence and no finite closure},\\
\mathrm{MC6}:&\mathsf I_{\rm clk}\hbox{ remains a selector/probe and not an RH theorem}.
\end{array}
\]
All six checks pass in the source-level scan:
\[
\boxed{
\mathrm{MC}^{r36}=(1,1,1,1,1,1),
\qquad
\mathrm{mapDebt}^{r36}=0.
}
\]

\subsection{No-Action Preservation}
\label{subsec:r36-no-action-preservation}

The no-action audit vector remains zero:
\[
\boxed{
\mathfrak A_{\rm noact}^{r36}=(0,0,0,0,0,0).
}
\]
No variational principle is introduced.  No Lagrangian is declared.  No
Euler--Lagrange equation is varied.  No physical field equation is derived.  No
all-readout promotion is made.  No C4, C5, C6/T2, RH, or global-readout guard is
spent.

\subsection{r36 Handoff}
\label{subsec:r36-handoff}

The r36 handoff object is
\[
\boxed{
\mathfrak H_{\rm int}^{r36}
=
\left(
\mathfrak H_{\rm int}^{r35},
\mathcal R_{\rm int}^{r36},
\mathrm{MC}^{r36},
\mathfrak A_{\rm noact}^{r36},
\Gamma_{\rm guard}^{r35}
\right).
}
\]
It is a guard-stable local interface consistency certificate and nothing more.
The safe next source-only step is
\[
\boxed{
\text{v4.7.0-r37: local MM/HC interface certificate and guard-stable residual packaging.}
}
\]


\section{v4.7.0-r37 Local MM/HC Interface Certificate and Guard-Stable Residual Packaging}
\label{sec:r37-local-interface-certificate}

This r37 source-only layer packages the r36 zero-residual interface scan as a
local certificate.  The purpose is deliberately narrow: the certificate records
that the six declared local MM/HC interface slots close their local residuals
inside the current guarded grammar.  It does not convert the local interface
map into a physical action, a Lagrangian, an Euler--Lagrange system, a derived
field equation, a T2 theorem, or an all-readout closure statement.

\subsection{Certificate Object}
\label{subsec:r37-certificate-object}

Let the ordered local interface list be
\[
\mathbb I^{r37}
=
(\mathsf I_{\rm bd},\mathsf I_{\rm src},\mathsf I_{\rm seam},
 \mathsf I_{\rm mix},\mathsf I_{\rm col},\mathsf I_{\rm clk}).
\]
The r36 local residual vector is inherited as
\[
\mathcal R_{\rm int}^{r36}=(0,0,0,0,0,0),
\]
and the r36 MM/HC map consistency vector is inherited as
\[
\mathrm{MC}^{r36}=(1,1,1,1,1,1),
\qquad
\mathrm{mapDebt}^{r36}=0.
\]
The r37 local interface certificate is the tuple
\[
\boxed{
\mathfrak C_{\rm int}^{r37}
=
\left(
\mathbb I^{r37},
\mathcal R_{\rm int}^{r36},
\mathrm{MC}^{r36},
\mathfrak A_{\rm noact}^{r36},
\Gamma_{\rm guard}^{r36}
\right).
}
\]
It is valid only as a local interface certificate.  The inclusion of
\(\Gamma_{\rm guard}^{r36}\) is part of the certificate, not a debt-spending
operation.

\subsection{Slotwise Certificate Ledger}
\label{subsec:r37-slotwise-certificate-ledger}

The slotwise certificate ledger is
\[
\begin{array}{c|c|c|c}
\hbox{slot} & \hbox{local residual} & \hbox{certificate status} & \hbox{guard status}\\
\hline
\mathsf{I_{\rm bd}} & 0 & \hbox{boundary relay certified} & \hbox{local only}\\
\mathsf{I_{\rm src}} & 0 & \hbox{two-sided source/relay typed} & \hbox{no all-readout source}\\
\mathsf{I_{\rm seam}} & 0 & \hbox{transition seam exact} & \hbox{no seam debt}\\
\mathsf{I_{\rm mix}} & 0 & \hbox{mixed-kernel guard-lock certified} & \beta_4=1/256\hbox{ unspent}\\
\mathsf{I_{\rm col}} & 0 & \hbox{bounded-collar recurrence certified} & \hbox{no finite }C5\hbox{ closure}\\
\mathsf{I_{\rm clk}} & 0 & \hbox{photon-clock selector/probe certified} & \hbox{no RH theorem}
\end{array}
\]
Thus every declared local interface slot has zero local residual, while the
interpretive guard attached to that slot remains active.

\subsection{Guard-Stable Residual Packaging}
\label{subsec:r37-guard-stable-residual-packaging}

The r37 residual package separates the local zero-residual layer from the
nonlocal guard layer:
\[
\boxed{
\mathcal P_{\rm res}^{r37}
=
\mathcal P_{\rm loc}^{r37}
\oplus
\mathcal P_{\rm guard}^{r37}.
}
\]
The local component is
\[
\mathcal P_{\rm loc}^{r37}
=
\left(
\mathcal R_{\rm int}^{r36},
\mathrm{MC}^{r36},
\mathrm{mapDebt}^{r36}
\right)
=
\left((0,0,0,0,0,0),(1,1,1,1,1,1),0\right).
\]
The guard component is
\[
\mathcal P_{\rm guard}^{r37}
=
\left(
\beta_4={1\over256},
\Phi_5^{(n)}>0\ (n<\infty),
\mathsf B_{\rm all/T2},
\mathsf{probe}_{\rm RH},
\mathsf{unclaimed}_{\rm all}
\right).
\]
This packaging is the main r37 safety point: zero local residual does not imply
zero global guard.  The two components are recorded together but not identified.

\subsection{Certificate Consequence}
\label{subsec:r37-certificate-consequence}

The only certified consequence is the local statement
\[
\boxed{
\mathfrak C_{\rm int}^{r37}
\vdash_{\rm loc}
\mathcal R_{\rm int}^{r37}=(0,0,0,0,0,0),
\qquad
\mathrm{mapDebt}^{r37}=0.
}
\]
The guard vector remains
\[
\boxed{
\Gamma_{\rm guard}^{r37}=(1,1,1,1,1).
}
\]
Hence the following remain unproved and unclaimed:
\[
\boxed{
C4\hbox{ closure},\quad
C5\hbox{ finite closure},\quad
C6/T2,\quad
\mathrm{RH},\quad
\hbox{global all-readout closure}.
}
\]
In particular, the r37 certificate preserves
\[
\boxed{
\beta_4={1\over256}\quad\hbox{as an unspent route-stable mixed-kernel guard.}
}
\]

\subsection{r37 Handoff}
\label{subsec:r37-handoff}

The r37 handoff object is
\[
\boxed{
\mathfrak H_{\rm cert}^{r37}
=
\left(
\mathfrak C_{\rm int}^{r37},
\mathcal P_{\rm res}^{r37},
\Gamma_{\rm guard}^{r37}
\right).
}
\]
The safe next source-only step is
\[
\boxed{
\text{v4.7.0-r38: certificate-to-local-field-interface precheck and no-promotion audit.}
}
\]


\section{v4.7.0-r38 Certificate-to-Local-Field-Interface Precheck and No-Promotion Audit}
\label{sec:v470-r38-certificate-field-interface-precheck}

\subsection{Purpose of the r38 Precheck}
\label{subsec:r38-purpose}

The r37 result produced a local MM/HC interface certificate
\(\mathfrak C_{\rm int}^{r37}\) with zero local interface residual, zero map
Debt, and an explicitly active guard vector.  The r38 step does not promote this
certificate to a field equation.  Its purpose is narrower: to test whether the
certificate can be used as the input object for a local field-interface layer.
Thus r38 introduces a pre-field-interface object, not a physical field law.

We write the r38 candidate as
\[
\boxed{
\mathfrak F_{\rm loc}^{r38}
:=
\Pi_{\rm guard}^{\rm loc}
\left(
\mathfrak C_{\rm int}^{r37},
\mathcal P_{\rm loc}^{r37},
\mathcal P_{\rm guard}^{r37}
\right).
}
\]
The projector \(\Pi_{\rm guard}^{\rm loc}\) means that only the declared local
interface data may be passed forward.  The guard package is carried along as a
read-only obstruction ledger.

\subsection{Local Field-Interface Precheck}
\label{subsec:r38-local-field-interface-precheck}

The precheck consists of six clauses:
\[
\begin{array}{ll}
\mathrm{FIP1} & \mathfrak C_{\rm int}^{r37}\hbox{ is a declared local certificate},\\
\mathrm{FIP2} & \mathcal R_{\rm int}^{r37}=(0,0,0,0,0,0),\\
\mathrm{FIP3} & \mathrm{MC}^{r37}=(1,1,1,1,1,1),\\
\mathrm{FIP4} & \mathrm{mapDebt}^{r37}=0,\\
\mathrm{FIP5} & \mathfrak A_{\rm noact}^{r37}=(0,0,0,0,0,0),\\
\mathrm{FIP6} & \Gamma_{\rm guard}^{r37}=(1,1,1,1,1)\hbox{ remains active}.
\end{array}
\]
The corresponding precheck vector is
\[
\boxed{
\mathrm{FIP}^{r38}=(1,1,1,1,1,1).
}
\]
This says only that the certificate is syntactically and residually admissible
as a local field-interface input.  It does not say that a field equation has
been derived.

\subsection{Pre-Field Residual Ledger}
\label{subsec:r38-pre-field-residual-ledger}

For bookkeeping, we introduce the pre-field residual package
\[
\boxed{
\mathcal R_{\rm prefield}^{r38}
=
\left(
R_{\rm cert},R_{\rm int},R_{\rm map},R_{\rm seam},R_{\rm action},R_{\rm guard}
\right).
}
\]
The r38 values are
\[
\boxed{
\mathcal R_{\rm prefield}^{r38}=(0,0,0,0,0,0).
}
\]
The last zero, \(R_{\rm guard}=0\), means that the guard ledger was not changed.
It does not mean that the guarded problems vanished.

\subsection{No-Promotion Audit}
\label{subsec:r38-no-promotion-audit}

The precheck is accompanied by an explicit no-promotion audit.  The audit vector
records whether any forbidden promotion occurred:
\[
\boxed{
\Pi_{\rm np}^{r38}
=
\left(
\pi_{\rm action},
\pi_{\rm EL},
\pi_{\rm field},
\pi_{C4},
\pi_{C5},
\pi_{C6/T2},
\pi_{\rm all}
\right).
}
\]
Here the entries mean, respectively: promotion to a physical action, promotion
to Euler--Lagrange variation, promotion to a derived field equation, spending of
the \(C4\) guard, finite \(C5\) closure, \(C6/T2\) promotion, and global
all-readout promotion.

The r38 audit result is
\[
\boxed{
\Pi_{\rm np}^{r38}=(0,0,0,0,0,0,0).
}
\]
Thus no forbidden promotion occurred.  In particular,
\[
\boxed{
\mathfrak F_{\rm loc}^{r38}
\not\Rightarrow
\hbox{physical field equation},
\quad
\mathfrak F_{\rm loc}^{r38}
\not\Rightarrow
T2,
\quad
\mathfrak F_{\rm loc}^{r38}
\not\Rightarrow
\mathrm{RH}.
}
\]

\subsection{Guard-Preserving Consequence}
\label{subsec:r38-guard-preserving-consequence}

The only r38 consequence is the local pre-field-interface statement
\[
\boxed{
\mathfrak C_{\rm int}^{r37}
\vdash_{\rm loc,pre}
\mathfrak F_{\rm loc}^{r38}
\quad\hbox{with}\quad
\mathcal R_{\rm prefield}^{r38}=0.
}
\]
All global guards remain unchanged:
\[
\boxed{
\Gamma_{\rm guard}^{r38}=(1,1,1,1,1).
}
\]
Equivalently,
\[
\boxed{
\beta_4={1\over256}\hbox{ remains unspent},
\quad
\Phi_5^{(n)}>0\ (n<\infty),
\quad
C6/T2\hbox{ remains blocked},
\quad
\mathrm{RH}\hbox{ remains probe-only},
\quad
\hbox{global all-readout closure remains unclaimed}.
}
\]

\subsection{r38 Handoff}
\label{subsec:r38-handoff}

The r38 handoff object is
\[
\boxed{
\mathfrak H_{\rm prefield}^{r38}
=
\left(
\mathfrak F_{\rm loc}^{r38},
\mathrm{FIP}^{r38},
\mathcal R_{\rm prefield}^{r38},
\Pi_{\rm np}^{r38},
\Gamma_{\rm guard}^{r38}
\right).
}
\]
The safe next source-only step is
\[
\boxed{
\text{v4.7.0-r39: local field-interface skeleton and admissibility-only operator tags.}
}
\]



\section{v4.7.0-r39 Local Field-Interface Skeleton and Admissibility-Only Operator Tags}
\label{sec:v470-r39-local-field-interface-skeleton}

\subsection{Purpose of the r39 Skeleton}
\label{subsec:r39-purpose}

The r38 result admitted the local certificate only as a pre-field-interface
object
\[
\mathfrak F_{\rm loc}^{r38}.
\]
The r39 step introduces a local field-interface skeleton.  This is not a field
equation, not an action principle, and not a dynamical law.  It is only a typed
interface grammar which records which admissibility channels may later be used
when a true MM/HC field-equation program is attempted.

We define the r39 skeleton by
\[
\boxed{
\mathfrak S_{\rm loc}^{r39}
=
\left(
\mathfrak F_{\rm loc}^{r38},
\mathbb O_{\rm adm}^{r39},
\mathsf T_{\rm adm}^{r39},
\mathcal R_{\rm skel}^{r39},
\Gamma_{\rm guard}^{r39}
\right).
}
\]
The symbol \(\mathbb O_{\rm adm}^{r39}\) denotes admissibility-only operator
symbols.  They are tags for interface tests, not physical operators.

\subsection{Admissibility-Only Operator Tags}
\label{subsec:r39-admissibility-only-operator-tags}

The r39 local skeleton carries six operator tags:
\[
\boxed{
\mathbb O_{\rm adm}^{r39}
=
\left(
\mathcal O_{\rm bd},
\mathcal O_{\rm src},
\mathcal O_{\rm seam},
\mathcal O_{\rm mix},
\mathcal O_{\rm col},
\mathcal O_{\rm clk}
\right).
}
\]
Their intended local meanings are:
\[
\begin{array}{lll}
\mathcal O_{\rm bd}   &:& \hbox{boundary relay admissibility monitor},\\
\mathcal O_{\rm src}  &:& \hbox{two-sided source/relay typing monitor},\\
\mathcal O_{\rm seam} &:& \hbox{compressed-atlas seam exactness monitor},\\
\mathcal O_{\rm mix}  &:& C4\hbox{ mixed-kernel guard-lock monitor},\\
\mathcal O_{\rm col}  &:& C5\hbox{ bounded-collar recurrence monitor},\\
\mathcal O_{\rm clk}  &:& \hbox{photon-clock selector/probe monitor}.
\end{array}
\]
Each tag is constrained by the declaration
\[
\boxed{
\mathcal O_i\in\mathrm{Op}_{\rm adm}^{\rm loc}
\quad\Longrightarrow\quad
\mathcal O_i\notin
\left\{
\mathrm{physical\ action},
\mathrm{Euler\mbox{-}Lagrange\ operator},
\mathrm{derived\ field\ operator},
\mathrm{global\ readout\ operator}
\right\}.
}
\]
Thus the operator tags are grammatical handles for admissibility tests only.
They do not generate equations of motion.

\subsection{Two-Sided Tag Typing}
\label{subsec:r39-two-sided-tag-typing}

The tags are recorded in the two-sided structural notation as local
source/readout-facing tags on the left and local relay/interface-facing tags on
the right:
\[
\boxed{
{}^{\partial{\rm src}}{}_{\rm loc}
\mathcal O_{\rm bd}
{}^{V0}_{\rm relay},
\qquad
{}^{{\rm src}}{}_{\rm loc}
\mathcal O_{\rm src}
{}^{{\rm type}}_{\rm relay},
}
\]
\[
\boxed{
{}^{{\rm atlas}}{}_{\rm loc}
\mathcal O_{\rm seam}
{}^{{\rm seam}}_{\rm trans},
\qquad
{}^{C4}{}_{\rm guard}
\mathcal O_{\rm mix}
{}^{{\rm mixker}}_{\rm lock},
}
\]
\[
\boxed{
{}^{C5}{}_{\rm collar}
\mathcal O_{\rm col}
{}^{{\rm rec}}_{\rm loc},
\qquad
{}^{\gamma}{}_{\rm probe}
\mathcal O_{\rm clk}
{}^{{\rm selector}}_{\rm loc}.
}
\]
This typing preserves the r38 rule: the local interface can be prepared, but it
cannot spend a guard or claim a global readout.

\subsection{Skeleton Residual Vector}
\label{subsec:r39-skeleton-residual-vector}

For the r39 skeleton we define
\[
\boxed{
\mathcal R_{\rm skel}^{r39}
=
\left(
R_{\rm bd},R_{\rm src},R_{\rm seam},R_{\rm mix},R_{\rm col},R_{\rm clk}
\right).
}
\]
Using the r36--r38 local interface certificate data, the skeleton residual is
\[
\boxed{
\mathcal R_{\rm skel}^{r39}=(0,0,0,0,0,0).
}
\]
This means only that the local admissibility tags are internally consistent with
the pre-field-interface object.  It does not mean that a field equation has been
found.

\subsection{Admissibility Tag Vector}
\label{subsec:r39-admissibility-tag-vector}

The tag admissibility vector is
\[
\boxed{
\mathsf T_{\rm adm}^{r39}
=
(t_{\rm bd},t_{\rm src},t_{\rm seam},t_{\rm mix},t_{\rm col},t_{\rm clk}).
}
\]
The r39 values are
\[
\boxed{
\mathsf T_{\rm adm}^{r39}=(1,1,1,1,1,1).
}
\]
The entries say that the tags are allowed as local admissibility monitors.  They
are not allowed as physical operators.

\subsection{No-Action and No-Promotion Guard}
\label{subsec:r39-no-action-no-promotion-guard}

The no-action status from r38 is preserved:
\[
\boxed{
\mathfrak A_{\rm noact}^{r39}=(0,0,0,0,0,0).
}
\]
The no-promotion status is also preserved:
\[
\boxed{
\Pi_{\rm np}^{r39}=(0,0,0,0,0,0,0).
}
\]
Thus r39 does not introduce a physical action, Lagrangian, Euler--Lagrange
variation, derived field equation, global readout promotion, or guard spending.

\subsection{Guard-Preserving Consequence}
\label{subsec:r39-guard-preserving-consequence}

The guard vector remains fully active:
\[
\boxed{
\Gamma_{\rm guard}^{r39}=(1,1,1,1,1).
}
\]
Equivalently,
\[
\boxed{
\beta_4={1\over256}\hbox{ remains unspent},
\quad
C5\hbox{ has no finite closure},
\quad
C6/T2\hbox{ remains blocked},
\quad
\mathrm{RH}\hbox{ remains photon-clock probe only},
\quad
\hbox{global all-readout closure remains unclaimed}.
}
\]

\subsection{r39 Handoff}
\label{subsec:r39-handoff}

The r39 handoff object is
\[
\boxed{
\mathfrak H_{\rm skel}^{r39}
=
\left(
\mathfrak S_{\rm loc}^{r39},
\mathbb O_{\rm adm}^{r39},
\mathsf T_{\rm adm}^{r39},
\mathcal R_{\rm skel}^{r39},
\Pi_{\rm np}^{r39},
\Gamma_{\rm guard}^{r39}
\right).
}
\]
The safe next source-only step is
\[
\boxed{
\text{v4.7.0-r40: admissibility-only operator composition table and guard leak scan.}
}
\]



\section{v4.7.0-r40: Admissibility-Only Operator Composition Table and Guard Leak Scan}
\label{sec:v4-7-0-r40-operator-composition-guard-leak}

\subsection{Purpose and Scope}
\label{subsec:r40-purpose-scope}

This source-only step composes the r39 admissibility-only operator tags and scans
for guard leakage.  The point is deliberately narrow.  The tags
\[
\mathbb O_{\rm adm}^{r39}
=
\{
\mathcal O_{\rm bd},
\mathcal O_{\rm src},
\mathcal O_{\rm seam},
\mathcal O_{\rm mix},
\mathcal O_{\rm col},
\mathcal O_{\rm clk}
\}
\]
are allowed to compose only as local admissibility monitors.  They are not
promoted into physical operators, action operators, Lagrangian terms,
Euler--Lagrange operators, derived field equations, or global readout operators.

The r40 rule is therefore:
\[
\boxed{
\mathcal O_j\circ\mathcal O_i
\quad\hbox{means ordered admissibility monitoring only.}
}
\]
It never means an operator product in a physical dynamics unless a later version
adds a separate, explicitly guarded promotion theorem.  No such promotion is made
in r40.

\subsection{Composition Classes}
\label{subsec:r40-composition-classes}

For two admissibility-only tags we write
\[
\boxed{
\mathcal C_{ji}^{r40}
=
(\mathcal O_j\circ\mathcal O_i)_{\rm adm}.
}
\]
The composition is assigned one of three local monitor classes:
\[
\boxed{
\mathsf A_{\rm loc},\qquad
\mathsf G_{\rm lock},
\qquad
\mathsf P_{\rm probe}.
}
\]
Here \(\mathsf A_{\rm loc}\) means ordinary local admissibility monitoring,
\(\mathsf G_{\rm lock}\) means admissibility monitoring under an explicit guard
lock, and \(\mathsf P_{\rm probe}\) means selector/probe monitoring with no
claim of closure.

The class assignment used in r40 is:
\[
\begin{array}{c|cccccc}
\circ & \mathcal O_{\rm bd} & \mathcal O_{\rm src} & \mathcal O_{\rm seam} & \mathcal O_{\rm mix} & \mathcal O_{\rm col} & \mathcal O_{\rm clk}\\
\hline
\mathcal O_{\rm bd}   & \mathsf A_{\rm loc} & \mathsf A_{\rm loc} & \mathsf A_{\rm loc} & \mathsf G_{\rm lock} & \mathsf A_{\rm loc} & \mathsf P_{\rm probe}\\
\mathcal O_{\rm src}  & \mathsf A_{\rm loc} & \mathsf A_{\rm loc} & \mathsf A_{\rm loc} & \mathsf G_{\rm lock} & \mathsf A_{\rm loc} & \mathsf P_{\rm probe}\\
\mathcal O_{\rm seam} & \mathsf A_{\rm loc} & \mathsf A_{\rm loc} & \mathsf A_{\rm loc} & \mathsf G_{\rm lock} & \mathsf A_{\rm loc} & \mathsf P_{\rm probe}\\
\mathcal O_{\rm mix}  & \mathsf G_{\rm lock} & \mathsf G_{\rm lock} & \mathsf G_{\rm lock} & \mathsf G_{\rm lock} & \mathsf G_{\rm lock} & \mathsf G_{\rm lock}\\
\mathcal O_{\rm col}  & \mathsf A_{\rm loc} & \mathsf A_{\rm loc} & \mathsf A_{\rm loc} & \mathsf G_{\rm lock} & \mathsf A_{\rm loc} & \mathsf P_{\rm probe}\\
\mathcal O_{\rm clk}  & \mathsf P_{\rm probe} & \mathsf P_{\rm probe} & \mathsf P_{\rm probe} & \mathsf G_{\rm lock} & \mathsf P_{\rm probe} & \mathsf P_{\rm probe}
\end{array}
\]
The row tag is applied after the column tag.  This table is not a multiplication
law for physical observables.  It is a local monitor-composition table.

\subsection{Guard Requirements for Composed Tags}
\label{subsec:r40-guard-requirements-composed-tags}

The guard requirement attached to a composition is:
\[
\boxed{
\mathsf{Req}(\mathcal C_{ji}^{r40})
\in
\{\mathrm{local},\mathrm{guard\mbox{-}locked},\mathrm{probe\mbox{-}only}\}.
}
\]
The requirements are read as follows:
\begin{itemize}
    \item \(\mathrm{local}\): the composition is an ordinary local admissibility
    monitor and carries no guard-spending power;
    \item \(\mathrm{guard\mbox{-}locked}\): the composition may mention the
    mixed-kernel channel but must preserve \(\beta_4=1/256\) and may not absorb
    \(C4\);
    \item \(\mathrm{probe\mbox{-}only}\): the composition may use the photon-clock
    selector as a timing probe but may not assert RH, global spectral closure, or
    all-readout closure.
\end{itemize}

Thus every composed tag remains in the local admissibility grammar:
\[
\boxed{
\mathcal C_{ji}^{r40}\in\mathsf{Adm}_{\rm loc}^{r40}
\quad\hbox{for all declared r40 tag pairs.}
}
\]

\subsection{Guard Leak Scan}
\label{subsec:r40-guard-leak-scan}

For each composition we scan the possible leak channels
\[
\boxed{
\ell^{r40}_{ji}
=
(\ell_{\rm act},\ell_{\rm Lag},\ell_{\rm EL},\ell_{\rm fld},
\ell_{\rm glob},\ell_{C4},\ell_{C5},\ell_{C6},\ell_{\rm RH},\ell_{\rm all}).
}
\]
These entries respectively monitor accidental promotion to physical action,
Lagrangian, Euler--Lagrange variation, derived field equation, global readout,
\(C4\)-spending, finite \(C5\)-closure, \(C6/T2\)-unblocking, RH assertion, and
all-readout closure.

The r40 scan gives
\[
\boxed{
\ell^{r40}_{ji}=(0,0,0,0,0,0,0,0,0,0)
\quad\hbox{for every declared tag pair.}
}
\]
Equivalently, the total leak vector is
\[
\boxed{
\mathcal L_{\rm total}^{r40}=(0,0,0,0,0,0,0,0,0,0).
}
\]

\subsection{Composition Residual and Admissibility Vectors}
\label{subsec:r40-composition-residual-admissibility-vectors}

The composition residual vector is ordered as
\[
\boxed{
\mathcal R_{\rm comp}^{r40}
=
(R_{\rm bd},R_{\rm src},R_{\rm seam},R_{\rm mix},R_{\rm col},R_{\rm clk}).
}
\]
The r40 result is
\[
\boxed{
\mathcal R_{\rm comp}^{r40}=(0,0,0,0,0,0).
}
\]
The tag-composition admissibility vector is
\[
\boxed{
\mathsf T_{\rm comp}^{r40}=(1,1,1,1,1,1).
}
\]
This says only that the composed monitors are well-typed local admissibility
monitors.  It does not assert that they generate equations of motion.

\subsection{No-Action and No-Promotion Audit}
\label{subsec:r40-no-action-no-promotion-audit}

The no-action vector remains
\[
\boxed{
\mathfrak A_{\rm noact}^{r40}=(0,0,0,0,0,0).
}
\]
The no-promotion vector remains
\[
\boxed{
\Pi_{\rm np}^{r40}=(0,0,0,0,0,0,0).
}
\]
Thus r40 introduces no physical action, no Lagrangian, no Euler--Lagrange
variation, no derived field equation, no global readout operator, no guard
spending, and no all-readout promotion.

\subsection{Guard-Preserving Consequence}
\label{subsec:r40-guard-preserving-consequence}

The guard vector is unchanged:
\[
\boxed{
\Gamma_{\rm guard}^{r40}=(1,1,1,1,1).
}
\]
Equivalently,
\[
\boxed{
\beta_4={1\over256}\hbox{ remains unspent},
\quad
C5\hbox{ has no finite closure},
\quad
C6/T2\hbox{ remains blocked},
\quad
\mathrm{RH}\hbox{ remains photon-clock probe only},
\quad
\hbox{global all-readout closure remains unclaimed}.
}
\]

\subsection{r40 Handoff}
\label{subsec:r40-handoff}

The r40 handoff object is
\[
\boxed{
\mathfrak H_{\rm comp}^{r40}
=
\left(
\mathfrak S_{\rm loc}^{r39},
\mathcal C^{r40},
\mathcal R_{\rm comp}^{r40},
\mathcal L_{\rm total}^{r40},
\mathfrak A_{\rm noact}^{r40},
\Pi_{\rm np}^{r40},
\Gamma_{\rm guard}^{r40}
\right).
}
\]
The safe next source-only step is
\[
\boxed{
\text{v4.7.0-r41: guarded composition normal form and local interface closure ledger.}
}
\]





\section{v4.7.0-r41 Source-Only: Guarded Composition Normal Form and Local Interface Closure Ledger}
\label{sec:v470-r41-guarded-composition-normal-form}

This source-only checkpoint packages the r40 admissibility-only composition scan
into a guarded normal form and records exactly which interface closures are local
monitor closures.  No PDF is generated in this intermediate checkpoint.

\subsection{Scope and Non-Promotion Clause}
\label{subsec:r41-scope-non-promotion-clause}

The r41 operation is purely administrative and local.  It acts on the r40 table
of admissibility-only tags
\[
\mathbb O_{\rm adm}^{r39}
=
\{\mathcal O_{\rm bd},\mathcal O_{\rm src},\mathcal O_{\rm seam},
\mathcal O_{\rm mix},\mathcal O_{\rm col},\mathcal O_{\rm clk}\}.
\]
It does not turn those tags into physical operators.  In particular, r41 does
not introduce an action, Lagrangian, Euler--Lagrange variation, physical field
equation, global readout operator, or all-readout promotion.

\subsection{Guarded Composition Normal Form}
\label{subsec:r41-guarded-composition-normal-form}

Every ordered r40 tag pair is now normalized into one of three guarded normal
forms:
\[
\boxed{
\mathsf{NF}_{\rm loc},\qquad
\mathsf{NF}_{\rm guard},\qquad
\mathsf{NF}_{\rm probe}.
}
\]
The interpretation is:
\[
\mathsf{NF}_{\rm loc}:
\hbox{ordinary local admissibility-monitor composition},
\]
\[
\mathsf{NF}_{\rm guard}:
\hbox{composition allowed only with the active guard-lock retained},
\]
\[
\mathsf{NF}_{\rm probe}:
\hbox{composition allowed only as a selector/probe monitor.}
\]
Thus the r40 count
\[
36=16+11+9
\]
becomes the r41 normal-form count
\[
\boxed{
\#\mathsf{NF}_{\rm loc}=16,
\quad
\#\mathsf{NF}_{\rm guard}=11,
\quad
\#\mathsf{NF}_{\rm probe}=9.
}
\]

The normal-form map is denoted
\[
\boxed{
\mathfrak n^{r41}:
\mathbb O_{\rm adm}^{r39}\times\mathbb O_{\rm adm}^{r39}
\longrightarrow
\{\mathsf{NF}_{\rm loc},\mathsf{NF}_{\rm guard},\mathsf{NF}_{\rm probe}\}.
}
\]
This map is a compression of the composition table only; it carries no dynamic
content.

\subsection{Local Interface Closure Ledger}
\label{subsec:r41-local-interface-closure-ledger}

The local interfaces are ordered as
\[
(\mathsf I_{\rm bd},\mathsf I_{\rm src},\mathsf I_{\rm seam},
\mathsf I_{\rm mix},\mathsf I_{\rm col},\mathsf I_{\rm clk}).
\]
The r41 local interface-closure vector is
\[
\boxed{
\mathsf C_{\rm loc}^{r41}=(1,1,1,1,1,1).
}
\]
The entries mean:
\begin{itemize}
    \item \(\mathsf I_{\rm bd}\): closed as local boundary relay admissibility;
    \item \(\mathsf I_{\rm src}\): closed as two-sided local source/relay typing;
    \item \(\mathsf I_{\rm seam}\): closed as compressed-atlas transition seam exactness;
    \item \(\mathsf I_{\rm mix}\): closed only as a \(C4\) guard-lock with \(\beta_4=1/256\) unspent;
    \item \(\mathsf I_{\rm col}\): closed only as bounded-collar recurrence monitor, not finite \(C5\)-closure;
    \item \(\mathsf I_{\rm clk}\): closed only as photon-clock selector/probe, not an RH assertion.
\end{itemize}
Therefore the phrase ``local interface closure'' in r41 means only closure of
the declared admissibility-monitor ledger:
\[
\boxed{
\mathrm{Closure}_{\rm interface}^{r41}
=
\mathrm{Closure}_{\rm local\ monitor}
\neq
\mathrm{Closure}_{\rm global}.
}
\]

\subsection{Closure-Scope Vector}
\label{subsec:r41-closure-scope-vector}

To prevent semantic drift, r41 records the closure-scope vector
\[
\boxed{
\mathsf S_{\rm closure}^{r41}
=({\rm local\ monitor},{\rm action},{\rm field},{\rm global},{\rm all})
=(1,0,0,0,0).
}
\]
Thus the only positive entry is the local monitor entry.  No action-level,
field-equation-level, global, or all-readout closure is claimed.

\subsection{Residual, Leak, and Promotion Audit}
\label{subsec:r41-residual-leak-promotion-audit}

The r41 normal-form residual vector is
\[
\boxed{
\mathcal R_{\rm nf}^{r41}=(0,0,0,0,0,0).
}
\]
The local closure-ledger residual vector is
\[
\boxed{
\mathcal R_{\rm clos}^{r41}=(0,0,0,0,0,0).
}
\]
The total leak vector remains
\[
\boxed{
\mathcal L_{\rm total}^{r41}=(0,0,0,0,0,0,0,0,0,0).
}
\]
The no-action vector remains
\[
\boxed{
\mathfrak A_{\rm noact}^{r41}=(0,0,0,0,0,0).
}
\]
The no-promotion vector remains
\[
\boxed{
\Pi_{\rm np}^{r41}=(0,0,0,0,0,0,0).
}
\]

\subsection{Guard Preservation}
\label{subsec:r41-guard-preservation}

The guard vector remains active:
\[
\boxed{
\Gamma_{\rm guard}^{r41}=(1,1,1,1,1).
}
\]
Equivalently,
\[
\boxed{
\beta_4={1\over256}\hbox{ remains unspent},
\quad
C5\hbox{ has no finite closure},
\quad
C6/T2\hbox{ remains blocked},
\quad
\mathrm{RH}\hbox{ remains photon-clock probe only},
\quad
\hbox{global all-readout closure remains unclaimed}.
}
\]

\subsection{r41 Handoff}
\label{subsec:r41-handoff}

The r41 handoff object is
\[
\boxed{
\mathfrak H_{\rm nf}^{r41}
=
\left(
\mathfrak H_{\rm comp}^{r40},
\mathfrak n^{r41},
\mathsf C_{\rm loc}^{r41},
\mathsf S_{\rm closure}^{r41},
\mathcal R_{\rm nf}^{r41},
\mathcal R_{\rm clos}^{r41},
\mathcal L_{\rm total}^{r41},
\Gamma_{\rm guard}^{r41}
\right).
}
\]
The safe next source-only step is
\[
\boxed{
\text{v4.7.0-r42: local monitor certificate compression and release-candidate guard audit.}
}
\]


\section{v4.7.0-r42: Local Monitor Certificate Compression and Release-Candidate Guard Audit}
\label{sec:v470-r42-monitor-certificate-rc-guard}

\subsection{Source-Only Scope}
\label{subsec:r42-source-only-scope}

This r42 layer is source-only.  It produces no PDF, assigns no new DOI,
authorizes no public upload, and does not change the v4.6.0 release-final
baseline DOI record.  Its purpose is to compress the r36--r41 local monitor
results into a single local certificate object and to audit whether the current
internal working source is release-candidate-ready under the active guards.
The audit is deliberately weaker than a release-final declaration.

\subsection{Local Monitor Certificate Compression}
\label{subsec:r42-local-monitor-certificate-compression}

The r41 local closure ledger provided
\[
\mathsf C_{\rm loc}^{r41}=(1,1,1,1,1,1),
\qquad
\mathsf S_{\rm closure}^{r41}=(1,0,0,0,0).
\]
We compress the local monitor evidence into the certificate
\[
\boxed{
\mathfrak C_{\rm mon}^{r42}
=
\left(
\mathfrak C_{\rm int}^{r37},
\mathfrak S_{\rm loc}^{r39},
\mathfrak n^{r41},
\mathsf C_{\rm loc}^{r41},
\mathsf S_{\rm closure}^{r41},
\Gamma_{\rm guard}^{r41}
\right).
}
\]
The certificate is local and monitor-only.  It records the following
six active local monitor channels:
\[
\boxed{
\mathsf M_{\rm loc}^{r42}
=
(\mathsf I_{\rm bd},\mathsf I_{\rm src},\mathsf I_{\rm seam},
 \mathsf I_{\rm mix},\mathsf I_{\rm col},\mathsf I_{\rm clk}).
}
\]
Their compressed residual vector is
\[
\boxed{
\mathcal R_{\rm mon}^{r42}=(0,0,0,0,0,0).
}
\]
Equivalently, the r42 compression preserves exactly the r41 local monitor
closure and does not promote any monitor into a field, action, or global
readout object.

\subsection{Release-Candidate Guard Audit}
\label{subsec:r42-release-candidate-guard-audit}

The guarded release-candidate audit is evaluated as a vector of local source
conditions:
\[
\boxed{
\mathsf{RC}_{\rm guard}^{r42}
=
(1,1,1,1,1,1,1,1).
}
\]
The eight entries mean, respectively:
\begin{enumerate}[label=\textnormal{RC\arabic*.},leftmargin=2.3em]
  \item source-only intermediate policy respected;
  \item no new DOI assigned;
  \item no PDF/public upload authorized;
  \item local monitor residuals vanish;
  \item local composition residuals vanish;
  \item total leak vector remains zero;
  \item no-promotion vector remains zero;
  \item all five global guards remain active.
\end{enumerate}
Thus r42 is allowed to be treated as an internal release-candidate-precheck
for the local monitor layer only.  The audit does not promote the branch to a
public release, does not create a Zenodo-ready final packet, and does not
remove the source-only status of the current working line.

\subsection{Guard Preservation}
\label{subsec:r42-guard-preservation}

The r42 guard vector is unchanged:
\[
\boxed{
\Gamma_{\rm guard}^{r42}=(1,1,1,1,1).
}
\]
The entries retain their previous meanings:
\[
\boxed{
\beta_4={1\over256}\hbox{ remains unspent},
\quad
C5\hbox{ has no finite closure},
\quad
C6/T2\hbox{ remains blocked},
\quad
\mathrm{RH}\hbox{ remains photon-clock probe only},
\quad
\hbox{global all-readout closure remains unclaimed}.
}
\]
The no-action and no-promotion ledgers are
\[
\boxed{
\mathfrak A_{\rm noact}^{r42}=(0,0,0,0,0,0),
\qquad
\Pi_{\rm np}^{r42}=(0,0,0,0,0,0,0).
}
\]
The total leak vector is still
\[
\boxed{
\mathcal L_{\rm total}^{r42}=(0,0,0,0,0,0,0,0,0,0).
}
\]

\subsection{r42 Decision}
\label{subsec:r42-decision}

The decision ledger is
\[
\boxed{
\mathsf D_{\rm r42}
=
\{
\mathrm{local\ monitor\ certificate}=\mathrm{pass},
\mathrm{RC\ guard\ audit}=\mathrm{pass},
\mathrm{public\ release}=\mathrm{false},
\mathrm{PDF}=\mathrm{not\ built}
\}.
}
\]
Therefore the safe next source-only step is
\[
\boxed{
\text{v4.7.0-r43: source-only metadata freeze check and local certificate handoff packet.}
}
\]



\section{v4.7.0-r43 Source-Only Metadata Freeze Check and Local Certificate Handoff}
\label{sec:r43-metadata-freeze-handoff}

This r43 layer is a source-only metadata-freeze and handoff-packaging pass over
r42.  It does not introduce a new DOI, does not authorize public upload, and
does not build a PDF.  Its purpose is to freeze the current source metadata and
separate the local monitor certificate from every still-active guard.

\subsection{r43 Metadata Freeze Check}
\label{subsec:r43-metadata-freeze-check}

The active source metadata is frozen as
\[
\boxed{
\mathsf M_{\rm freeze}^{r43}
=
(\mathrm{version}=4.7.0\hbox{-}r43,
\mathrm{mode}=\mathrm{source\mbox{-}only},
\mathrm{baseline\ DOI}=10.5281/zenodo.20117361,
\mathrm{new\ DOI}=\varnothing,
\mathrm{PDF}=0).
}
\]
The concept DOI remains the stable Companion-line pointer, while the v4.6.0
version DOI remains a baseline DOI only.  The r43 layer therefore changes the
working-candidate metadata state, not the public release record.

The freeze ledger is
\[
\boxed{
\mathsf F_{\rm meta}^{r43}
=(1,1,1,1,1,1,1,1).
}
\]
The entries mean: active version string present; baseline DOI present; no new
v4.7.0 DOI assigned; source-only mode preserved; PDF build omitted; public
release authorization absent; title-page DOI separation preserved; and guard
language preserved.

\subsection{r43 Local Certificate Handoff Object}
\label{subsec:r43-local-certificate-handoff}

The r42 monitor certificate is handed forward as
\[
\boxed{
\mathfrak H_{\rm cert}^{r43}
=
(\mathfrak C_{\rm mon}^{r42},
\mathcal R_{\rm mon}^{r42},
\mathsf{RC}_{\rm guard}^{r42},
\Gamma_{\rm guard}^{r42},
\mathsf F_{\rm meta}^{r43}).
}
\]
Using the r42 values, this is equivalently the local handoff packet
\[
\boxed{
\mathfrak H_{\rm cert}^{r43}
=
(\mathfrak C_{\rm mon}^{r42},
(0,0,0,0,0,0),
(1,1,1,1,1,1,1,1),
(1,1,1,1,1),
(1,1,1,1,1,1,1,1)).
}
\]
This handoff is local-monitor only.  It is not a field-equation handoff, not a
public release packet, and not a DOI-bearing object.

\subsection{r43 Guard Preservation Ledger}
\label{subsec:r43-guard-preservation-ledger}

The r43 guard vector remains
\[
\boxed{
\Gamma_{\rm guard}^{r43}=(1,1,1,1,1).
}
\]
The entries retain their previous meanings:
\[
\boxed{
\beta_4={1\over256}\hbox{ remains unspent},
\quad
C5\hbox{ has no finite closure},
\quad
C6/T2\hbox{ remains blocked},
\quad
\mathrm{RH}\hbox{ remains photon-clock probe only},
\quad
\hbox{global all-readout closure remains unclaimed}.
}
\]
Thus r43 freezes and packages the local monitor certificate without spending any
guard, without converting the skeleton into a physical action, and without
promoting the local monitor closure to a global all-readout closure.

\subsection{r43 Decision}
\label{subsec:r43-decision}

The decision ledger is
\[
\boxed{
\mathsf D_{\rm r43}
=
\{
\mathrm{metadata\ freeze}=\mathrm{pass},
\mathrm{local\ certificate\ handoff}=\mathrm{pass},
\mathrm{public\ release}=\mathrm{false},
\mathrm{new\ DOI}=\varnothing,
\mathrm{PDF}=\mathrm{not\ built}
\}.
}
\]
Therefore the safe next source-only step is
\[
\boxed{
\text{v4.7.0-r44: handoff integrity scan and local certificate index.}
}
\]



\section{v4.7.0-r44 Source-Only Handoff Integrity Scan and Local Certificate Index}
\label{sec:r44-handoff-integrity-certificate-index}

This r44 layer is a source-only handoff-integrity and local-certificate-index
pass over r43.  It builds no PDF, assigns no new DOI, and authorizes no public
release.  Its purpose is to check that the r43 handoff object is complete and to
index the local certificates needed by the next working layer.

\subsection{r44 Handoff Integrity Scan}
\label{subsec:r44-handoff-integrity-scan}

The incoming handoff packet is
\[
\mathfrak H_{\rm cert}^{r43}
=
(\mathfrak C_{\rm mon}^{r42},
\mathcal R_{\rm mon}^{r42},
\mathsf{RC}_{\rm guard}^{r42},
\Gamma_{\rm guard}^{r42},
\mathsf F_{\rm meta}^{r43}).
\]
The r44 integrity scan records the component checks as
\[
\boxed{
\mathsf I_{\rm hand}^{r44}
=(1,1,1,1,1,1,1,1,1).
}
\]
The entries mean: monitor certificate present; monitor residual present; RC
precheck present; guard vector present; metadata-freeze vector present;
source-only mode preserved; baseline DOI preserved; no new v4.7.0 DOI assigned;
and no PDF build requested or generated.

The integrity residual is therefore
\[
\boxed{
\Delta_{\rm hand}^{r44}=0.
}
\]
Equivalently, no component of the r43 handoff packet is missing, retyped as a
physical field equation, or promoted beyond the declared local-monitor scope.

\subsection{r44 Local Certificate Index}
\label{subsec:r44-local-certificate-index}

The local certificate index is
\[
\boxed{
\mathfrak I_{\rm cert}^{r44}
=
\{I_0,I_1,I_2,I_3,I_4,I_5,I_6,I_7\}.
}
\]
The indexed entries are
\[
\begin{array}{c|l|l}
\text{index} & \text{object} & \text{scope}\\
\hline
I_0 & \mathfrak H_{\rm seed}^{r34} & \text{comfort-calibrated local seed}\\
I_1 & \mathfrak F_{\rm loc}^{r38} & \text{local pre-field interface only}\\
I_2 & \mathfrak S_{\rm loc}^{r39} & \text{admissibility-only operator-tag skeleton}\\
I_3 & \mathfrak n^{r41} & \text{guarded composition normal form}\\
I_4 & \mathfrak C_{\rm int}^{r37} & \text{local MM/HC interface certificate}\\
I_5 & \mathfrak C_{\rm mon}^{r42} & \text{compressed local monitor certificate}\\
I_6 & \mathfrak H_{\rm cert}^{r43} & \text{source-only local certificate handoff}\\
I_7 & \mathsf F_{\rm meta}^{r43} & \text{metadata-freeze support vector}
\end{array}
\]
All entries are local-monitor or local-interface objects.  None is a physical
action, a Lagrangian, an Euler--Lagrange variation, a derived field equation, a
T2 theorem, or a global all-readout certificate.

\subsection{r44 Certificate Index Guard Ledger}
\label{subsec:r44-certificate-index-guard-ledger}

The certificate index has the guard ledger
\[
\boxed{
\mathsf{CI}_{\rm guard}^{r44}
=(1,1,1,1,1,1,1,1).
}
\]
The entries mean: local scope preserved; source-only scope preserved; no DOI
promotion; no PDF build; no action promotion; no field-equation promotion; no
all-readout promotion; and all explicit guards remain active.

The guard vector remains
\[
\boxed{
\Gamma_{\rm guard}^{r44}=(1,1,1,1,1).
}
\]
Thus
\[
\boxed{
\beta_4={1\over256}\text{ remains unspent},\quad
C5\text{ has no finite closure},\quad
C6/T2\text{ remains blocked},\quad
\mathrm{RH}\text{ remains photon-clock probe only},\quad
\text{global all-readout closure remains unclaimed}.
}
\]

\subsection{r44 Decision}
\label{subsec:r44-decision}

The decision ledger is
\[
\boxed{
\mathsf D_{\rm r44}
=
\{
\mathrm{handoff\ integrity}=\mathrm{pass},
\mathrm{certificate\ index}=\mathrm{pass},
\mathrm{public\ release}=\mathrm{false},
\mathrm{new\ DOI}=\varnothing,
\mathrm{PDF}=\mathrm{not\ built}
\}.
}
\]
Therefore the safe next source-only step is
\[
\boxed{
\text{v4.7.0-r45: local certificate reader-routing map and next-branch queue.}
}
\]



\section{v4.7.0-r45 Source-Only Local Certificate Reader-Routing Map and Next-Branch Queue}
\label{sec:r45-reader-routing-next-branch}

This r45 layer is a source-only reader-routing and next-branch-queue pass over
r44.  It builds no PDF, assigns no new DOI, and authorizes no public release.
Its purpose is to make the local certificate index readable by assigning each
certificate object to a declared reader route and by recording the next safe
source-only working queues.

\subsection{r45 Reader-Routing Map}
\label{subsec:r45-reader-routing-map}

The incoming local certificate index is
\[
\mathfrak I_{\rm cert}^{r44}
=
\{I_0,I_1,I_2,I_3,I_4,I_5,I_6,I_7\}.
\]
The reader-routing map is
\[
\boxed{
\mathsf{RR}^{r45}:
\mathfrak I_{\rm cert}^{r44}
\longrightarrow
\mathsf{Route}_{\rm loc}^{r45}.
}
\]
It assigns the following local routes:
\[
\begin{array}{c|l|l|l}
\text{index} & \text{object} & \text{reader route} & \text{use}\\
\hline
I_0 & \mathfrak H_{\rm seed}^{r34} & R_{\rm seed} & \text{comfort seed lookup}\\
I_1 & \mathfrak F_{\rm loc}^{r38} & R_{\rm prefield} & \text{local pre-field interface lookup}\\
I_2 & \mathfrak S_{\rm loc}^{r39} & R_{\rm skel} & \text{admissibility-only skeleton lookup}\\
I_3 & \mathfrak n^{r41} & R_{\rm nf} & \text{guarded normal-form lookup}\\
I_4 & \mathfrak C_{\rm int}^{r37} & R_{\rm int} & \text{local interface certificate lookup}\\
I_5 & \mathfrak C_{\rm mon}^{r42} & R_{\rm mon} & \text{compressed monitor certificate lookup}\\
I_6 & \mathfrak H_{\rm cert}^{r43} & R_{\rm hand} & \text{handoff packet lookup}\\
I_7 & \mathsf F_{\rm meta}^{r43} & R_{\rm meta} & \text{metadata-freeze lookup}
\end{array}
\]
Every route is local and read-only.  No route is allowed to reinterpret its
object as a physical action, a Lagrangian, an Euler--Lagrange variation, a
derived field equation, a T2 theorem, or a global all-readout certificate.

\subsection{r45 Reader-Routing Integrity Vector}
\label{subsec:r45-reader-routing-integrity-vector}

The r45 routing integrity vector is
\[
\boxed{
\mathsf I_{\rm route}^{r45}
=(1,1,1,1,1,1,1,1).
}
\]
The entries mean: every certificate-index entry has a route; every route has a
scope; every route has a use; every route remains source-only; every route
preserves local scope; every route preserves the no-action guard; every route
preserves the no-field-equation guard; and every route preserves the no-global
all-readout guard.

The routing residual is therefore
\[
\boxed{
\Delta_{\rm route}^{r45}=0.
}
\]
Thus the local certificate index is now not only complete, but also reader
routable without changing its mathematical status.

\subsection{r45 Next-Branch Queue}
\label{subsec:r45-next-branch-queue}

The next-branch queue is recorded as
\[
\boxed{
\mathsf Q_{\rm next}^{r45}
=
(Q_1,Q_2,Q_3,Q_4,Q_5).
}
\]
The entries are
\[
\begin{array}{c|l|l|l}
\text{queue} & \text{working title} & \text{mode} & \text{guard status}\\
\hline
Q_1 & \text{local certificate dependency graph} & \text{source-only} & \text{guarded}\\
Q_2 & \text{reader-route cross-reference hygiene} & \text{source-only} & \text{guarded}\\
Q_3 & \text{local pre-field interface theoremlet candidate} & \text{source-only} & \text{no action, no field equation}\\
Q_4 & \text{MM/HC interface handoff to next research branch} & \text{source-only} & \text{local only}\\
Q_5 & \text{release-candidate hygiene preflight, if requested} & \text{source-only unless PDF requested} & \text{no DOI by default}
\end{array}
\]
The default next source-only step is
\[
\boxed{
Q_1:
\text{v4.7.0-r46 local certificate dependency graph and reader-route cross-check.}
}
\]
This default is intentionally still a local certificate/readout step.  It does
not start a public release, it does not assign a DOI, and it does not build a
PDF.

\subsection{r45 Guard Preservation Ledger}
\label{subsec:r45-guard-preservation-ledger}

The r45 reader-routing layer preserves the guard vector
\[
\boxed{
\Gamma_{\rm guard}^{r45}=(1,1,1,1,1).
}
\]
The entries retain their previous meanings:
\[
\boxed{
\beta_4={1\over256}\text{ remains unspent},\quad
C5\text{ has no finite closure},\quad
C6/T2\text{ remains blocked},\quad
\mathrm{RH}\text{ remains photon-clock probe only},\quad
\text{global all-readout closure remains unclaimed}.
}
\]
The r45 reader-routing map therefore changes only the navigability of the local
certificate layer.  It does not spend any guard, does not promote a monitor to a
field equation, and does not convert local monitor closure into global closure.

\subsection{r45 Decision}
\label{subsec:r45-decision}

The r45 decision ledger is
\[
\boxed{
\mathsf D_{\rm r45}
=
\{
\mathrm{reader\ routing}=\mathrm{pass},
\mathrm{routing\ residual}=0,
\mathrm{next\ branch\ queue}=\mathrm{pass},
\mathrm{public\ release}=\mathrm{false},
\mathrm{new\ DOI}=\varnothing,
\mathrm{PDF}=\mathrm{not\ built}
\}.
}
\]
Therefore the safe next source-only step is
\[
\boxed{
\text{v4.7.0-r46: local certificate dependency graph and reader-route cross-check.}
}
\]




\section{v4.7.0-r46 Source-Only Local Certificate Dependency Graph and Reader-Route Cross-Check}
\label{sec:r46-dependency-graph-route-crosscheck}

This r46 layer is a source-only dependency-graph and reader-route cross-check over
\(\mathfrak I_{\rm cert}^{r44}\) and the r45 reader-routing map.  It does not
build a PDF, does not assign a new DOI, and does not promote the local monitor
certificate to an action, a Lagrangian, a derived field equation, or global
all-readout closure.

\subsection{r46 Local Certificate Dependency Graph}
\label{subsec:r46-local-certificate-dependency-graph}

The local certificate index contains the eight nodes
\[
\boxed{
\mathfrak I_{\rm cert}^{r44}
=
\{I_0,I_1,I_2,I_3,I_4,I_5,I_6,I_7\}.
}
\]
For r46 the nodes are routed through the following directed acyclic dependency
skeleton:
\[
\boxed{
\mathsf{DG}_{\rm cert}^{r46}
=
(\mathsf V_{\rm cert}^{r46},\mathsf E_{\rm cert}^{r46}).
}
\]
The vertex set is
\[
\mathsf V_{\rm cert}^{r46}
=
\{
\mathfrak H_{\rm seed}^{r34},
\mathfrak C_{\rm int}^{r37},
\mathfrak F_{\rm loc}^{r38},
\mathfrak S_{\rm loc}^{r39},
\mathfrak n^{r41},
\mathfrak C_{\rm mon}^{r42},
\mathfrak H_{\rm cert}^{r43},
\mathsf F_{\rm meta}^{r43}
\}.
\]
The dependency edges are
\[
\begin{array}{c|c|l}
\text{edge} & \text{dependency} & \text{meaning}\\
\hline
E_1  & \mathfrak H_{\rm seed}^{r34}\to\mathfrak C_{\rm int}^{r37}  & \text{seed supports local interface certificate}\\
E_2  & \mathfrak C_{\rm int}^{r37}\to\mathfrak F_{\rm loc}^{r38}  & \text{certificate admitted as pre-field interface}\\
E_3  & \mathfrak F_{\rm loc}^{r38}\to\mathfrak S_{\rm loc}^{r39}  & \text{pre-field interface gives admissibility-only skeleton}\\
E_4  & \mathfrak S_{\rm loc}^{r39}\to\mathfrak n^{r41}  & \text{skeleton supports guarded composition normal form}\\
E_5  & \mathfrak C_{\rm int}^{r37}\to\mathfrak C_{\rm mon}^{r42}  & \text{interface certificate enters monitor compression}\\
E_6  & \mathfrak S_{\rm loc}^{r39}\to\mathfrak C_{\rm mon}^{r42}  & \text{skeleton enters monitor compression}\\
E_7  & \mathfrak n^{r41}\to\mathfrak C_{\rm mon}^{r42}  & \text{normal form enters monitor compression}\\
E_8  & \mathfrak C_{\rm mon}^{r42}\to\mathfrak H_{\rm cert}^{r43}  & \text{monitor certificate enters handoff packet}\\
E_9  & \mathsf F_{\rm meta}^{r43}\to\mathfrak H_{\rm cert}^{r43}  & \text{metadata-freeze vector enters handoff packet}.
\end{array}
\]
The graph is acyclic.  One compatible topological order is
\[
\boxed{
\mathfrak H_{\rm seed}^{r34},
\mathfrak C_{\rm int}^{r37},
\mathfrak F_{\rm loc}^{r38},
\mathfrak S_{\rm loc}^{r39},
\mathfrak n^{r41},
\mathsf F_{\rm meta}^{r43},
\mathfrak C_{\rm mon}^{r42},
\mathfrak H_{\rm cert}^{r43}.
}
\]
This order is a reader dependency order only.  It is not a proof of any global
closure and it does not change the mathematical status of any guard.

\subsection{r46 Reader-Route Cross-Check}
\label{subsec:r46-reader-route-crosscheck}

The r45 reader routes are cross-checked against the dependency graph by the map
\[
\boxed{
\mathsf{XC}_{\rm route}^{r46}:
\mathsf{Route}_{\rm loc}^{r45}\to\mathsf{DG}_{\rm cert}^{r46}.
}
\]
The cross-check table is
\[
\begin{array}{c|c|c|c}
\text{route} & \text{node} & \text{dependency status} & \text{cross-check}\\
\hline
R_{\rm seed} & \mathfrak H_{\rm seed}^{r34} & \text{root seed} & 1\\
R_{\rm int} & \mathfrak C_{\rm int}^{r37} & \text{seed-supported} & 1\\
R_{\rm prefield} & \mathfrak F_{\rm loc}^{r38} & \text{certificate-supported} & 1\\
R_{\rm skel} & \mathfrak S_{\rm loc}^{r39} & \text{pre-field-supported} & 1\\
R_{\rm nf} & \mathfrak n^{r41} & \text{skeleton-supported} & 1\\
R_{\rm mon} & \mathfrak C_{\rm mon}^{r42} & \text{compressed monitor-supported} & 1\\
R_{\rm meta} & \mathsf F_{\rm meta}^{r43} & \text{metadata root} & 1\\
R_{\rm hand} & \mathfrak H_{\rm cert}^{r43} & \text{monitor+metadata-supported} & 1.
\end{array}
\]
Thus the reader-route cross-check vector is
\[
\boxed{
\mathsf{XC}_{\rm route}^{r46}
=
(1,1,1,1,1,1,1,1).
}
\]
The unresolved-route residual is
\[
\boxed{
\Delta_{\rm route}^{r46}=0.
}
\]

\subsection{r46 Dependency Integrity and Guard Separation}
\label{subsec:r46-dependency-integrity-guard-separation}

The dependency integrity vector is
\[
\boxed{
\mathsf I_{\rm dep}^{r46}
=
(
\mathrm{nodes},
\mathrm{edges},
\mathrm{acyclic},
\mathrm{topological\ order},
\mathrm{route\ coverage},
\mathrm{no\ dangling\ node},
\mathrm{no\ dangling\ route},
\mathrm{guard\ separation}
)
=(1,1,1,1,1,1,1,1).
}
\]
The dependency residual is
\[
\boxed{
\Delta_{\rm dep}^{r46}=0.
}
\]
The graph deliberately separates local certificate dependencies from guard
claims.  The local certificate graph may be complete while the guard layer stays
active.  Formally,
\[
\boxed{
\mathsf{complete}(\mathsf{DG}_{\rm cert}^{r46})
\not\Rightarrow
\mathrm{closure}_{\rm global}.
}
\]
This is the same separation already enforced by the r41 closure-scope vector and
by the r42--r45 handoff ledgers.

\subsection{r46 Guard Preservation Ledger}
\label{subsec:r46-guard-preservation-ledger}

The r46 dependency graph and reader-route cross-check preserve the guard vector
\[
\boxed{
\Gamma_{\rm guard}^{r46}=(1,1,1,1,1).
}
\]
The entries retain their previous meanings:
\[
\boxed{
\beta_4={1\over256}\text{ remains unspent},\quad
C5\text{ has no finite closure},\quad
C6/T2\text{ remains blocked},\quad
\mathrm{RH}\text{ remains photon-clock probe only},\quad
\text{global all-readout closure remains unclaimed}.
}
\]
No r46 route, dependency edge, or topological ordering spends any of these
guards.  In particular, \(\mathfrak H_{\rm cert}^{r43}\) is a local handoff
packet, not a public release certificate and not a global theorem.

\subsection{r46 Decision and Next Queue}
\label{subsec:r46-decision-next-queue}

The r46 decision ledger is
\[
\boxed{
\mathsf D_{\rm r46}
=
\{
\mathrm{dependency\ graph}=\mathrm{pass},
\mathrm{route\ cross\ check}=\mathrm{pass},
\Delta_{\rm dep}=0,
\Delta_{\rm route}=0,
\mathrm{PDF}=\mathrm{not\ built},
\mathrm{new\ DOI}=\varnothing,
\mathrm{public\ release}=\mathrm{false}
\}.
}
\]
The safe next source-only step is
\[
\boxed{
\text{v4.7.0-r47: dependency-order theoremlet and branch-launch guard preflight.}
}
\]



\section{v4.7.0-r47 Source-Only Dependency-Order Theoremlet and Branch-Launch Guard Preflight}
\label{sec:r47-dependency-order-theoremlet-branch-preflight}

This r47 layer is source-only.  It converts the r46 directed acyclic certificate
index into a small dependency-order theoremlet and then performs a branch-launch
preflight for the next local work block.  It does not build a PDF, assign a new
DOI, authorize a public release, promote the local monitor certificate into a
field equation, or spend any global guard.

\subsection{r47 Dependency-Order Theoremlet}
\label{subsec:r47-dependency-order-theoremlet}

Let
\[
\mathsf{DG}_{\rm cert}^{r46}
=
(\mathsf V_{\rm cert}^{r46},\mathsf E_{\rm cert}^{r46})
\]
be the r46 local certificate dependency graph with eight certificate nodes and
nine dependency edges.  The r46 cross-check established
\[
\Delta_{\rm dep}^{r46}=0,
\qquad
\Delta_{\rm route}^{r46}=0,
\qquad
\mathsf{XC}_{\rm route}^{r46}=(1,1,1,1,1,1,1,1).
\]
The dependency-order theoremlet recorded here is the following strictly local
statement:
\[
\boxed{
\mathsf{Acyclic}(\mathsf{DG}_{\rm cert}^{r46})
\wedge
\Delta_{\rm dep}^{r46}=0
\wedge
\Delta_{\rm route}^{r46}=0
\Rightarrow
\exists\,\prec_{\rm cert}^{r47}
\text{ compatible with every reader route.}
}
\]
One compatible order is
\[
\boxed{
\mathfrak H_{\rm seed}^{r34}
\prec
\mathfrak C_{\rm int}^{r37}
\prec
\mathfrak F_{\rm loc}^{r38}
\prec
\mathfrak S_{\rm loc}^{r39}
\prec
\mathfrak n^{r41}
\prec
\mathfrak C_{\rm mon}^{r42}
\prec
\mathfrak H_{\rm cert}^{r43},
}
\]
with the metadata vector \(\mathsf F_{\rm meta}^{r43}\) treated as a separate
metadata root for the source-only handoff packet.  The theoremlet is only a
reader-routing and dependency-order statement.  It does not assert that local
certificate completion implies physical, global, or all-readout closure.

\subsection{r47 Branch-Launch Guard Preflight}
\label{subsec:r47-branch-launch-guard-preflight}

The next local branch may be launched only through a guard-preserving queue.
The preflight checks are
\[
\boxed{
\mathsf{BL}_{\rm pre}^{r47}
=
(
\mathrm{source\mbox{-}only},
\mathrm{metadata\ freeze},
\mathrm{dependency\ order},
\mathrm{reader\ routes},
\mathrm{guard\ vector},
\mathrm{no\ PDF},
\mathrm{no\ DOI},
\mathrm{no\ public\ release},
\mathrm{no\ promotion}
)
=(1,1,1,1,1,1,1,1,1).
}
\]
Thus the branch-launch residual is
\[
\boxed{
\Delta_{\rm BL}^{r47}=0.
}
\]
The preflight authorizes only another internal source-only layer.  It does not
turn the r42 monitor certificate into a release certificate and does not turn the
r38--r41 local interface stack into a field-theoretic operator calculus.

\subsection{r47 Guard Preservation Ledger}
\label{subsec:r47-guard-preservation-ledger}

The r47 dependency-order theoremlet preserves the guard vector
\[
\boxed{
\Gamma_{\rm guard}^{r47}=(1,1,1,1,1).
}
\]
The five entries retain the standing meanings:
\[
\boxed{
\beta_4={1\over256}\text{ remains unspent},
\quad C5\text{ has no finite closure},
\quad C6/T2\text{ remains blocked},
\quad \mathrm{RH}\text{ remains photon-clock probe only},
\quad \text{global all-readout closure remains unclaimed}.
}
\]
Equivalently,
\[
\boxed{
\mathrm{LocalDependencyOrder}^{r47}
\not\Rightarrow
\mathrm{GlobalClosure}.
}
\]

\subsection{r47 Decision and Next Queue}
\label{subsec:r47-decision-next-queue}

The r47 decision ledger is
\[
\boxed{
\mathsf D_{\rm r47}
=
\{
\mathrm{dependency\ theoremlet}=\mathrm{pass},
\mathrm{branch\ preflight}=\mathrm{pass},
\Delta_{\rm BL}=0,
\mathrm{PDF}=\mathrm{not\ built},
\mathrm{new\ DOI}=\varnothing,
\mathrm{public\ release}=\mathrm{false}
\}.
}
\]
The safe next source-only step is
\[
\boxed{
\text{v4.7.0-r48: branch queue selection and source-only work-package split.}
}
\]


% -----------------------------------------------------------------------------
\section{v4.7.0-r48 Source-Only Branch Queue Selection and Work-Package Split}
\label{sec:v470-r48-branch-queue-work-package-split}

This source-only layer consumes the r47 branch-launch preflight and selects a
next internal queue without changing release status.  No PDF is built, no new
version DOI is assigned, and no public release is authorized.  The purpose is to
split the next work into guard-stable source packages that can be executed in
small local steps.

\subsection{r48 Queue Selection Rule}
\label{subsec:r48-queue-selection-rule}

Let
\[
\mathsf{BL}_{\rm pre}^{r47}=(1,1,1,1,1,1,1,1,1),
\qquad
\Delta_{\rm BL}^{r47}=0.
\]
A branch candidate \(Q\) is admissible for the next source-only layer only if it
satisfies
\[
\boxed{
\mathsf{Adm}_{\rm branch}^{r48}(Q)
=
\mathrm{sourceOnly}(Q)
\wedge
\mathrm{dependencyRouted}(Q)
\wedge
\mathrm{guardPreserving}(Q)
\wedge
\neg\mathrm{claimsClosure}(Q)
\wedge
\neg\mathrm{requestsDOI}(Q).
}
\]
This rule selects branch work by admissibility and guard preservation, not by
mathematical ambition.  In particular, it forbids promoting the local monitor
certificate to a field equation, action principle, T2 closure, RH statement, or
global all-readout theorem.

\subsection{r48 Candidate Branch Queue}
\label{subsec:r48-candidate-branch-queue}

The queue is split into five candidate work packages:
\[
\boxed{
\mathsf Q_{\rm branch}^{r48}
=
(Q_1,Q_2,Q_3,Q_4,Q_5).
}
\]
The candidates are:
\[
\begin{array}{lll}
Q_1&=&\text{guarded local field-interface elaboration},\\
Q_2&=&\text{boundary-source robustness extension},\\
Q_3&=&\text{photon-clock selector-window calibration},\\
Q_4&=&\text{C4 mixed-kernel localization refinement},\\
Q_5&=&\text{C5 asymptotic-collar bookkeeping consolidation}.
\end{array}
\]
All five are source-only candidates.  Their admissibility statuses are
\[
\boxed{
\mathsf{Adm}_{\rm branch}^{r48}(Q_i)=1
\quad(1\leq i\leq5),
}
\]
provided that they keep the standing guard vector active.

The selected default branch is
\[
\boxed{
Q_1:
\text{guarded local field-interface elaboration}.
}
\]
The reason for selecting \(Q_1\) first is structural: r38--r41 already produced
a local pre-field/interface stack, while r42--r47 compressed and routed its
certificate.  Thus \(Q_1\) can extend an already routed object without touching
C4, C5 finite closure, C6/T2, RH, or global readout claims.

\subsection{r48 Source-Only Work-Package Split}
\label{subsec:r48-source-only-work-package-split}

The selected branch \(Q_1\) is split into four source-only work packages:
\[
\boxed{
Q_1=(W_1^{r48},W_2^{r48},W_3^{r48},W_4^{r48}).
}
\]
They are:
\[
\begin{array}{lll}
W_1^{r48}&=&\text{interface-slot glossary and two-sided typing recap},\\
W_2^{r48}&=&\text{admissibility-only operator-domain refinement},\\
W_3^{r48}&=&\text{local pre-field interface diagram in text form},\\
W_4^{r48}&=&\text{no-action/no-promotion audit refresh}.
\end{array}
\]
Each work package inherits the source-only constraint
\[
\boxed{
\mathrm{PDF}(W_j^{r48})=0,
\quad
\mathrm{newDOI}(W_j^{r48})=\varnothing,
\quad
\mathrm{publicRelease}(W_j^{r48})=0.
}
\]
The dependency order is
\[
\boxed{
W_1^{r48}\prec W_2^{r48}\prec W_3^{r48}\prec W_4^{r48}.
}
\]
This order is compatible with the r47 certificate order, because it starts from
\(\mathfrak H_{\rm cert}^{r43}\) and only refines the local interface reader
path.

\subsection{r48 Queue-Selection Audit}
\label{subsec:r48-queue-selection-audit}

The branch queue selection vector is
\[
\boxed{
\mathsf{QS}^{r48}
=
(1,1,1,1,1,1,1,1).
}
\]
The entries mean, respectively: source-only mode, dependency-routed queue,
guard-preserving queue, no DOI, no public release, no PDF, no closure claim, and
selected default branch.  Thus
\[
\boxed{
\Delta_{\rm QS}^{r48}=0.
}
\]
The work-package split vector is
\[
\boxed{
\mathsf{WP}^{r48}=(1,1,1,1),
\qquad
\Delta_{\rm WP}^{r48}=0.
}
\]
Therefore r48 authorizes the next source-only step but does not authorize a
release candidate build or any global claim.

\subsection{r48 Guard Preservation Ledger}
\label{subsec:r48-guard-preservation-ledger}

The guard vector remains
\[
\boxed{
\Gamma_{\rm guard}^{r48}=(1,1,1,1,1).
}
\]
Its entries retain the standing meanings:
\[
\boxed{
\beta_4={1\over256}\text{ remains unspent},
\quad C5\text{ has no finite closure},
\quad C6/T2\text{ remains blocked},
\quad \mathrm{RH}\text{ remains photon-clock probe only},
\quad \text{global all-readout closure remains unclaimed}.
}
\]
Equivalently,
\[
\boxed{
\mathrm{BranchQueue}^{r48}
\not\Rightarrow
\mathrm{FieldEquation}
\not\Rightarrow
\mathrm{GlobalClosure}.
}
\]

\subsection{r48 Decision and Next Queue}
\label{subsec:r48-decision-next-queue}

The r48 decision ledger is
\[
\boxed{
\mathsf D_{r48}
=
\{
\mathrm{queue\ selected}=Q_1,
\mathrm{work\ packages}=4,
\Delta_{\rm QS}=0,
\Delta_{\rm WP}=0,
\mathrm{PDF}=\mathrm{not\ built},
\mathrm{new\ DOI}=\varnothing,
\mathrm{public\ release}=\mathrm{false}
\}.
}
\]
The safe next source-only step is
\[
\boxed{
\text{v4.7.0-r49: interface-slot glossary and two-sided typing recap.}
}
\]



\section{v4.7.0-r49 Source-Only Interface-Slot Glossary and Two-Sided Typing Recap}
\label{sec:v470-r49-interface-slot-glossary-typing-recap}

This r49 layer executes the first work package selected in r48,
\(W_1^{r48}\): interface-slot glossary and two-sided typing recap.  It is
source-only.  It introduces no new DOI, no PDF build, no public release, no
action functional, no Lagrangian, no Euler--Lagrange variation, no derived field
equation, and no global readout promotion.

\subsection{r49 Interface-Slot Glossary}
\label{subsec:r49-interface-slot-glossary}

The local field-interface stack uses the following six admissibility-only slots:
\[
\boxed{
\mathbb I_{\rm loc}^{r49}
=
\{\mathsf I_{\rm bd},\mathsf I_{\rm src},\mathsf I_{\rm seam},
\mathsf I_{\rm mix},\mathsf I_{\rm col},\mathsf I_{\rm clk}\}.
}
\]
They are glossary entries for local monitoring only:
\[
\begin{array}{lll}
\mathsf{I_{\rm bd}} &:& \text{boundary relay admissibility, measured by }E_{\rm rel}^{-}+E_{\rm src}^{-},\\
\mathsf{I_{\rm src}} &:& \text{two-sided local source/relay typing},\\
\mathsf{I_{\rm seam}} &:& \text{compressed-atlas transition seam exactness},\\
\mathsf{I_{\rm mix}} &:& \text{mixed-kernel guard lock preserving }\beta_4=1/256,\\
\mathsf{I_{\rm col}} &:& \text{bounded-collar recurrence monitor for }C5,\\
\mathsf{I_{\rm clk}} &:& \text{photon-clock selector/probe window.}
\end{array}
\]
The residual readings remain
\[
\boxed{
\mathcal R_{\rm slot}^{r49}=(0,0,0,0,0,0).
}
\]
This is a local slot statement only: it does not infer a physical operator,
physical equation, or global readout theorem from the slot names.

\subsection{r49 Two-Sided Typing Recap}
\label{subsec:r49-two-sided-typing-recap}

The two-sided notation used for these slots keeps a left/front readout side and a
right/back structural side.  A schematic local interface object is written as
\[
\boxed{
{}^{L}{}_{R}\mathcal T^{a}{}_{b},
}
\]
where the front indices \({}^{L}{}_{R}\) mark the local readout/source side, and
\({}^{a}{}_{b}\) marks the relay/structure side.  Upper--lower contraction keeps
its ordinary action/readout meaning.  Same-variance comparison remains forbidden
unless an explicit comparison tensor is declared.

For r49 the canonical local typing is
\[
\boxed{
{}^{\partial{\rm src}}{}_{\rm loc}
\mathcal B_0^{\rm res}
{}^{V0}_{\rm relay}
\quad\leadsto\quad
{}^{\rm loc}{}_{\rm guard}
\mathfrak F_{\rm loc}^{r38}
{}^{\rm relay}_{\rm mon}.
}
\]
The notation means: a local boundary source is allowed to feed a local relay
monitor only after the guard projection has separated it from \(C4\), \(C6/T2\),
and the all-readout channel.

\subsection{r49 Slot-to-Operator-Tag Recap}
\label{subsec:r49-slot-to-operator-tag-recap}

Each slot has a corresponding admissibility-only tag from r39:
\[
\boxed{
\mathsf I_{\rm bd}\leftrightarrow\mathcal O_{\rm bd},\quad
\mathsf I_{\rm src}\leftrightarrow\mathcal O_{\rm src},\quad
\mathsf I_{\rm seam}\leftrightarrow\mathcal O_{\rm seam},\quad
\mathsf I_{\rm mix}\leftrightarrow\mathcal O_{\rm mix},\quad
\mathsf I_{\rm col}\leftrightarrow\mathcal O_{\rm col},\quad
\mathsf I_{\rm clk}\leftrightarrow\mathcal O_{\rm clk}.
}
\]
These tags are not physical operators.  They are monitor tags whose composition
was normalized in r41.  Their scope vector remains
\[
\boxed{
\mathsf S_{\rm closure}^{r49}=(1,0,0,0,0),
}
\]
meaning: local monitor closure only; no action closure, no field-equation
closure, no global readout closure, and no all-readout closure.

\subsection{r49 Local Reader Recap}
\label{subsec:r49-local-reader-recap}

For a reader entering the local field-interface branch at r49, the minimum safe
reading route is
\[
\boxed{
\mathfrak H_{\rm seed}^{r34}
\prec
\mathfrak F_{\rm loc}^{r38}
\prec
\mathfrak S_{\rm loc}^{r39}
\prec
\mathfrak n^{r41}
\prec
\mathfrak C_{\rm mon}^{r42}
\prec
\mathfrak H_{\rm cert}^{r43}.
}
\]
The route may also pass through \(\mathfrak C_{\rm int}^{r37}\) for the full
interface certificate.  The metadata-freeze support vector \(\mathsf F_{\rm
meta}^{r43}\) remains a separate source-only root.

\subsection{r49 Audit and Guard Preservation}
\label{subsec:r49-audit-guard-preservation}

The r49 glossary and typing recap satisfy the source-only work-package vector
\[
\boxed{
\mathsf W_1^{r49}=(1,1,1,1),
\qquad
\Delta_{W_1}^{r49}=0.
}
\]
The entries mean: glossary inserted, two-sided typing recapped, operator-tag
scope preserved, and no-promotion audit preserved.  The guard vector remains
\[
\boxed{
\Gamma_{\rm guard}^{r49}=(1,1,1,1,1).
}
\]
Thus
\[
\boxed{
\beta_4={1\over256}\text{ remains unspent},
\quad C5\text{ has no finite closure},
\quad C6/T2\text{ remains blocked},
\quad \mathrm{RH}\text{ remains photon-clock probe only},
\quad \text{global all-readout closure remains unclaimed}.
}
\]
No new DOI, PDF, release packet, public upload, or global theorem is authorized.

\subsection{r49 Decision and Next Step}
\label{subsec:r49-decision-next-step}

The r49 decision ledger is
\[
\boxed{
\mathsf D_{r49}
=
\{
W_1\text{ complete},
\Delta_{W_1}=0,
\mathrm{PDF}=\mathrm{not\ built},
\mathrm{new\ DOI}=\varnothing,
\mathrm{public\ release}=\mathrm{false},
\Gamma_{\rm guard}=(1,1,1,1,1)
\}.
}
\]
The safe next source-only step is
\[
\boxed{
\text{v4.7.0-r51: local pre-field interface diagram in text form.}
}
\]



\section{v4.7.0-r50 Source-Only Addendum: Admissibility-Only Operator-Domain Refinement}
\label{sec:r50-operator-domain-refinement}

\subsection{Purpose and Scope}
\label{subsec:r50-purpose-scope}

This source-only addendum executes the second work package from r48:
\[
\boxed{W_2^{r48}:\quad \text{admissibility-only operator-domain refinement}.}
\]
The purpose is to specify the local domains and codomains of the six r39/r49
operator tags without promoting them to physical operators.  The refinement is
strictly local, monitor-facing, and guarded.  It does not introduce an action,
a Lagrangian, an Euler--Lagrange variation, a derived field equation, a global
readout operator, or a closure theorem.

\subsection{Admissibility Domains}
\label{subsec:r50-admissibility-domains}

Let
\[
\mathbb D_{\rm loc}^{r50}
=
\mathbb D_{\rm bd}
\times
\mathbb D_{\rm src}
\times
\mathbb D_{\rm seam}
\times
\mathbb D_{\rm mix}
\times
\mathbb D_{\rm col}
\times
\mathbb D_{\rm clk}
\]
be the local product domain of admissibility monitors.  Each component domain is
read only as an interface-domain declaration:
\[
\begin{array}{lll}
\mathbb D_{\rm bd}   &:=& \{(N,\kappa,G):H=G(N+\kappa N(N-1))\ge 16\},\\[2mm]
\mathbb D_{\rm src}  &:=& \{\text{two-sided local source/relay typed charts}\},\\[2mm]
\mathbb D_{\rm seam} &:=& \{\Sigma_{\rm trans}=0\},\\[2mm]
\mathbb D_{\rm mix}  &:=& \{P_{\partial}S_{4,\mu}^{\rm mixker}=0,\\[-1mm]
&&\hspace{21mm}\beta_4=1/256\},\\[2mm]
\mathbb D_{\rm col}  &:=& \{\Phi_5^{(n+1)}=(3/4)\Phi_5^{(n)},\ n<\infty\},\\[2mm]
\mathbb D_{\rm clk}  &:=& \{\mathcal C_{\gamma}(\tau_v)\ge 16/S(N,\kappa)\}.
\end{array}
\]
Here
\[
S(N,\kappa)=N+\kappa N(N-1).
\]
The domain product is not a phase space and not a space of physical field
configurations.  It is only the declared local domain on which the six
admissibility-only monitors are well-typed.

\subsection{Codomain and Residual Convention}
\label{subsec:r50-codomain-residual-convention}

Each tag maps into a monitor codomain, not into a dynamical equation space:
\[
\mathcal O_x:
\mathbb D_x\longrightarrow \mathbb M_x,
\qquad
x\in\{\mathrm{bd},\mathrm{src},\mathrm{seam},\mathrm{mix},\mathrm{col},\mathrm{clk}\}.
\]
The codomains are declared as
\[
\begin{array}{lll}
\mathbb M_{\rm bd}   &:=& \mathbb R_{\ge0}\times\mathbb R_{\ge0},\\
\mathbb M_{\rm src}  &:=& \{0,1\}_{\rm type},\\
\mathbb M_{\rm seam} &:=& \mathbb R_{\ge0},\\
\mathbb M_{\rm mix}  &:=& \mathbb R_{\ge0}\times\mathbb R_{\ge0},\\
\mathbb M_{\rm col}  &:=& \mathbb R,\\
\mathbb M_{\rm clk}  &:=& \mathbb R_{\ge0}.
\end{array}
\]
The boundary-relay monitor records
\[
(E_{\rm rel}^{-},E_{\rm src}^{-}),
\qquad
E_{\rm rel}^{-}=\max(0,16-H),
\qquad
E_{\rm src}^{-}=\max(0,16-H_{\partial}).
\]
The signed quantity
\[
E_{\rm rel}=H-16
\]
remains a frontier monitor only.  It is not a debt term and must not be squared
as a minimization target.

\subsection{Operator-Domain Table}
\label{subsec:r50-operator-domain-table}

The r50 operator-domain table is
\[
\begin{array}{c|c|c|c}
\text{tag} & \text{domain} & \text{codomain} & \text{status}\\
\hline
\mathcal O_{\rm bd} & \mathbb D_{\rm bd} & \mathbb M_{\rm bd} & \text{local monitor}\\
\mathcal O_{\rm src} & \mathbb D_{\rm src} & \mathbb M_{\rm src} & \text{type monitor}\\
\mathcal O_{\rm seam} & \mathbb D_{\rm seam} & \mathbb M_{\rm seam} & \text{hard seam guard}\\
\mathcal O_{\rm mix} & \mathbb D_{\rm mix} & \mathbb M_{\rm mix} & \text{C4 guard-lock}\\
\mathcal O_{\rm col} & \mathbb D_{\rm col} & \mathbb M_{\rm col} & \text{finite recurrence monitor}\\
\mathcal O_{\rm clk} & \mathbb D_{\rm clk} & \mathbb M_{\rm clk} & \text{selector/probe monitor}
\end{array}
\]
The table is an admissibility-domain refinement only.  None of these tags is a
physical differential operator, Hamiltonian, action-gradient, stress tensor,
field equation, or global readout map.

\subsection{Two-Sided Domain Typing}
\label{subsec:r50-two-sided-domain-typing}

The two-sided typing convention from r49 is preserved.  Schematically,
\[
{}^{L_x}{}_{G_x}\mathcal O_x{}^{R_x}_{M_x}.
\]
This takes a local readout/source/guard-facing domain on the left/front side and a
relay/structure/monitor-facing domain on the right/back side.  The intended
readings are
\[
\begin{array}{lll}
L_x&=&\text{local readout/source side},\\
G_x&=&\text{guard or selector side},\\
R_x&=&\text{relay/structure side},\\
M_x&=&\text{monitor codomain side}.
\end{array}
\]
Upper--lower contractions retain the classical action/readout meaning.  Same
variance comparisons remain forbidden unless an explicit comparison tensor is
declared.  Thus the r50 domain refinement does not loosen the tensor grammar.

\subsection{Domain Residual Vector}
\label{subsec:r50-domain-residual-vector}

The local operator-domain residual vector is
\[
\mathcal R_{\rm dom}^{r50}
=
(r_{\rm bd},r_{\rm src},r_{\rm seam},r_{\rm mix},r_{\rm col},r_{\rm clk}).
\]
For the comfort-calibrated local seed
\[
(N,\kappa,G)=\left(5,{3\over4},1\right)
\]
one obtains
\[
S=20,
\qquad
H=20,
\qquad
E_{\rm rel}^{-}=0,
\qquad
E_{\rm src}^{-}=0.
\]
Together with the already certified seam, mix, collar, and clock monitors this
gives
\[
\boxed{
\mathcal R_{\rm dom}^{r50}=(0,0,0,0,0,0).
}
\]
The associated domain-admissibility vector is
\[
\boxed{
\mathsf D_{\rm adm}^{r50}=(1,1,1,1,1,1).
}
\]

\subsection{No-Promotion and Guard Preservation}
\label{subsec:r50-no-promotion-guard}

The no-promotion vector remains
\[
\boxed{
\Pi_{\rm np}^{r50}=(0,0,0,0,0,0,0).
}
\]
The entries mean: no action, no Lagrangian, no Euler--Lagrange variation, no
physical field equation, no global readout promotion, no C4/C5/C6 spending, and
no RH promotion.

The guard vector remains
\[
\boxed{
\Gamma_{\rm guard}^{r50}=(1,1,1,1,1).
}
\]
Thus
\[
\boxed{
\beta_4={1\over256}\text{ remains unspent},
\quad C5\text{ has no finite closure},
\quad C6/T2\text{ remains blocked},
\quad \mathrm{RH}\text{ remains photon-clock probe only},
\quad \text{global all-readout closure remains unclaimed}.
}
\]
No new DOI, PDF, release packet, public upload, physical action, or field
equation is authorized.

\subsection{r50 Decision and Next Step}
\label{subsec:r50-decision-next-step}

The r50 decision ledger is
\[
\boxed{
\mathsf D_{r50}
=
\{
W_2\text{ complete},
\mathcal R_{\rm dom}=(0,0,0,0,0,0),
\mathsf D_{\rm adm}=(1,1,1,1,1,1),
\mathrm{PDF}=\mathrm{not\ built},
\mathrm{new\ DOI}=\varnothing,
\mathrm{public\ release}=\mathrm{false},
\Gamma_{\rm guard}=(1,1,1,1,1)
\}.
}
\]
The safe next source-only step is
\[
\boxed{
\text{v4.7.0-r51: local pre-field interface diagram in text form.}
}
\]


\section{v4.7.0-r51 Source-Only Local Pre-Field Interface Diagram}
\label{sec:r51-prefield-interface-diagram}

\subsection{Source-Only Status and Scope}
\label{subsec:r51-source-only-status}

This source-only layer executes
\[
\boxed{W_3^{r48}:\ \text{local pre-field interface diagram in text form}.}
\]
No PDF is built, no new DOI is assigned, no public upload is authorized, and no
physical action, Lagrangian, Euler--Lagrange variation, derived field equation,
T2 bridge, RH theorem, or global all-readout closure is claimed.

The r50 admissibility-only operator-domain refinement remains the immediate
input.  The present layer only gives the reader a textual wiring diagram for the
already certified local monitor skeleton.  It is a routing and typing diagram,
not a dynamical diagram.

\subsection{Textual Pre-Field Interface Diagram}
\label{subsec:r51-textual-diagram}

The local pre-field interface is recorded as the following guarded pipeline:
\[
\boxed{
\begin{array}{ccccccc}
\mathfrak H_{\rm seed}^{r34}
&\xrightarrow{\;\mathsf I_{\rm bd}\;}&
\mathfrak F_{\rm loc}^{r38}
&\xrightarrow{\;\mathsf I_{\rm src}\;}&
\mathfrak S_{\rm loc}^{r39}
&\xrightarrow{\;\mathsf I_{\rm seam}\;}&
\mathfrak n^{r41}
\end{array}}
\]
followed by the guard-locked monitor tail
\[
\boxed{
\begin{array}{ccccc}
\mathfrak n^{r41}
&\xrightarrow{\;\mathsf I_{\rm mix}\;}&
\mathfrak C_{\rm int}^{r37}
&\xrightarrow{\;\mathsf I_{\rm col}\;}&
\mathfrak C_{\rm mon}^{r42}
\xrightarrow{\;\mathsf I_{\rm clk}\;}
\mathfrak H_{\rm cert}^{r43}.
\end{array}}
\]
The two displayed rows are not two independent proofs.  They are a textual
reader diagram for one local admissibility-monitor path:
\[
\boxed{
\mathfrak H_{\rm seed}^{r34}
\longrightarrow
\mathfrak F_{\rm loc}^{r38}
\longrightarrow
\mathfrak S_{\rm loc}^{r39}
\longrightarrow
\mathfrak n^{r41}
\longrightarrow
\mathfrak C_{\rm int}^{r37}
\longrightarrow
\mathfrak C_{\rm mon}^{r42}
\longrightarrow
\mathfrak H_{\rm cert}^{r43}.
}
\]
The metadata-freeze support vector \(\mathsf F_{\rm meta}^{r43}\) remains a side
root rather than a dynamical input:
\[
\boxed{
\mathsf F_{\rm meta}^{r43}
\dashrightarrow
\mathfrak H_{\rm cert}^{r43}
\quad\text{as source-only metadata support only}.}
\]

\subsection{Diagram Node Table}
\label{subsec:r51-node-table}

\begin{longtable}{P{0.20\linewidth}P{0.27\linewidth}P{0.43\linewidth}}
\toprule
Node & Role & Guarded meaning \\
\midrule
\endhead
\(\mathfrak H_{\rm seed}^{r34}\) & comfort-calibrated local seed & Contains the selected local anchor \((N,\kappa,G)=(5,3/4,1)\), with \(H=20\) and positive margin. \\
\(\mathfrak F_{\rm loc}^{r38}\) & local pre-field interface & The certificate is admitted only after local guard projection; not a field equation. \\
\(\mathfrak S_{\rm loc}^{r39}\) & admissibility-only skeleton & Six monitor tags exist as local admissibility tags only. \\
\(\mathfrak n^{r41}\) & guarded composition normal form & Ordered tag compositions collapse to \(\mathsf{NF}_{\rm loc}\), \(\mathsf{NF}_{\rm guard}\), or \(\mathsf{NF}_{\rm probe}\). \\
\(\mathfrak C_{\rm int}^{r37}\) & local MM/HC interface certificate & Local residuals vanish, while guard residuals remain separated. \\
\(\mathfrak C_{\rm mon}^{r42}\) & compressed local monitor certificate & Monitor evidence is compressed without promoting to action or field dynamics. \\
\(\mathfrak H_{\rm cert}^{r43}\) & source-only handoff packet & The local certificate is handed forward with frozen source-only metadata. \\
\(\mathsf F_{\rm meta}^{r43}\) & metadata side root & Records baseline DOI, no new DOI, no PDF, and no public release. \\
\bottomrule
\end{longtable}

\subsection{Diagram Edge Table}
\label{subsec:r51-edge-table}

The local diagram edges are typed as monitor transports rather than as physical
propagators:
\begin{longtable}{P{0.18\linewidth}P{0.26\linewidth}P{0.46\linewidth}}
\toprule
Edge & Tag & Meaning \\
\midrule
\endhead
\(e_{01}\) & \(\mathsf I_{\rm bd}\) & boundary relay admissibility: \(H\ge16\), debt only through \(E_{\rm rel}^{-}\). \\
\(e_{12}\) & \(\mathsf I_{\rm src}\) & two-sided local source/relay typing: left source/readout, right relay/structure. \\
\(e_{23}\) & \(\mathsf I_{\rm seam}\) & compressed-atlas seam exactness: \(\Sigma_{\rm trans}=0\). \\
\(e_{34}\) & \(\mathsf I_{\rm mix}\) & C4 mixed-kernel guard-lock: \(P_{\partial}S_{4,\mu}^{\rm mixker}=0\), \(\beta_4=1/256\). \\
\(e_{45}\) & \(\mathsf I_{\rm col}\) & finite C5 collar recurrence monitor: \(\Phi_5^{(n+1)}=3\Phi_5^{(n)}/4\). \\
\(e_{56}\) & \(\mathsf I_{\rm clk}\) & photon-clock selector/probe window: \(\mathcal C_{\gamma}(\tau_v)\ge16/S(N,\kappa)\). \\
\(e_{m}\) & metadata support & source-only metadata support: no DOI, no PDF, no public release. \\
\bottomrule
\end{longtable}

\subsection{Pre-Field Diagram Residuals}
\label{subsec:r51-prefield-diagram-residuals}

The textual diagram is valid exactly because every local monitor edge is already
admissible in the r50 domain refinement.  The diagram residual vector is
\[
\boxed{
\mathcal R_{\rm diag}^{r51}=(0,0,0,0,0,0,0).
}
\]
The first six entries correspond to \(e_{01},\ldots,e_{56}\).  The seventh entry
is the metadata support edge.  The diagram-admissibility vector is
\[
\boxed{
\mathsf D_{\rm diag}^{r51}=(1,1,1,1,1,1,1).
}
\]
The local diagram therefore supports reader navigation and local interface
handoff, but it does not create a field equation, a variational principle, or a
global readout theorem.

\subsection{No-Action, No-Promotion, and Guard Preservation}
\label{subsec:r51-no-action-guard}

The no-action vector remains
\[
\boxed{
\mathfrak A_{\rm noact}^{r51}=(0,0,0,0,0,0).
}
\]
The no-promotion vector remains
\[
\boxed{
\Pi_{\rm np}^{r51}=(0,0,0,0,0,0,0).
}
\]
The guard vector remains
\[
\boxed{
\Gamma_{\rm guard}^{r51}=(1,1,1,1,1).
}
\]
Thus
\[
\boxed{
\beta_4={1\over256}\text{ remains unspent},
\quad C5\text{ has no finite closure},
\quad C6/T2\text{ remains blocked},
\quad \mathrm{RH}\text{ remains photon-clock probe only},
\quad \text{global all-readout closure remains unclaimed}.
}
\]

\subsection{r51 Decision and Next Step}
\label{subsec:r51-decision-next-step}

The r51 decision ledger is
\[
\boxed{
\mathsf D_{r51}
=
\{
W_3\text{ complete},
\mathcal R_{\rm diag}=(0,0,0,0,0,0,0),
\mathsf D_{\rm diag}=(1,1,1,1,1,1,1),
\mathrm{PDF}=\mathrm{not\ built},
\mathrm{new\ DOI}=\varnothing,
\mathrm{public\ release}=\mathrm{false},
\Gamma_{\rm guard}=(1,1,1,1,1)
\}.
}
\]
The safe next source-only step is
\[
\boxed{
\text{v4.7.0-r53: no-action/no-promotion audit refresh.}
}
\]



\section{v4.7.0-r52 Source-Only No-Action/No-Promotion Audit Refresh}
\label{sec:r52-no-action-no-promotion-audit-refresh}

\subsection{Source-Only Status}
\label{subsec:r52-source-only-status}

This r52 layer executes \(W_4^{r48}\): the no-action and no-promotion audit refresh after the r51 textual local pre-field interface diagram.  The layer is source-only.  It creates no PDF, assigns no new DOI, authorizes no public release, and does not alter the v4.6.0 release-final baseline DOI
\[
\boxed{10.5281/zenodo.20117361.}
\]
The purpose is only to refresh the audit boundary between the admissibility-only local monitor stack and any later field-interface work.

\subsection{Audit Refresh Objects}
\label{subsec:r52-audit-refresh-objects}

The refreshed input stack is
\[
\boxed{
\mathfrak X_{\rm audit}^{r52}
=
\left(
\mathfrak H_{\rm seed}^{r34},
\mathfrak F_{\rm loc}^{r38},
\mathfrak S_{\rm loc}^{r39},
\mathfrak n^{r41},
\mathfrak C_{\rm int}^{r37},
\mathfrak C_{\rm mon}^{r42},
\mathfrak H_{\rm cert}^{r43},
\Gamma_{\rm guard}^{r51}
\right).
}
\]
The refreshed no-action vector is
\[
\boxed{
\mathfrak A_{\rm noact}^{r52}
=(0,0,0,0,0,0).
}
\]
Its entries mean respectively: no physical action, no Lagrangian, no Euler--Lagrange variation, no derived field equation, no global readout promotion, and no guard spending.

The refreshed no-promotion vector is
\[
\boxed{
\Pi_{\rm np}^{r52}
=(0,0,0,0,0,0,0).
}
\]
Its entries mean respectively: no local-monitor-to-action promotion, no action-to-field-equation promotion, no local-to-global readout promotion, no \(C4\) closure promotion, no finite \(C5\) closure promotion, no \(C6/T2\) promotion, and no RH theorem promotion.

\subsection{Guard Preservation Refresh}
\label{subsec:r52-guard-preservation-refresh}

The guard vector remains
\[
\boxed{
\Gamma_{\rm guard}^{r52}=(1,1,1,1,1).
}
\]
Equivalently,
\[
\boxed{
\beta_4={1\over256}\text{ remains unspent},
\quad C5\text{ has no finite closure},
\quad C6/T2\text{ remains blocked},
\quad \mathrm{RH}\text{ remains photon-clock probe only},
\quad \text{global all-readout closure remains unclaimed}.
}
\]
The local pre-field interface diagram may therefore be used as a reader-routing and admissibility-monitor object, but it still does not license a physical field equation.

\subsection{Refreshed Audit Ledger}
\label{subsec:r52-refreshed-audit-ledger}

The refreshed audit ledger is
\[
\boxed{
\mathsf A_{\rm refresh}^{r52}
=
\left(
\mathfrak A_{\rm noact}^{r52},
\Pi_{\rm np}^{r52},
\Gamma_{\rm guard}^{r52},
\mathcal R_{\rm diag}^{r51},
\mathsf D_{\rm diag}^{r51}
\right).
}
\]
The refresh residual is
\[
\boxed{
\Delta_{\rm audit}^{r52}=0.
}
\]
Thus r52 completes \(W_4^{r48}\) without changing the closure status of any global or guarded object.

\subsection{r52 Decision and Next Step}
\label{subsec:r52-decision-next-step}

The r52 decision ledger is
\[
\boxed{
\mathsf D_{r52}
=
\{
W_4\text{ complete},
\mathfrak A_{\rm noact}=(0,0,0,0,0,0),
\Pi_{\rm np}=(0,0,0,0,0,0,0),
\Gamma_{\rm guard}=(1,1,1,1,1),
\Delta_{\rm audit}=0,
\mathrm{PDF}=\mathrm{not\ built},
\mathrm{new\ DOI}=\varnothing,
\mathrm{public\ release}=\mathrm{false}
\}.
}
\]
The safe next source-only step is
\[
\boxed{
\text{v4.7.0-r53: guarded local-interface consolidation and post-W1--W4 branch checkpoint.}
}
\]



\clearpage
\section{v4.7.0-r53 Source-Only Guarded Local-Interface Consolidation and Post-W1--W4 Branch Checkpoint}
\label{sec:r53-local-interface-consolidation-checkpoint}

\subsection{Source-Only Scope and Baseline}
\label{subsec:r53-source-only-scope-baseline}

This section records an internal source-only consolidation layer.  It does not build a PDF, does not assign a new DOI, does not authorize a public release, and does not change the v4.6.0 baseline DOI
\[
\boxed{10.5281/zenodo.20117361.}
\]
The purpose of r53 is to consolidate the completed r48 work-package sequence
\[
W_1\prec W_2\prec W_3\prec W_4
\]
after r49--r52, and to freeze the local-interface branch state before opening the next internal source-only task.

\subsection{Completed Work-Package Stack}
\label{subsec:r53-completed-work-package-stack}

The r48 branch queue selected the guarded local field-interface elaboration branch.  The completed branch stack is
\[
\boxed{
\mathsf W_{Q1}^{r53}
=
\left(
W_1^{r49},
W_2^{r50},
W_3^{r51},
W_4^{r52}
\right).
}
\]
Here
\[
\begin{array}{ll}
W_1^{r49}:& \text{interface-slot glossary and two-sided typing recap},\\
W_2^{r50}:& \text{admissibility-only operator-domain refinement},\\
W_3^{r51}:& \text{local pre-field interface diagram in text form},\\
W_4^{r52}:& \text{no-action/no-promotion audit refresh}.
\end{array}
\]
The completion vector is
\[
\boxed{
\mathsf W_{Q1}^{r53}=(1,1,1,1),
\qquad
\Delta_{Q1}^{r53}=0.
}
\]
Thus the first r48 branch package is complete as a local, source-only, guard-preserving interface layer.

\subsection{Consolidated Local-Interface Object}
\label{subsec:r53-consolidated-local-interface-object}

The consolidated local-interface object is
\[
\boxed{
\mathfrak L_{\rm int}^{r53}
=
\left(
\mathcal R_{\rm slot}^{r49},
\mathsf D_{\rm adm}^{r50},
\mathcal R_{\rm diag}^{r51},
\mathfrak A_{\rm noact}^{r52},
\Pi_{\rm np}^{r52},
\Gamma_{\rm guard}^{r52}
\right).
}
\]
Its local residual package is
\[
\boxed{
\mathcal R_{\rm cons}^{r53}
=(0,0,0,0,0,0).
}
\]
The entries record respectively: slot residual zero, domain residual zero, diagram residual zero, no-action residual zero, no-promotion residual zero, and guard residual zero.

\subsection{Guarded Branch Checkpoint}
\label{subsec:r53-guarded-branch-checkpoint}

The post-W1--W4 branch checkpoint is
\[
\boxed{
\mathsf{BC}_{Q1}^{r53}
=
\left(
\mathsf W_{Q1}^{r53},
\mathfrak L_{\rm int}^{r53},
\mathcal R_{\rm cons}^{r53},
\Gamma_{\rm guard}^{r52}
\right).
}
\]
The guard vector remains
\[
\boxed{
\Gamma_{\rm guard}^{r53}=(1,1,1,1,1).
}
\]
Equivalently,
\[
\boxed{
\beta_4={1\over256}\text{ remains unspent},\quad
C5\text{ has no finite closure},\quad
C6/T2\text{ remains blocked},\quad
\mathrm{RH}\text{ remains photon-clock probe only},\quad
\text{global all-readout closure remains unclaimed}.
}
\]
No guard entry is consumed by the local-interface consolidation.

\subsection{No-Action and No-Promotion Confirmation}
\label{subsec:r53-no-action-no-promotion-confirmation}

The r53 confirmation vectors are
\[
\boxed{
\mathfrak A_{\rm noact}^{r53}=(0,0,0,0,0,0),
\qquad
\Pi_{\rm np}^{r53}=(0,0,0,0,0,0,0).
}
\]
Thus r53 still asserts no physical action, no Lagrangian, no Euler--Lagrange variation, no derived field equation, no global readout promotion, and no guard spending.  The local-interface object is an admissibility and reader-routing layer only.

\subsection{r53 Decision and Next Step}
\label{subsec:r53-decision-next-step}

The r53 decision ledger is
\[
\boxed{
\mathsf D_{r53}
=
\{
Q1\text{ consolidated},
W_1\text{--}W_4\text{ complete},
\mathcal R_{\rm cons}=0,
\Gamma_{\rm guard}=(1,1,1,1,1),
\mathrm{PDF}=\mathrm{not\ built},
\mathrm{new\ DOI}=\varnothing,
\mathrm{public\ release}=\mathrm{false}
\}.
}
\]
The safe next source-only step is
\[
\boxed{
\text{v4.7.0-r54: local-interface theoremlet extraction and next-branch selector preflight.}
}
\]


\clearpage
\section{v4.7.0-r54 Source-Only Local-Interface Theoremlet Extraction and Next-Branch Selector Preflight}
\label{sec:r54-local-interface-theoremlet-selector-preflight}

\subsection{Source-Only Scope and Non-Release Clause}
\label{subsec:r54-source-only-scope-non-release}

This section records an internal source-only theoremlet-extraction layer after the r53 guarded local-interface consolidation.  It does not build a PDF, does not assign a new DOI, does not authorize a public release, and does not change the v4.6.0 baseline DOI
\[
\boxed{10.5281/zenodo.20117361.}
\]
The task is to decide whether the completed local-interface stack supports a bounded theoremlet, and to select the next internal source-only branch without promoting the result to a physical field equation or a global closure theorem.

\subsection{Extracted Local-Interface Theoremlet}
\label{subsec:r54-extracted-local-interface-theoremlet}

The r53 branch checkpoint provides
\[
\mathsf W_{Q1}^{r53}=(1,1,1,1),
\qquad
\Delta_{Q1}^{r53}=0,
\qquad
\mathcal R_{\rm cons}^{r53}=(0,0,0,0,0,0),
\qquad
\Gamma_{\rm guard}^{r53}=(1,1,1,1,1).
\]
Hence the following bounded theoremlet is admissible.

\paragraph{Theoremlet r54.1: Guarded Local-Interface Admissibility.}
If the r49 interface-slot glossary, the r50 admissibility-only operator-domain refinement, the r51 local pre-field interface diagram, and the r52 no-action/no-promotion audit refresh all pass with zero local residual, then the consolidated object
\[
\mathfrak L_{\rm int}^{r53}
=\left(
\mathcal R_{\rm slot}^{r49},
\mathsf D_{\rm adm}^{r50},
\mathcal R_{\rm diag}^{r51},
\mathfrak A_{\rm noact}^{r52},
\Pi_{\rm np}^{r52},
\Gamma_{\rm guard}^{r52}
\right)
\]
is a local admissibility interface object.  Equivalently,
\[
\boxed{
\mathfrak L_{\rm int}^{r53}\in
\mathsf{Adm}_{\rm loc}^{\rm interface}
}
\]
with
\[
\boxed{
\mathcal R_{\rm thm}^{r54}=(0,0,0,0,0,0),
\qquad
\Gamma_{\rm guard}^{r54}=(1,1,1,1,1).
}
\]
The theoremlet is local, source-only, and interface-level.  It is not an action principle, not a Lagrangian, not an Euler--Lagrange statement, not a derived field equation, and not a global readout theorem.

\subsection{Boundary of the Theoremlet}
\label{subsec:r54-boundary-of-theoremlet}

The theoremlet closes only the local interface admissibility question.  Its closure-scope vector is
\[
\boxed{
\mathsf S_{\rm thm}^{r54}=(1,0,0,0,0,0).
}
\]
The entries mean:
\[
\begin{array}{c|c|c}
\text{slot} & \text{meaning} & \text{status}\\
\hline
1 & \text{local interface admissibility} & 1\\
2 & \text{physical action / Lagrangian} & 0\\
3 & \text{derived field equation} & 0\\
4 & \text{global readout} & 0\\
5 & C6/T2 & 0\\
6 & \text{RH / all-readout closure} & 0
\end{array}
\]
Thus r54 extracts a real local theoremlet, but it does not solve the global problem.

\subsection{Next-Branch Selector Preflight}
\label{subsec:r54-next-branch-selector-preflight}

The r48 branch queue remains available after Q1 completion.  The selector evaluates the remaining branch families as follows:
\[
\begin{array}{c|l|c|l}
\text{branch} & \text{candidate} & \text{admissible} & \text{reason}\\
\hline
Q_2 & \text{local theoremlet-to-lemma packaging} & 1 & \text{uses r54 theoremlet, no guard spending}\\
Q_3 & \text{field-equation ansatz promotion} & 0 & \text{would violate no-action/no-field guard}\\
Q_4 & \text{C4 mixed-kernel spending attempt} & 0 & \beta_4=1/256\text{ remains unspent}\\
Q_5 & \text{all-readout/T2 branch} & 0 & C6/T2\text{ remains blocked}
\end{array}
\]
The selector vector is therefore
\[
\boxed{
\mathsf{SEL}_{\rm branch}^{r54}=(1,0,0,0),
\qquad
\Delta_{\rm sel}^{r54}=0.
}
\]
The selected next branch is
\[
\boxed{
Q_2:\ \text{local theoremlet-to-lemma packaging and guard-stable lemma registry.}
}
\]

\subsection{No-Action, No-Promotion, and Guard Audit}
\label{subsec:r54-no-action-no-promotion-guard-audit}

The r54 audit vectors are
\[
\boxed{
\mathfrak A_{\rm noact}^{r54}=(0,0,0,0,0,0),
\qquad
\Pi_{\rm np}^{r54}=(0,0,0,0,0,0,0).
}
\]
No physical action, no Lagrangian, no Euler--Lagrange variation, no derived field equation, no global readout promotion, and no guard spending occurs in r54.

The guard vector remains
\[
\boxed{
\Gamma_{\rm guard}^{r54}=(1,1,1,1,1).
}
\]
Equivalently,
\[
\boxed{
\beta_4={1\over256}\text{ remains unspent},\quad
C5\text{ has no finite closure},\quad
C6/T2\text{ remains blocked},\quad
\mathrm{RH}\text{ remains photon-clock probe only},\quad
\text{global all-readout closure remains unclaimed}.
}
\]

\subsection{r54 Decision}
\label{subsec:r54-decision}

The r54 decision is
\[
\boxed{
\mathsf D_{r54}
=
\{
\text{local-interface theoremlet found},
\text{global solution not found},
Q_2\text{ selected},
\Gamma_{\rm guard}=(1,1,1,1,1),
\mathrm{PDF}=\mathrm{not\ built},
\mathrm{new\ DOI}=\varnothing,
\mathrm{public\ release}=\mathrm{false}
\}.
}
\]
The safe next source-only step is
\[
\boxed{
\text{v4.7.0-r55: local theoremlet-to-lemma packaging and guard-stable lemma registry.}
}
\]



% =============================================================
% v4.7.0 post-r54 guarded solve chain appended for RC9 full monolith recovery
% =============================================================


% ---- appended source checkpoint r64 from Quantum_Sphaera_Companion_v4_7_0_r64_detector_threshold_sufficiency.tex ----
% Quantum Sphaera Companion -- v4.7.0-r62 source-only
% Admissible terminal-rule search for finite C5 closure.
% Extends r61 finite C5 closure-replacement search under r60+ conservative relative-seam extension.
% No PDF generated. No new DOI. No public release.



\clearpage
